% from ega3-1-fr.tex
\oldpage[0\textsubscript{III}]{79}
(déduits de
$\mathfrak J\otimes_S\mathrm F^p(M)\to\mathrm F^{p+1}(M)$, compte tenu de ce
que l'image de l'homomorphisme composé
\[
\mathfrak J\otimes_S\mathrm F^{p+1}(M)\longrightarrow
\mathfrak J\otimes_S\mathrm F^p(M)\longrightarrow\mathrm F^{p+1}(M)
\]
est $\mathfrak J\mathrm F^{p+1}(M)\subset\mathrm F^{p+2}(M)$) sont
surjectifs.
\end{enumerate}
Pour la démonstration, voir Bourbaki, \emph{Alg. comm.}, chap.~III, \S~3,
n\textsuperscript{o}~1, prop.~3.
\end{lemma}

\begin{env}[13.7.7.4]
\label{0.13.7.7.4-fr}
Pour appliquer le lemme (\hyperref[0.13.7.7.2-fr]{13.7.7.2}), il reste à
observer que la topologie définie sur
$\mathrm R'^{\,n}\mathrm T(\mathbf A)$ par la filtration considérée fait de
$\mathrm R'^{\,n}\mathrm T(\mathbf A)$ un $S$-module séparé et complet,
cette topologie étant celle de la limite projective des groupes discrets
$\mathrm R^n\mathrm T(A_k)$~; d'autre part, si $A_k=0$ pour $k<k_0$, on a
aussi $\mathrm R^n\mathrm T(A_k)=0$ pour $k<k_0$, donc
$\mathrm F^{k_0}(\mathrm R'^{\,n}\mathrm T(\mathbf A))=
\mathrm R'^{\,n}\mathrm T(\mathbf A)$, et on est bien dans les conditions
d'application du lemme.
\end{env}

\begin{corollary}[13.7.8]
\label{0.13.7.8-fr}
Si l'hypothèse $(\mathrm F_n)$ est vérifiée, on a, pour tout élément $f\in S$,
un isomorphisme canonique
\begin{equation}
\label{0.13.7.8.1-fr}
\tag{13.7.8.1}
\varprojlim_k\bigl((\mathrm R^n\mathrm T(A_k))_f\bigr)
\xrightarrow{\sim}\mathrm R^n\mathrm T(\mathbf A)\otimes_S S_{\{f\}}.
\end{equation}
\end{corollary}

\begin{proof}
En effet, $\mathrm R^n\mathrm T(\mathbf A)$ est un $S$-module de type fini,
$S_{\{f\}}$ une $S$-algèbre adique noethérienne
(\hyperref[0.7.6.11-fr]{0\textsubscript{I}, 7.6.11}), séparée complétée de
$S_f$ pour la topologie $\mathfrak J$-préadique
(\hyperref[0.7.6.2-fr]{0\textsubscript{I}, 7.6.2}). On conclut de
(\hyperref[0.7.7.8-fr]{0\textsubscript{I}, 7.7.8}) et
(\hyperref[0.7.7.1-fr]{0\textsubscript{I}, 7.7.1}) que
$\mathrm R^n\mathrm T(\mathbf A)\otimes_S S_{\{f\}}$ est isomorphe au séparé
complété de $\mathrm R^n\mathrm T(\mathbf A)\otimes_S S_f$ pour la topologie
$\mathfrak J$-préadique~; un système fondamental de voisinages de $0$ pour
cette topologie est
$(\mathfrak J^p\mathrm R^n\mathrm T(\mathbf A))\otimes_S S_f$, donc
$\mathrm F^p(\mathrm R^n\mathrm T(\mathbf A))\otimes_S S_f$ est aussi un tel
système~; ce dernier est le noyau de l'application canonique
$(\mathrm R^n\mathrm T(\mathbf A))_f\to
(\mathrm R^n\mathrm T(A_{p-1}))_f$ et par suite le groupe séparé associé à
$\mathrm R^n\mathrm T(\mathbf A)\otimes_S S_f$ s'identifie à un sous-groupe
$G$ de $\varprojlim_k((\mathrm R^n\mathrm T(A_k))_f)$. Mais le système
projectif $((\mathrm R^n\mathrm T(A_k))_f)$ vérifie évidemment la condition
(ML), et l'image de $(\mathrm R^n\mathrm T(\mathbf A))_f$ dans chacun des
$(\mathrm R^n\mathrm T(A_k))_f$ est égale à l'image commune des
$(\mathrm R^n\mathrm T(A_h))_f$ pour $h\geq k$ assez grand. On en conclut
aussitôt que $G$ est \emph{partout dense} dans
$\varprojlim_k((\mathrm R^n\mathrm T(A_k))_f)$, et comme ce dernier groupe est
séparé et complet, le corollaire est démontré.
\end{proof}

\begin{flushright}
\textit{(A suivre.)}
\end{flushright}

\clearpage
\thispagestyle{empty}
\null
\clearpage
\thispagestyle{plain}
\markboth{ÉTUDE COHOMOLOGIQUE DES FAISCEAUX COHÉRENTS}{CHAPITRE III}
\oldpage[III]{81}
\phantomsection
\label{chapter:III-fr}

\begin{center}
{\large CHAPITRE III\par}
\vspace{1.5em}
{\Large\bfseries ÉTUDE COHOMOLOGIQUE\par}
{\Large\bfseries DES FAISCEAUX COHÉRENTS\par}
\vspace{1.8em}
{\bfseries Sommaire\par}
\end{center}

\medskip
\begin{tabular}{@{}r@{\ }p{0.82\textwidth}@{}}
\S~1. & Cohomologie des schémas affines.\\
\S~2. & Étude cohomologique des morphismes projectifs.\\
\S~3. & Le théorème de finitude pour les morphismes propres.\\
\S~4. & Le théorème fondamental des morphismes propres~; applications.\\
\S~5. & Un théorème d'existence de faisceaux algébriques cohérents.\\
\S~6. & Foncteurs Tor et Ext locaux et globaux~; formule de Künneth.\\
\S~7. & Étude du changement de base dans les foncteurs cohomologiques
covariants de Modules.\\
\S~8. & Le théorème de dualité sur les fibrés projectifs.\\
\S~9. & Cohomologie relative et cohomologie locale~; dualité locale.\\
\S~10. & Relations entre cohomologie projective et cohomologie locale.
Technique de complétion formelle le long d'un diviseur.\\
\S~11. & Groupes de Picard globaux et locaux\footnotemark.
\end{tabular}

\medskip
Ce chapitre donne les théorèmes fondamentaux sur la cohomologie des faisceaux
algébriques cohérents, à l'exception de ceux découlant de la théorie des
résidus (théorèmes de dualité), qui feront l'objet d'un chapitre ultérieur.
Parmi les premiers, il y a essentiellement six théorèmes fondamentaux,
faisant l'objet des six premiers paragraphes du présent chapitre. Ces
résultats seront dans la suite des outils essentiels, même dans des questions
qui ne sont pas de nature proprement cohomologique~; le lecteur en verra les
premiers exemples dès le \S~4. Le \S~7 donne des résultats de nature plus
technique, mais d'un usage constant dans les applications. Enfin, dans les
\S\S~8 à 11, nous développerons certains résultats, liés à la dualité des
faisceaux cohérents, particulièrement importants pour les applications, et
qui peuvent s'exposer antérieurement à la théorie générale des résidus.

\footnotetext{Le chapitre IV ne dépend pas des \S\S~8 à 11, et sera sans doute
publié avant ces derniers.}

\oldpage[III]{82}
Le contenu des \S\S~1 et 2 est dû à J.-P. Serre, et le lecteur constatera que
nous n'avons eu qu'à suivre l'exposé de (FAC). Les \S\S~8 et 9 sont également
inspirés par (FAC) (les transpositions nécessitées par les contextes
différents étant toutefois moins évidentes). Enfin, comme nous l'avons dit
dans l'Introduction, le \S~4 doit être considéré comme la mise en forme, en
langage moderne, du «~théorème d'invariance~» fondamental de la «~théorie des
fonctions holomorphes~» de Zariski.

Signalons enfin que les résultats du n\textsuperscript{o}~3.4 (et les
propositions préliminaires de (0, 13.4 à 13.7)) ne seront pas utilisés dans la
suite du chap.~III et peuvent donc être omis en première lecture.

\setcounter{section}{0}
\section{Cohomologie des schémas affines}
\label{section:III.1-fr}

\subsection{Rappels sur le complexe de l'algèbre extérieure}
\label{subsection:III.1.1-fr}

\begin{env}[1.1.1]
\label{III.1.1.1-fr}
Soient $A$ un anneau, $\mathbf f=(f_i)_{1\leq i\leq r}$ un système de $r$
éléments de $A$. Le \emph{complexe de l'algèbre extérieure
$K_\bullet(\mathbf f)$} correspondant à $\mathbf f$ est un complexe de chaînes
(G, I, 2.2) se définissant de la façon suivante~: le $A$-module gradué
$K_\bullet(\mathbf f)$ est égal à l'\emph{algèbre extérieure
$\bigwedge(A^r)$}, graduée de la façon usuelle, et l'opérateur bord est la
\emph{multiplication intérieure $i_{\mathbf f}$} par $\mathbf f$ considéré
comme élément du dual $(A^r)^\vee$~; on rappelle que $i_{\mathbf f}$ est une
\emph{antidérivation} de degré $-1$ de $\bigwedge(A^r)$, et que si
$(\mathbf e_i)_{1\leq i\leq r}$ est la base canonique de $A^r$, on a
$i_{\mathbf f}(\mathbf e_i)=f_i$~; la vérification de la condition
$i_{\mathbf f}\circ i_{\mathbf f}=0$ est immédiate.

Une définition équivalente est la suivante~: pour chaque $i$, on considère
un complexe de chaînes $K_\bullet(f_i)$ défini comme suit~:
$K_0(f_i)=K_1(f_i)=A$, $K_n(f_i)=0$ pour $n\neq0,1$~; l'opérateur bord est
défini par la condition que $d_1:A\to A$ est la \emph{multiplication par
$f_i$}. On prend alors pour $K_\bullet(\mathbf f)$ le \emph{produit tensoriel
$K_\bullet(f_1)\otimes K_\bullet(f_2)\otimes\cdots\otimes K_\bullet(f_r)$}
(G, I, 2.7) muni de son degré total~; la vérification de l'isomorphisme de ce
complexe et du complexe défini plus haut est immédiate.
\end{env}

\begin{env}[1.1.2]
\label{III.1.1.2-fr}
Pour tout $A$-module $M$, on définit le \emph{complexe de chaînes}
\begin{equation}
\label{III.1.1.2.1-fr}
\tag{1.1.2.1}
K_\bullet(\mathbf f,M)=K_\bullet(\mathbf f)\otimes_A M
\end{equation}
et le \emph{complexe de cochaînes} (G, I, 2.2)
\begin{equation}
\label{III.1.1.2.2-fr}
\tag{1.1.2.2}
K^\bullet(\mathbf f,M)=\operatorname{Hom}_A(K_\bullet(\mathbf f),M).
\end{equation}
Si $g$ est une $k$-cochaîne de ce dernier complexe, et si on pose
\[
g(i_1,\ldots,i_k)=g(\mathbf e_{i_1}\wedge\cdots\wedge\mathbf e_{i_k}),
\]
$g$ s'identifie à une application \emph{alternée} de $[1,r]^k$ dans $M$, et
il résulte des définitions précédentes que l'on a
\begin{equation}
\label{III.1.1.2.3-fr}
\tag{1.1.2.3}
d^kg(i_1,i_2,\ldots,i_{k+1})=
\sum_{h=1}^{k+1}(-1)^{h-1}f_{i_h}
g(i_1,\ldots,\widehat{i_h},\ldots,i_{k+1}).
\end{equation}
\end{env}

\oldpage[III]{83}
\begin{env}[1.1.3]
\label{III.1.1.3-fr}
Des complexes précédents, on déduit comme d'ordinaire les \emph{$A$-modules
d'homologie et de cohomologie} (G, I, 2.2)
\begin{align}
\label{III.1.1.3.1-fr}
\tag{1.1.3.1}
H_\bullet(\mathbf f,M)&=H_\bullet(K_\bullet(\mathbf f,M)),\\
\label{III.1.1.3.2-fr}
\tag{1.1.3.2}
H^\bullet(\mathbf f,M)&=H^\bullet(K^\bullet(\mathbf f,M)).
\end{align}
On définit d'ailleurs un $A$-isomorphisme
$K_i(\mathbf f,M)\xrightarrow{\sim}K^{r-i}(\mathbf f,M)$ en faisant
correspondre à toute chaîne
$z=\sum(\mathbf e_{i_1}\wedge\cdots\wedge\mathbf e_{i_k})\otimes
z_{i_1\ldots i_k}$ la cochaîne $g_z$ telle que
$g_z(j_1,\ldots,j_{r-k})=\varepsilon z_{i_1\ldots i_k}$, où
$(j_h)_{1\leq h\leq r-k}$ est la suite strictement croissante complémentaire
de la suite strictement croissante $(i_h)_{1\leq h\leq k}$ dans $[1,r]$ et
$\varepsilon=(-1)^\nu$, $\nu$ étant le nombre d'inversions de la permutation
$i_1,\ldots,i_k,j_1,\ldots,j_{r-k}$ de $[1,r]$. On vérifie que
$g_{dz}=d(g_z)$, ce qui donne un isomorphisme
\begin{equation}
\label{III.1.1.3.3-fr}
\tag{1.1.3.3}
H^i(\mathbf f,M)\xrightarrow{\sim}H_{r-i}(\mathbf f,M)
\qquad\text{pour }0\leq i\leq r.
\end{equation}
Dans ce chapitre, nous aurons surtout à considérer les modules de cohomologie
$H^\bullet(\mathbf f,M)$.

Pour un $\mathbf f$ donné, il est immédiat (G, I, 2.1) que
$M\mapsto H^\bullet(\mathbf f,M)$ est un \emph{foncteur cohomologique}
(T, II, 2.1), de la catégorie des $A$-modules dans celle des $A$-modules
gradués, nul pour les degrés $<0$ et $>r$. En outre, on a
\begin{equation}
\label{III.1.1.3.4-fr}
\tag{1.1.3.4}
H^0(\mathbf f,M)=\operatorname{Hom}_A(A/(\mathbf f),M)
\end{equation}
en désignant par $(\mathbf f)$ l'idéal de $A$ engendré par
$f_1,\ldots,f_r$~; cela résulte aussitôt de
(\hyperref[III.1.1.2.3-fr]{1.1.2.3}), et il est clair que
$H^0(\mathbf f,M)$ s'identifie au sous-module de $M$ \emph{annulé par
$(\mathbf f)$}. De même, on a d'après
(\hyperref[III.1.1.2.3-fr]{1.1.2.3})
\begin{equation}
\label{III.1.1.3.5-fr}
\tag{1.1.3.5}
H^r(\mathbf f,M)=M/\Bigl(\sum_{i=1}^r f_iM\Bigr)
=(A/(\mathbf f))\otimes_A M.
\end{equation}
Nous utiliserons le résultat connu suivant, dont nous rappellerons une
démonstration pour être complet~:
\end{env}

\begin{proposition}[1.1.4]
\label{III.1.1.4-fr}
Soient $A$ un anneau, $\mathbf f=(f_i)_{1\leq i\leq r}$ une famille finie
d'éléments de $A$, $M$ un $A$-module. Si, pour $1\leq i\leq r$, l'homothétie
$z\mapsto f_i z$ dans $M_{i-1}=M/(f_1M+\cdots+f_{i-1}M)$ est injective, on a
$H^i(\mathbf f,M)=0$ pour $i\neq r$.
\end{proposition}

\begin{proof}
Il suffit en effet de prouver que $H_i(\mathbf f,M)=0$ pour tout $i>0$ en
vertu de (\hyperref[III.1.1.3.3-fr]{1.1.3.3}). Raisonnons par récurrence sur
$r$, le cas $r=0$ étant trivial. Posons
$\mathbf f'=(f_i)_{1\leq i\leq r-1}$~; cette famille vérifie les conditions de
l'énoncé, donc, si on pose $L_\bullet=K_\bullet(\mathbf f',M)$, on a
$H_i(L_\bullet)=0$ pour $i>0$ par hypothèse, et
$H_0(L_\bullet)=M_{r-1}$ en vertu de
(\hyperref[III.1.1.3.3-fr]{1.1.3.3}) et
(\hyperref[III.1.1.3.5-fr]{1.1.3.5}). Posons pour abréger
$K_\bullet=K_\bullet(f_r)=K_0\oplus K_1$, avec $K_0=K_1=A$,
$d_1:K_1\to K_0$ étant la multiplication par $f_r$~; on a par définition
(\hyperref[III.1.1.1-fr]{1.1.1})
$K_\bullet(\mathbf f,M)=K_\bullet\otimes_A L_\bullet$. Or, on a le lemme
suivant~:

\begin{lemma}[1.1.4.1]
\label{III.1.1.4.1-fr}
Soit $K_\bullet$ un complexe de chaînes formé de $A$-modules libres, nuls sauf
en dimensions $0$ et $1$. Pour tout complexe de chaînes $L_\bullet$ formé de
$A$-modules, on a une suite exacte
\[
0\longrightarrow H_0(K_\bullet\otimes H_p(L_\bullet))
\longrightarrow H_p(K_\bullet\otimes L_\bullet)
\longrightarrow H_1(K_\bullet\otimes H_{p-1}(L_\bullet))
\longrightarrow0
\]
pour tout indice $p$.
\end{lemma}

\oldpage[III]{84}
C'est un cas particulier d'une suite exacte de termes de bas degré de la
suite spectrale de Künneth (M, XVII, 5.2 a) et G, I, 5.5.2)~; on peut le
démontrer directement de la façon suivante. Considérons $K_0$ et $K_1$ comme
des complexes de chaînes (nuls en dimensions $\neq0$ et $\neq1$
respectivement)~; on a alors une suite exacte de complexes
\[
0\longrightarrow K_0\otimes L_\bullet
\longrightarrow K_\bullet\otimes L_\bullet
\longrightarrow K_1\otimes L_\bullet\longrightarrow0
\]
à laquelle nous pouvons appliquer la suite exacte d'homologie
\[
\cdots\longrightarrow H_{p+1}(K_1\otimes L_\bullet)
\xrightarrow{\partial}H_p(K_0\otimes L_\bullet)
\longrightarrow H_p(K_\bullet\otimes L_\bullet)
\longrightarrow H_p(K_1\otimes L_\bullet)
\xrightarrow{\partial}H_{p-1}(K_0\otimes L_\bullet)
\longrightarrow\cdots.
\]
Mais il est évident que
$H_p(K_0\otimes L_\bullet)=K_0\otimes H_p(L_\bullet)$ et
$H_p(K_1\otimes L_\bullet)=K_1\otimes H_{p-1}(L_\bullet)$ pour tout $p$~; en
outre, on vérifie aussitôt que l'opérateur
$\partial:K_1\otimes H_p(L_\bullet)\to K_0\otimes H_p(L_\bullet)$ n'est autre
que $d_1\otimes1$~; le lemme résulte donc de la suite exacte précédente et de
la définition de $H_0(K_\bullet\otimes H_p(L_\bullet))$ et de
$H_1(K_\bullet\otimes H_{p-1}(L_\bullet))$.

Ce lemme étant établi, la fin de la démonstration de
(\hyperref[III.1.1.4-fr]{1.1.4}) est immédiate~: l'hypothèse de récurrence et
ce lemme (\hyperref[III.1.1.4.1-fr]{1.1.4.1}) donnent
$H_p(K_\bullet\otimes L_\bullet)=0$ pour $p\geq2$~; en outre, si on prouve que
$H_1(K_\bullet,H_0(L_\bullet))=0$, on déduira aussi
$H_1(K_\bullet\otimes L_\bullet)=0$ du lemme
(\hyperref[III.1.1.4.1-fr]{1.1.4.1})~; mais par définition,
$H_1(K_\bullet,H_0(L_\bullet))$ n'est autre que le noyau de l'homothétie
$z\mapsto f_rz$ dans $M_{r-1}$, et comme par hypothèse ce noyau est nul, cela
achève la démonstration.
\end{proof}

\begin{env}[1.1.5]
\label{III.1.1.5-fr}
Soit $\mathbf g=(g_i)_{1\leq i\leq r}$ une seconde suite de $r$ éléments de
$A$, et posons $\mathbf f\mathbf g=(f_ig_i)_{1\leq i\leq r}$. On peut définir
un homomorphisme canonique de complexes
\begin{equation}
\label{III.1.1.5.1-fr}
\tag{1.1.5.1}
\varphi_{\mathbf g}:K_\bullet(\mathbf f\mathbf g)\longrightarrow
K_\bullet(\mathbf f)
\end{equation}
comme l'extension canonique à l'algèbre extérieure $\bigwedge(A^r)$ de
l'application $A$-linéaire
$(x_1,\ldots,x_r)\mapsto(g_1x_1,\ldots,g_rx_r)$ de $A^r$ dans lui-même. Pour
voir qu'on a bien un homomorphisme de complexes, il suffit de remarquer, de
façon générale, que si $u:E\to F$ est une application $A$-linéaire,
$\mathbf x\in F^\vee$ et $\mathbf y={}^tu(\mathbf x)\in E^\vee$, alors on a
la formule
\begin{equation}
\label{III.1.1.5.2-fr}
\tag{1.1.5.2}
(\bigwedge u)\circ i_{\mathbf y}=i_{\mathbf x}\circ(\bigwedge u)~;
\end{equation}
en effet, les deux membres sont des antidérivations de $\bigwedge F$, et il
suffit de vérifier qu'ils coïncident dans $F$, ce qui résulte aussitôt des
définitions.

Lorsqu'on identifie $K_\bullet(\mathbf f)$ au produit tensoriel des
$K_\bullet(f_i)$ (\hyperref[III.1.1.1-fr]{1.1.1}), $\varphi_{\mathbf g}$ est
le produit tensoriel des $\varphi_{g_i}$, où $\varphi_{g_i}$ se réduit à
l'identité en degré $0$ et à la multiplication par $g_i$ en degré $1$.
\end{env}

\begin{env}[1.1.6]
\label{III.1.1.6-fr}
En particulier, pour tout couple d'entiers $m,n$ tels que $0\leq n\leq m$,
on a des homomorphismes de complexes
\begin{equation}
\label{III.1.1.6.1-fr}
\tag{1.1.6.1}
\varphi_{\mathbf f^{m-n}}:K_\bullet(\mathbf f^m)\longrightarrow
K_\bullet(\mathbf f^n)
\end{equation}
et par suite des homomorphismes
\begin{align}
\label{III.1.1.6.2-fr}
\tag{1.1.6.2}
\varphi_{\mathbf f^{m-n}} &:K^\bullet(\mathbf f^n,M)\longrightarrow
K^\bullet(\mathbf f^m,M),\\
\label{III.1.1.6.3-fr}
\tag{1.1.6.3}
\varphi_{\mathbf f^{m-n}} &:H^\bullet(\mathbf f^n,M)\longrightarrow
H^\bullet(\mathbf f^m,M).
\end{align}

\oldpage[III]{85}
Ces derniers vérifient évidemment la condition de transitivité
$\varphi_{\mathbf f^{m-p}}=
\varphi_{\mathbf f^{m-n}}\circ\varphi_{\mathbf f^{n-p}}$ pour $p\leq n\leq m$~;
ils définissent donc deux \emph{systèmes inductifs} de $A$-modules~; on posera
\begin{align}
\label{III.1.1.6.4-fr}
\tag{1.1.6.4}
C^\bullet((\mathbf f),M)&=\varinjlim_n K^\bullet(\mathbf f^n,M),\\
\label{III.1.1.6.5-fr}
\tag{1.1.6.5}
H^\bullet((\mathbf f),M)&=H^\bullet(C^\bullet((\mathbf f),M))
=\varinjlim_n H^\bullet(\mathbf f^n,M),
\end{align}
la dernière égalité provenant du fait que le passage à la limite inductive
permute au foncteur $H^\bullet$ (G, I, 2.1). On verra ultérieurement
(\hyperref[III.1.4.3-fr]{1.4.3}) que $H^\bullet((\mathbf f),M)$ ne dépend en
fait que de l'\emph{idéal $(\mathbf f)$} dans $A$ (et même de la topologie
$(\mathbf f)$-préadique sur $A$), ce qui justifie les notations.

Il est clair que $M\mapsto C^\bullet((\mathbf f),M)$ est un foncteur
$A$-linéaire exact, et $M\mapsto H^\bullet((\mathbf f),M)$ un foncteur
cohomologique.
\end{env}

\begin{env}[1.1.7]
\label{III.1.1.7-fr}
Soient $\mathbf f=(f_i)\in A^r$, $\mathbf g=(g_i)\in A^r$~; désignons par
$e_{\mathbf g}$ la multiplication à gauche par le vecteur
$\mathbf g\in A^r$ dans l'algèbre extérieure $\bigwedge(A^r)$~; on sait que
l'on a la \emph{formule d'homotopie}
\begin{equation}
\label{III.1.1.7.1-fr}
\tag{1.1.7.1}
i_{\mathbf f}e_{\mathbf g}+e_{\mathbf g}i_{\mathbf f}
=\langle\mathbf g,\mathbf f\rangle 1
\end{equation}
dans le $A$-module $A^r$ ($1$ désignant l'automorphisme identique de $A^r$)~;
cette relation signifie aussi que, \emph{dans le complexe
$K_\bullet(\mathbf f)$}, on a
\begin{equation}
\label{III.1.1.7.2-fr}
\tag{1.1.7.2}
de_{\mathbf g}+e_{\mathbf g}d=\langle\mathbf g,\mathbf f\rangle 1.
\end{equation}
Si l'idéal $(\mathbf f)$ est égal à $A$, il existe $\mathbf g\in A^r$ tel que
$\langle\mathbf g,\mathbf f\rangle=\sum_{i=1}^r g_if_i=1$. Par suite
(G, I, 2.4)~:
\end{env}

\begin{proposition}[1.1.8]
\label{III.1.1.8-fr}
Supposons que l'idéal $(\mathbf f)$ engendré par les $f_i$ soit égal à $A$.
Alors le complexe $K_\bullet(\mathbf f)$ est homotopiquement trivial, et il en
est donc de même des complexes $K_\bullet(\mathbf f,M)$ et
$K^\bullet(\mathbf f,M)$ pour tout $A$-module $M$.
\end{proposition}

\begin{corollary}[1.1.9]
\label{III.1.1.9-fr}
Si $(\mathbf f)=A$, on a $H^\bullet(\mathbf f,M)=0$ et
$H^\bullet((\mathbf f),M)=0$ pour tout $A$-module $M$.
\end{corollary}

\begin{proof}
En effet, on a alors $(\mathbf f^n)=A$ pour tout $n$.
\end{proof}

\begin{remark}[1.1.10]
\label{III.1.1.10-fr}
Avec les mêmes notations que ci-dessus, soient $X=\operatorname{Spec}(A)$,
$Y$ le sous-schéma fermé de $X$ défini par l'idéal $(\mathbf f)$. Nous
prouverons au \S~9 que $H^\bullet((\mathbf f),M)$ est isomorphe à la
cohomologie $H_Y^\bullet(X,\widetilde M)$ correspondant à l'antifiltre $\Phi$
des parties fermées de $Y$ (T, 3.2). Nous montrerons aussi que la prop.
(\hyperref[III.1.2.3-fr]{1.2.3}) appliquée à $X$ et à
$\mathcal F=\widetilde M$, est un cas particulier d'une suite exacte de
cohomologie
\[
\cdots\longrightarrow H_Y^p(X,\mathcal F)\longrightarrow H^p(X,\mathcal F)
\longrightarrow H^p(X-Y,\mathcal F)\longrightarrow H_Y^{p+1}(X,\mathcal F)
\longrightarrow\cdots.
\]
\end{remark}

\subsection{Cohomologie de Čech d'un recouvrement ouvert}
\label{subsection:III.1.2-fr}

\begin{notation}[1.2.1]
\label{III.1.2.1-fr}
Dans ce numéro, on notera~:
\begin{itemize}
\item $X$ un préschéma~;
\item $\mathcal F$ un $\mathcal O_X$-Module quasi-cohérent~;
\oldpage[III]{86}
\item $A=\Gamma(X,\mathcal O_X)$, $M=\Gamma(X,\mathcal F)$~;
\item $\mathbf f=(f_i)_{1\leq i\leq r}$ un système fini d'éléments de $A$~;
\item $U_i=X_{f_i}$ l'ensemble ouvert
(\hyperref[0.5.5.2-fr]{0\textsubscript{I}, 5.5.2}) des $x\in X$ tels que
$f_i(x)\neq0$~;
\item $U=\bigcup_{i=1}^r U_i$~;
\item $\mathfrak U$ le recouvrement $(U_i)_{1\leq i\leq r}$ de $U$.
\end{itemize}
\end{notation}

\begin{env}[1.2.2]
\label{III.1.2.2-fr}
Supposons que $X$ soit, ou bien un préschéma dont l'espace sous-jacent est
\emph{noethérien}, ou bien un \emph{schéma} dont l'espace sous-jacent soit
\emph{quasi-compact}. On sait alors (\hyperref[I.9.3.3-fr]{I, 9.3.3}) que
l'on a $\Gamma(U_i,\mathcal F)=M_{f_i}$. Nous poserons
\[
U_{i_0i_1\ldots i_p}=\bigcap_{k=0}^p U_{i_k}
=X_{f_{i_0}f_{i_1}\cdots f_{i_p}}
\]
(\hyperref[0.5.5.3-fr]{0\textsubscript{I}, 5.5.3})~; on a donc aussi
\begin{equation}
\label{III.1.2.2.1-fr}
\tag{1.2.2.1}
\Gamma(U_{i_0i_1\ldots i_p},\mathcal F)=M_{f_{i_0}f_{i_1}\cdots f_{i_p}}.
\end{equation}
Or (\hyperref[0.1.6.1-fr]{0\textsubscript{I}, 1.6.1})
$M_{f_{i_0}f_{i_1}\cdots f_{i_p}}$ s'identifie à la limite inductive
$\varinjlim_n M_{i_0i_1\ldots i_p}^{(n)}$, où le système inductif est formé
des $M_{i_0i_1\ldots i_p}^{(n)}=M$, l'homomorphisme
$\varphi_{nm}:M_{i_0i_1\ldots i_p}^{(m)}\to
M_{i_0i_1\ldots i_p}^{(n)}$ étant la multiplication par
$(f_{i_0}f_{i_1}\cdots f_{i_p})^{n-m}$ pour $m\leq n$. Désignons par
$C_n^p(M)$ l'ensemble des applications \emph{alternées} de $[1,r]^{p+1}$
dans $M$ (pour tout $n$)~; ces $A$-modules forment encore un système inductif
pour les $\varphi_{nm}$. Si $C^p(\mathfrak U,\mathcal F)$ est le groupe des
$p$-cochaînes \emph{alternées} de Čech relatif au recouvrement $\mathfrak U$,
à coefficients dans $\mathcal F$ (G, II, 5.1), il résulte de ce qui précède
que l'on peut écrire
\begin{equation}
\label{III.1.2.2.2-fr}
\tag{1.2.2.2}
C^p(\mathfrak U,\mathcal F)=\varinjlim_n C_n^p(M).
\end{equation}
Or, avec les notations de (\hyperref[III.1.1.2-fr]{1.1.2}), $C_n^p(M)$
s'identifie à $K^{p+1}(\mathbf f^n,M)$, et l'application $\varphi_{nm}$
s'identifie à l'application $\varphi_{\mathbf f^{n-m}}$ définie dans
(\hyperref[III.1.1.6-fr]{1.1.6}). On a donc, pour tout $p\geq0$, un
isomorphisme canonique fonctoriel en $\mathcal F$
\begin{equation}
\label{III.1.2.2.3-fr}
\tag{1.2.2.3}
C^p(\mathfrak U,\mathcal F)\xrightarrow{\sim}C^{p+1}((\mathbf f),M).
\end{equation}
En outre, la formule (\hyperref[III.1.1.2.3-fr]{1.1.2.3}) et la définition de
la cohomologie d'un recouvrement (G, II, 5.1) montrent que les isomorphismes
(\hyperref[III.1.2.2.3-fr]{1.2.2.3}) sont compatibles avec les opérateurs
cobords.
\end{env}

\begin{proposition}[1.2.3]
\label{III.1.2.3-fr}
Si $X$ est un préschéma dont l'espace sous-jacent est noethérien, ou un
schéma dont l'espace sous-jacent est quasi-compact, il existe un isomorphisme
canonique fonctoriel en $\mathcal F$
\begin{equation}
\label{III.1.2.3.1-fr}
\tag{1.2.3.1}
H^p(\mathfrak U,\mathcal F)\xrightarrow{\sim}H^{p+1}((\mathbf f),M)
\qquad\text{pour tout }p\geq1.
\end{equation}
On a en outre une suite exacte fonctorielle en $\mathcal F$
\begin{equation}
\label{III.1.2.3.2-fr}
\tag{1.2.3.2}
0\longrightarrow H^0((\mathbf f),M)\longrightarrow M
\longrightarrow H^0(\mathfrak U,\mathcal F)
\longrightarrow H^1((\mathbf f),M)\longrightarrow0.
\end{equation}
\end{proposition}

\oldpage[III]{87}
\begin{proof}
Les relations (\hyperref[III.1.2.3.1-fr]{1.2.3.1}) sont en effet conséquences
immédiates de ce qu'on a vu dans (\hyperref[III.1.2.2-fr]{1.2.2}). D'autre
part, on a
$C^0(\mathfrak U,\mathcal F)=C^1((\mathbf f),M)$~;
$H^0(\mathfrak U,\mathcal F)$ s'identifie par suite au sous-groupe des
$1$-cocycles de $C^1((\mathbf f),M)$~; comme
$M=C^0((\mathbf f),M)$, la suite exacte
(\hyperref[III.1.2.3.2-fr]{1.2.3.2}) n'est autre que celle qui résulte de la
définition des groupes de cohomologie $H^0((\mathbf f),M)$ et
$H^1((\mathbf f),M)$.
\end{proof}

\begin{corollary}[1.2.4]
\label{III.1.2.4-fr}
Supposons que les $X_{f_i}$ soient quasi-compacts et qu'il existe des
$g_i\in\Gamma(U,\mathcal F)$ tels que
$\sum_i g_i(f_i\vert U)=1\vert U$. Alors, pour tout
$(\mathcal O_X\vert U)$-Module quasi-cohérent $\mathcal G$, on a
$H^p(\mathfrak U,\mathcal G)=0$ pour $p>0$~; si en outre $U=X$,
l'homomorphisme canonique (\hyperref[III.1.2.3.2-fr]{1.2.3.2})
$M\to H^0(\mathfrak U,\mathcal F)$ est bijectif.
\end{corollary}

\begin{proof}
Comme par hypothèse les $U_i=X_{f_i}$ sont quasi-compacts, il en est de même
de $U$, et on peut donc se borner au cas où $U=X$~; l'hypothèse entraîne alors
$H^p((\mathbf f),M)=0$ pour tout $p\geq0$
(\hyperref[III.1.1.9-fr]{1.1.9}). Le corollaire résulte alors aussitôt de
(\hyperref[III.1.2.3.1-fr]{1.2.3.1}) et
(\hyperref[III.1.2.3.2-fr]{1.2.3.2}).
\end{proof}

On observera que puisque
$H^0(\mathfrak U,\mathcal F)=H^0(U,\mathcal F)$ (G, II, 5.2.2), on a ainsi
démontré à nouveau (I, 1.3.7) comme cas particulier.

\begin{remark}[1.2.5]
\label{III.1.2.5-fr}
Supposons que $X$ soit un \emph{schéma affine}~; alors les
$U_i=X_{f_i}=D(f_i)$ sont des ouverts affines, ainsi que les
$U_{i_0i_1\ldots i_p}$ (mais $U$ n'est pas nécessairement affine). Dans ce
cas, les foncteurs $\Gamma(X,\mathcal F)$ et
$\Gamma(U_{i_0i_1\ldots i_p},\mathcal F)$ sont exacts en $\mathcal F$
(I, 1.3.11). Si on a une suite exacte
$0\to\mathcal F'\to\mathcal F\to\mathcal F''\to0$ de
$\mathcal O_X$-Modules quasi-cohérents, la suite des complexes
\[
0\longrightarrow C^\bullet(\mathfrak U,\mathcal F')
\longrightarrow C^\bullet(\mathfrak U,\mathcal F)
\longrightarrow C^\bullet(\mathfrak U,\mathcal F'')
\longrightarrow0
\]
est exacte, et donne donc une suite exacte de cohomologie
\[
\cdots\longrightarrow H^p(\mathfrak U,\mathcal F')
\longrightarrow H^p(\mathfrak U,\mathcal F)
\longrightarrow H^p(\mathfrak U,\mathcal F'')
\xrightarrow{\partial}H^{p+1}(\mathfrak U,\mathcal F')
\longrightarrow\cdots.
\]
D'autre part, si on pose $M'=\Gamma(X,\mathcal F')$,
$M''=\Gamma(X,\mathcal F'')$, la suite
$0\to M'\to M\to M''\to0$ est exacte~; comme
$C^\bullet((\mathbf f),M)$ est un foncteur exact en $M$, on a aussi la suite
exacte de cohomologie
\[
\cdots\longrightarrow H^p((\mathbf f),M')
\longrightarrow H^p((\mathbf f),M)
\longrightarrow H^p((\mathbf f),M'')
\xrightarrow{\partial}H^{p+1}((\mathbf f),M')
\longrightarrow\cdots.
\]
Cela étant, comme le diagramme
\[
\xymatrix{
0\ar[r]&C^\bullet(\mathfrak U,\mathcal F')\ar[r]\ar[d]
&C^\bullet(\mathfrak U,\mathcal F)\ar[r]\ar[d]
&C^\bullet(\mathfrak U,\mathcal F'')\ar[r]\ar[d]&0\\
0\ar[r]&C^\bullet((\mathbf f),M')\ar[r]
&C^\bullet((\mathbf f),M)\ar[r]
&C^\bullet((\mathbf f),M'')\ar[r]&0
}
\]
\oldpage[III]{88}
est commutatif, on en conclut que les diagrammes
\begin{equation}
\label{III.1.2.5.1-fr}
\tag{1.2.5.1}
\xymatrix{
H^p(\mathfrak U,\mathcal F'')\ar[r]^{\partial}\ar[d]
&H^{p+1}(\mathfrak U,\mathcal F')\ar[d]\\
H^{p+1}((\mathbf f),M'')\ar[r]^{\partial}
&H^{p+2}((\mathbf f),M')
}
\end{equation}
sont commutatifs pour tout $p$ (G, I, 2.1.1).
\end{remark}

\subsection{Cohomologie d'un schéma affine}
\label{subsection:III.1.3-fr}

\begin{theorem}[1.3.1]
\label{III.1.3.1-fr}
Soit $X$ un schéma affine. Pour tout $\mathcal O_X$-Module quasi-cohérent
$\mathcal F$, on a $H^p(X,\mathcal F)=0$ pour tout $p>0$.
\end{theorem}

\begin{proof}
Soit $\mathfrak U$ un recouvrement fini de $X$ par des ouverts affines
$X_{f_i}=D(f_i)$ ($1\leq i\leq r$)~; on sait qu'alors l'idéal engendré dans
$A=\Gamma(X,\mathcal O_X)$ par les $f_i$ est égal à $A$. On conclut donc de
(\hyperref[III.1.2.4-fr]{1.2.4}) que l'on a
$H^p(\mathfrak U,\mathcal F)=0$ pour $p>0$. Comme il y a des recouvrements
finis de $X$ par des ouverts affines qui sont arbitrairement fins (I, 1.1.10),
la définition de la cohomologie de Čech (G, II, 5.8) montre que l'on a aussi
$\check H^p(X,\mathcal F)=0$ pour $p>0$. Mais ceci s'applique aussi à tout
préschéma $X_f$ pour $f\in A$ (I, 1.3.6), donc
$\check H^p(X_f,\mathcal F)=0$ pour $p>0$. Comme l'on a
$X_f\cap X_g=X_{fg}$, on en déduit que l'on a aussi
$H^p(X,\mathcal F)=0$ pour tout $p>0$, en vertu de (G, II, 5.9.2).
\end{proof}

\begin{corollary}[1.3.2]
\label{III.1.3.2-fr}
Soient $Y$ un préschéma, $f:X\to Y$ un morphisme affine (II, 1.6.1). Pour
tout $\mathcal O_X$-Module quasi-cohérent $\mathcal F$, on a
$R^qf_*(\mathcal F)=0$ pour $q>0$.
\end{corollary}

\begin{proof}
En effet, par définition $R^qf_*(\mathcal F)$ est le $\mathcal O_Y$-Module
associé au préfaisceau $U\mapsto H^q(f^{-1}(U),\mathcal F)$, $U$ parcourant
les ouverts de $Y$. Mais les ouverts affines forment une base de $Y$, et pour
un tel ouvert $U$, $f^{-1}(U)$ est affine (II, 1.3.2), donc
$H^q(f^{-1}(U),\mathcal F)=0$ par
(\hyperref[III.1.3.1-fr]{1.3.1}), ce qui démontre le corollaire.
\end{proof}

\begin{corollary}[1.3.3]
\label{III.1.3.3-fr}
Soient $Y$ un préschéma, $f:X\to Y$ un morphisme affine. Pour tout
$\mathcal O_X$-Module quasi-cohérent $\mathcal F$, l'homomorphisme canonique
$H^p(Y,f_*(\mathcal F))\to H^p(X,\mathcal F)$ (0, 12.1.3.1) est bijectif
pour tout $p$.
\end{corollary}

\begin{proof}
En effet, il suffit (en vertu de (0, 12.1.7)) de montrer que les
edge-homomorphismes
$''E_2^{p0}=H^p(Y,f_*(\mathcal F))\to H^p(X,\mathcal F)$ de la seconde suite
spectrale du foncteur composé $\Gamma f_*$ sont bijectifs. Mais le terme
$E_2$ de cette suite est donné par
$''E_2^{pq}=H^p(Y,R^qf_*(\mathcal F))$ (G, II, 4.17.1), donc il résulte de
(\hyperref[III.1.3.2-fr]{1.3.2}) que $''E_2^{pq}=0$ pour $q>0$, et la suite
spectrale dégénère~; d'où notre assertion (0, 11.1.6).
\end{proof}

\begin{corollary}[1.3.4]
\label{III.1.3.4-fr}
Soient $f:X\to Y$ un morphisme affine, $g:Y\to Z$ un morphisme.
\oldpage[III]{89}
Pour tout $\mathcal O_X$-Module quasi-cohérent $\mathcal F$,
l'homomorphisme canonique
$R^pg_*(f_*(\mathcal F))\to R^p(g\circ f)_*(\mathcal F)$ (0, 12.2.5.1)
est bijectif pour tout $p$.
\end{corollary}

\begin{proof}
Il suffit de remarquer que, d'après
(\hyperref[III.1.3.3-fr]{1.3.3}), pour tout ouvert affine $W$ de $Z$,
l'homomorphisme canonique
$H^p(g^{-1}(W),f_*(\mathcal F))\to
H^p(f^{-1}(g^{-1}(W)),\mathcal F)$ est bijectif~; cela prouve que
l'homomorphisme de préfaisceaux définissant l'homomorphisme canonique
$R^pg_*(f_*(\mathcal F))\to R^p(g\circ f)_*(\mathcal F)$ est bijectif
(0, 12.2.5).
\end{proof}

\subsection{Application à la cohomologie des préschémas quelconques}
\label{subsection:III.1.4-fr}

\begin{proposition}[1.4.1]
\label{III.1.4.1-fr}
Soient $X$ un schéma, $\mathfrak U=(U_\alpha)$ un recouvrement de $X$ par
des ouverts affines. Pour tout $\mathcal O_X$-Module quasi-cohérent
$\mathcal F$, les modules de cohomologie $H^\bullet(X,\mathcal F)$ et
$H^\bullet(\mathfrak U,\mathcal F)$ (sur $\Gamma(X,\mathcal O_X)$) sont
canoniquement isomorphes.
\end{proposition}

\begin{proof}
En effet, comme $X$ est un schéma, toute intersection finie $V$ d'ouverts du
recouvrement $\mathfrak U$ est affine (I, 5.5.6), donc
$H^q(V,\mathcal F)=0$ pour $q\geq1$ en vertu de
(\hyperref[III.1.3.1-fr]{1.3.1}). La proposition résulte alors du th. de
Leray (G, II, 5.4.1).
\end{proof}

\begin{remark}[1.4.2]
\label{III.1.4.2-fr}
On notera que la conclusion de (\hyperref[III.1.4.1-fr]{1.4.1}) est encore
valable lorsque les intersections finies des ensembles $U_\alpha$ sont
affines, même lorsque l'on ne suppose pas nécessairement que $X$ soit un
schéma.
\end{remark}

\begin{corollary}[1.4.3]
\label{III.1.4.3-fr}
Soient $X$ un schéma dont l'espace sous-jacent est quasi-compact,
$A=\Gamma(X,\mathcal O_X)$, $\mathbf f=(f_i)_{1\leq i\leq r}$ une suite
finie d'éléments de $A$ tels que les $X_{f_i}$ (notations de
(\hyperref[III.1.2.1-fr]{1.2.1})) soient affines. Alors (avec les notations
de (\hyperref[III.1.2.1-fr]{1.2.1})), pour tout $\mathcal O_X$-Module
quasi-cohérent $\mathcal F$, on a un isomorphisme canonique fonctoriel en
$\mathcal F$
\begin{equation}
\label{III.1.4.3.1-fr}
\tag{1.4.3.1}
H^q(U,\mathcal F)\xrightarrow{\sim}H^{q+1}((\mathbf f),M)
\qquad\text{pour }q\geq1
\end{equation}
et une suite exacte fonctorielle en $\mathcal F$
\begin{equation}
\label{III.1.4.3.2-fr}
\tag{1.4.3.2}
0\longrightarrow H^0((\mathbf f),M)\longrightarrow M
\longrightarrow H^0(U,\mathcal F)
\longrightarrow H^1((\mathbf f),M)\longrightarrow0.
\end{equation}
\end{corollary}

\begin{proof}
Cela résulte en effet aussitôt de
(\hyperref[III.1.4.1-fr]{1.4.1}) et
(\hyperref[III.1.2.3-fr]{1.2.3}).
\end{proof}

\begin{env}[1.4.4]
\label{III.1.4.4-fr}
Si $X$ est un \emph{schéma affine}, il résulte de
(\hyperref[III.1.2.5-fr]{1.2.5}) et
(\hyperref[III.1.4.1-fr]{1.4.1}) que pour tout $q\geq0$, les diagrammes
\begin{equation}
\label{III.1.4.4.1-fr}
\tag{1.4.4.1}
\xymatrix{
H^q(U,\mathcal F'')\ar[r]^{\partial}\ar[d]
&H^{q+1}(U,\mathcal F')\ar[d]\\
H^{q+1}((\mathbf f),M'')\ar[r]^{\partial}
&H^{q+2}((\mathbf f),M')
}
\end{equation}
correspondant à une suite exacte
$0\to\mathcal F'\to\mathcal F\to\mathcal F''\to0$ de
$\mathcal O_X$-Modules quasi-cohérents (avec les notations de
(\hyperref[III.1.2.5-fr]{1.2.5})), sont commutatifs.
\end{env}

\oldpage[III]{90}
\begin{proposition}[1.4.5]
\label{III.1.4.5-fr}
Soient $X$ un schéma quasi-compact, $\mathcal L$ un $\mathcal O_X$-Module
inversible, et considérons l'anneau gradué
$A_\bullet=\Gamma_\bullet(\mathcal L)$ (0\textsubscript{I}, 5.4.6)~; alors
\[
H^\bullet(\mathcal F,\mathcal L)
=\bigoplus_{n\in\mathbb Z}H^\bullet(X,\mathcal F\otimes\mathcal L^{\otimes n})
\]
est un $A_\bullet$-module gradué, et pour tout $f\in A_n$, on a un
isomorphisme canonique
\begin{equation}
\label{III.1.4.5.1-fr}
\tag{1.4.5.1}
H^\bullet(X_f,\mathcal F)
\xrightarrow{\sim}(H^\bullet(\mathcal F,\mathcal L))_{(f)}
\end{equation}
de $(A_\bullet)_{(f)}$-modules.
\end{proposition}

\begin{proof}
Comme $X$ est un schéma quasi-compact, on peut calculer la cohomologie de
tous les $\mathcal O_X$-Modules
$\mathcal F\otimes\mathcal L^{\otimes n}$ à l'aide d'un même recouvrement
fini $\mathfrak U=(U_i)$ par des ouverts affines tels que la restriction
$\mathcal L\vert U_i$ soit isomorphe à $\mathcal O_X\vert U_i$ pour chaque
$i$ (\hyperref[III.1.4.1-fr]{1.4.1}). Il est alors immédiat que les
$U_i\cap X_f$ sont des ouverts affines (I, 1.3.6), et l'on peut donc aussi
calculer la cohomologie
$H^\bullet(X_f,\mathcal F\otimes\mathcal L^{\otimes n})$ à l'aide du
recouvrement $\mathfrak U\vert X_f=(U_i\cap X_f)$
(\hyperref[III.1.4.1-fr]{1.4.1}). Il est immédiat que pour tout $f\in A_n$,
la multiplication par $f$ définit un homomorphisme
\[
C^\bullet(\mathfrak U,\mathcal F\otimes\mathcal L^{\otimes m})
\longrightarrow
C^\bullet(\mathfrak U,\mathcal F\otimes\mathcal L^{\otimes(m+n)}),
\]
d'où un homomorphisme
\[
H^\bullet(\mathfrak U,\mathcal F\otimes\mathcal L^{\otimes m})
\longrightarrow
H^\bullet(\mathfrak U,\mathcal F\otimes\mathcal L^{\otimes(m+n)}),
\]
ce qui établit la première assertion. D'autre part, pour $f\in A_n$ donné,
il résulte de (I, 9.3.2) que l'on a un isomorphisme de complexes de
$(A_\bullet)_{(f)}$-modules
\[
C^\bullet(\mathfrak U\vert X_f,\mathcal F)
\xrightarrow{\sim}
\left(C^\bullet\left(\mathfrak U,
\bigoplus_{n\in\mathbb Z}\mathcal F\otimes\mathcal L^{\otimes n}
\right)\right)_{(f)}
\]
compte tenu de (I, 1.3.9, (ii)). Passant à la cohomologie de ces deux
complexes, on en déduit l'isomorphisme
(\hyperref[III.1.4.5.1-fr]{1.4.5.1}), en se souvenant que le foncteur
$M\mapsto M_{(f)}$ est exact dans la catégorie des $A_\bullet$-modules
gradués.
\end{proof}

\begin{corollary}[1.4.6]
\label{III.1.4.6-fr}
On suppose vérifiées les hypothèses de
(\hyperref[III.1.4.5-fr]{1.4.5}) et on suppose en outre que
$\mathcal L=\mathcal O_X$. Si on pose $A=\Gamma(X,\mathcal O_X)$, alors,
pour tout $f\in A$, on a un isomorphisme canonique
\[
H^\bullet(X_f,\mathcal F)
\xrightarrow{\sim}(H^\bullet(X,\mathcal F))_f
\]
de $A_f$-modules.
\end{corollary}

\begin{corollary}[1.4.7]
\label{III.1.4.7-fr}
Soient $X$ un schéma quasi-compact, $f$ un élément de
$\Gamma(X,\mathcal O_X)$.
\begin{enumerate}
\item[(i)] Supposons l'ouvert $X_f$ affine. Alors, pour tout
$\mathcal O_X$-Module quasi-cohérent $\mathcal F$, tout $i>0$ et tout
$\xi\in H^i(X,\mathcal F)$, il existe un entier $n>0$ tel que $f^n\xi=0$.
\item[(ii)] Inversement, supposons que $X_f$ soit quasi-compact et que pour
tout faisceau quasi-cohérent $\mathcal J$ d'idéaux de $\mathcal O_X$ et tout
$\zeta\in H^1(X,\mathcal J)$, il existe $n>0$ tel que $f^n\zeta=0$. Alors
$X_f$ est affine.
\end{enumerate}
\end{corollary}

\begin{proof}
\begin{enumerate}
\item[(i)] Si $X_f$ est affine, on a $H^i(X_f,\mathcal F)=0$ pour tout
$i>0$ (\hyperref[III.1.3.1-fr]{1.3.1}), donc l'assertion résulte directement
de (\hyperref[III.1.4.6-fr]{1.4.6}).
\item[(ii)] En vertu du critère de Serre (II, 5.2.1), il suffit de prouver
que pour tout Idéal quasi-cohérent $\mathcal K$ de
$\mathcal O_X\vert X_f$, on a $H^1(X_f,\mathcal K)=0$. Comme $X_f$ est un
ouvert quasi-compact dans un schéma quasi-compact $X$, il existe un Idéal
quasi-cohérent $\mathcal J$ de $\mathcal O_X$ tel que
$\mathcal K=\mathcal J\vert X_f$ (I, 9.4.2). En vertu de
(\hyperref[III.1.4.6-fr]{1.4.6}), on a
$H^1(X_f,\mathcal K)=(H^1(X,\mathcal J))_f$, et l'hypothèse entraîne que le
second membre est nul, d'où la conclusion.
\end{enumerate}
\end{proof}

\begin{remark}[1.4.8]
\label{III.1.4.8-fr}
On notera que (\hyperref[III.1.4.7-fr]{1.4.7}, (i)) redonne une
démonstration plus simple de la relation (II, 4.5.13.2).
\end{remark}

\oldpage[III]{91}
\begin{lemma}[1.4.9]
\label{III.1.4.9-fr}
Soient $X$ un schéma quasi-compact,
$\mathfrak U=(U_i)_{1\leq i\leq n}$ un recouvrement fini de $X$ par des
ouverts affines, $\mathcal F$ un $\mathcal O_X$-Module quasi-cohérent. Le
complexe de faisceaux $\mathcal C^\bullet(\mathfrak U,\mathcal F)$ défini
par le recouvrement $\mathfrak U$ (G, II, 5.2) est alors un
$\mathcal O_X$-Module quasi-cohérent.
\end{lemma}

\begin{proof}
Il résulte des définitions (G, II, 5.2) que
$\mathcal C^p(\mathfrak U,\mathcal F)$ est somme directe des faisceaux
images directes des $\mathcal F\vert U_{i_0\ldots i_p}$ par l'injection
canonique $U_{i_0\ldots i_p}\to X$. Or, l'hypothèse que $X$ est un schéma
entraîne que ces injections sont des morphismes affines (I, 5.5.6), donc les
$\mathcal C^p(\mathfrak U,\mathcal F)$ sont quasi-cohérents (II, 1.2.6).
\end{proof}

\begin{proposition}[1.4.10]
\label{III.1.4.10-fr}
Soit $u:X\to Y$ un morphisme séparé et quasi-compact. Pour tout
$\mathcal O_X$-Module quasi-cohérent $\mathcal F$, les
$R^qu_*(\mathcal F)$ sont des $\mathcal O_Y$-Modules quasi-cohérents.
\end{proposition}

\begin{proof}
La question étant locale sur $Y$, on peut supposer $Y$ affine. Alors $X$ est
réunion finie d'ouverts affines $U_i$ ($1\leq i\leq n$)~; soit
$\mathfrak U$ le recouvrement $(U_i)$. En outre, comme $Y$ est un schéma,
il résulte de (I, 5.5.10) que pour tout ouvert affine $V\subset Y$,
l'injection canonique $u^{-1}(V)\to X$ est un morphisme affine~; on en
conclut ((\hyperref[III.1.4.1-fr]{1.4.1}) et (G, II, 5.2)) que l'on a un
isomorphisme canonique
\begin{equation}
\label{III.1.4.10.1-fr}
\tag{1.4.10.1}
H^\bullet(u^{-1}(V),\mathcal F)
\xrightarrow{\sim}H^\bullet(\Gamma(V,\mathcal K^\bullet))
\end{equation}
où l'on a posé
$\mathcal K^\bullet=u_*(\mathcal C^\bullet(\mathfrak U,\mathcal F))$. En
vertu de (\hyperref[III.1.4.9-fr]{1.4.9}) et (I, 9.2.2),
$\mathcal K^\bullet$ est un $\mathcal O_Y$-Module quasi-cohérent~; par
ailleurs, il constitue un \emph{complexe de faisceaux} puisqu'il en est
ainsi de $\mathcal C^\bullet(\mathfrak U,\mathcal F)$. Il résulte alors de la
définition de la cohomologie
$\mathcal H^\bullet(\mathcal K^\bullet)$ (G, II, 4.1) que cette dernière est
formée de $\mathcal O_Y$-Modules quasi-cohérents (I, 4.1.1). Comme (pour $V$
affine dans $Y$) le foncteur $\Gamma(V,\mathcal G)$ est exact en
$\mathcal G$ dans la catégorie des $\mathcal O_Y$-Modules quasi-cohérents,
on a (G, II, 4.1)
\begin{equation}
\label{III.1.4.10.2-fr}
\tag{1.4.10.2}
H^\bullet(\Gamma(V,\mathcal K^\bullet))
=\Gamma(V,\mathcal H^\bullet(\mathcal K^\bullet)).
\end{equation}
Notons enfin qu'il résulte de la définition de l'homomorphisme canonique
\[
H^\bullet(\mathfrak U,\mathcal F)\longrightarrow H^\bullet(X,\mathcal F)
\]
donnée dans (G, II, 5.2), que si $V'\subset V$ est un second ouvert affine
de $Y$, le diagramme
\[
\xymatrix{
H^\bullet(u^{-1}(V),\mathcal F)\ar[r]^{\sim}\ar[d]
&H^\bullet(\Gamma(V,\mathcal K^\bullet))\ar[d]\\
H^\bullet(u^{-1}(V'),\mathcal F)\ar[r]^{\sim}
&H^\bullet(\Gamma(V',\mathcal K^\bullet))
}
\]
est commutatif. On conclut donc de ce qui précède que les isomorphismes
(\hyperref[III.1.4.10.1-fr]{1.4.10.1}) définissent un isomorphisme de
$\mathcal O_Y$-Modules
\begin{equation}
\label{III.1.4.10.3-fr}
\tag{1.4.10.3}
R^\bullet u_*(\mathcal F)
\xrightarrow{\sim}\mathcal H^\bullet(\mathcal K^\bullet)
\end{equation}
\oldpage[III]{92}
et par suite $R^\bullet u_*(\mathcal F)$ est quasi-cohérent.
\end{proof}

Il résulte en outre de (\hyperref[III.1.4.10.3-fr]{1.4.10.3}),
(\hyperref[III.1.4.10.2-fr]{1.4.10.2}) et
(\hyperref[III.1.4.10.1-fr]{1.4.10.1}) que~:

\begin{corollary}[1.4.11]
\label{III.1.4.11-fr}
Sous les hypothèses de (\hyperref[III.1.4.10-fr]{1.4.10}), pour tout ouvert
affine $V$ de $Y$, l'homomorphisme canonique
\begin{equation}
\label{III.1.4.11.1-fr}
\tag{1.4.11.1}
H^q(u^{-1}(V),\mathcal F)\longrightarrow
\Gamma(V,R^qu_*(\mathcal F))
\end{equation}
est un isomorphisme pour tout $q\geq0$.
\end{corollary}

\begin{corollary}[1.4.12]
\label{III.1.4.12-fr}
Supposons vérifiées les hypothèses de
(\hyperref[III.1.4.10-fr]{1.4.10}) et supposons en outre que $Y$ soit
quasi-compact. Alors il existe un entier $r>0$ tel que pour tout
$\mathcal O_X$-Module quasi-cohérent $\mathcal F$ et tout entier $q>r$, on
ait $R^qu_*(\mathcal F)=0$. Si $Y$ est affine, on peut prendre pour $r$ un
entier tel qu'il existe un recouvrement de $X$ formé de $r$ ouverts affines.
\end{corollary}

\begin{proof}
Comme on peut recouvrir $Y$ par un nombre fini d'ouverts affines, on est
ramené à démontrer la seconde assertion, en vertu de
(\hyperref[III.1.4.11-fr]{1.4.11}). Or, si $\mathfrak U$ est un
recouvrement de $X$ par $r$ ouverts affines, on a
$H^q(\mathfrak U,\mathcal F)=0$ pour $q>r$, puisque les cochaînes de
$C^q(\mathfrak U,\mathcal F)$ sont alternées~; la conclusion résulte donc de
(\hyperref[III.1.4.1-fr]{1.4.1}).
\end{proof}

\begin{corollary}[1.4.13]
\label{III.1.4.13-fr}
Supposons vérifiées les hypothèses de
(\hyperref[III.1.4.10-fr]{1.4.10}) et en outre supposons que
$Y=\operatorname{Spec}(A)$ soit affine. Alors, pour tout
$\mathcal O_X$-Module quasi-cohérent $\mathcal F$ et tout $f\in A$, on a
\[
\Gamma(Y_f,R^qu_*(\mathcal F))
=(\Gamma(Y,R^qu_*(\mathcal F)))_f
\]
à un isomorphisme canonique près.
\end{corollary}

\begin{proof}
En effet, cela résulte de ce que $R^qu_*(\mathcal F)$ est un
$\mathcal O_Y$-Module quasi-cohérent (I, 1.3.7).
\end{proof}

\begin{proposition}[1.4.14]
\label{III.1.4.14-fr}
Soient $f:X\to Y$ un morphisme séparé quasi-compact, $g:Y\to Z$ un
morphisme affine. Pour tout $\mathcal O_X$-Module quasi-cohérent
$\mathcal F$, l'homomorphisme canonique
$R^p(g\circ f)_*(\mathcal F)\to g_*(R^pf_*(\mathcal F))$ (0, 12.2.5.2) est
bijectif pour tout $p$.
\end{proposition}

\begin{proof}
En effet, pour tout ouvert affine $W$ de $Z$, $g^{-1}(W)$ est un ouvert
affine de $Y$. L'homomorphisme de préfaisceaux définissant l'homomorphisme
canonique
\[
R^p(g\circ f)_*(\mathcal F)\longrightarrow g_*(R^pf_*(\mathcal F))
\]
(0, 12.2.5) est donc bijectif en vertu de
(\hyperref[III.1.4.11-fr]{1.4.11}).
\end{proof}

\begin{proposition}[1.4.15]
\label{III.1.4.15-fr}
Soient $u:X\to Y$ un morphisme séparé de type fini, $v:Y'\to Y$ un
morphisme plat de préschémas (0\textsubscript{I}, 6.7.1)~; soit
$u'=u_{(Y')}$, de sorte qu'on a le diagramme commutatif
\begin{equation}
\label{III.1.4.15.1-fr}
\tag{1.4.15.1}
\xymatrix{
X\ar[d]_u & X'=X_{(Y')}\ar[l]_{v'}\ar[d]^{u'}\\
Y & Y'\ar[l]_v
}
\end{equation}
Alors, pour tout $\mathcal O_X$-Module quasi-cohérent $\mathcal F$,
$R^qu_*'(\mathcal F')$, où
$\mathcal F'=v'^*(\mathcal F)=\mathcal F\otimes_{\mathcal O_Y}\mathcal O_{Y'}$,
est canoniquement isomorphe à
$R^qu_*(\mathcal F)\otimes_{\mathcal O_Y}\mathcal O_{Y'}
=v^*(R^qu_*(\mathcal F))$ pour tout $q\geq0$.
\end{proposition}

\oldpage[III]{93}
\begin{proof}
L'homomorphisme canonique
$\rho:\mathcal F\to v_*'(v'^*(\mathcal F))$ (0\textsubscript{I}, 4.4.3.2)
définit par fonctorialité un homomorphisme
\begin{equation}
\label{III.1.4.15.2-fr}
\tag{1.4.15.2}
R^qu_*(\mathcal F)\longrightarrow R^qu_*(v_*'(\mathcal F')).
\end{equation}
D'autre part, on a, en posant $w=u\circ v'=v\circ u'$, les homomorphismes
canoniques (0, 12.2.5.1 et 12.2.5.2)
\begin{equation}
\label{III.1.4.15.3-fr}
\tag{1.4.15.3}
R^qu_*(v_*'(\mathcal F'))\longrightarrow R^qw_*(\mathcal F')
\longrightarrow v_*(R^qu_*'(\mathcal F')).
\end{equation}
Composant (\hyperref[III.1.4.15.3-fr]{1.4.15.3}) et
(\hyperref[III.1.4.15.2-fr]{1.4.15.2}), on a un homomorphisme
\[
\psi:R^qu_*(\mathcal F)\longrightarrow v_*(R^qu_*'(\mathcal F'))
\]
et finalement on en déduit l'homomorphisme canonique (dont la définition ne
fait intervenir \emph{aucune hypothèse sur $v$})
\begin{equation}
\label{III.1.4.15.4-fr}
\tag{1.4.15.4}
\psi^\sharp:v^*(R^qu_*(\mathcal F))\longrightarrow R^qu_*'(\mathcal F')
\end{equation}
dont il s'agit de prouver que c'est un isomorphisme lorsque $v$ est
\emph{plat}. Il est clair que la question est locale sur $Y$ et $Y'$ et l'on
peut donc supposer que $Y=\operatorname{Spec}(A)$,
$Y'=\operatorname{Spec}(B)$~; nous utiliserons en outre le

\begin{lemma}[1.4.15.5]
\label{III.1.4.15.5-fr}
Soient $\varphi:A\to B$ un homomorphisme d'anneaux,
$Y=\operatorname{Spec}(A)$, $X=\operatorname{Spec}(B)$,
$f:X\to Y$ le morphisme correspondant à $\varphi$, $M$ un $B$-module. Pour
que le $\mathcal O_X$-Module $\widetilde M$ soit $f$-plat
(0\textsubscript{I}, 6.7.1), il faut et il suffit que $M$ soit un $A$-module
plat. En particulier, pour que le morphisme $f$ soit plat, il faut et il
suffit que $B$ soit un $A$-module plat.
\end{lemma}

En effet, cela résulte de la définition (0\textsubscript{I}, 6.7.1) et de
(0\textsubscript{I}, 6.3.3), compte tenu de (I, 1.3.4).

Cela étant, il résulte de
(\hyperref[III.1.4.11.1-fr]{1.4.11.1}) et des définitions des
homomorphismes (\hyperref[III.1.4.15.3-fr]{1.4.15.3})
(cf. 0, 12.2.5) que $\psi$ correspond alors au morphisme composé
\[
H^q(X,\mathcal F)\xrightarrow{\rho_q}
H^q(X,v_*'(v'^*(\mathcal F)))\xrightarrow{\theta_q}
H^q(X',v'^*(v_*'(v'^*(\mathcal F))))\xrightarrow{\sigma_q}
H^q(X',v'^*(\mathcal F))
\]
où $\rho_q$ et $\sigma_q$ sont les homomorphismes correspondant dans la
cohomologie aux morphismes canoniques $\rho$ et
$\sigma:v'^*(v_*'(\mathcal G'))\to\mathcal G'$, et $\theta_q$ est le
$\varphi$-morphisme (0, 12.1.3.1) relatif au $\mathcal O_X$-Module
$v_*'(v'^*(\mathcal F))$. Mais par fonctorialité de $\theta_q$, on a le
diagramme commutatif
\[
\xymatrix{
H^q(X,\mathcal F)\ar[rr]^{\rho_q}\ar[dd]_{\theta_q}
&&H^q(X,v_*'(v'^*(\mathcal F)))\ar[dd]^{\theta_q}\\\\
H^q(X',v'^*(\mathcal F))\ar[rr]^{v'^*(\rho_q)}
&&H^q(X',v'^*(v_*'(v'^*(\mathcal F))))
}
\]
\oldpage[III]{94}
et comme par définition (0\textsubscript{I}, 4.4.3), $v'^*(\rho)$ est
l'inverse de $\sigma$, on voit que le morphisme composé considéré plus haut
n'est autre finalement que $\theta_q$~; $\psi^\sharp$ est par suite le
$B$-homomorphisme associé
\[
H^q(X,\mathcal F)\otimes_A B\longrightarrow H^q(X',\mathcal F').
\]
Comme $u$ est de type fini, $X$ est réunion finie d'ouverts affines $U_i$
($1\leq i\leq r$)~; soit $\mathfrak U$ le recouvrement $(U_i)$. D'ailleurs
$v$ est un morphisme affine, donc il en est de même de $v'$ (II, 1.6.2,
(iii)), et les $U_i'=v'^{-1}(U_i)$ forment par suite un recouvrement ouvert
affine $\mathfrak U'$ de $X'$. On sait alors (0, 12.1.4.2) que le diagramme
\[
\xymatrix{
H^q(\mathfrak U,\mathcal F)\ar[r]^{\theta_q}\ar[d]
&H^q(\mathfrak U',\mathcal F')\ar[d]\\
H^q(X,\mathcal F)\ar[r]^{\theta_q}&H^q(X',\mathcal F')
}
\]
est commutatif, et les flèches verticales sont des isomorphismes puisque
$X$ et $X'$ sont des schémas (\hyperref[III.1.4.1-fr]{1.4.1}). Il suffit par
suite de prouver que le $\varphi$-morphisme canonique
$\theta_q:H^q(\mathfrak U,\mathcal F)\to H^q(\mathfrak U',\mathcal F')$
est tel que le $B$-homomorphisme associé
\[
H^q(\mathfrak U,\mathcal F)\otimes_A B
\longrightarrow H^q(\mathfrak U',\mathcal F')
\]
soit un isomorphisme. Pour toute suite
$\mathbf s=(i_k)_{0\leq k\leq p}$ de $p+1$ indices de $[1,r]$, posons
\[
U_{\mathbf s}=\bigcap_{k=0}^pU_{i_k},\qquad
U_{\mathbf s}'=\bigcap_{k=0}^pU_{i_k}'=v'^{-1}(U_{\mathbf s}),\qquad
M_{\mathbf s}=\Gamma(U_{\mathbf s},\mathcal F),\qquad
M_{\mathbf s}'=\Gamma(U_{\mathbf s}',\mathcal F').
\]
L'application canonique
$M_{\mathbf s}\otimes_A B\to M_{\mathbf s}'$ est un isomorphisme
(I, 1.6.5), donc l'application canonique
$C^p(\mathfrak U,\mathcal F)\otimes_A B\to
C^p(\mathfrak U',\mathcal F')$ est un isomorphisme, par lequel
$d\otimes1$ s'identifie à l'opérateur cobord
$C^p(\mathfrak U',\mathcal F')\to C^{p+1}(\mathfrak U',\mathcal F')$.
Comme $B$ est un $A$-module \emph{plat}, il résulte aussitôt de la définition
des modules de cohomologie que l'application canonique
$H^q(\mathfrak U,\mathcal F)\otimes_A B\to
H^q(\mathfrak U',\mathcal F')$ est un isomorphisme
(0\textsubscript{I}, 6.1.1). Ce résultat sera généralisé plus tard (\S~6).
\end{proof}

\begin{corollary}[1.4.16]
\label{III.1.4.16-fr}
Soient $A$ un anneau, $X$ un $A$-schéma de type fini, $B$ une $A$-algèbre
fidèlement plate sur $A$. Pour que $X$ soit affine, il faut et il suffit que
$X\otimes_A B$ le soit.
\end{corollary}

\begin{proof}
La condition est évidemment nécessaire (I, 3.2.2)~; montrons qu'elle est
suffisante. Comme $X$ est séparé sur $A$ et que le morphisme
$\operatorname{Spec}(B)\to\operatorname{Spec}(A)$ est plat, il résulte de
(1.4.1) que l'on a
\begin{equation}
\label{III.1.4.16.1-fr}
\tag{1.4.16.1}
H^i(X\otimes_A B,\mathcal F\otimes_A B)
=H^i(X,\mathcal F)\otimes_A B
\end{equation}
pour tout $i\geq0$ et tout $\mathcal O_X$-Module quasi-cohérent $\mathcal F$.
Si $X\otimes_A B$ est affine, le premier membre de
(\hyperref[III.1.4.16.1-fr]{1.4.16.1}) est nul pour $i=1$, donc il en est de
même de $H^1(X,\mathcal F)$ puisque $B$ est un $A$-module fidèlement plat.
Comme $X$ est un schéma quasi-compact, on conclut par le critère de Serre
(II, 5.2.1).
\end{proof}

\begin{proposition}[1.4.17]
\label{III.1.4.17-fr}
Soient $X$ un préschéma,
$0\to\mathcal F\xrightarrow{u}\mathcal G\xrightarrow{v}\mathcal H\to0$
une suite exacte de $\mathcal O_X$-Modules. Si $\mathcal F$ et $\mathcal H$
sont quasi-cohérents, il en est de même de $\mathcal G$.
\end{proposition}

\oldpage[III]{95}
\begin{proof}
La question étant locale sur $X$, on peut supposer
$X=\operatorname{Spec}(A)$ affine et il suffit alors de prouver que
$\mathcal G$ vérifie les conditions d~1) et d~2) de (I, 1.4.1) (avec
$V=X$). La vérification de d~2) est immédiate, car si
$t\in\Gamma(X,\mathcal G)$ a une restriction nulle à $D(f)$, il en est de
même de son image $v(t)\in\Gamma(X,\mathcal H)$~; il y a donc $m>0$ tel que
$f^mv(t)=v(f^mt)=0$ (I, 1.4.1), et comme $\Gamma$ est exact à gauche,
$f^mt=u(s)$, où $s\in\Gamma(X,\mathcal F)$~; comme $u$ est injectif, la
restriction de $s$ à $D(f)$ est nulle, d'où (I, 1.4.1) l'existence d'un
entier $n>0$ tel que $f^ns=0$~; on en déduit finalement
$f^{m+n}t=u(f^ns)=0$.

Vérifions maintenant d~2), et soit $t'\in\Gamma(D(f),\mathcal G)$~; comme
$\mathcal H$ est quasi-cohérent, il existe un entier $m$ tel que
$f^mv(t')=v(f^mt')$ se prolonge en une section
$z\in\Gamma(X,\mathcal H)$ (I, 1.4.1). Mais en vertu de
(\hyperref[III.1.3.1-fr]{1.3.1}) (ou de (I, 5.1.9.2)) appliqué au
$\mathcal O_X$-Module quasi-cohérent $\mathcal F$, la suite
$\Gamma(X,\mathcal G)\to\Gamma(X,\mathcal H)\to0$ est exacte, donc il existe
$t\in\Gamma(X,\mathcal G)$ tel que $z=v(t)$~; on voit donc que
$v(f^mt'-t'')=0$, en désignant par $t''$ la restriction de $t$ à $D(f)$~;
on a donc $f^mt'-t''=u(s')$, où $s'\in\Gamma(D(f),\mathcal F)$. Mais comme
$\mathcal F$ est quasi-cohérent, il existe un entier $n>0$ tel que $f^ns'$
se prolonge en une section $s\in\Gamma(X,\mathcal F)$~; comme
$f^{m+n}t'-f^nt''=u(f^ns')$, on voit que $f^{m+n}t'$ est la restriction à
$D(f)$ de la section $f^nt+u(f^ns)\in\Gamma(X,\mathcal G)$, ce qui achève
la démonstration.
\end{proof}

\section{Étude cohomologique des morphismes projectifs}
\label{section:III.2-fr}

\subsection{Calculs explicites de certains groupes de cohomologie}
\label{subsection:III.2.1-fr}

\begin{env}[2.1.1]
\label{III.2.1.1-fr}
Soient $X$ un préschéma, $\mathcal L$ un $\mathcal O_X$-Module inversible~;
considérons l'anneau gradué (0\textsubscript{I}, 5.4.6)
\begin{equation}
\label{III.2.1.1.1-fr}
\tag{2.1.1.1}
S=\Gamma_\bullet(X,\mathcal L)
=\bigoplus_{n\in\mathbb Z}\Gamma(X,\mathcal L^{\otimes n}).
\end{equation}
Soit $(f_i)_{1\leq i\leq r}$ une famille finie d'éléments \emph{homogènes}
de $S$, soit $f_i\in S_{d_i}$~; posons
$U_i=X_{f_i}$, $U=\bigcup_iU_i$, et désignons par $\mathfrak U$ le
recouvrement $(U_i)$ de $U$. Pour tout $\mathcal O_X$-Module quasi-cohérent
$\mathcal F$, on posera
\begin{align}
\label{III.2.1.1.2-fr}
\tag{2.1.1.2}
H^\bullet(U,\mathcal F(*))
&=\bigoplus_{n\in\mathbb Z}
H^\bullet(U,\mathcal F\otimes\mathcal L^{\otimes n}),\\
\label{III.2.1.1.3-fr}
\tag{2.1.1.3}
H^\bullet(\mathfrak U,\mathcal F(*))
&=\bigoplus_{n\in\mathbb Z}
H^\bullet(\mathfrak U,\mathcal F\otimes\mathcal L^{\otimes n}).
\end{align}
Les groupes abéliens (\hyperref[III.2.1.1.2-fr]{2.1.1.2}) et
(\hyperref[III.2.1.1.3-fr]{2.1.1.3}) sont \emph{bigradués}, en prenant
\[
(H^\bullet(U,\mathcal F(*)))_{mn}
=H^m(U,\mathcal F\otimes\mathcal L^{\otimes n})
\]
et une définition analogue pour
(\hyperref[III.2.1.1.3-fr]{2.1.1.3}). Pour le deuxième degré, il est clair
que ces groupes sont des $S$-modules gradués, comme il résulte par exemple du
fait que $\mathcal F\mapsto H^m(U,\mathcal F)$ et
$\mathcal F\mapsto H^m(\mathfrak U,\mathcal F)$ sont des foncteurs.
\end{env}

\begin{env}[2.1.2]
\label{III.2.1.2-fr}
Considérons maintenant le $S$-module gradué (0\textsubscript{I}, 5.4.6)
\begin{equation}
\label{III.2.1.2.1-fr}
\tag{2.1.2.1}
M=\Gamma_\bullet(\mathcal L,\mathcal F)
=H^0(X,\mathcal F(*))
=\bigoplus_{n\in\mathbb Z}
\Gamma(X,\mathcal F\otimes\mathcal L^{\otimes n}).
\end{equation}
\oldpage[III]{96}
Si $X$ est un préschéma dont l'espace sous-jacent est noethérien, ou un
schéma quasi-compact, il résulte de (I, 9.3.1) qu'en posant comme d'ordinaire
$U_{i_0i_1\ldots i_p}=\bigcap_{k=0}^pU_{i_k}$, on a, à un isomorphisme
canonique près,
\[
\Gamma(U_{i_0i_1\ldots i_p},\mathcal F(*))
=H^0(U_{i_0i_1\ldots i_p},\mathcal F(*))
=M_{f_{i_0}f_{i_1}\cdots f_{i_p}}.
\]
On peut encore, avec les notations de
(\hyperref[III.1.2.2-fr]{1.2.2}), identifier
$M_{f_{i_0}f_{i_1}\cdots f_{i_p}}$ à
$\varinjlim_n M_{i_0i_1\ldots i_p}^{(n)}$. Cette identification est un
isomorphisme de $S$-modules gradués, si l'on définit le degré d'un élément
homogène $z\in\varinjlim_n M_{i_0i_1\ldots i_p}^{(n)}$ de la façon
suivante~: $z$ est l'image canonique d'un élément homogène
$x\in M_{i_0i_1\ldots i_p}^{(n)}=M$, de degré $m$~; on prend alors pour
degré de $z$ le nombre
$m-n(d_{i_0}+d_{i_1}+\cdots+d_{i_p})$. Tenant compte de la définition des
homomorphismes
$\varphi_{kn}:M_{i_0i_1\ldots i_p}^{(n)}\to
M_{i_0i_1\ldots i_p}^{(k)}$ (\hyperref[III.1.2.2-fr]{1.2.2}), on voit
aussitôt que cette définition ne dépend pas du «~représentant~» $x$ de $z$
que l'on a considéré. Désignant comme dans
(\hyperref[III.1.2.2-fr]{1.2.2}) par $C_n^p(M)$ l'ensemble des applications
alternées de $[1,r]^{p+1}$ dans $M$ (pour tout $n$), on définit de la même
manière que ci-dessus une structure de $S$-module gradué sur
$\varinjlim_n C_n^p(M)$~; on a encore comme dans
(\hyperref[III.1.2.2-fr]{1.2.2})
\begin{equation}
\label{III.2.1.2.2-fr}
\tag{2.1.2.2}
C^p(\mathfrak U,\mathcal F(*))=\varinjlim_n C_n^p(M)
\end{equation}
l'isomorphisme des deux membres \emph{respectant les degrés}. On a alors,
comme dans (\hyperref[III.1.2.2-fr]{1.2.2})
\begin{equation}
\label{III.2.1.2.3-fr}
\tag{2.1.2.3}
C^p(\mathfrak U,\mathcal F(*))
=C^{p+1}((\mathbf f),M)
=\varinjlim_n K^{p+1}(\mathbf f^n,M)
\end{equation}
l'isomorphisme conservant les degrés~: le degré d'un élément de
$\varinjlim_n K^{p+1}(\mathbf f^n,M)$, image canonique d'une cochaîne
$g\in K^{p+1}(\mathbf f^n,M)$ dont les valeurs
$g(i_0,\ldots,i_p)$ sont dans une même composante homogène $M_m$ de $M$,
est $m-n(d_{i_0}+\cdots+d_{i_p})$, et il est indépendant du choix de cette
cochaîne comme représentant de l'élément considéré.
\end{env}

Comme les isomorphismes précédents sont compatibles avec les opérateurs
cobords, on en conclut, comme dans
(\hyperref[III.1.2.2-fr]{1.2.2}) que l'on a~:

\begin{proposition}[2.1.3]
\label{III.2.1.3-fr}
Soit $X$ un préschéma dont l'espace sous-jacent est noethérien, ou un schéma
quasi-compact. Il existe un isomorphisme canonique fonctoriel en $\mathcal F$
\begin{equation}
\label{III.2.1.3.1-fr}
\tag{2.1.3.1}
H^p(\mathfrak U,\mathcal F(*))
\xrightarrow{\sim}H^{p+1}((\mathbf f),M)
\qquad\text{pour tout }p\geq1.
\end{equation}
On a en outre une suite exacte fonctorielle en $\mathcal F$
\begin{equation}
\label{III.2.1.3.2-fr}
\tag{2.1.3.2}
0\longrightarrow H^0((\mathbf f),M)\longrightarrow M
\longrightarrow H^0(\mathfrak U,\mathcal F(*))
\longrightarrow H^1((\mathbf f),M)\longrightarrow0.
\end{equation}
De plus, tous les homomorphismes introduits sont de degré $0$ pour les
structures de $S$-modules gradués ($S$ étant l'anneau
(\hyperref[III.2.1.1.1-fr]{2.1.1.1})).
\end{proposition}
\providecommand{\bbGamma}{\boldsymbol{\Gamma}}
\providecommand{\shProj}{\mathcal{P}\mathrm{roj}}

\oldpage[III]{97}
\begin{corollary}[2.1.4]
\label{III.2.1.4-fr}
Si $X$ est un schéma quasi-compact et si les $U_i=X_{f_i}$ sont affines,
il existe un isomorphisme canonique fonctoriel en $\mathcal F$, de degré $0$,
\begin{equation}
\label{III.2.1.4.1-fr}
\tag{2.1.4.1}
H^p(U,\mathcal F(*))\xrightarrow{\sim}H^{p+1}((\mathbf f),M)
\qquad\text{pour }p\geq1,
\end{equation}
et une suite exacte fonctorielle en $\mathcal F$
\begin{equation}
\label{III.2.1.4.2-fr}
\tag{2.1.4.2}
0\longrightarrow H^0((\mathbf f),M)\longrightarrow M
\longrightarrow H^0(U,\mathcal F(*))
\longrightarrow H^1((\mathbf f),M)\longrightarrow0,
\end{equation}
où tous les homomorphismes sont de degré $0$.
\end{corollary}

\begin{proof}
Il suffit en effet d'appliquer (1.4.1) au résultat de
(\hyperref[III.2.1.3-fr]{2.1.3}).
\end{proof}

La proposition «~locale~» analogue à
(\hyperref[III.2.1.3-fr]{2.1.3}) est la suivante~:

\begin{proposition}[2.1.5]
\label{III.2.1.5-fr}
Soient $S$ un anneau gradué à degrés positifs,
$f_i$ $(1\leq i\leq r)$ un élément homogène de $S_+$ de degré $d_i$,
$M$ un $S$-module gradué. Soit $X=\operatorname{Proj}(S)$ le spectre
premier homogène de $S$ et posons $U_i=D_+(f_i)$,
$U=\bigcup_iU_i$,
\[
H^\bullet(U,\widetilde M(*))
=\bigoplus_{n\in\mathbb Z}H^\bullet(U,(M(n))^\sim).
\]
Il existe alors des isomorphismes canoniques fonctoriels en $M$, de degré
$0$ pour les structures de $S$-module gradué
\begin{equation}
\label{III.2.1.5.1-fr}
\tag{2.1.5.1}
H^p(U,\widetilde M(*))\xrightarrow{\sim}H^{p+1}((\mathbf f),M)
\qquad\text{pour }p\geq1,
\end{equation}
et une suite exacte fonctorielle en $M$
\begin{equation}
\label{III.2.1.5.2-fr}
\tag{2.1.5.2}
0\longrightarrow H^0((\mathbf f),M)\longrightarrow M
\longrightarrow H^0(U,\widetilde M(*))
\longrightarrow H^1((\mathbf f),M)\longrightarrow0,
\end{equation}
où tous les homomorphismes sont de degré $0$.
\end{proposition}

\begin{proof}
En effet, on a
\[
\Gamma(U_{i_0i_1\ldots i_p},(M(n))^\sim)
=(M_{f_{i_0}\cdots f_{i_p}})_n
\]
par définition (II, 2.5.2), donc
$\Gamma(U_{i_0i_1\ldots i_p},\widetilde M(*))
=M_{f_{i_0}\cdots f_{i_p}}$. Le reste du raisonnement est alors le même
que pour démontrer (\hyperref[III.2.1.4-fr]{2.1.4}), tenant compte de ce
que $X$ est un schéma.
\end{proof}

\begin{env}[2.1.6]
\label{III.2.1.6-fr}
\emph{Remarques.} ---
(i) Sous les conditions de (\hyperref[III.2.1.5-fr]{2.1.5}), les foncteurs
$\Gamma(U_{i_0i_1\ldots i_p},\widetilde M(*))$ sont exacts en $M$, en
vertu de (0\textsubscript{I}, 1.3.2)~; le même raisonnement que dans
(1.2.5) montre alors que si
$0\to M'\to M\to M''\to0$ est une suite exacte de $S$-modules gradués
(où les homomorphismes sont de degré $0$), on a des diagrammes
commutatifs pour tout $p\geq0$
\begin{equation}
\label{III.2.1.6.1-fr}
\tag{2.1.6.1}
\xymatrix{
H^p(U,\widetilde {M''}(*))\ar[r]^-\partial\ar[d] &
H^{p+1}(U,\widetilde {M'}(*))\ar[d]\\
H^{p+1}((\mathbf f),M'')\ar[r]^-\partial &
H^{p+2}((\mathbf f),M').
}
\end{equation}

(ii) La prop. (\hyperref[III.2.1.5-fr]{2.1.5}) sera surtout intéressante
lorsque $S$ sera une $A$-algèbre engendrée par un nombre fini d'éléments de
degré $1$, $A$ étant supposé noethérien~; en
\oldpage[III]{98}
effet, lorsqu'il en est ainsi, tout $\mathcal O_X$-Module quasi-cohérent est
de la forme $\widetilde M$ (II, 2.7.7).
\end{env}

\begin{env}[2.1.7]
\label{III.2.1.7-fr}
Nous allons appliquer (\hyperref[III.2.1.5-fr]{2.1.5}) dans le cas
$S=A[T_0,\ldots,T_r]$, où $A$ est un anneau quelconque, les $T_i$ sont des
indéterminées, avec $M=S$ et $f_i=T_i$. On est donc essentiellement ramené
à calculer $H^\bullet((\mathbf T),S)$, où
$\mathbf T=(T_i)_{0\leq i\leq r}$.
\end{env}

\begin{lemma}[2.1.8]
\label{III.2.1.8-fr}
Si $S=A[T_0,\ldots,T_r]$, on a, avec
$\mathbf T=(T_i)_{0\leq i\leq r}$,
\begin{align}
\label{III.2.1.8.1-fr}
\tag{2.1.8.1}
H^i(\mathbf T^n,S)&=0 &&\text{si }i\neq r+1,\\
\label{III.2.1.8.2-fr}
\tag{2.1.8.2}
H^{r+1}(\mathbf T^n,S)&=S/(\mathbf T^n).
\end{align}
Le $A$-module $H^{r+1}(\mathbf T^n,S)$ a donc une base sur $A$ formée des
classes mod. $(\mathbf T^n)$ des monômes
\[
T_{\mathbf p}=T_0^{p_0}\cdots T_r^{p_r},
\qquad \mathbf p=(p_0,\ldots,p_r),\quad 0\leq p_i<n
\quad\text{pour tout }i.
\]
\end{lemma}

\begin{proof}
C'est une conséquence immédiate de (1.1.3.5) et de la prop. (1.1.4), dont
les hypothèses sont trivialement vérifiées.
\end{proof}

\begin{env}[2.1.9]
\label{III.2.1.9-fr}
Passons à la limite inductive sur $n$~; les relations
(\hyperref[III.2.1.8.1-fr]{2.1.8.1}) donnent
$H^i((\mathbf T),S)=0$ pour $i\neq r+1$. Pour $i=r+1$, le système inductif
est formé des $S/(\mathbf T^n)$, l'homomorphisme
\[
\varphi_{nm}:S/(\mathbf T^n)\longrightarrow S/(\mathbf T^m)
\qquad(0\leq n\leq m)
\]
étant la multiplication par $(T_0\cdots T_r)^{m-n}$. Pour
$n\geq\sup(p_i)_{0\leq i\leq r}$, désignons par
$\xi_{\mathbf p}^{(n)}=\xi_{p_0,\ldots,p_r}^{(n)}$ la classe de
$T_0^{n-p_0}\cdots T_r^{n-p_r}\pmod{(\mathbf T^n)}$~; on a alors
$\varphi_{nm}(\xi_{\mathbf p}^{(n)})=\xi_{\mathbf p}^{(m)}$, et ces éléments
ont donc même image canonique
$\xi_{\mathbf p}=\xi_{p_0,\ldots,p_r}$ dans la limite inductive
$H^{r+1}((\mathbf T),S)$~; en vertu de la définition du degré donnée dans
(\hyperref[III.2.1.2-fr]{2.1.2}), le degré de $\xi_{\mathbf p}$ est donc
égal à $-|\mathbf p|=-(p_0+p_1+\cdots+p_r)$. Il est clair que les
$\xi_{\mathbf p}^{(n)}$ pour $0<p_i\leq n$ et $0\leq i\leq r$ forment une
base de $S/(\mathbf T^n)$. On déduit donc aussitôt de
(\hyperref[III.2.1.8-fr]{2.1.8})~:
\end{env}

\begin{corollary}[2.1.10]
\label{III.2.1.10-fr}
Avec les notations de (\hyperref[III.2.1.8-fr]{2.1.8}), on a
\begin{equation}
\label{III.2.1.10.1-fr}
\tag{2.1.10.1}
H^i((\mathbf T),S)=0\qquad\text{pour }i\neq r+1,
\end{equation}
et $H^{r+1}((\mathbf T),S)$ est un $A$-module libre dont une base est formée
des éléments $\xi_{p_0,\ldots,p_r}$ tels que $p_i>0$ pour $0\leq i\leq r$.
\end{corollary}

\begin{env}[2.1.11]
\label{III.2.1.11-fr}
\emph{Remarque.} --- Soit $N$ un $A$-module quelconque et soit
$M=S\otimes_A N$~; le raisonnement de
(\hyperref[III.2.1.8-fr]{2.1.8}) montre que l'on a plus généralement
\begin{align}
\label{III.2.1.11.1-fr}
\tag{2.1.11.1}
H^i(\mathbf T^n,M)&=0 &&\text{si }i\neq r+1,\\
\label{III.2.1.11.2-fr}
\tag{2.1.11.2}
H^{r+1}(\mathbf T^n,M)&=(S/(\mathbf T^n))\otimes_A N,
\end{align}
car la dernière formule se déduit directement de (1.1.3.5), et d'autre
part il est clair que
\[
M/(T_0^nM+\cdots+T_{i-1}^nM)
=\bigl(S/(T_0^nS+\cdots+T_{i-1}^nS)\bigr)\otimes_A N,
\]
l'idéal $T_0^nS+\cdots+T_{i-1}^nS$ étant facteur direct dans le
$A$-module $S$~; cela permet d'appliquer (1.1.4) à $M$, et on obtient ainsi
(\hyperref[III.2.1.11.1-fr]{2.1.11.1}).
\end{env}

Combinant (\hyperref[III.2.1.10-fr]{2.1.10}) et
(\hyperref[III.2.1.5-fr]{2.1.5}), on obtient~:

\begin{proposition}[2.1.12]
\label{III.2.1.12-fr}
Soient $A$ un anneau, $r$ un entier $>0$, et
$X=\mathbf P_A^r$ (II, 4.1.1). Alors~:
\begin{enumerate}
\item[{\rm(i)}] On a $H^i(X,\mathcal O_X(*))=0$ pour $i\neq0,r$.
\item[{\rm(ii)}] L'homomorphisme canonique
$\alpha:S\to H^0(X,\mathcal O_X(*))$ (II, 2.6.2) est bijectif.
\oldpage[III]{99}
\item[{\rm(iii)}] $H^r(X,\mathcal O_X(*))$ est un $A$-module libre ayant
une base formée d'éléments $\xi_{p_0,\ldots,p_r}$, où $p_i>0$ pour
$0\leq i\leq r$, $\xi_{p_0,\ldots,p_r}$ étant de degré
$-|\mathbf p|=-(p_0+\cdots+p_r)$, et le produit
$T_i\xi_{p_0,\ldots,p_r}$ étant
$\xi_{p_0,\ldots,p_i-1,\ldots,p_r}$.
\end{enumerate}
\end{proposition}

\begin{proof}
Remarquons en effet que, dans la suite exacte
(\hyperref[III.2.1.5.2-fr]{2.1.5.2}) appliquée à
$M=S=A[T_0,\ldots,T_r]$, on a
$H^0((\mathbf T),S)=0$ et $H^1((\mathbf T),S)=0$ d'après
(\hyperref[III.2.1.10.1-fr]{2.1.10.1}), et que la prop.
(\hyperref[III.2.1.5-fr]{2.1.5}) s'applique à $U=X$, puisque $X$ est
réunion des $D_+(T_i)$ (II, 2.3.14). Il reste à identifier l'application
$S\to H^0(X,\mathcal O_X(*))$ de la suite exacte
(\hyperref[III.2.1.5.2-fr]{2.1.5.2}) et l'application canonique $\alpha$~;
mais cela résulte de l'identification canonique de
$H^0(U,\mathcal O_X(*))$ et de $H^0(\mathfrak U,\mathcal O_X(*))$.
\end{proof}

\begin{corollary}[2.1.13]
\label{III.2.1.13-fr}
Les seules valeurs de $(i,n)$ pour lesquelles on puisse avoir
$H^i(X,\mathcal O_X(n))\neq0$ sont les suivantes~:
$i=0$ et $n\geq0$, $i=r$ et $n\leq-(r+1)$.
\end{corollary}

On notera que si $A\neq0$, on a effectivement
$H^i(X,\mathcal O_X(n))\neq0$ pour les couples énumérés dans
(\hyperref[III.2.1.13-fr]{2.1.13})~; cela résulte de
(\hyperref[III.2.1.12-fr]{2.1.12}), puisque $S_n$ est alors $\neq0$ pour
tous les degrés $n\geq0$.

Dans les applications qui seront faites dans ce chapitre, nous utiliserons
surtout le résultat moins précis~:

\begin{corollary}[2.1.14]
\label{III.2.1.14-fr}
Les $A$-modules $H^i(X,\mathcal O_X(n))$ sont libres de type fini~; si
$i>0$, ils sont nuls pour $n>0$.
\end{corollary}

\begin{proposition}[2.1.15]
\label{III.2.1.15-fr}
Soient $Y$ un préschéma, $\mathcal E$ un $\mathcal O_Y$-Module localement
libre de rang $r+1$, $X=\mathbf P(\mathcal E)$ le fibré projectif défini
par $\mathcal E$, $f:X\to Y$ le morphisme structural. Les seules valeurs de
$i$ et $n$ pour lesquelles $R^if_*(\mathcal O_X(n))\neq0$ sont
$i=0$ et $n\geq0$, $i=r$ et $n\leq-(r+1)$~; en outre, l'homomorphisme
canonique (II, 3.3.2)
\[
\alpha:S_{\mathcal O_Y}(\mathcal E)
\longrightarrow\Gamma_*(\mathcal O_X)
=R^0f_*(\mathcal O_X(*))
=\bigoplus_{n\in\mathbb Z}f_*(\mathcal O_X(n))
\]
est un isomorphisme.
\end{proposition}

\begin{proof}
La question étant locale sur $Y$, on peut supposer $Y$ affine d'anneau
$A$ et $\mathcal E=\widetilde E$, où $E=A^{r+1}$~; on est alors
immédiatement ramené à (\hyperref[III.2.1.12-fr]{2.1.12}), compte tenu de
(1.4.11).
\end{proof}

\begin{env}[2.1.16]
\label{III.2.1.16-fr}
\emph{Remarque.} --- Nous compléterons plus tard les résultats de
(\hyperref[III.2.1.15-fr]{2.1.15}) en démontrant les propositions
suivantes~: posons
\[
\boldsymbol\omega=f^*\left(\bigwedge^{r+1}\mathcal E\right)(-r-1),
\]
qui est un $\mathcal O_X$-Module inversible. Alors~:
\begin{enumerate}
\item[{\rm(i)}] On a un isomorphisme canonique
\begin{equation}
\label{III.2.1.16.1-fr}
\tag{2.1.16.1}
\rho:R^rf_*(\boldsymbol\omega)\xleftarrow{\sim}\mathcal O_Y.
\end{equation}
\item[{\rm(ii)}] L'accouplement par cup-produit (0, 12.2.2)
\begin{equation}
\label{III.2.1.16.2-fr}
\tag{2.1.16.2}
R^rf_*(\mathcal O_X(n))\times R^0f_*(\boldsymbol\omega(-n))
\longrightarrow R^rf_*(\boldsymbol\omega)
\end{equation}
composé avec l'isomorphisme $\rho^{-1}$, définit un \emph{isomorphisme}
de $R^rf_*(\mathcal O_X(n))$ sur le \emph{dual} du $\mathcal O_Y$-Module
localement libre
\[
R^0f_*(\boldsymbol\omega(-n))
=\left(\bigwedge^{r+1}\mathcal E\right)
\otimes_{\mathcal O_Y}\bigl(S_{\mathcal O_Y}(\mathcal E)\bigr)_{-n}.
\]
\end{enumerate}
\end{env}
\subsection{Le théorème fondamental des morphismes projectifs}
\label{subsection:III.2.2-fr}
\oldpage[III]{100}

\begin{theorem}[2.2.1]
\label{III.2.2.1-fr}
(Serre). Soient $Y$ un préschéma noethérien, $f:X\to Y$ un morphisme
propre, $\mathcal L$ un $\mathcal O_X$-Module inversible ample pour $f$.
Pour tout $\mathcal O_X$-Module $\mathcal F$, posons
\[
\mathcal F(n)=\mathcal F\otimes_{\mathcal O_X}\mathcal L^{\otimes n}
\qquad\text{pour tout }n\in\mathbb Z.
\]
Alors, pour tout $\mathcal O_X$-Module cohérent $\mathcal F$~:
\begin{enumerate}
\item[{\rm(i)}] Les $R^qf_*(\mathcal F)$ sont des
$\mathcal O_Y$-Modules cohérents.
\item[{\rm(ii)}] Il existe un entier $N$ tel que pour $n\geq N$, on ait
\[
R^qf_*(\mathcal F(n))=0\qquad\text{pour tout }q>0.
\]
\item[{\rm(iii)}] Il existe un entier $N$ tel que, pour $n\geq N$,
l'homomorphisme canonique
\[
f^*(f_*(\mathcal F(n)))\longrightarrow\mathcal F(n)
\]
soit surjectif.
\end{enumerate}
\end{theorem}

\begin{proof}
Notons d'abord que si le théorème est vrai quand on y remplace
$\mathcal L$ par $\mathcal L^{\otimes d}$ $(d>0)$, il est vrai sous sa
forme initiale. En effet, on peut alors écrire
\[
\mathcal F(n)
=(\mathcal F\otimes\mathcal L^{\otimes r})
\otimes\mathcal L^{\otimes hd}
\]
avec un $h>0$ et $0\leq r<d$, et par hypothèse pour chaque $r$ il y a
un entier $N_r$ tel que pour $h\geq N_r$ les propriétés (ii) et (iii)
aient lieu pour le $\mathcal O_X$-Module
$\mathcal F\otimes\mathcal L^{\otimes r}$~; prenant pour $N$ le plus grand
des $dN_r$, (ii) et (iii) auront lieu pour $n\geq N$. On peut donc
supposer $\mathcal L$ très ample relativement à $f$ (II, 4.6.11)~; il
existe par suite une $Y$-immersion ouverte dominante
$i:X\to P$, où $P=\operatorname{Proj}(\mathcal S)$, où $\mathcal S$ est
une $\mathcal O_Y$-Algèbre quasi-cohérente graduée à degrés positifs, dans
laquelle $\mathcal S_1$ est de type fini et engendre $\mathcal S$~; en
outre, $\mathcal L$ est isomorphe à $i^*(\mathcal O_P(1))$ (II, 4.4.7).
Mais comme $f$ est propre, il en est de même de $i$ (II, 5.4.4), donc
$i$ est un isomorphisme $X\xrightarrow{\sim}P$. On peut donc se borner au
cas où $X=\operatorname{Proj}(\mathcal S)$ et
$\mathcal L=\mathcal O_X(1)$. Le th. (2.2.1) est alors conséquence de la
\end{proof}

\begin{proposition}[2.2.2]
\label{III.2.2.2-fr}
Soient $A$ un anneau noethérien, $S$ une $A$-algèbre graduée à degrés
positifs, dans laquelle $S_1$ est un $A$-module ayant un système de $r+1$
générateurs, et qui engendre l'algèbre $S$. Soit
$X=\operatorname{Proj}(S)$. Pour tout $\mathcal O_X$-Module cohérent
$\mathcal F$~:
\begin{enumerate}
\item[{\rm(i)}] Les $A$-modules $H^q(X,\mathcal F)$ sont de type fini.
\item[{\rm(ii)}] On a $H^q(X,\mathcal F)=0$ pour $q>r$.
\item[{\rm(iii)}] Il existe un entier $N$ tel que pour $n\geq N$, on ait
$H^q(X,\mathcal F(n))=0$ pour tout $q>0$.
\item[{\rm(iv)}] Il existe un entier $N$ tel que pour $n\geq N$,
$\mathcal F(n)$ soit engendré par ses sections au-dessus de $X$.
\end{enumerate}
\end{proposition}

\begin{proof}
Montrons d'abord comment (2.2.2) entraîne (2.2.1)~: dans (2.2.1)
(ramené au cas particulier $X=\operatorname{Proj}(\mathcal S)$ considéré
ci-dessus), $Y$ est quasi-compact, donc peut être recouvert par un nombre
fini d'ouverts affines, d'anneaux noethériens, tels que la restriction de
$\mathcal S_1$ à chacun de ces ouverts $U_\alpha$ soit engendrée par un
nombre fini de sections de $\mathcal S_1$ au-dessus de $U_\alpha$. Si on
suppose (2.2.2) démontré, il suffira alors de prendre pour $N$ dans les
parties (ii) et (iii) de (2.2.1) le plus grand des entiers analogues
correspondant aux $U_\alpha$ (tenant compte de (1.4.11) et de
(II, 3.4.7)).

Pour prouver (2.2.2), remarquons que $X$ s'identifie à un sous-schéma
fermé de $P=\mathbf P_A^r$ (II, 3.6.2)~; en outre, si $j:X\to P$ est
l'injection canonique, $j_*(\mathcal F)$ est un
$\mathcal O_P$-Module cohérent et on a
$j_*(\mathcal F(n))=(j_*(\mathcal F))(n)$ (II, 3.4.5 et 3.5.2).
Compte tenu de (G, II, cor. du th. 4.9.1), on est donc ramené à prouver
(2.2.2) dans le cas particulier où
\[
X=\mathbf P_A^r\quad\text{et}\quad S=A[T_0,\ldots,T_r].
\]

\oldpage[III]{101}
Comme $X$ est recouvert par les ouverts affines $D_+(T_i)$ en nombre
$r+1$, (ii) résulte de (1.4.12). Notons d'autre part que (iv) a déjà été
démontré (II, 2.7.9).

Nous allons prouver simultanément (i) et (iii). Notons que ces assertions
sont vraies pour $\mathcal F=\mathcal O_X(m)$
(\hyperref[III.2.1.13-fr]{2.1.13})~; elles le sont donc aussi lorsque
$\mathcal F$ est somme directe d'un nombre fini de
$\mathcal O_X$-Modules de la forme $\mathcal O_X(m_j)$. D'autre part,
(i) et (iii) sont vraies trivialement pour $q>r$ en vertu de (ii). Nous
allons procéder par \emph{récurrence descendante} sur $q$. On sait que
$\mathcal F$ est isomorphe à un quotient d'une somme directe $\mathcal E$
d'un nombre fini de faisceaux $\mathcal O_X(m_j)$ (II, 2.7.10)~;
autrement dit, on a une suite exacte
\[
0\longrightarrow\mathcal R\longrightarrow\mathcal E
\longrightarrow\mathcal F\longrightarrow0,
\]
où $\mathcal R$ est cohérent (0\textsubscript{I}, 5.3.3) et où
$\mathcal E$ vérifie (i) et (iii). Comme $\mathcal F(n)$ est un foncteur
exact en $\mathcal F$, on a aussi la suite exacte
\[
0\longrightarrow\mathcal R(n)\longrightarrow\mathcal E(n)
\longrightarrow\mathcal F(n)\longrightarrow0
\]
pour tout $n\in\mathbb Z$. On en déduit la suite exacte de cohomologie
\[
H^{q-1}(X,\mathcal E(n))\longrightarrow
H^{q-1}(X,\mathcal F(n))\longrightarrow
H^q(X,\mathcal R(n)).
\]
Comme $\mathcal E(n)$ est somme directe des $\mathcal O_X(n+m_j)$
(II, 2.5.14), $H^{q-1}(X,\mathcal E(n))$ est de type fini, et il en est
de même de $H^q(X,\mathcal R(n))$ par l'hypothèse de récurrence~; comme
$A$ est noethérien, on en conclut que
$H^{q-1}(X,\mathcal F(n))$ est de type fini pour tout $n\in\mathbb Z$, et
en particulier pour $n=0$. D'autre part, par l'hypothèse de récurrence,
il existe un entier $N$ tel que pour $n\geq N$ on ait
$H^q(X,\mathcal R(n))=0$~; par ailleurs, on peut supposer aussi $N$ choisi
tel que $H^{q-1}(X,\mathcal E(n))=0$ pour $n\geq N$, puisque $\mathcal E$
vérifie (iii)~; on en conclut que
$H^{q-1}(X,\mathcal F(n))=0$ pour $n\geq N$, ce qui achève la
démonstration.
\end{proof}

\begin{corollary}[2.2.3]
\label{III.2.2.3-fr}
Sous les hypothèses de (\hyperref[III.2.2.1-fr]{2.2.1}), soit
$\mathcal F\to\mathcal G\to\mathcal H$ une suite exacte de
$\mathcal O_X$-Modules cohérents. Il existe alors un entier $N$ tel que
pour $n\geq N$, la suite
\[
f_*(\mathcal F(n))\longrightarrow f_*(\mathcal G(n))
\longrightarrow f_*(\mathcal H(n))
\]
soit exacte.
\end{corollary}

\begin{proof}
Soient $\mathcal F'$, $\mathcal G'$, $\mathcal G''$ le noyau, l'image et
le conoyau de $\mathcal F\to\mathcal G$~; $\mathcal G'$ est le noyau et
$\mathcal G''$ l'image de $\mathcal G\to\mathcal H$, soit $\mathcal H''$
le conoyau de cet homomorphisme~; tous ces $\mathcal O_X$-Modules sont
cohérents (0\textsubscript{I}, 5.3.4). Comme $\mathcal F(n)$ est un
foncteur exact en $\mathcal F$, il suffit de prouver que pour $n$ assez
grand, chacune des suites
\begin{align*}
0&\longrightarrow f_*(\mathcal F'(n))
\longrightarrow f_*(\mathcal F(n))
\longrightarrow f_*(\mathcal G'(n))\longrightarrow0,\\
0&\longrightarrow f_*(\mathcal G'(n))
\longrightarrow f_*(\mathcal G(n))
\longrightarrow f_*(\mathcal G''(n))\longrightarrow0,\\
0&\longrightarrow f_*(\mathcal G''(n))
\longrightarrow f_*(\mathcal H(n))
\longrightarrow f_*(\mathcal H''(n))\longrightarrow0
\end{align*}
est exacte~; par suite, on peut supposer que
$0\to\mathcal F\to\mathcal G\to\mathcal H\to0$ est exacte. On a alors
la suite exacte de cohomologie
\[
0\longrightarrow f_*(\mathcal F(n))
\longrightarrow f_*(\mathcal G(n))
\longrightarrow f_*(\mathcal H(n))
\longrightarrow R^1f_*(\mathcal F(n))\longrightarrow\cdots
\]
et la conclusion résulte de (\hyperref[III.2.2.1-fr]{2.2.1}, (ii)).
\end{proof}

\begin{corollary}[2.2.4]
\label{III.2.2.4-fr}
Soient $Y$ un préschéma noethérien, $f:X\to Y$ un morphisme de type fini,
$\mathcal L$ un $\mathcal O_X$-Module inversible ample pour $f$~; pour
tout $\mathcal O_X$-Module $\mathcal F$, on pose
$\mathcal F(n)=\mathcal F\otimes_{\mathcal O_X}\mathcal L^{\otimes n}$
(pour $n\in\mathbb Z$). Soit
$\mathcal F\to\mathcal G\to\mathcal H$ une suite exacte de
$\mathcal O_X$-Modules cohérents, telle que les supports de $\mathcal F$
et de $\mathcal H$ soient propres sur $Y$ (II, 5.4.10). Il existe alors
un entier $N$ tel que, pour $n\geq N$, la suite
\[
f_*(\mathcal F(n))\longrightarrow f_*(\mathcal G(n))
\longrightarrow f_*(\mathcal H(n))
\]
soit exacte.
\end{corollary}

\begin{proof}
Le même raisonnement qu'au début de
(\hyperref[III.2.2.1-fr]{2.2.1}) montre que si le corollaire est vrai
pour $\mathcal L^{\otimes d}$ $(d>0)$, il l'est aussi pour $\mathcal L$~;
on peut donc se borner au cas où $\mathcal L$ est très ample pour $f$
(II, 4.6.11), et par suite on peut identifier $X$ à un ouvert dans un
$Y$-schéma $Z=\operatorname{Proj}(\mathcal S)$, où $\mathcal S$ est une
$\mathcal O_Y$-Algèbre quasi-cohérente graduée à degrés positifs, dans
laquelle $\mathcal S_1$ est de type fini et engendre $\mathcal S$, de
sorte que $\mathcal L=i^*(\mathcal O_Z(1))$, où $i$ est l'immersion
canonique $X\to Z$ (II, 4.4.7). Cela étant, comme
$\operatorname{Supp}(\mathcal G)$ est fermé dans $X$ et contenu dans
$\operatorname{Supp}(\mathcal F)\cap\operatorname{Supp}(\mathcal H)$, il
est propre sur $Y$~; les supports de $\mathcal F$, $\mathcal G$,
$\mathcal H$ sont donc \emph{fermés dans $Z$} (II, 5.4.10). Les faisceaux
$\mathcal F'=i_*(\mathcal F)$, $\mathcal G'=i_*(\mathcal G)$,
$\mathcal H'=i_*(\mathcal H)$ sont donc des $\mathcal O_Z$-Modules
cohérents, et la suite $\mathcal F'\to\mathcal G'\to\mathcal H'$ est
exacte~; en outre, si $g:Z\to Y$ est le morphisme structural, on a
$f=g\circ i$, et il est clair que
$\mathcal F'(n)=i_*(\mathcal F(n))$, et de même pour $\mathcal G'$ et
$\mathcal H'$~; la conclusion résulte donc de
(\hyperref[III.2.2.3-fr]{2.2.3}) appliqué à
$\mathcal F'$, $\mathcal G'$, $\mathcal H'$.
\end{proof}

\begin{env}[2.2.5]
\label{III.2.2.5-fr}
\emph{Remarques.} ---
(i) L'assertion (i) de (\hyperref[III.2.2.1-fr]{2.2.1}) est encore vraie
lorsqu'on suppose seulement que $Y$ est \emph{localement noethérien}~; en
effet, la propriété est évidemment locale sur $Y$~; d'autre part, les
hypothèses de (\hyperref[III.2.2.1-fr]{2.2.1}) impliquent que pour tout
ouvert $U\subset Y$, la restriction de $f$ à $f^{-1}(U)$ est un morphisme
projectif $f^{-1}(U)\to U$ (II, 5.5.5, (iii)) et
$\mathcal L|f^{-1}(U)$ est ample pour ce morphisme (II, 4.6.4).

(ii) L'assertion (iii) de (\hyperref[III.2.2.1-fr]{2.2.1}) est encore
valable, comme on l'a vu, lorsque l'on suppose seulement que $X$ est un
schéma quasi-compact ou un préschéma dont l'espace sous-jacent est
noethérien, et $f:X\to Y$ un morphisme quasi-compact (II, 4.6.8). Mais il
faut noter que même lorsqu'on suppose que $Y$ est le \emph{spectre d'un
corps} $K$ et que $f$ est \emph{quasi-projectif}, l'assertion (ii) de
(\hyperref[III.2.2.1-fr]{2.2.1}) n'est plus nécessairement vérifiée. Par
exemple, soit $X'=\operatorname{Spec}(K[T_0,\ldots,T_r])$ et soit $X$ la
réunion des ouverts affines $D(T_i)$ de $X'$ $(0\leq i\leq r)$~; comme
l'immersion $X\to X'$ est quasi-compacte, le morphisme structural
$f:X\to Y$ est quasi-affine (II, 5.1.10), donc $\mathcal O_X$ est très
ample pour $f$ (II, 5.1.6). Mais l'anneau
$\Gamma(X,\mathcal O_X)$ s'identifie à l'intersection des anneaux de
fractions $(K[T_0,\ldots,T_r])_{T_i}$ pour $0\leq i\leq r$
(I, 8.2.1.1), c'est-à-dire à $K[T_0,\ldots,T_r]$. Par suite, il résulte
des formules (1.4.3.1) et (1.1.3.5) que l'on a
$H^r(X,\mathcal O_X^{\otimes n})=H^r(X,\mathcal O_X)=A\neq0$ pour tout
$n$.
\end{env}

\subsection{Application aux faisceaux gradués d'algèbres et de modules}
\label{subsection:III.2.3-fr}

\begin{theorem}[2.3.1]
\label{III.2.3.1-fr}
Soient $Y$ un préschéma noethérien, $\mathcal S$ une
$\mathcal O_Y$-Algèbre graduée à degrés positifs, quasi-cohérente et de
type fini, $X=\operatorname{Proj}(\mathcal S)$,
$q:X\to Y$ le morphisme structural, $\mathcal M$ un $\mathcal S$-Module
gradué quasi-cohérent vérifiant la condition \textbf{(TF)}. Alors il existe
un entier $N$ tel que, pour $n\geq N$, l'homomorphisme canonique
(II, 8.14.5.1)
\[
\alpha_n:\mathcal M_n\longrightarrow
q_*(\shProj_0(\mathcal M(n)))
=q_*((\shProj(\mathcal M))_n)
\]
\oldpage[III]{103}
soit bijectif. En d'autres termes, l'homomorphisme canonique
\[
\alpha:\mathcal M\longrightarrow\bbGamma_*(\shProj(\mathcal M))
\]
est un \textbf{(TN)}-isomorphisme.
\end{theorem}

\begin{proof}
On peut se borner au cas où $\mathcal M$ est un $\mathcal S$-module de
type fini (II, 3.4.2). Comme $Y$ est quasi-compact, il existe un entier
$d>0$ tel que $\mathcal S^{(d)}$ soit engendrée par le
$\mathcal O_Y$-Module quasi-cohérent $\mathcal S_d$, ce dernier étant de
type fini (II, 3.1.10), donc cohérent puisque $Y$ est noethérien.
Remarquons maintenant que $\mathcal M$ est somme directe des
$\mathcal M^{(d,k)}$ pour $0\leq k<d$ et que chacun des
$\mathcal M^{(d,k)}$ est un $\mathcal S^{(d)}$-Module quasi-cohérent de
type fini, ainsi qu'il résulte de (II, 2.1.6, (iii)), la question étant
locale sur $Y$. Or, il suffit évidemment de prouver que chacun des
homomorphismes canoniques
\[
\alpha:\mathcal M^{(d,k)}
\longrightarrow\bbGamma_*((\shProj(\mathcal M))^{(d,k)})
\]
est un \textbf{(TN)}-isomorphisme. Compte tenu de (II, 8.14.13), et
notamment du diagramme (8.14.13.4), on voit qu'on est ramené à prouver le
théorème lorsque $\mathcal S$ est engendrée par $\mathcal S_1$ et que
$\mathcal S_1$ est un $\mathcal O_Y$-Module cohérent.

Comme $Y$ est noethérien, le même raisonnement qu'au début de
(\hyperref[III.2.2.2-fr]{2.2.2}) montre qu'on peut se borner au cas où
$Y=\operatorname{Spec}(A)$, $\mathcal S=\widetilde S$,
$\mathcal M=\widetilde M$, $A$ étant un anneau noethérien, $S_1$ un
$A$-module de type fini et $M$ un $S$-module gradué de type fini.
Montrons qu'il suffit alors de prouver le théorème lorsque $M=S$. En effet,
dans le cas général, on a une suite exacte
\[
L'\longrightarrow L\longrightarrow M\longrightarrow0,
\]
où $L$ et $L'$ sont des sommes directes de modules gradués de la forme
$S(m)$. Si le résultat est vrai pour $M=S$, il l'est aussi pour $M=S(m)$,
donc pour $L$ et $L'$. Considérons alors le diagramme commutatif
\[
\xymatrix{
\widetilde {L'_n}\ar[r]\ar[d]_{\alpha_n} &
\widetilde {L_n}\ar[r]\ar[d]^{\alpha_n} &
\widetilde {M_n}\ar[r]\ar[d]^{\alpha_n} & 0\\
q_*(\widetilde {L'}(n))\ar[r] &
q_*(\widetilde L(n))\ar[r] &
q_*(\widetilde M(n))\ar[r] & 0 .
}
\]
La deuxième ligne est exacte en vertu de
(\hyperref[III.2.2.3-fr]{2.2.3}), dès que $n$ est assez grand~; comme il
en est de même de la première et que les deux flèches verticales de gauche
sont des isomorphismes, il en est de même de la troisième.

Cela étant, pour prouver le théorème lorsque $M=S$, supposons d'abord que
$S=A[T_0,\ldots,T_r]$ ($T_i$ indéterminées)~; dans ce cas, notre assertion
n'est autre que (\hyperref[III.2.1.11-fr]{2.1.11}, (ii)). Dans le cas
général, $S$ s'identifie à un quotient d'un anneau
$S'=A[T_0,\ldots,T_r]$ par un idéal gradué, donc $X$ à un sous-schéma
fermé de $X'=\mathbf P_A^r$ (II, 2.9.2). Si $j$ est l'injection canonique
$X\to X'$, $j_*(\widetilde S(n))$ n'est autre que le
$\mathcal O_{X'}$-Module $(\shProj(\widetilde S))(n)$ où $S$ est considéré
comme un $S'$-module gradué~; cela résulte en effet aussitôt de
(II, 2.8.7). Comme $j_*(\widetilde S(n))$ est un
$\mathcal O_{X'}$-Module vérifiant \textbf{(TF)}, l'homomorphisme canonique
$\alpha_n:S_n\to\Gamma(X',j_*(\widetilde S(n)))$ est bijectif pour $n$
assez grand, en vertu de ce qui précède~; cela achève la démonstration,
puisque
$\Gamma(X',j_*(\widetilde S(n)))=\Gamma(X,\widetilde S(n))$.
\end{proof}

\oldpage[III]{104}
\begin{corollary}[2.3.2]
\label{III.2.3.2-fr}
Sous les hypothèses de (\hyperref[III.2.3.1-fr]{2.3.1}), soit
\[
\mathcal S_X=\bigoplus_{n\in\mathbb Z}\mathcal O_X(n),
\]
et soit $\mathcal F$ un $\mathcal S_X$-Module gradué quasi-cohérent de
type fini. Alors $\bbGamma_*(\mathcal F)$ vérifie la condition
\textbf{(TF)}.
\end{corollary}

\begin{proof}
On a vu dans la démonstration de
(\hyperref[III.2.3.1-fr]{2.3.1}) que $X$, qui est isomorphe à
$\operatorname{Proj}(\mathcal S^{(d)})$ (II, 3.1.8), est de type fini sur
$Y$ (II, 3.4.1). Il résulte alors de (II, 8.14.9) que $\mathcal F$ est
isomorphe à un $\mathcal S_X$-Module gradué de la forme
$\shProj(\mathcal M)$, où $\mathcal M$ est un $\mathcal S$-Module gradué
quasi-cohérent de type fini. En vertu de
(\hyperref[III.2.3.1-fr]{2.3.1}), $\bbGamma_*(\mathcal F)$ est
\textbf{(TN)}-isomorphe à $\mathcal M$ et par suite vérifie
\textbf{(TF)}.
\end{proof}

\begin{env}[2.3.3]
\label{III.2.3.3-fr}
\emph{Scholie.} --- Soient $Y$ un préschéma noethérien, $\mathcal S$ une
$\mathcal O_Y$-Algèbre graduée vérifiant les conditions de
(\hyperref[III.2.3.1-fr]{2.3.1}) et $X=\operatorname{Proj}(\mathcal S)$.
Soient $K_{\mathcal S}$ la catégorie abélienne des $\mathcal S$-Modules
gradués quasi-cohérents vérifiant \textbf{(TF)}, $K'_{\mathcal S}$ la
sous-catégorie de $K_{\mathcal S}$ formée des $\mathcal S$-Modules
vérifiant \textbf{(TN)}~; enfin, soit $K_X$ la catégorie des
$\mathcal S_X$-Modules gradués quasi-cohérents de type fini $\mathcal F$
(ce qui revient à dire, puisque $\mathcal S_X$ est périodique
(II, 8.14.4 et 8.14.12), que les $\mathcal F_i$ sont des
$\mathcal O_X$-Modules cohérents). Alors les foncteurs
\[
\mathcal M\longmapsto\shProj(\mathcal M)
\quad\text{dans }K_{\mathcal S},
\qquad
\mathcal F\longmapsto\bbGamma_*(\mathcal F)
\quad\text{dans }K_X
\]
définissent, en vertu de (II, 8.14.8 et 8.14.10) et de
(\hyperref[III.2.3.2-fr]{2.3.2}), une équivalence (T, I, 1.2) de la
catégorie quotient $K_{\mathcal S}/K'_{\mathcal S}$ (T, I, 1.11) avec la
catégorie $K_X$. Lorsque $\mathcal S$ est engendrée par $\mathcal S_1$,
on peut remplacer $K_X$ par la catégorie des $\mathcal O_X$-Modules
cohérents (II, 8.14.12).
\end{env}

\begin{proposition}[2.3.4]
\label{III.2.3.4-fr}
Soit $Y$ un préschéma noethérien.
\begin{enumerate}
\item[{\rm(i)}] Soit $\mathcal S$ une $\mathcal O_Y$-Algèbre graduée à
degrés positifs, quasi-cohérente et de type fini. Soient
$X=\operatorname{Proj}(\mathcal S)$, et
\[
\mathcal S_X=\shProj(\mathcal S)
=\bigoplus_{n\in\mathbb Z}\mathcal O_X(n).
\]
Alors $\mathcal S_X$ est une $\mathcal O_X$-Algèbre graduée périodique
(II, 8.14.12) dont les composants homogènes
$(\mathcal S_X)_n=\mathcal O_X(n)$ sont des $\mathcal O_X$-Modules
cohérents, et si $d>0$ est une période de $\mathcal S_X$,
$(\mathcal S_X)_d=\mathcal O_X(d)$ est un $\mathcal O_X$-Module
inversible $Y$-ample. En outre, l'homomorphisme canonique
\[
\alpha:\mathcal S\longrightarrow\bbGamma_*(\mathcal S_X)
\]
est un \textbf{(TN)}-isomorphisme.
\item[{\rm(ii)}] Inversement, soit $q:X\to Y$ un morphisme projectif, et
soit $\mathcal S'$ une $\mathcal O_X$-Algèbre graduée, dont les composants
homogènes $\mathcal S'_n$ $(n\in\mathbb Z)$ sont des
$\mathcal O_X$-Modules cohérents, et qui admet une période $d>0$ telle que
$\mathcal S'_d$ soit un $\mathcal O_X$-Module inversible ample pour $q$.
Alors
\[
\mathcal S=\bigoplus_{n\geq0}q_*(\mathcal S'_n)
\]
est une $\mathcal O_Y$-Algèbre graduée à degrés positifs quasi-cohérente
et de type fini, et il existe un $Y$-isomorphisme
\[
r:X\xrightarrow{\sim}\operatorname{Proj}(\mathcal S)
\]
tel que $r^*(\shProj(\mathcal S))$ soit isomorphe (en tant que
$\mathcal O_X$-Algèbre graduée) à $\mathcal S'$.
\end{enumerate}
\end{proposition}

\begin{proof}
(i) Toutes les assertions ont pratiquement déjà été démontrées, la
dernière n'étant autre qu'un cas particulier de
(\hyperref[III.2.3.2-fr]{2.3.2}). Le fait que $\mathcal S_X$ est
périodique a été vu en (II, 8.14.14) et le fait qu'il y a une période
$d>0$ telle que $\mathcal O_X(d)$ soit inversible et $Y$-ample n'est
autre que (II, 4.6.18). Enfin, pour $0\leq k<d$,
$(\mathcal S_X)^{(d,k)}$ est un $(\mathcal S_X)^{(d)}$-Module de type fini
(II, 8.14.14), donc chacun des $(\mathcal S_X)_n$ est un
$\mathcal O_X$-Module quasi-cohérent de type fini, en vertu de
(II, 2.1.6, (ii))~; la question étant locale, comme $\mathcal O_X$ est
cohérent, il en est de même des $(\mathcal S_X)_n$.

(ii) Quitte à remplacer la période $d$ par un de ses multiples, on peut
supposer que $\mathcal L=\mathcal S'_d$ est un $\mathcal O_X$-Module
\emph{très ample} relativement à $q$ (II, 4.6.11). On a en outre
\[
\mathcal S'^{(d)}=\bigoplus_{n\in\mathbb Z}\mathcal L^{\otimes n}
\quad\text{par hypothèse, donc}\quad
\mathcal S^{(d)}=\bigoplus_{n\geq0}q_*(\mathcal L^{\otimes n})~;
\]
on sait (II, 3.1.8 et 3.2.9) qu'il y a un $Y$-isomorphisme $s$ de
$X'=\operatorname{Proj}(\mathcal S)$ sur
$X''=\operatorname{Proj}(\mathcal S^{(d)})$ tel que
\oldpage[III]{105}
\[
s^*(\mathcal O_{X''}(n))=\mathcal O_{X'}(nd).
\]
On établira donc l'existence d'un $Y$-isomorphisme $X\xrightarrow{\sim}X'$
si l'on prouve la

\begin{env}[2.3.4.1]
\label{III.2.3.4.1-fr}
Soient $Y$ un préschéma noethérien, $q:X\to Y$ un morphisme projectif,
$\mathcal L$ un $\mathcal O_X$-Module inversible très ample pour $q$.
Alors $\mathcal S=\bigoplus_{n\geq0}q_*(\mathcal L^{\otimes n})$ est une
$\mathcal O_Y$-Algèbre graduée quasi-cohérente de type fini, telle que
$\mathcal S_n=\mathcal S_1^n$ pour $n$ assez grand, et il existe un
$Y$-isomorphisme $r:X\xrightarrow{\sim}P=\operatorname{Proj}(\mathcal S)$
tel que $\mathcal L=r^*(\mathcal O_P(1))$.
\end{env}

Comme $q$ est un morphisme projectif, il résulte de (II, 5.4.4 et 4.4.7)
qu'il existe un $Y$-isomorphisme
$r':X\xrightarrow{\sim}P'=\operatorname{Proj}(\mathcal T)$, où
$\mathcal T$ est une $\mathcal O_Y$-Algèbre quasi-cohérente,
$\mathcal T_1$ est de type fini et engendre $\mathcal T$, et
$\mathcal L=r'^*(\mathcal O_{P'}(1))$. Si $q':P'\to Y$ est le morphisme
structural, on a
\[
\mathcal S=\bigoplus_{n\geq0}q'_*(\mathcal O_{P'}(n)).
\]
En vertu de (\hyperref[III.2.3.1-fr]{2.3.1}), l'homomorphisme canonique
$\alpha_n:\mathcal T_n\to\mathcal S_n=q'_*(\mathcal O_{P'}(n))$ est
bijectif pour $n$ assez grand~; comme $\mathcal T_n=\mathcal T_1^n$, on a
\emph{a fortiori} $\mathcal S_n=\mathcal S_1^n$ dès que $n$ est assez
grand. En outre, l'homomorphisme canonique
$\alpha:\mathcal T\to\mathcal S$ de $\mathcal O_Y$-Algèbres graduées est
un \textbf{(TN)}-isomorphisme~; par suite
\phantomsection\label{III.2.3.5.1-fr}
\[
\Phi=\operatorname{Proj}(\alpha):
\operatorname{Proj}(\mathcal S)\longrightarrow
\operatorname{Proj}(\mathcal T)
\]
est un isomorphisme (II, 3.6.1). En outre, il résulte de (II, 3.5.2) et
du fait que les $\mathcal T$-Modules gradués obtenus à partir des
$\mathcal S(n)$ par restriction des scalaires sont \textbf{(TN)}-isomorphes
aux $\mathcal T(n)$ que
\[
\Phi_*(\mathcal O_P(n))=\mathcal O_{P'}(n)
\]
pour tout $n$ (II, 3.4.2). Pour achever de prouver
(\hyperref[III.2.3.4.1-fr]{2.3.4.1}), il reste à montrer que $\mathcal S$
est une $\mathcal O_Y$-Algèbre de type fini~; or les
$\mathcal S_n=q'_*(\mathcal O_{P'}(n))$ sont des
$\mathcal O_Y$-Modules cohérents en vertu de
(\hyperref[III.2.2.1-fr]{2.2.1}), et si
$\mathcal S_n=\mathcal S_1^n$ pour $n\geq n_0$, $\mathcal S$ est
engendrée par $\bigoplus_{i\leq n_0}\mathcal S_i$, qui est cohérent~;
d'où notre assertion (I, 9.6.2).

Revenons à la démonstration de
(\hyperref[III.2.3.4-fr]{2.3.4}), dont nous reprenons les notations. Nous
avons démontré l'existence d'un $Y$-isomorphisme $r'':X\xrightarrow{\sim}X''$
tel que $r''_*(\mathcal L^{\otimes n})=\mathcal O_{X''}(n)$ pour tout
$n\in\mathbb Z$~; nous désignerons par $q''$ le morphisme structural
$X''\to Y$. Notons maintenant que $\mathcal S'$ est somme directe des
$\mathcal S'^{(d)}$-Modules gradués $\mathcal S'^{(d,k)}$~; chacun de ces
derniers est un $\mathcal S'^{(d)}$-Module quasi-cohérent de type fini, en
vertu de la périodicité de $\mathcal S'$ et de l'hypothèse que les
$\mathcal S'_n$ sont des $\mathcal O_X$-Modules de type fini
(II, 8.14.12). Posons
$\mathcal F^{(k)}=r''_*(\mathcal S'^{(d,k)})$, de sorte que les
$\mathcal F^{(k)}$ sont des $\mathcal S_{X''}$-Modules gradués
quasi-cohérents de type fini~; par suite (II, 8.14.8), l'homomorphisme
canonique
\[
\beta:\shProj(\bbGamma_*(\mathcal F^{(k)}))
\longrightarrow\mathcal F^{(k)}
\]
est un isomorphisme de $\mathcal S_{X''}$-Modules. Mais on a
$q''_*((\mathcal F^{(k)})_n)=q_*((\mathcal S'^{(d,k)})_n)$, et pour
$n\geq0$ ce dernier $\mathcal O_Y$-Module est par définition égal à
$(\mathcal S^{(d,k)})_n$. Autrement dit, l'injection canonique
$\mathcal S^{(d,k)}\to\bbGamma_*(\mathcal F^{(k)})$ est un
\textbf{(TN)}-isomorphisme~; donc
\[
\shProj(\mathcal S^{(d,k)})
=\shProj(\bbGamma_*(\mathcal F^{(k)})),
\]
et par suite
$r''{}^*(\shProj(\mathcal S^{(d,k)}))\simeq\mathcal S'^{(d,k)}$.
Il reste à remarquer que, à un isomorphisme canonique près,
\[
\shProj(\mathcal S^{(d,k)})
=s_*((\shProj(\mathcal S))^{(d,k)})
\]
(II, 8.14.13.1) pour avoir démontré l'isomorphisme de
$r^*(\shProj(\mathcal S))$ et de $\mathcal S'$.

Enfin, en vertu de (\hyperref[III.2.3.2-fr]{2.3.2}), chacun des
$\bbGamma_*(\mathcal F^{(k)})$ vérifie la condition \textbf{(TF)}, donc il
en est de même de chacun des $\mathcal S^{(d,k)}$~; en outre, les
$\mathcal S_n=q_*(\mathcal S'_n)$ sont cohérents en vertu de
(\hyperref[III.2.2.1-fr]{2.2.1}), puisque les $\mathcal S'_n$ sont
cohérents. On en conclut aussitôt que les $\mathcal S^{(d,k)}$ sont des
$\mathcal S^{(d)}$-Modules de type fini~; comme
(\hyperref[III.2.3.4.1-fr]{2.3.4.1}) montre que $\mathcal S^{(d)}$ est une
$\mathcal O_Y$-Algèbre de type fini, il en est de même de $\mathcal S$.
\end{proof}

\oldpage[III]{106}
\begin{proposition}[2.3.5]
\label{III.2.3.5-fr}
Soient $Y$ un préschéma intègre noethérien, $X$ un préschéma intègre,
$f:X\to Y$ un morphisme projectif birationnel. Il existe alors un Idéal
fractionnaire cohérent $\mathcal J\subset\mathcal R(Y)$ (II, 8.1.2) tel
que $X$ soit $Y$-isomorphe au préschéma obtenu en faisant éclater
$\mathcal J$ (II, 8.1.3). En outre, il existe un ouvert $U$ de $Y$ tel que
la restriction de $f$ à $f^{-1}(U)$ soit un isomorphisme de $f^{-1}(U)$
sur $U$ (cf. I, 6.5.5), et que $\mathcal J|U$ soit inversible.
\end{proposition}

\begin{proof}
Comme il existe un $\mathcal O_X$-Module inversible $\mathcal L$ très
ample pour $f$ (II, 4.4.2 et 5.3.2), on peut appliquer
(\hyperref[III.2.3.4.1-fr]{2.3.4.1}), et on voit que $X$ s'identifie à
$\operatorname{Proj}(\mathcal S)$, où
$\mathcal S=\bigoplus_{n\geq0}f_*(\mathcal L^{\otimes n})$. On sait en
outre que les $f_*(\mathcal L^{\otimes n})$ sont des
$\mathcal O_Y$-Modules sans torsion (I, 7.4.5), donc il en est de même du
$\mathcal O_Y$-Module $\mathcal S$, et par suite $\mathcal S$ s'identifie
canoniquement à un sous-$\mathcal O_Y$-Module de
$\mathcal S\otimes_{\mathcal O_Y}\mathcal R(Y)$ (I, 7.4.1)~; ce dernier
est un faisceau \emph{simple} (I, 7.3.6) qui est connu lorsqu'on connaît
sa restriction à un ouvert non vide, par exemple à un ouvert non vide
$U'\subset U$ tel que $\mathcal L|f^{-1}(U')$ soit isomorphe à
$\mathcal O_X|f^{-1}(U')$. Comme par hypothèse les
$f_*(\mathcal L^{\otimes n})|U'$ sont alors isomorphes à
$\mathcal O_Y|U'$, on voit que
$\mathcal S\otimes_{\mathcal O_Y}\mathcal R(Y)$ est un
$\mathcal R(Y)$-Module isomorphe à $\mathcal R(Y)[T]$, où $T$ est une
indéterminée, et $\mathcal S$ est \textbf{(TN)}-isomorphe à la
sous-$\mathcal O_Y$-Algèbre engendrée par l'image canonique de
$f_*(\mathcal L)$ dans $\mathcal S\otimes_{\mathcal O_Y}\mathcal R(Y)$
(\hyperref[III.2.3.4.1-fr]{2.3.4.1})~; mais si on identifie ce dernier
faisceau à $\mathcal R(Y)[T]$, l'image de $f_*(\mathcal L)$ s'identifie à
$\mathcal J T$, où $\mathcal J$ est un sous-$\mathcal O_Y$-Module
cohérent (\hyperref[III.2.2.1-fr]{2.2.1}) de $\mathcal R(Y)$, dont la
restriction à $U'$ est isomorphe à $\mathcal O_Y|U'$, et qui par suite
est tel que $\mathcal J|U$ soit inversible. On voit alors que $\mathcal S$
est \textbf{(TN)}-isomorphe à $\bigoplus_{n\geq0}\mathcal J^n$, ce qui
achève la démonstration.
\end{proof}

\begin{corollary}[2.3.6]
\label{III.2.3.6-fr}
Sous les hypothèses de (\hyperref[III.2.3.5-fr]{2.3.5}), supposons en
outre que, pour tout sous-$\mathcal O_Y$-Module cohérent
$\mathcal J\neq0$ de $\mathcal R(Y)$, il existe un
$\mathcal O_Y$-Module inversible $\mathcal L$ tel que
\[
\Gamma(Y,\mathcal L\otimes_{\mathcal O_Y}
\mathcal{H}om(\mathcal J,\mathcal O_Y))\neq0~;
\]
alors, dans l'énoncé de (\hyperref[III.2.3.5-fr]{2.3.5}), on peut supposer
que $\mathcal J$ est un Idéal de $\mathcal O_Y$. Cette condition
supplémentaire est toujours vérifiée s'il existe un
$\mathcal O_Y$-Module ample.
\end{corollary}

\begin{proof}
En effet, on a (0\textsubscript{I}, 5.4.2)
\[
\mathcal L\otimes\mathcal{H}om(\mathcal J,\mathcal O_Y)
=\mathcal{H}om(\mathcal L^{-1},\mathcal{H}om(\mathcal J,\mathcal O_Y))
=\mathcal{H}om(\mathcal J\otimes\mathcal L^{-1},\mathcal O_Y)~;
\]
l'hypothèse signifie donc qu'il y a un homomorphisme non nul $u$ de
$\mathcal J\otimes\mathcal L^{-1}$ dans $\mathcal O_Y$. Comme, pour tout
$y\in Y$, $(\mathcal J\otimes\mathcal L^{-1})_y$ s'identifie à un
sous-$\mathcal O_y$-Module du corps des fractions $\mathcal R(Y)_y$ de
$\mathcal O_y$ (I, 7.1.5), $u_y$ est nécessairement injectif, donc $u$
est un isomorphisme de $\mathcal J\otimes\mathcal L^{-1}$ sur un Idéal
$\mathcal J'$ de $\mathcal O_Y$. Mais comme
\[
\operatorname{Proj}\left(\bigoplus_{n\geq0}\mathcal J^n\right)
\quad\text{et}\quad
\operatorname{Proj}\left(
\bigoplus_{n\geq0}(\mathcal J\otimes\mathcal L^{-1})^n\right)
\]
sont $Y$-isomorphes (II, 3.1.8), cela prouve la première assertion du
corollaire. Pour démontrer la seconde, notons que
$\mathcal F=\mathcal{H}om_{\mathcal O_Y}(\mathcal J,\mathcal O_Y)$ est
cohérent et $\neq0$, puisqu'il existe un ouvert $U$ de $Y$ tel que
$\mathcal J|U$ soit inversible. Si $\mathcal L$ est un
$\mathcal O_Y$-Module ample, il existe un entier $n$ tel que
$\mathcal F(n)=\mathcal F\otimes\mathcal L^{\otimes n}$ soit engendré par
ses sections au-dessus de $Y$ (II, 4.5.5)~; \emph{a fortiori}, on a
$\Gamma(Y,\mathcal F(n))\neq0$, d'où la conclusion.
\end{proof}

\begin{corollary}[2.3.7]
\label{III.2.3.7-fr}
Soient $X$ et $Y$ deux schémas intègres, projectifs sur un corps $k$, et
soit $f:X\to Y$ un $k$-morphisme birationnel. Alors $X$ est $k$-isomorphe
à un $Y$-schéma obtenu en faisant éclater un sous-schéma fermé $Y'$
(non nécessairement réduit) de $Y$.
\end{corollary}

\oldpage[III]{107}
\begin{proof}
En effet, $f$ est projectif (II, 5.5.5, (v)) et comme $Y$ est projectif
sur $k$, la condition supplémentaire de
(\hyperref[III.2.3.6-fr]{2.3.6}) est vérifiée~; il suffit alors de
considérer le sous-schéma fermé $Y'$ de $Y$ défini par l'Idéal cohérent
$\mathcal J$ du cor. (\hyperref[III.2.3.6-fr]{2.3.6}).
\end{proof}

\begin{env}[2.3.8]
\label{III.2.3.8-fr}
\emph{Remarque.} Au chap. IV, en étudiant la notion de diviseur, nous
verrons que si, dans l'énoncé de
(\hyperref[III.2.3.5-fr]{2.3.5}), on suppose que les anneaux
$\mathcal O_y$ $(y\in Y)$ sont factoriels (ce qui est le cas par exemple
si $Y$ est non singulier), alors $X$ peut se déduire de $Y$ en faisant
éclater un sous-préschéma fermé $Y'$ de $Y$ dont l'espace sous-jacent est
contenu dans $Y-U$.
\end{env}

\subsection{Une généralisation du théorème fondamental}
\label{subsection:III.2.4-fr}

\begin{theorem}[2.4.1]
\label{III.2.4.1-fr}
Soient $Y$ un préschéma noethérien, $\mathcal S$ une
$\mathcal O_Y$-Algèbre quasi-cohérente de type fini. Soient
$f:X\to Y$ un morphisme projectif, $\mathcal S'=f^*(\mathcal S)$,
$\mathcal M$ un $\mathcal S'$-Module quasi-cohérent de type fini.
Alors~:
\begin{enumerate}
\item[(i)] Pour tout $p\in\mathbb Z$, $R^pf_*(\mathcal M)$ est un
$\mathcal S$-Module de type fini.
\item[(ii)] Soit de plus $\mathcal L$ un $\mathcal O_X$-Module inversible
ample pour $f$, et posons
$\mathcal M(n)=\mathcal M\otimes\mathcal L^{\otimes n}$ pour tout
$n\in\mathbb Z$. Il existe un entier $N$ tel que, pour $n\geq N$, on ait
\[
R^pf_*(\mathcal M(n))=0
\tag{2.4.1.1}\label{III.2.4.1.1-fr}
\]
pour tout $p>0$, et que l'homomorphisme canonique
\[
f^*(f_*(\mathcal M(n)))\longrightarrow\mathcal M(n)
\]
(0\textsubscript{I}, 4.4.3) soit surjectif.
\end{enumerate}
\end{theorem}

\begin{proof}
Posons $Y'=\operatorname{Spec}(\mathcal S)$,
$X'=\operatorname{Spec}(\mathcal S')$ de sorte que
$X'=X\times_Y Y'$ (II, 1.5.5)~; soient $g:Y'\to Y$,
$g':X'\to X$ les morphismes structuraux, qui sont affines par
définition, et $f'=f_{(Y')}:X'\to Y'$~; on a donc un diagramme
commutatif
\[
\xymatrix{
X\ar[d]_f & X'\ar[l]_{g'}\ar[d]^{f'}\ar[dl]^{h}\\
Y & Y'\ar[l]^{g}
}
\]
et le morphisme $f'$ est projectif (II, 5.5.5, (iii))~; posons
$h=f\circ g'=g\circ f'$.

(i) Soit $\widetilde{\mathcal M}$ le $\mathcal O_{X'}$-Module associé au
$\mathcal S'$-Module quasi-cohérent $\mathcal M$, quand $X'$ est
considéré comme un $X$-schéma affine (II, 1.4.3), de sorte que l'on a
$\mathcal M=g'_*(\widetilde{\mathcal M})$~; comme $\mathcal M$ est un
$\mathcal S'$-Module de type fini, $\widetilde{\mathcal M}$ est un
$\mathcal O_{X'}$-Module de type fini (II, 1.4.5)~; comme $h$ est de type
fini, puisque $g$ et $f'$ le sont (II, 1.3.7 et I, 6.3.4, (ii)), $X'$ est
noethérien (I, 6.3.7) et $\widetilde{\mathcal M}$ est par suite cohérent.
Cela étant, comme $g'$ est affine, l'homomorphisme canonique
$R^pf_*(\mathcal M)\to R^ph_*(\widetilde{\mathcal M})$ est bijectif
(\hyperref[III.1.3.4-fr]{1.3.4}). En outre, cet homomorphisme est un
homomorphisme de $\mathcal S$-Modules~; en effet, de l'homomorphisme
canonique
\[
g'_*(\mathcal O_{X'})\otimes_{\mathcal O_X}
g'_*(\widetilde{\mathcal M})\longrightarrow
g'_*(\widetilde{\mathcal M})
\tag{2.4.1.2}\label{III.2.4.1.2-fr}
\]
qui définit la structure de $\mathcal S'$-Module de $\mathcal M$ (en se
rappelant que $\mathcal S'=g'_*(\mathcal O_{X'})$), on déduit
canoniquement un homomorphisme
\[
f_*(g'_*(\mathcal O_{X'}))\otimes
R^pf_*(g'_*(\widetilde{\mathcal M}))\longrightarrow
R^pf_*(g'_*(\widetilde{\mathcal M})).
\]

\oldpage[III]{108}
(0, 12.2.2), et comme (\hyperref[III.2.4.1.2-fr]{2.4.1.2}) provient
lui-même (par application de (0\textsubscript{I}, 4.2.2.1)) de
l'homomorphisme
\[
\mathcal O_{X'}\otimes\widetilde{\mathcal M}
\longrightarrow\widetilde{\mathcal M}
\]
définissant la structure de $\mathcal O_{X'}$-Module de
$\widetilde{\mathcal M}$, le diagramme
\[
\xymatrix{
f_*(g'_*(\mathcal O_{X'}))\otimes
R^pf_*(g'_*(\widetilde{\mathcal M}))
\ar[r]\ar[d] &
R^pf_*(g'_*(\widetilde{\mathcal M}))\ar[d]\\
h_*(\mathcal O_{X'})\otimes R^ph_*(\widetilde{\mathcal M})
\ar[r] & R^ph_*(\widetilde{\mathcal M})
}
\]
est commutatif (0, 12.2.6)~; composant les flèches horizontales avec
l'homomorphisme provenant de l'homomorphisme canonique
\[
\mathcal S\longrightarrow f_*(f^*(\mathcal S))=f_*(\mathcal S')
=f_*(g'_*(\mathcal O_{X'}))=h_*(\mathcal O_{X'}),
\]
on obtient notre assertion. D'autre part, puisque $g$ est affine et que
$f'$ est séparé et quasi-compact, l'homomorphisme canonique
$R^ph_*(\widetilde{\mathcal M})\to
g_*(R^pf'_*(\widetilde{\mathcal M}))$ est bijectif
(\hyperref[III.1.4.14-fr]{1.4.14}), et on démontre comme ci-dessus que
c'est un isomorphisme de $\mathcal S$-Modules (en utilisant cette fois la
commutativité de (0, 12.2.6.2)). Or, $f'$ étant projectif et
$\widetilde{\mathcal M}$ cohérent,
$R^pf'_*(\widetilde{\mathcal M})$ est un
$\mathcal O_{Y'}$-Module cohérent en vertu de
(\hyperref[III.2.2.1-fr]{2.2.1})~; on en conclut que
$g_*(R^pf'_*(\widetilde{\mathcal M}))$ est un $\mathcal S$-Module de type
fini (II, 1.4.5).

(ii) Soit $\mathcal L'=g'^*(\mathcal L)$, qui est un
$\mathcal O_{X'}$-Module inversible~; pour tout $n\in\mathbb Z$, on a
\[
g'_*(\widetilde{\mathcal M}\otimes\mathcal L'^{\otimes n})
=g'_*(\widetilde{\mathcal M})\otimes\mathcal L^{\otimes n}
=\mathcal M(n)
\]
(0\textsubscript{I}, 5.4.10) à un isomorphisme près~; on peut appliquer
à $\widetilde{\mathcal M}\otimes\mathcal L'^{\otimes n}$ le
raisonnement fait dans (i) pour $\widetilde{\mathcal M}$, qui prouve que
\[
R^pf_*(g'_*(\widetilde{\mathcal M}\otimes\mathcal L'^{\otimes n}))
\quad\hbox{est isomorphe à}\quad
g_*(R^pf'_*(\widetilde{\mathcal M}\otimes\mathcal L'^{\otimes n})).
\]
Or $\mathcal L'$ est ample pour $f'$ (II, 4.6.13, (iii)), donc il résulte
de (\hyperref[III.2.2.1-fr]{2.2.1}) qu'il existe un entier $N$ tel que
\[
R^pf'_*(\widetilde{\mathcal M}\otimes\mathcal L'^{\otimes n})=0
\]
pour tout $p$ et tout $n\geq N$, ce qui prouve
(\hyperref[III.2.4.1.1-fr]{2.4.1.1}). Enfin, il résulte encore de
(\hyperref[III.2.2.1-fr]{2.2.1}) qu'on peut supposer $N$ tel que pour
$n\geq N$, l'homomorphisme canonique
\[
f'^*(f'_*(\widetilde{\mathcal M}\otimes\mathcal L'^{\otimes n}))
\longrightarrow\widetilde{\mathcal M}\otimes\mathcal L'^{\otimes n}
\]
soit surjectif~; comme $g'_*$ est un foncteur exact (II, 1.4.4),
l'homomorphisme correspondant
\[
g'_*(f'^*(f'_*(\widetilde{\mathcal M}\otimes
\mathcal L'^{\otimes n})))
\longrightarrow
g'_*(\widetilde{\mathcal M}\otimes\mathcal L'^{\otimes n})
=\mathcal M(n)
\]
est surjectif. Or, on a
\[
g'_*(f'^*(f'_*(\widetilde{\mathcal M}\otimes
\mathcal L'^{\otimes n})))
=f^*(g_*(f'_*(\widetilde{\mathcal M}\otimes
\mathcal L'^{\otimes n})))
\]
(II, 1.5.2) et comme $g_*\circ f'_* = f_*\circ g'_*$, on voit
finalement que l'on a
\[
g'_*(f'^*(f'_*(\widetilde{\mathcal M}\otimes\mathcal L'^{\otimes n})))
=f^*(f_*(g'_*(\widetilde{\mathcal M}\otimes
\mathcal L'^{\otimes n})))=f^*(f_*(\mathcal M(n))),
\]
ce qui achève la démonstration.
\end{proof}

\begin{env}[2.4.2]
\label{III.2.4.2-fr}
Nous aurons en particulier à appliquer
(\hyperref[III.2.4.1-fr]{2.4.1}) lorsque $\mathcal S$ est une
$\mathcal O_Y$-Algèbre graduée à degrés positifs,
$\mathcal M=\bigoplus_{k\in\mathbb Z}\mathcal M_k$ un
$\mathcal S'$-Module gradué. Alors (avec les

\oldpage[III]{109}
mêmes hypothèses de finitude sur $\mathcal S$ et $\mathcal M$) on conclut
de (\hyperref[III.2.4.1-fr]{2.4.1}) que
\[
\bigoplus_{k\in\mathbb Z}R^pf_*(\mathcal M_k)
\]
est un $\mathcal S$-Module de type fini pour tout $p$, et (sous les
hypothèses supplémentaires de (\hyperref[III.2.4.1-fr]{2.4.1}), (ii))
qu'il existe $N$ tel que pour $n\geq N$, on ait
$R^pf_*(\mathcal M_k(n))=0$ pour tout $p>0$ et tout $k\in\mathbb Z$, et
que l'homomorphisme canonique
$f^*(f_*(\mathcal M_k(n)))\to\mathcal M_k(n)$ soit surjectif pour tout
$k\in\mathbb Z$.
\end{env}

\subsection{Caractéristique d'Euler-Poincaré et polynôme de Hilbert}
\label{subsection:III.2.5-fr}

\begin{env}[2.5.1]
\label{III.2.5.1-fr}
Soient $A$ un anneau artinien, $X$ un $A$-schéma projectif sur
$Y=\operatorname{Spec}(A)$. Pour tout $\mathcal O_X$-Module cohérent
$\mathcal F$, les $H^i(X,\mathcal F)$ $(i\geq0)$ sont des $A$-modules de
type fini (\hyperref[III.2.2.1-fr]{2.2.1}), donc ici de longueur finie
puisque $A$ est artinien. On sait en outre
(\hyperref[III.2.2.1-fr]{2.2.1}) que $H^i(X,\mathcal F)=0$ sauf pour un
nombre fini de valeurs de $i\geq0$~; le nombre entier
\[
\chi_A(\mathcal F)=\sum_{i=0}^{\infty}(-1)^i
\operatorname{long}(H^i(X,\mathcal F))
\tag{2.5.1.1}\label{III.2.5.1.1-fr}
\]
est donc défini pour tout $\mathcal O_X$-Module cohérent $\mathcal F$.
Lorsque $A$ est un anneau local artinien, on dit que $\chi_A(\mathcal F)$
est la caractéristique d'Euler-Poincaré de $\mathcal F$ (par rapport à
l'anneau $A$). Pour $\mathcal F=\mathcal O_X$, on dit que
$\chi_A(\mathcal O_X)$ est le genre arithmétique de $X$ (par rapport à
$A$).
\end{env}

\begin{proposition}[2.5.2]
\label{III.2.5.2-fr}
Soit
\[
0\longrightarrow\mathcal F'\longrightarrow\mathcal F
\longrightarrow\mathcal F''\longrightarrow0
\]
une suite exacte de $\mathcal O_X$-Modules cohérents~; on a alors
\[
\chi_A(\mathcal F)=\chi_A(\mathcal F')+\chi_A(\mathcal F'').
\tag{2.5.2.1}\label{III.2.5.2.1-fr}
\]
\end{proposition}

\begin{proof}
Comme les modules de cohomologie de $\mathcal F$, $\mathcal F'$,
$\mathcal F''$ sont nuls sauf un nombre fini d'entre eux, il y a un
entier $r>0$ tel que la suite exacte de cohomologie s'écrive
\[
0\to H^0(X,\mathcal F')\to H^0(X,\mathcal F)
\to H^0(X,\mathcal F'')\to H^1(X,\mathcal F')\to\cdots
\]
\[
\cdots\to H^r(X,\mathcal F')\to H^r(X,\mathcal F)
\to H^r(X,\mathcal F'')\to0.
\]
Or, on sait que dans une suite exacte de $A$-modules de longueur finie,
ayant des 0 aux deux extrémités, la somme alternée des longueurs est
nulle (0, 11.10.1)~; appliquant ce résultat, on trouve immédiatement la
formule (\hyperref[III.2.5.2.1-fr]{2.5.2.1}).
\end{proof}

On notera que le résultat de (\hyperref[III.2.5.2-fr]{2.5.2}) s'applique
toutes les fois que l'on sait que $X$ est un $A$-schéma quasi-compact et
que les $A$-modules $H^i(X,\mathcal F)$ sont de type fini pour tout
$\mathcal O_X$-Module cohérent $\mathcal F$
(\hyperref[III.1.4.12-fr]{1.4.12}).

\begin{theorem}[2.5.3]
\label{III.2.5.3-fr}
Soient $A$ un anneau local artinien, $X$ un schéma projectif sur
$Y=\operatorname{Spec}(A)$, $\mathcal L$ un $\mathcal O_X$-Module
inversible très ample relativement à $Y$, $\mathcal F$ un
$\mathcal O_X$-Module cohérent~; on pose
$\mathcal F(n)=\mathcal F\otimes_{\mathcal O_X}\mathcal L^{\otimes n}$
pour tout $n\in\mathbb Z$.
\begin{enumerate}
\item[(i)] Il existe un polynôme unique $P\in\mathbb Q[T]$ tel que
$\chi_A(\mathcal F(n))=P(n)$ pour tout $n\in\mathbb Z$ (on dit que $P$
est le polynôme de Hilbert de $\mathcal F$ par rapport à $A$).
\item[(ii)] Pour $n$ assez grand, on a
$\chi_A(\mathcal F(n))=\operatorname{long}_A
\Gamma(X,\mathcal F(n))$.
\item[(iii)] Le terme de plus haut degré de $\chi_A(\mathcal F(n))$ a un
coefficient $\geq0$.
\end{enumerate}
\end{theorem}

Ajoutons qu'au chap. IV, dans le paragraphe consacré à la notion de
dimension, nous démontrerons en outre que le degré de
$\chi_A(\mathcal F(n))$ est égal à la dimension du support de
$\mathcal F$.

\begin{proof}
Comme on a $H^i(X,\mathcal F(n))=0$ pour tout $i>0$ dès que $n$ est
assez grand (\hyperref[III.2.2.1-fr]{2.2.1}),

\oldpage[III]{110}
on a
\[
\chi_A(\mathcal F(n))=\operatorname{long}H^0(X,\mathcal F(n))
=\operatorname{long}\Gamma(X,\mathcal F(n))
\]
pour $n$ assez grand, d'où (ii)~; cela entraîne
$\chi_A(\mathcal F(n))\geq0$ pour $n$ assez grand, et (iii) résulte donc
de (i)~; comme d'ailleurs l'assertion d'unicité de (i) est immédiate, il
reste à prouver l'existence du polynôme $P$.

Montrons d'abord qu'on peut supposer que $\mathfrak m\mathcal F=0$, où
$\mathfrak m$ est l'idéal maximal de $A$. En effet, il existe un entier
$s>0$ tel que $\mathfrak m^s=0$, et $\mathcal F(n)$ admet donc une
filtration finie
\[
\mathcal F(n)\supset\mathfrak m\mathcal F(n)\supset\cdots
\supset\mathfrak m^{s-1}\mathcal F(n)\supset0.
\]
Par récurrence, on déduit de
(\hyperref[III.2.5.2.1-fr]{2.5.2.1}) que
\[
\chi_A(\mathcal F(n))=\sum_{k=1}^{s}
\chi_A(\mathfrak m^{k-1}\mathcal F(n)/\mathfrak m^k\mathcal F(n))~;
\]
comme
$\mathfrak m^{k-1}\mathcal F(n)/\mathfrak m^k\mathcal F(n)
=\mathcal F'_k(n)$, où
$\mathcal F'_k=\mathfrak m^{k-1}\mathcal F/\mathfrak m^k\mathcal F$,
cela prouve notre assertion.

Supposons donc $\mathfrak m\mathcal F=0$~; si $X'$ est le sous-schéma
fermé de $X$, image réciproque par le morphisme structural
$X\to\operatorname{Spec}(A)$ de l'unique point fermé de
$\operatorname{Spec}(A)$, et $j:X'\to X$ l'injection canonique, on a
$\mathcal F=j_*(\mathcal F')$, où $\mathcal F'$ est un
$\mathcal O_{X'}$-Module cohérent~; $X'$ est un schéma projectif sur
$\operatorname{Spec}(K)$, où $K=A/\mathfrak m$. Si
$\mathcal L'=j^*(\mathcal L)$, $\mathcal L'$ est très ample relativement
à $\operatorname{Spec}(K)$ (II, 4.4.10), et on a
$\mathcal F(n)=j_*(\mathcal F'(n))$, où
$\mathcal F'(n)=\mathcal F'\otimes_{\mathcal O_{X'}}
\mathcal L'^{\otimes n}$ (0\textsubscript{I}, 5.4.10). On en conclut que
$\chi_A(\mathcal F(n))=\chi_K(\mathcal F'(n))$ (G, II, 4.9.1), et on est
donc ramené au cas où $A$ est un corps.

Notons maintenant que $X$ peut être considéré comme un sous-schéma
fermé de $P=\mathbf P_A^r$ pour un $r$ convenable (II, 5.5.4, (ii))~;
si $i:X\to P$ est l'injection canonique, on voit comme ci-dessus que l'on
a $\chi_A(\mathcal F(n))=\chi_A(i_*(\mathcal F)(n))$, de sorte que l'on
peut se borner au cas où $X=\mathbf P_A^r=\operatorname{Proj}(S)$ avec
$S=A[T_0,\ldots,T_r]$, $A$ étant un corps.

Cela étant, on a $\mathcal F=\widetilde M$, où $M$ est un $S$-module
gradué de type fini (II, 2.7.8)~; il existe par suite une résolution finie
de $M$ par des $S$-modules gradués libres de type fini
\[
0\to L_q\to L_{q-1}\to\cdots\to L_1\to M\to0
\]
en vertu du th. des syzygies de Hilbert (M, VIII, 6.5)~; comme
$M\mapsto\widetilde M$ est un foncteur exact en $M$ (II, 2.5.4), on a
aussi une suite exacte
\[
0\to\widetilde L_q\to\widetilde L_{q-1}\to\cdots
\to\widetilde L_1\to\widetilde M\to0
\]
et par suite, pour tout $n\in\mathbb Z$, la suite
\[
0\to\widetilde L_q(n)\to\widetilde L_{q-1}(n)\to\cdots
\to\widetilde L_1(n)\to\widetilde M(n)\to0
\]
est exacte~; appliquant par récurrence sur $q$ la prop. (2.5.1), il vient
\[
\chi_A(\widetilde M(n))=\sum_{j=1}^{q}(-1)^{j+1}
\chi_A(\widetilde L_j(n))
\]
et pour prouver (i), on est donc ramené au cas où $M$ est libre et gradué
de type fini, donc au cas où $M=S(h)$ pour un $h\in\mathbb Z$. Comme
alors $\widetilde M(n)=(M(n))^\sim=(S(n+h))^\sim$ (II, 2.5.15), on voit
finalement que le théorème résultera du
\end{proof}

\oldpage[III]{111}
\begin{lemma}[2.5.3.1]
\label{III.2.5.3.1-fr}
Soient $A$ un corps, $r$ un entier $>0$, et $X=\mathbf P_A^r$~; on a
alors
\[
\chi_A(\mathcal O_X(n))=\binom{n+r}{r}
\]
pour tout $n\in\mathbb Z$.
\end{lemma}

\begin{proof}
En effet, pour $n\geq0$, on a
$\chi_A(\mathcal O_X(n))=\operatorname{long}H^0(X,\mathcal O_X(n))$, qui
est le nombre de monômes en les $T_i$ de degré total $n$, c'est-à-dire
$\binom{n+r}{r}$ (\hyperref[III.2.1.12-fr]{2.1.12}). Pour
$n\leq-r-1$, on a de même
\[
\chi_A(\mathcal O_X(n))=(-1)^r
\operatorname{long}H^r(X,\mathcal O_X(n))~;
\]
si $n=-r-h$, la dimension de $H^r(X,\mathcal O_X(n))$ sur $A$ est le
nombre des suites $(p_i)_{0\leq i\leq r}$ d'entiers $p_i>0$ tels que
$\sum_{i=0}^{r}p_i=r+h$ (\hyperref[III.2.1.12-fr]{2.1.12}), ou encore le
nombre des suites d'entiers $q_i\geq0$ $(0\leq i\leq r)$ tels que
$\sum_{i=0}^{r}q_i=h-1$~; c'est donc le nombre
\[
\binom{h+r-1}{r}=(-1)^r\binom{n+r}{r}.
\]
Enfin, pour $-r\leq n<0$, on a $\binom{n+r}{r}=0$ et d'autre part
$H^i(X,\mathcal O_X(n))=0$ pour tout $i\geq0$
(\hyperref[III.2.1.12-fr]{2.1.12}), ce qui prouve le lemme.
\end{proof}

\begin{corollary}[2.5.4]
\label{III.2.5.4-fr}
Soient $A$ un anneau local artinien, $S$ une $A$-algèbre graduée de type
fini engendrée par $S_1$, $M$ un $S$-module gradué de type fini,
$X=\operatorname{Proj}(S)$. On a alors
\[
\chi_A(\widetilde M(n))=\operatorname{long}M_n
\]
pour $n$ assez grand.
\end{corollary}

\begin{proof}
Cela résulte de ce que $M_n$ et $\Gamma(X,\widetilde M(n))$ sont
isomorphes pour $n$ assez grand
(\hyperref[III.2.3.1-fr]{2.3.1}).
\end{proof}

\subsection{Application : critères d'amplitude}
\label{subsection:III.2.6-fr}

\begin{proposition}[2.6.1]
\label{III.2.6.1-fr}
Soient $Y$ un préschéma noethérien, $f:X\to Y$ un morphisme propre,
$\mathcal L$ un $\mathcal O_X$-Module inversible. Les conditions
suivantes sont équivalentes~:
\begin{enumerate}
\item[a)] $\mathcal L$ est ample pour $f$.
\item[b)] Pour tout $\mathcal O_X$-Module cohérent $\mathcal F$, il
existe un entier $N$ tel que pour $n\geq N$, on ait
\[
R^qf_*(\mathcal F\otimes\mathcal L^{\otimes n})=0
\]
pour tout $q>0$.
\item[c)] Pour tout Idéal cohérent $\mathcal J$ de $\mathcal O_X$, il
existe un entier $N$ tel que pour $n\geq N$ on ait
\[
R^1f_*(\mathcal J\otimes\mathcal L^{\otimes n})=0.
\]
\end{enumerate}
\end{proposition}

\begin{proof}
On a vu que a) entraîne b) (\hyperref[III.2.2.1-fr]{2.2.1}, (ii)). Il
est trivial que b) entraîne c), et il reste à prouver que c) implique
a). On peut se borner au cas où $Y$ est affine (II, 4.6.4), et prouver
dans ce cas que $\mathcal L$ est ample~; il suffira de montrer que lorsque
$h$ parcourt l'ensemble des sections des $\mathcal L^{\otimes n}$
$(n>0)$ au-dessus de $X$, ceux des $X_h$ qui sont affines forment un
recouvrement de $X$ (II, 4.5.2). Pour cela, montrons que pour tout point
fermé $x$ de $X$ et tout voisinage ouvert affine $U$ de $x$, il existe un
$n$ et un $h\in\Gamma(X,\mathcal L^{\otimes n})$ tels que
$x\in X_h\subset U$~; $X_h$ sera nécessairement affine (I, 1.3.6) et la
réunion de ces $X_h$ sera un ouvert de $X$ contenant tous les points
fermés de $X$, et par suite $X$ lui-même puisque $X$ est noethérien
(I, 6.3.7 et 0\textsubscript{I}, 2.1.3). Soit $\mathcal J$ (resp.
$\mathcal J'$) le faisceau quasi-cohérent d'idéaux de $\mathcal O_X$
définissant le sous-préschéma fermé réduit de $X$ ayant pour espace
sous-jacent $X-U$ (resp. $(X-U)\cup\{x\}$) (I, 5.2.1)~; il est clair que
$\mathcal J$ et $\mathcal J'$ sont cohérents (I, 6.1.1), que
$\mathcal J'\subset\mathcal J$ et que
$\mathcal J''=\mathcal J/\mathcal J'$ est un $\mathcal O_X$-Module
cohérent (0\textsubscript{I}, 5.3.3) ayant pour support $\{x\}$ et tel
que $\mathcal J''_x=k(x)$. Comme $\mathcal L$ est localement libre, la
suite
\[
0\to\mathcal J'\otimes\mathcal L^{\otimes n}
\to\mathcal J\otimes\mathcal L^{\otimes n}
\to\mathcal J''\otimes\mathcal L^{\otimes n}\to0
\]
est exacte pour tout $n$, et par hypothèse il existe $n$ assez grand tel
que $H^1(X,\mathcal J'\otimes\mathcal L^{\otimes n})=0$~; la suite
exacte de cohomologie

\oldpage[III]{112}
prouve donc que l'homomorphisme
\[
\Gamma(X,\mathcal J\otimes\mathcal L^{\otimes n})
\longrightarrow
\Gamma(X,\mathcal J''\otimes\mathcal L^{\otimes n})
\]
est surjectif. Une section $g$ de
$\mathcal J''\otimes\mathcal L^{\otimes n}$ au-dessus de $X$ telle que
$g(x)\neq0$ est donc l'image d'une section
$h\in\Gamma(X,\mathcal J\otimes\mathcal L^{\otimes n})
\subset\Gamma(X,\mathcal L^{\otimes n})$ (car en vertu de
(0\textsubscript{I}, 5.4.1),
$\mathcal J\otimes\mathcal L^{\otimes n}$ est un sous-$\mathcal O_X$-Module
de $\mathcal L^{\otimes n}$)~; on a par définition $h(x)\neq0$ et
$h(z)=0$ pour $z\notin U$, ce qui achève la démonstration.
\end{proof}

\begin{proposition}[2.6.2]
\label{III.2.6.2-fr}
Soient $Y$ un préschéma noethérien, $f:X\to Y$ un morphisme de type
fini, $g:X'\to X$ un morphisme fini surjectif, $\mathcal L$ un
$\mathcal O_X$-Module inversible et $\mathcal L'=g^*(\mathcal L)$. On
suppose vérifiée la condition suivante~: Il existe une partie $Z$ de
$X$, propre sur $Y$ (II, 5.4.10) telle que pour tout $x\in X-Z$, ou bien
$X$ est normal au point $x$, ou bien $(g_*(\mathcal O_{X'}))_x$ est un
$\mathcal O_x$-module libre. Dans ces conditions, pour que $\mathcal L$
soit ample pour $f$, il faut et il suffit que $\mathcal L'$ soit ample
pour $f\circ g$.
\end{proposition}

\begin{proof}
\begin{env}[2.6.2.1]
\label{III.2.6.2.1-fr}
Puisque $g$ est affine, la condition est nécessaire (II, 5.1.12). Pour
voir qu'elle est suffisante, on peut supposer $Y$ affine (II, 4.6.4).
Montrons en outre qu'on peut se borner au cas où $X$ est réduit. En
effet, soit $j:X_{\mathrm{red}}\to X$ l'injection canonique, et posons
$X_1=X_{\mathrm{red}}$, $X'_1=X'\times_X X_1$, de sorte que l'on a le
diagramme commutatif
\[
\xymatrix{
X'\ar[d]_g & X'_1\ar[l]_{j'}\ar[d]^{g_1\, .}\\
X & X_1\ar[l]^{j}
}
\tag{2.6.2.2}\label{III.2.6.2.2-fr}
\]
Le morphisme $f\circ j$ est alors de type fini (I, 6.3.4) et $g_1$ est
un morphisme fini (II, 6.1.5, (iii))~; si $\mathcal L'$ est ample pour
$f\circ g$, $j'^*(\mathcal L')$ est ample pour $f\circ g\circ j'$ puisque
$j'$ est une immersion fermée (II, 5.1.12 et I, 4.3.2). Si on pose
$Z_1=j^{-1}(Z)$, $Z_1$ est propre sur $Y$ (II, 5.4.10)~; d'autre part,
si $X$ est normal en un point $x$, il en est évidemment de même de
$X_{\mathrm{red}}$~; enfin, si $(g_*(\mathcal O_{X'}))_x$ est un
$\mathcal O_x$-module libre, il résulte aussitôt de (II, 1.5.2) que
$((g_1)_*(\mathcal O_{X'_1}))_x$ est un $\mathcal O_{X_1,x}$-module
libre. Enfin, comme $X$ est noethérien (I, 6.3.7), si
$j^*(\mathcal L)$ est ample, $\mathcal L$ est ample (II, 4.5.14), et
comme $j'^*(\mathcal L')=g_1^*(j^*(\mathcal L))$, cela achève la
réduction annoncée. Nous supposons donc désormais $Y$ affine et $X$
réduit.
\end{env}

Les hypothèses de (II, 6.6.11) étant alors vérifiées, il existe un
$Y$-préschéma réduit $X_2$ et un $Y$-morphisme $h:X_2\to X$ fini et
birationnel tel que la restriction de $h$ à $h^{-1}(X-Z)$ soit un
isomorphisme sur $X-Z$ et que $h^*(\mathcal L)$ soit ample. Remplaçant
$X'$ par $X_2$, on voit qu'on est ramené à démontrer la proposition en
supposant en outre que $g$ possède les propriétés qui viennent d'être
énumérées pour $h$. Nous désignerons encore par $Z$ un sous-préschéma de
$X$ ayant $Z$ pour espace sous-jacent, qui est propre sur $Y$
(II, 5.4.10).

\begin{env}[2.6.2.3]
\label{III.2.6.2.3-fr}
Soient maintenant $X_1$ un sous-préschéma fermé de $X$,
$j:X_1\to X$ l'injection canonique,
$X'_1=g^{-1}(X_1)=X'\times_X X_1$ son image réciproque,
$j':X'_1\to X'$ l'injection canonique, de sorte que l'on a le diagramme
commutatif (\hyperref[III.2.6.2.2-fr]{2.6.2.2})~; posons
$\mathcal L_1=j^*(\mathcal L)$,
$\mathcal L'_1=j'^*(\mathcal L')=g_1^*(\mathcal L_1)$, de sorte que
$\mathcal L'_1$ est ample pour $f\circ g\circ j'$ (II, 5.1.12). Si on
pose $Z_1=j^{-1}(Z)$, le sous-préschéma fermé $Z_1$ de $X_1$ est propre
sur $Y$ (II, 5.4.2, (ii)). En d'autres termes, les hypothèses de
(\hyperref[III.2.6.2-fr]{2.6.2}) sont vérifiées pour $X_1$,
$\mathcal L_1$, $g_1$ et $Z_1$.
\end{env}

\oldpage[III]{113}
Ceci va nous permettre de démontrer
(\hyperref[III.2.6.2-fr]{2.6.2}) par récurrence noethérienne
(0\textsubscript{I}, 2.2.2) dans le cas où la restriction de $g$ à
$g^{-1}(X-Z)$ est un isomorphisme sur $X-Z$ (ce qui est suffisant pour
notre propos, comme on l'a vu en
(\hyperref[III.2.6.2.1-fr]{2.6.2.1}))~: il suffira d'établir que si, pour
tout sous-préschéma fermé $X_1$ de $X$, dont l'espace sous-jacent est
$\neq X$, la conclusion de (\hyperref[III.2.6.2-fr]{2.6.2}) est vraie
pour le faisceau $\mathcal L_1$, alors elle est vraie aussi pour le
faisceau $\mathcal L$.

\begin{env}[2.6.2.4]
\label{III.2.6.2.4-fr}
Soient alors $\mathcal A=\mathcal O_X$, $\mathcal B=g_*(\mathcal O_{X'})$,
de sorte que $\mathcal B$ est une sous-$\mathcal A$-Algèbre de
$\mathcal R(X)$, qui est un $\mathcal A$-Module cohérent~; en outre, la
restriction $\mathcal B|(X-Z)$ est égale à $\mathcal A|(X-Z)$. Soit
$\mathcal K$ le conducteur de $\mathcal B$ sur $\mathcal A$, c'est-à-dire
le plus grand sous-$\mathcal A$-Module quasi-cohérent de $\mathcal A$ tel
que $\mathcal B\mathcal K\subset\mathcal A$ (ou encore l'annulateur du
$\mathcal A$-Module $\mathcal B/\mathcal A$
(0\textsubscript{I}, 5.3.7)), ce qui entraîne
$\mathcal B\mathcal K=\mathcal K$. Il est clair que
$\mathcal K_x=\mathcal A_x$ en tous les points admettant un voisinage
$W_x$ tel que $g$ soit un isomorphisme de $g^{-1}(W_x)$ sur $W_x$, et en
particulier en tous les points de $X-Z$ et dans un voisinage de tout
point générique d'une composante irréductible de $X$. Considérons alors
le sous-préschéma fermé $Z_1=\operatorname{Spec}(\mathcal A/\mathcal K)$
de $X$ défini par $\mathcal K$~; il est encore propre sur $Y$ car le
sous-espace $Z_1$ est fermé dans $Z$ (II, 5.4.10). De plus, la définition
de $\mathcal K$ montre que
$\mathcal B|(X-Z_1)=\mathcal A|(X-Z_1)$~; on voit donc qu'on peut toujours
se ramener au cas où
$Z=\operatorname{Spec}(\mathcal A/\mathcal K)$, et comme on a vu que
$X-Z_1$ est un ouvert non vide de $X$, on peut toujours supposer que
l'espace $Z$ est distinct de $X$.
\end{env}

\begin{env}[2.6.2.5]
\label{III.2.6.2.5-fr}
Considérons $X'$ comme égal à $\operatorname{Spec}(\mathcal B)$ (puisque
$g$ est affine) et soit $\mathcal K'=\widetilde{\mathcal K}$, Idéal
cohérent de $\mathcal O_{X'}$ tel que $g_*(\mathcal K')=\mathcal K$
(II, 1.4.1)~; le sous-préschéma fermé
$Z'=g^{-1}(Z)=Z\times_X X'$ de $X'$ est défini par $\mathcal K'$ et égal
à $\operatorname{Spec}(\mathcal B/\mathcal K)$ (II, 1.4.10)~; comme
$h:Z'\to Z$ est un morphisme fini (II, 6.1.5, (iii)), $Z'$ est propre
sur $Y$ (II, 6.1.11 et 5.4.2, (ii)).

Cela posé, il nous faut prouver que pour tout $x\in X$ et tout voisinage
ouvert $U$ de $x$, il existe une section $s$ d'un
$\mathcal L^{\otimes n}$ $(n>0)$ au-dessus de $X$ telle que
$x\in X_s\subset U$ (II, 4.5.2)~; nous distinguerons deux cas~:

$1^\circ$ On a $x\in X-Z$~; on peut évidemment supposer alors que l'on a
aussi $U\subset X-Z$, donc l'ouvert $U'=g^{-1}(U)$ ne rencontre pas
$Z'$. Comme $\mathcal L'$ est ample par hypothèse, il existe un $n>0$ et
une section $s'$ de $\mathcal L'^{\otimes n}$ au-dessus de $X'$ telle
que $x'=g^{-1}(x)\in X'_{s'}\subset g^{-1}(U)$ (II, 4.5.2). En outre, on
peut supposer que $\mathcal K'\otimes\mathcal L'^{\otimes n}$ soit
engendré par ses sections au-dessus de $X'$ (II, 4.5.5), donc, comme
$\mathcal K'_{x'}=\mathcal O_{x'}$, il y a une de ces sections $s''$
telle que $s''(x')\neq0$~; en la multipliant par $s'$ (ce qui revient à
remplacer $n$ par $2n$), on voit qu'on peut supposer aussi que
$x'\in X'_{s''}\subset g^{-1}(U)$. Cela étant, il résulte de
(0\textsubscript{I}, 5.4.10) que l'on a un isomorphisme canonique
\[
\Gamma(X,\mathcal K\otimes\mathcal L^{\otimes n})
\xrightarrow{\sim}
\Gamma(X',\mathcal K'\otimes\mathcal L'^{\otimes n}).
\]
La section $s$ de $\mathcal K\otimes\mathcal L^{\otimes n}$ qui
correspond à $s''$ par cet isomorphisme a évidemment les propriétés
voulues.

$2^\circ$ On a $x\in Z$. Soit $\mathcal J$ l'Idéal cohérent de
$\mathcal O_X$ définissant le sous-préschéma fermé réduit de $X$ ayant
pour espace sous-jacent $X-U$, et considérons dans $\mathcal B$ les deux

\oldpage[III]{114}
Idéaux cohérents $\mathcal J\mathcal B$ et
$\mathcal J\mathcal K=\mathcal J(\mathcal K\mathcal B)
=\mathcal K(\mathcal J\mathcal B)$, de sorte que l'on a le diagramme
d'inclusions
\[
\xymatrix@R=1.5em{
\mathcal J\mathcal B\ar@{^{(}->}[r] & \mathcal B\\
\mathcal J\ar@{^{(}->}[r]\ar@{^{(}->}[u] &
\mathcal A\ar@{^{(}->}[u]\\
\mathcal J\mathcal K\mathcal B=\mathcal J\mathcal K
\ar@{^{(}->}[r]\ar@{^{(}->}[u] &
\mathcal K\ar@{^{(}->}[u]
}
\tag{2.6.2.6}\label{III.2.6.2.6-fr}
\]
Soit $\mathcal J'$ l'Idéal cohérent $(\mathcal J\mathcal B)^\sim$ de
$\mathcal O_{X'}$, de sorte que $\mathcal J\mathcal B=g_*(\mathcal J')$,
$\mathcal J'\mathcal K'=(\mathcal J\mathcal K\mathcal B)^\sim$, et par
suite
\[
\mathcal J'/\mathcal J'\mathcal K'
=(\mathcal J\mathcal B/\mathcal J\mathcal K\mathcal B)^\sim
\]
(II, 1.4.4). Comme $\mathcal J|V=\mathcal J\mathcal K|V$ pour tout
ouvert $V$ ne rencontrant pas $Z$, on voit que le support de
$\mathcal J'/\mathcal J'\mathcal K'$ est contenu dans $Z'$. Comme $Z'$ est
propre sur $Y$, on peut appliquer (\hyperref[III.2.2.4-fr]{2.2.4}) et on
voit que pour $n$ assez grand, l'application canonique
\[
\Gamma(X',\mathcal J'\otimes\mathcal L'^{\otimes n})
\longrightarrow
\Gamma(X',(\mathcal J'/\mathcal J'\mathcal K')
\otimes\mathcal L'^{\otimes n})
\]
est surjective. Mais en vertu de (0\textsubscript{I}, 5.4.10), on en
conclut que l'application canonique
\[
\Gamma(X,\mathcal J\mathcal B\otimes\mathcal L^{\otimes n})
\longrightarrow
\Gamma(X,(\mathcal J\mathcal B/\mathcal J\mathcal K\mathcal B)
\otimes\mathcal L^{\otimes n})
\]
est surjective.

Cela étant, soient $i:Z\to X$ l'injection canonique,
$i':Z'\to X'$ l'injection canonique, de sorte qu'on a le diagramme
commutatif
\[
\xymatrix{
X'\ar[d]_g & Z'\ar[l]_{i'}\ar[d]^{h}\\
X & Z\ar[l]^{i}
}
\]
Soient $\mathcal M=i^*(\mathcal L)$, $\mathcal M'=i'^*(\mathcal L')$~;
comme $\mathcal L'$ est ample, $\mathcal M'$ est ample (II, 5.1.12), et
d'autre part $\mathcal M'=h^*(\mathcal M)$~; on conclut donc de
l'hypothèse de récurrence noethérienne (puisque $Z\neq X$) que
$\mathcal M$ est ample. Par suite
$i^*(\mathcal J/\mathcal J\mathcal K)\otimes\mathcal M^{\otimes n}$ est
engendré par ses sections au-dessus de $Z$ pour $n$ assez grand
(II, 4.5.5). Comme
$\mathcal J/\mathcal J\mathcal K=i_*(i^*(\mathcal J/\mathcal J\mathcal K))$,
on déduit encore de (0\textsubscript{I}, 5.4.10) qu'il existe une section
$s$ de $(\mathcal J/\mathcal J\mathcal K)\otimes
\mathcal L^{\otimes n}$ au-dessus de $X$ (pour un $n$ assez grand) telle
que $s(x)\neq0$, puisque l'on a $\mathcal J_x=\mathcal A_x$ par
définition de $\mathcal J$ et $\mathcal K_x\neq\mathcal A_x$ par
hypothèse. Le diagramme (\hyperref[III.2.6.2.6-fr]{2.6.2.6}) montre que
$s$ est aussi une section de
$(\mathcal J\mathcal B/\mathcal J\mathcal K\mathcal B)
\otimes\mathcal L^{\otimes n}$ au-dessus de $X$, donc $s$ est l'image
canonique d'une section $t$ de
$(\mathcal J\mathcal B)\otimes\mathcal L^{\otimes n}$ au-dessus de $X$.
Mais par définition, l'image canonique $s$ de $t$ mod.
$(\mathcal J\mathcal K\mathcal B)\otimes\mathcal L^{\otimes n}$ est
dans
\[
(\mathcal J\otimes\mathcal L^{\otimes n})/
(\mathcal J\mathcal K\mathcal B\otimes\mathcal L^{\otimes n})
\]
donc, en vertu de (\hyperref[III.2.6.2.6-fr]{2.6.2.6}), cela implique que
$t$ est une section de $\mathcal J\otimes\mathcal L^{\otimes n}$
au-dessus de $X$, et \emph{a fortiori} une section de
$\mathcal L^{\otimes n}$. On a vu ci-dessus que $t(x)\neq0$, donc
$x\in X_t$, et par définition de $\mathcal J$, $t(y)=0$ dans $X-U$ qui
est le support de $\mathcal O_X/\mathcal J$~; donc $X_t\subset U$, ce qui
achève la démonstration.
\end{env}
\end{proof}

\oldpage[III]{115}
\begin{env}[2.6.3]
\label{III.2.6.3-fr}
\emph{Remarque.} Lorsque $X$ est propre sur $Y$, on peut démontrer
(\hyperref[III.2.6.2-fr]{2.6.2}) plus simplement, en raisonnant comme dans
le th. de Chevalley (II, 6.7.1), à l'aide de
(\hyperref[III.2.6.1-fr]{2.6.1}) et du lemme (II, 6.7.1.1).
\end{env}

\section{Le théorème de finitude pour les morphismes propres}
\label{section:III.3-fr}

\subsection{Le lemme de dévissage}
\label{subsection:III.3.1-fr}

\begin{definition}[3.1.1]
\label{III.3.1.1-fr}
Soit $K$ une catégorie abélienne. On dit qu'un sous-ensemble $K'$ de
l'ensemble des objets de $K$ est \emph{exact} si $0\in K'$ et si, pour toute
suite exacte
\[
0\longrightarrow A'\longrightarrow A\longrightarrow A''\longrightarrow0
\]
dans $K$ telle que deux des objets $A$, $A'$, $A''$ soit dans $K'$, alors le
troisième est aussi dans $K'$.
\end{definition}

\begin{theorem}[3.1.2]
\label{III.3.1.2-fr}
Soit $X$ un préschéma noethérien~; on désigne par $K$ la catégorie abélienne
des $\mathcal O_X$-Modules cohérents. Soient $K'$ un sous-ensemble exact de
$K$, $X'$ une partie fermée de l'espace sous-jacent à $X$. On suppose que
pour toute partie fermée irréductible $Y$ de $X'$, de point générique $y$, il
existe un $\mathcal O_X$-Module $\mathcal G\in K'$ tel que $\mathcal G_y$ soit
un $k(y)$-espace vectoriel de dimension $1$. Alors tout $\mathcal O_X$-Module
cohérent de support contenu dans $X'$ appartient à $K'$ (et en particulier,
si $X'=X$, on a $K'=K$).
\end{theorem}

\begin{proof}
Considérons la propriété suivante $P(Y)$ d'une partie fermée $Y$ de $X'$~:
tout $\mathcal O_X$-Module cohérent de support contenu dans $Y$ appartient à
$K'$. En vertu du principe de récurrence noethérienne (0\textsubscript{I},
2.2.2), on voit qu'on est ramené à démontrer que \emph{si $Y$ est une partie
fermée de $X'$ telle que la propriété $P(Y')$ soit vraie pour toute partie
fermée $Y'$ de $Y$, distincte de $Y$, alors $P(Y)$ est vraie}.

Soit donc $\mathcal F\in K$ à support contenu dans $Y$ et prouvons que
$\mathcal F\in K'$. Désignons encore par $Y$ le sous-préschéma fermé réduit de
$X$ ayant $Y$ pour espace sous-jacent (I, 5.2.1)~; il est défini par un Idéal
cohérent $\mathcal J$ de $\mathcal O_X$. On sait (I, 9.3.4) qu'il existe un
entier $n>0$ tel que $\mathcal J^n\mathcal F=0$~; pour $1\leq k\leq n$, on a
donc une suite exacte
\[
0\longrightarrow
\mathcal J^{k-1}\mathcal F/\mathcal J^k\mathcal F
\longrightarrow \mathcal F/\mathcal J^k\mathcal F
\longrightarrow \mathcal F/\mathcal J^{k-1}\mathcal F
\longrightarrow0
\]
de $\mathcal O_X$-Modules cohérents (0\textsubscript{I}, 5.3.6 et 5.3.3)~;
comme $K'$ est exact, on voit, par récurrence sur $k$, qu'il suffit de montrer
que chacun des
$\mathcal F_k=\mathcal J^{k-1}\mathcal F/\mathcal J^k\mathcal F$ est dans
$K'$. On est donc ramené à prouver que $\mathcal F\in K'$ sous l'hypothèse
supplémentaire que $\mathcal J\mathcal F=0$~; il revient au même de dire que
$\mathcal F=j_*(j^*(\mathcal F))$, où $j$ est l'injection canonique
$Y\to X$. Distinguons maintenant deux cas~:
\begin{enumerate}
\item[a)] $Y$ est \emph{réductible}. Soient $Y=Y'\cup Y''$, $Y'$ et $Y''$
étant des parties fermées de $Y$, distinctes de $Y$~; soient encore $Y'$,
$Y''$ les sous-préschémas fermés réduits de $X$ ayant pour espaces
sous-jacents $Y'$, $Y''$ respectivement, qui sont définis respectivement par
les Idéaux cohérents $\mathcal J'$, $\mathcal J''$ de $\mathcal O_X$. Posons
\[
\mathcal F'=\mathcal F\otimes_{\mathcal O_X}
(\mathcal O_X/\mathcal J'),\qquad
\mathcal F''=\mathcal F\otimes_{\mathcal O_X}
(\mathcal O_X/\mathcal J'').
\]
Les homomorphismes canoniques $\mathcal F\to\mathcal F'$,
$\mathcal F\to\mathcal F''$ définissent donc un homomorphisme
$u:\mathcal F\to\mathcal F'\oplus\mathcal F''$. Montrons que pour tout
$z\notin Y'\cap Y''$, l'homomorphisme
$u_z:\mathcal F_z\to\mathcal F'_z\oplus\mathcal F''_z$ est \emph{bijectif}.
En effet, on a $\mathcal J'\cap\mathcal J''=\mathcal J$, car la question est
locale et

\oldpage[III]{116}
l'égalité précédente résulte de (I, 5.2.1 et 1.1.5)~; si $z\notin Y''$, on a
donc $\mathcal J'_z=\mathcal J_z$, d'où $\mathcal F'_z=\mathcal F_z$ et
$\mathcal F''_z=0$, ce qui établit notre assertion dans ce cas~; on raisonne
de même si $z\notin Y'$. Par suite, le noyau et le conoyau de $u$, qui sont
dans $K$ (0\textsubscript{I}, 5.3.4) ont leurs supports dans $Y'\cap Y''$,
et sont donc dans $K'$ par hypothèse~; pour la même raison, $\mathcal F'$ et
$\mathcal F''$ sont dans $K'$, donc aussi $\mathcal F'\oplus\mathcal F''$,
puisque $K'$ est exact. La conclusion résulte alors de la considération des
deux suites exactes
\[
0\longrightarrow\operatorname{Im}u\longrightarrow
\mathcal F'\oplus\mathcal F''\longrightarrow
\operatorname{Coker}u\longrightarrow0,
\]
\[
0\longrightarrow\operatorname{Ker}u\longrightarrow\mathcal F
\longrightarrow\operatorname{Im}u\longrightarrow0
\]
et de l'hypothèse que $K'$ est exact.

\item[b)] $Y$ est irréductible, et par suite le sous-préschéma $Y$ de $X$ est
\emph{intègre}. Si $y$ est son point générique, on a
$(\mathcal O_Y)_y=k(y)$, et comme $j^*(\mathcal F)$ est un
$\mathcal O_Y$-Module cohérent, $\mathcal F_y=(j^*(\mathcal F))_y$ est un
$k(y)$-espace vectoriel de dimension finie $m$. Par hypothèse, il y a un
$\mathcal O_X$-Module cohérent $\mathcal G$ (nécessairement de support $Y$)
tel que $\mathcal G_y$ soit un $k(y)$-espace vectoriel de dimension $1$. Par
suite, il y a un $k(y)$-isomorphisme
$(\mathcal G_y)^m\xrightarrow{\sim}\mathcal F_y$, qui est aussi un
$\mathcal O_Y$-isomorphisme, et comme $\mathcal G^m$ et $\mathcal F$ sont
cohérents, il existe un voisinage ouvert $W$ de $y$ dans $X$ et un
isomorphisme $\mathcal G^m|W\xrightarrow{\sim}\mathcal F|W$
(0\textsubscript{I}, 5.2.7). Soit $\mathcal H$ le graphe de cet isomorphisme,
qui est un sous-$(\mathcal O_X|W)$-Module cohérent de
$(\mathcal G^m\oplus\mathcal F)|W$, canoniquement isomorphe à
$\mathcal G^m|W$ et à $\mathcal F|W$~; il existe donc un sous-$\mathcal O_X$-
Module cohérent $\mathcal H_0$ de $\mathcal G^m\oplus\mathcal F$, induisant
$\mathcal H$ sur $W$ et $0$ sur $X-Y$ puisque $\mathcal G^m$ et
$\mathcal F$ ont pour support $Y$ (I, 9.4.7). Les restrictions
$v:\mathcal H_0\to\mathcal G^m$ et $w:\mathcal H_0\to\mathcal F$ des
projections canoniques de $\mathcal G^m\oplus\mathcal F$ sont alors des
homomorphismes de $\mathcal O_X$-Modules cohérents, qui, dans $W$ et dans
$X-Y$, se réduisent à des isomorphismes~; autrement dit, les noyaux et
conoyaux de $v$ et $w$ ont leur support dans l'ensemble fermé
$Y-(Y\cap W)$, distinct de $Y$. Ils appartiennent donc à $K'$~; d'autre
part, on a $\mathcal G^m\in K'$ puisque $\mathcal G\in K'$ et que $K'$ est
exact. On en conclut successivement, par l'exactitude de $K'$, que
$\mathcal H_0\in K'$ puis que $\mathcal F\in K'$. C.Q.F.D.
\end{enumerate}
\end{proof}

\begin{corollary}[3.1.3]
\label{III.3.1.3-fr}
Supposons que le sous-ensemble exact $K'$ de $K$ ait en outre la propriété
que tout facteur direct cohérent d'un $\mathcal O_X$-Module cohérent
$\mathcal M\in K'$ appartienne encore à $K'$. Dans ces conditions, la
conclusion de (\hyperref[III.3.1.2-fr]{3.1.2}) subsiste encore lorsque la
condition «~$\mathcal G_y$ est un $k(y)$-espace vectoriel de dimension $1$~»
est remplacée par $\mathcal G_y\neq0$ (ce qui équivaut à
$\operatorname{Supp}(\mathcal G)=Y$).
\end{corollary}

\begin{proof}
En effet, le raisonnement de (\hyperref[III.3.1.2-fr]{3.1.2}) ne doit être
modifié que dans le cas b)~; cette fois $\mathcal G_y$ est un $k(y)$-espace
vectoriel de dimension $q>0$, et on a par suite un $\mathcal O_Y$-
isomorphisme $(\mathcal G_y)^m\xrightarrow{\sim}(\mathcal F_y)^q$~; la fin
du raisonnement de (\hyperref[III.3.1.2-fr]{3.1.2}) prouve alors que
$\mathcal F^q\in K'$, et l'hypothèse additionnelle sur $K'$ implique que
$\mathcal F\in K'$.
\end{proof}

\subsection{Le théorème de finitude : cas des schémas usuels}
\label{subsection:III.3.2-fr}

\begin{theorem}[3.2.1]
\label{III.3.2.1-fr}
Soient $Y$ un préschéma localement noethérien, $f:X\to Y$ un morphisme
propre. Pour tout $\mathcal O_X$-Module cohérent $\mathcal F$, les
$\mathcal O_Y$-Modules $R^qf_*(\mathcal F)$ sont cohérents pour $q\geq0$.
\end{theorem}

\begin{proof}
La question étant locale sur $Y$, on peut supposer $Y$ noethérien, donc $X$
noethérien (I, 6.3.7). Les $\mathcal O_X$-Modules cohérents $\mathcal F$ pour
lesquels la conclusion du th. (\hyperref[III.3.2.1-fr]{3.2.1}) est vraie
forment un sous-ensemble exact $K'$ de la catégorie $K$ des
$\mathcal O_X$-Modules cohérents.

\oldpage[III]{117}
En effet, soit
$0\to\mathcal F'\to\mathcal F\to\mathcal F''\to0$ une suite exacte de
$\mathcal O_X$-Modules cohérents~; supposons par exemple que $\mathcal F'$
et $\mathcal F''$ appartiennent à $K'$~; on a la suite exacte de cohomologie
\[
R^{q-1}f_*(\mathcal F'')\xrightarrow{\partial}R^qf_*(\mathcal F')
\longrightarrow R^qf_*(\mathcal F)\longrightarrow R^qf_*(\mathcal F'')
\xrightarrow{\partial}R^{q+1}f_*(\mathcal F')
\]
dans laquelle par hypothèse les quatre termes extrêmes sont cohérents~; il
en est donc de même du terme médian $R^qf_*(\mathcal F)$ par
(0\textsubscript{I}, 5.3.4 et 5.3.3). On montre de la même manière que
lorsque $\mathcal F$ et $\mathcal F'$ (resp. $\mathcal F$ et
$\mathcal F''$) sont dans $K'$, il en est de même de $\mathcal F''$ (resp.
$\mathcal F'$). De plus, tout \emph{facteur direct} cohérent $\mathcal F'$
d'un $\mathcal O_X$-Module $\mathcal F\in K'$ appartient aussi à $K'$~: en
effet, $R^qf_*(\mathcal F')$ est alors facteur direct de
$R^qf_*(\mathcal F)$ (G, II, 4.4.4), donc est de type fini, et comme il est
quasi-cohérent (\hyperref[III.1.4.10-fr]{1.4.10}), il est cohérent, $Y$
étant noethérien. En vertu de (\hyperref[III.3.1.3-fr]{3.1.3}), on est
ramené à prouver que lorsque $X$ est \emph{irréductible} de point générique
$x$, il existe un $\mathcal O_X$-Module cohérent $\mathcal F$ appartenant à
$K'$, tel que $\mathcal F_x\neq0$~: en effet, si ce point est établi, on
pourra l'appliquer à tout sous-préschéma fermé irréductible $Y$ de $X$, car
si $j:Y\to X$ est l'injection canonique, $f\circ j$ est propre (II, 5.4.2),
et si $\mathcal G$ est un $\mathcal O_Y$-Module cohérent de support $Y$,
$j_*(\mathcal G)$ est un $\mathcal O_X$-Module cohérent tel que
\[
R^q(f\circ j)_*(\mathcal G)=R^qf_*(j_*(\mathcal G))
\]
(G, II, 4.9.1), donc on est bien dans les conditions d'application de
(\hyperref[III.3.1.3-fr]{3.1.3}).

Or, en vertu du lemme de Chow (II, 5.6.2), il existe un préschéma
irréductible $X'$ et un morphisme \emph{projectif} et surjectif
$g:X'\to X$ tel que $f\circ g:X'\to Y$ soit \emph{projectif}. Il existe un
$\mathcal O_{X'}$-Module $\mathcal L$ ample pour $g$ (II, 5.3.1)~;
appliquons le th. fondamental des morphismes projectifs
(\hyperref[III.2.2.1-fr]{2.2.1}) à $g:X'\to X$ et à $\mathcal L$~: il
existe donc un entier $n$ tel que
$\mathcal F=g_*(\mathcal O_{X'}(n))$ soit un $\mathcal O_X$-Module cohérent
et $R^qg_*(\mathcal O_{X'}(n))=0$ pour tout $q>0$~; en outre, comme
\[
g^*(g_*(\mathcal O_{X'}(n)))\longrightarrow\mathcal O_{X'}(n)
\]
est surjectif pour $n$ assez grand (\hyperref[III.2.2.1-fr]{2.2.1}), on
voit qu'on peut supposer que, au point générique $x$ de $X$, on a
$\mathcal F_x\neq0$ (II, 3.4.7). D'autre part, comme $f\circ g$ est
projectif et $Y$ noethérien, les
$R^p(f\circ g)_*(\mathcal O_{X'}(n))$ sont \emph{cohérents}
(\hyperref[III.2.2.1-fr]{2.2.1}). Cela étant,
$R^*(f\circ g)_*(\mathcal O_{X'}(n))$ est l'aboutissement d'une suite
spectrale de Leray, dont le terme $E_2$ est donné par
\[
E_2^{pq}=R^pf_*(R^qg_*(\mathcal O_{X'}(n)))~;
\]
ce qui précède montre que cette suite spectrale est dégénérée, et on sait
alors (0, 11.1.6) que $E_2^{p0}=R^pf_*(\mathcal F)$ est isomorphe à
$R^p(f\circ g)_*(\mathcal O_{X'}(n))$, ce qui achève la démonstration.
\end{proof}

\begin{corollary}[3.2.2]
\label{III.3.2.2-fr}
Soit $Y$ un préschéma localement noethérien. Pour tout morphisme propre
$f:X\to Y$, l'image directe par $f$ de tout $\mathcal O_X$-Module cohérent
est un $\mathcal O_Y$-Module cohérent.
\end{corollary}

\begin{corollary}[3.2.3]
\label{III.3.2.3-fr}
Soient $A$ un anneau noethérien, $X$ un schéma propre sur $A$~; pour tout
$\mathcal O_X$-Module cohérent $\mathcal F$, les $H^p(X,\mathcal F)$ sont
des $A$-modules de type fini, et il existe un entier $r>0$ tel que pour tout
$\mathcal O_X$-Module cohérent $\mathcal F$ et tout $p>r$,
$H^p(X,\mathcal F)=0$.
\end{corollary}

\begin{proof}
La seconde assertion a déjà été démontrée
(\hyperref[III.1.4.12-fr]{1.4.12})~; la première résulte du th. de finitude
(\hyperref[III.3.2.1-fr]{3.2.1}), compte tenu de
(\hyperref[III.1.4.11-fr]{1.4.11}).
\end{proof}

En particulier, si $X$ est un \emph{schéma algébrique propre} sur un corps
$k$, alors, pour tout $\mathcal O_X$-Module cohérent $\mathcal F$, les
$H^p(X,\mathcal F)$ sont des $k$-espaces vectoriels de \emph{dimension
finie}.

\begin{corollary}[3.2.4]
\label{III.3.2.4-fr}
Soient $Y$ un préschéma localement noethérien, $f:X\to Y$ un morphisme de
type fini. Pour tout $\mathcal O_X$-Module cohérent $\mathcal F$ dont le
support est propre sur $Y$ (II, 5.4.10), les $\mathcal O_Y$-Modules
$R^qf_*(\mathcal F)$ sont cohérents.
\end{corollary}

\oldpage[III]{118}
\begin{proof}
La question étant locale sur $Y$, on peut supposer $Y$ noethérien, et il en
est donc de même de $X$ (I, 6.3.7). Par hypothèse, tout sous-préschéma fermé
$Z$ de $X$ dont l'espace sous-jacent est $\operatorname{Supp}(\mathcal F)$
est propre sur $Y$, autrement dit, si $j:Z\to X$ est l'injection canonique,
$f\circ j:Z\to Y$ est propre. Or, on peut supposer $Z$ tel que
$\mathcal F=j_*(\mathcal G)$, où $\mathcal G=j^*(\mathcal F)$ est un
$\mathcal O_Z$-Module cohérent (I, 9.3.5)~; comme on a
$R^qf_*(\mathcal F)=R^q(f\circ j)_*(\mathcal G)$ par
(\hyperref[III.1.3.4-fr]{1.3.4}), la conclusion résulte aussitôt de
(\hyperref[III.3.2.1-fr]{3.2.1}).
\end{proof}

\subsection{Généralisation du théorème de finitude (schémas usuels)}
\label{subsection:III.3.3-fr}

\begin{proposition}[3.3.1]
\label{III.3.3.1-fr}
Soient $Y$ un préschéma noethérien, $\mathcal S$ une
$\mathcal O_Y$-Algèbre graduée à degrés positifs, quasi-cohérente et de type
fini, $Y'=\operatorname{Proj}(\mathcal S)$ et $g:Y'\to Y$ le morphisme
structural. Soient $f:X\to Y$ un morphisme propre,
$\mathcal S'=f^*(\mathcal S)$,
$\mathcal M=\bigoplus_{k\in\mathbb Z}\mathcal M_k$ un $\mathcal S'$-Module
gradué quasi-cohérent et de type fini. Alors les
\[
R^pf_*(\mathcal M)=\bigoplus_{k\in\mathbb Z}R^pf_*(\mathcal M_k)
\]
sont des $\mathcal S$-Modules gradués de type fini pour tout $p$. Supposons
en outre que $\mathcal S$ soit engendrée par $\mathcal S_1$~; alors, pour
chaque $p\in\mathbb Z$, il existe un entier $k_p$ tel que pour tout
$k\geq k_p$ et tout $r\geq0$, on ait
\[
R^pf_*(\mathcal M_{k+r})=\mathcal S_rR^pf_*(\mathcal M_k).
\tag{3.3.1.1}\label{III.3.3.1.1-fr}
\]
\end{proposition}

\begin{proof}
La première assertion est identique à l'énoncé de
(\hyperref[III.2.4.1-fr]{2.4.1}, (i)), où on a simplement remplacé
«~morphisme projectif~» par «~morphisme propre~». Or, dans la démonstration
de (\hyperref[III.2.4.1-fr]{2.4.1}, (i)), l'hypothèse sur $f$ a été utilisée
uniquement pour montrer (avec les notations de cette démonstration) que
$R^pf'_*(\widetilde{\mathcal M})$ est un $\mathcal O_{Y'}$-Module cohérent.
Avec les hypothèses de (\hyperref[III.3.3.1-fr]{3.3.1}), $f'$ est propre
(II, 5.4.2, (iii)), donc on peut reprendre sans changement toute la
démonstration de (\hyperref[III.2.4.1-fr]{2.4.1}, (i)), grâce au théorème de
finitude (\hyperref[III.3.2.1-fr]{3.2.1}).

Quant à la seconde assertion, il suffit de remarquer qu'il y a un
recouvrement ouvert affine fini $(U_i)$ de $Y$ tel que les restrictions à
$U_i$ des deux membres de
(\hyperref[III.3.3.1.1-fr]{3.3.1.1}) soient égales pour tout
$k\geq k_{p,i}$ (II, 2.1.6, (ii))~; il suffit de prendre pour $k_p$ le plus
grand des $k_{p,i}$.
\end{proof}

\begin{corollary}[3.3.2]
\label{III.3.3.2-fr}
Soient $A$ un anneau noethérien, $\mathfrak m$ un idéal de $A$, $X$ un
$A$-schéma propre, $\mathcal F$ un $\mathcal O_X$-Module cohérent. Alors,
pour tout $p\geq0$, la somme directe
$\bigoplus_{k\geq0}H^p(X,\mathfrak m^k\mathcal F)$ est un module de type fini
sur l'anneau $S=\bigoplus_{k\geq0}\mathfrak m^k$~; en particulier, il existe
un entier $k_p\geq0$ tel que pour tout $k\geq k_p$, et tout $r\geq0$, on ait
\[
H^p(X,\mathfrak m^{k+r}\mathcal F)
=\mathfrak m^rH^p(X,\mathfrak m^k\mathcal F).
\tag{3.3.3.1}\label{III.3.3.2.1-fr}
\]
\end{corollary}

\begin{proof}
Il suffit d'appliquer (\hyperref[III.3.3.2-fr]{3.3.2}) avec
$Y=\operatorname{Spec}(A)$, $\mathcal S=\widetilde S$,
$\mathcal M_k=\mathfrak m^k\mathcal F$, compte tenu de
(\hyperref[III.1.4.11-fr]{1.4.11}).
\end{proof}

Il convient de rappeler que la structure de $S$-module de
$\bigoplus_{k\geq0}H^p(X,\mathfrak m^k\mathcal F)$ s'obtient en considérant,
pour tout $a\in\mathfrak m^r$, l'application
\[
H^p(X,\mathfrak m^k\mathcal F)
\longrightarrow H^p(X,\mathfrak m^{k+r}\mathcal F)
\]
qui provient par passage à la cohomologie de la multiplication
$\mathfrak m^k\mathcal F\to\mathfrak m^{k+r}\mathcal F$ définie par $a$
(\hyperref[III.2.4.1-fr]{2.4.1}).

\oldpage[III]{119}
\subsection{Le théorème de finitude : cas des schémas formels}
\label{subsection:III.3.4-fr}

Les résultats de cette section (sauf la définition
(\hyperref[III.3.4.1-fr]{3.4.1})) ne seront pas utilisés dans la suite de ce
chapitre.

\begin{env}[3.4.1]
\label{III.3.4.1-fr}
Soient $\mathfrak X$ et $\mathfrak S$ deux préschémas formels localement
noethériens (I, 10.4.2), $f:\mathfrak X\to\mathfrak S$ un morphisme de
préschémas formels. Nous dirons que $f$ est un morphisme \emph{propre} s'il
vérifie les conditions suivantes~:
\begin{enumerate}
\item[$1^\circ$] $f$ est un morphisme de type fini (I, 10.13.3).
\item[$2^\circ$] Si $\mathcal K$ est un Idéal de définition de $\mathfrak S$
et si l'on pose
$\mathcal J=f^*(\mathcal K)\mathcal O_{\mathfrak X}$,
$X_0=(\mathfrak X,\mathcal O_{\mathfrak X}/\mathcal J)$,
$S_0=(\mathfrak S,\mathcal O_{\mathfrak S}/\mathcal K)$, le morphisme
$f_0:X_0\to Y_0$ déduit de $f$ (I, 10.5.6) est propre.
\end{enumerate}

Il est immédiat que cette définition ne dépend pas de l'Idéal de définition
$\mathcal K$ de $\mathfrak S$ considéré~; en effet, si $\mathcal K'$ est un
second Idéal de définition tel que $\mathcal K'\subset\mathcal K$, et si on
pose $\mathcal J'=f^*(\mathcal K')\mathcal O_{\mathfrak X}$,
$X'_0=(\mathfrak X,\mathcal O_{\mathfrak X}/\mathcal J')$,
$S'_0=(\mathfrak S,\mathcal O_{\mathfrak S}/\mathcal K')$, le morphisme
$f'_0:X'_0\to S'_0$ déduit de $f$ est tel que le diagramme
\[
\xymatrix{
X_0\ar[r]^{f_0}\ar[d]_i & S_0\ar[d]^j\\
X'_0\ar[r]_{f'_0} & S'_0
}
\]
est commutatif, $i$ et $j$ étant des immersions surjectives~; il revient donc
au même de dire que $f_0$ ou $f'_0$ est propre, en vertu de (II, 5.4.5).

On notera que, pour tout $n\geq0$, si on pose
$X_n=(\mathfrak X,\mathcal O_{\mathfrak X}/\mathcal J^{n+1})$,
$S_n=(\mathfrak S,\mathcal O_{\mathfrak S}/\mathcal K^{n+1})$, le morphisme
$f_n:X_n\to Y_n$ déduit de $f$ (I, 10.5.6) est propre pour tout $n$ dès qu'il
l'est pour $n=0$ (II, 5.4.6).

Si $g:Y\to Z$ est un morphisme propre de préschémas usuels, localement
noethériens, $Z'$ une partie fermée de $Z$, $Y'$ une partie fermée de $Y$
telle que $g(Y')\subset Z'$, le prolongement
$\widehat g:Y_{/Y'}\to Z_{/Z'}$ de $g$ aux complétés (I, 10.9.1) est un
morphisme propre de préschémas formels, comme il résulte de la définition et
de (II, 5.4.5).

Soient $\mathfrak X$ et $\mathfrak S$ deux préschémas formels localement
noethériens, $f:\mathfrak X\to\mathfrak S$ un morphisme \emph{de type fini}
(I, 10.13.3)~; les notations étant les mêmes que ci-dessus, on dit qu'une
partie $Z$ de l'espace sous-jacent à $\mathfrak X$ est \emph{propre sur}
$\mathfrak S$ (ou propre pour $f$) si, considérée comme partie de $X_0$, $Z$
est \emph{propre sur $S_0$} (II, 5.4.10). Toutes les propriétés des parties
propres de préschémas usuels énoncées dans (II, 5.4.10) sont encore valables
pour les parties propres de préschémas formels comme il résulte aussitôt des
définitions.
\end{env}

\begin{theorem}[3.4.2]
\label{III.3.4.2-fr}
Soient $\mathfrak X$, $\mathfrak Y$ deux préschémas formels localement
noethériens, $f:\mathfrak X\to\mathfrak Y$ un morphisme propre. Pour tout
$\mathcal O_{\mathfrak X}$-Module cohérent $\mathcal F$, les
$\mathcal O_{\mathfrak Y}$-Modules $R^qf_*(\mathcal F)$ sont cohérents pour
tout $q\geq0$.
\end{theorem}

Soient $\mathcal J$ un Idéal de définition de $\mathfrak Y$,
$\mathcal K=f^*(\mathcal J)\mathcal O_{\mathfrak X}$, et considérons les
$\mathcal O_{\mathfrak X}$-Modules
\[
\mathcal F_k=\mathcal F\otimes_{\mathcal O_{\mathfrak Y}}
(\mathcal O_{\mathfrak Y}/\mathcal J^{k+1})
=\mathcal F/\mathcal K^{k+1}\mathcal F\qquad(k\geq0)
\tag{3.4.2.1}\label{III.3.4.2.1-fr}
\]
qui forment évidemment un \emph{système projectif} de
$\mathcal O_{\mathfrak X}$-Modules topologiques, tel que

\oldpage[III]{120}
\[
\mathcal F=\varprojlim_k\mathcal F_k
\]
par (I, 10.11.3). D'autre part, il résultera de
(\hyperref[III.3.4.2-fr]{3.4.2}) que chacun des $R^qf_*(\mathcal F)$, étant
cohérent, est naturellement muni d'une structure de
$\mathcal O_{\mathfrak Y}$-Module topologique (I, 10.11.6), et il en est de
même des $R^qf_*(\mathcal F_k)$. Aux homomorphismes canoniques
$\mathcal F\to\mathcal F_k=\mathcal F/\mathcal K^{k+1}\mathcal F$
correspondent canoniquement des homomorphismes
\[
R^qf_*(\mathcal F)\longrightarrow R^qf_*(\mathcal F_k)
\]
qui sont nécessairement continus pour les structures de
$\mathcal O_{\mathfrak Y}$-Module topologique précédentes (I, 10.11.6), et
forment un système projectif, donnant à la limite un homomorphisme canonique
fonctoriel
\[
R^qf_*(\mathcal F)\longrightarrow\varprojlim_kR^qf_*(\mathcal F_k)
\tag{3.4.2.2}\label{III.3.4.2.2-fr}
\]
qui sera un homomorphisme continu de $\mathcal O_{\mathfrak Y}$-Modules
topologiques. Nous allons démontrer en même temps que
(\hyperref[III.3.4.2-fr]{3.4.2}) le

\begin{corollary}[3.4.3]
\label{III.3.4.3-fr}
Chacun des homomorphismes (\hyperref[III.3.4.2.2-fr]{3.4.2.2}) est un
isomorphisme topologique. En outre, si $\mathfrak Y$ est noethérien, le
système projectif
$(R^qf_*(\mathcal F/\mathcal K^{k+1}\mathcal F))_{k\geq0}$ satisfait à la
condition (ML) (0, 13.1.1).
\end{corollary}

Nous commencerons par établir (\hyperref[III.3.4.2-fr]{3.4.2}) et
(\hyperref[III.3.4.3-fr]{3.4.3}) lorsque $Y$ est un schéma affine
formel noethérien (I, 10.4.1)~:

\begin{corollary}[3.4.4]
\label{III.3.4.4-fr}
Sous les hypothèses de (\hyperref[III.3.4.2-fr]{3.4.2}), supposons en outre
que $\mathfrak Y=\operatorname{Spf}(A)$, où $A$ est un anneau adique
noethérien. Soit $\mathfrak J$ un idéal de définition de $A$, et posons
$\mathcal F_k=\mathcal F/\mathfrak J^{k+1}\mathcal F$ pour $k\geq0$. Alors
les $H^n(\mathfrak X,\mathcal F)$ sont des $A$-modules de type fini~; le
système projectif $(H^n(\mathfrak X,\mathcal F_k))_{k\geq0}$ satisfait à la
condition (ML) pour tout $n$~; si on pose
\[
N_{n,k}=\operatorname{Ker}\bigl(H^n(\mathfrak X,\mathcal F)
\longrightarrow H^n(\mathfrak X,\mathcal F_k)\bigr)
\tag{3.4.4.1}\label{III.3.4.4.1-fr}
\]
(égal aussi à
$\operatorname{Im}(H^n(\mathfrak X,\mathfrak J^{k+1}\mathcal F)
\to H^n(\mathfrak X,\mathcal F))$ par la suite exacte de cohomologie), les
$N_{k,n}$ définissent sur $H^n(\mathfrak X,\mathcal F)$ une filtration
$\mathfrak J$-bonne (0, 13.7.7)~; enfin, l'homomorphisme canonique
\[
H^n(\mathfrak X,\mathcal F)
\longrightarrow\varprojlim_kH^n(\mathfrak X,\mathcal F_k)
\tag{3.4.4.2}\label{III.3.4.4.2-fr}
\]
est un isomorphisme topologique pour tout $n$ (le premier membre étant muni
de la topologie $\mathfrak J$-adique, les $H^n(\mathfrak X,\mathcal F_k)$ de
la topologie discrète).
\end{corollary}

Posons
\[
S=\operatorname{gr}(A)=\bigoplus_{k\geq0}
\mathfrak J^k/\mathfrak J^{k+1},\qquad
\mathcal M=\operatorname{gr}(\mathcal F)=\bigoplus_{k\geq0}
\mathfrak J^k\mathcal F/\mathfrak J^{k+1}\mathcal F.
\tag{3.4.4.3}\label{III.3.4.4.3-fr}
\]
On sait que $\mathfrak J^\Delta$ est un Idéal de définition de $\mathfrak Y$
(I, 10.3.1)~; soient
$\mathcal K=f^*(\mathfrak J^\Delta)\mathcal O_{\mathfrak X}$,
$X_0=(\mathfrak X,\mathcal O_{\mathfrak X}/\mathcal K)$,
$Y_0=(\mathfrak Y,\mathcal O_{\mathfrak Y}/\mathfrak J^\Delta)
=\operatorname{Spec}(A_0)$ avec $A_0=A/\mathfrak J$. Il est clair que les
$\mathcal M_k=\mathfrak J^k\mathcal F/\mathfrak J^{k+1}\mathcal F$ sont des
$\mathcal O_{X_0}$-Modules cohérents (I, 10.11.3). Considérons d'autre part la
$\mathcal O_{X_0}$-Algèbre graduée quasi-cohérente
\[
\mathcal S=\mathcal O_{X_0}\otimes_{A_0}S
=\operatorname{gr}(\mathcal O_{\mathfrak X})
=\bigoplus_{k\geq0}\mathcal K^k/\mathcal K^{k+1}.
\tag{3.4.4.4}\label{III.3.4.4.4-fr}
\]

L'hypothèse que $\mathcal F$ est un $\mathcal O_{\mathfrak X}$-Module de type
fini entraîne d'abord que $\mathcal M$ est

\oldpage[III]{121}
un $\mathcal S$-Module gradué \emph{de type fini}. En effet, la question est
locale sur $\mathfrak X$, et on peut donc supposer pour la traiter que
$\mathfrak X=\operatorname{Spf}(B)$, où $B$ est un anneau adique noethérien,
et $\mathcal F=N^\Delta$, où $N$ est un $B$-module de type fini
(I, 10.10.5)~; on a en outre $X_0=\operatorname{Spec}(B_0)$ où
$B_0=B/\mathfrak J B$, et les $\mathcal O_{X_0}$-Modules quasi-cohérents
$\mathcal S$ et $\mathcal M$ sont respectivement égaux à $\widetilde{S'}$ et
$\widetilde{M'}$, où
\[
S'=\bigoplus_{k\geq0}
((\mathfrak J^k/\mathfrak J^{k+1})\otimes_{A_0}B_0),\qquad
M'=\bigoplus_{k\geq0}
((\mathfrak J^k/\mathfrak J^{k+1})\otimes_{A_0}N_0),
\]
avec $N_0=N/\mathfrak JN$~; on a donc évidemment
$M'=S'\otimes_{B_0}N_0$, et comme $N_0$ est un $B_0$-module de type fini,
$M'$ est un $S'$-module de type fini, d'où notre assertion (I, 1.3.13).

Comme le morphisme $f_0:X_0\to Y_0$ est \emph{propre} par hypothèse, on peut
appliquer (\hyperref[III.3.3.2-fr]{3.3.2}) à $\mathcal S$, $\mathcal M$ et au
morphisme $f_0$~; compte tenu de (\hyperref[III.1.4.11-fr]{1.4.11}), on en
conclut que pour \emph{tout $n\geq0$},
$\bigoplus_{k\geq0}H^n(X_0,\mathcal M_k)$ est un $S$-module gradué \emph{de
type fini}. Cela prouve que la condition $(\mathrm F_n)$ de (0, 13.7.7) est
vérifiée pour \emph{tout $n\geq0$}, lorsqu'on considère le système projectif
strict $(\mathcal F/\mathfrak J^k\mathcal F)_{k\geq0}$ de faisceaux de
groupes abéliens sur $X_0$, munis chacun de sa structure naturelle de
«~$A$-module filtré~». On peut donc appliquer (0, 13.7.7), qui prouve que~:
\begin{enumerate}
\item[$1^\circ$] Le système projectif
$(H^n(\mathfrak X,\mathcal F_k))_{k\geq0}$ vérifie la condition (ML).
\item[$2^\circ$] Si
$H'^{,n}=\varprojlim_kH^n(\mathfrak X,\mathcal F_k)$, $H'^{,n}$ est un
$A$-module de type fini.
\item[$3^\circ$] La filtration définie sur $H'^{,n}$ par les noyaux des
homomorphismes canoniques $H'^{,n}\to H^n(\mathfrak X,\mathcal F_k)$ est
$\mathfrak J$-bonne.
\end{enumerate}

Notons d'autre part que si on pose
$X_k=(\mathfrak X,\mathcal O_{\mathfrak X}/\mathcal K^{k+1})$,
$\mathcal F_k$ est un $\mathcal O_{X_k}$-Module cohérent (I, 10.11.3) et si
$U$ est un ouvert affine dans $X_0$, $U$ est aussi un ouvert affine dans
chacun des $X_k$ (I, 5.1.9), donc $H^n(U,\mathcal F_k)=0$ pour tout $n>0$ et
tout $k$ (\hyperref[III.1.3.1-fr]{1.3.1}) et
$H^0(U,\mathcal F_k)\to H^0(U,\mathcal F_h)$ est surjectif pour $h\leq k$
(I, 1.3.9). On est donc dans les conditions de (0, 13.3.2) et l'application
de (0, 13.3.1) prouve que $H'^{,n}$ s'identifie canoniquement à
$H^n(\mathfrak X,\varprojlim_k\mathcal F_k)=H^n(\mathfrak X,\mathcal F)$~;
cela achève de prouver (\hyperref[III.3.4.4-fr]{3.4.4}).

\begin{env}[3.4.5]
\label{III.3.4.5-fr}
Revenons maintenant à la démonstration de (\hyperref[III.3.4.2-fr]{3.4.2})
et (\hyperref[III.3.4.3-fr]{3.4.3}). Prouvons d'abord ces propositions dans
le cas $\mathfrak Y=\operatorname{Spf}(A)$ envisagé dans
(\hyperref[III.3.4.4-fr]{3.4.4})~; pour cela, pour tout $g\in A$, appliquons
(\hyperref[III.3.4.4-fr]{3.4.4}) au schéma formel affine noethérien induit
sur l'ouvert $\mathfrak Y_g=\mathfrak D(g)$ de $\mathfrak Y$, qui est égal à
$\operatorname{Spf}(A_{\{g\}})$, et au préschéma formel induit par
$\mathfrak X$ sur $f^{-1}(\mathfrak Y_g)$~; notons que $\mathfrak Y_g$ est
aussi un ouvert affine dans le préschéma
$Y_k=(\mathfrak Y,\mathcal O_{\mathfrak Y}/(\mathfrak J^\Delta)^{k+1})$,
et comme $\mathcal F_k$ est un $\mathcal O_{X_k}$-Module cohérent, on a
\[
H^n(f^{-1}(\mathfrak Y_g),\mathcal F_k)
=\Gamma(\mathfrak Y_g,R^nf_*(\mathcal F_k))
\]
pour tout $k\geq0$ en vertu de
(\hyperref[III.1.4.11-fr]{1.4.11}). L'homomorphisme canonique
\[
H^n(f^{-1}(\mathfrak Y_g),\mathcal F)
\longrightarrow\varprojlim_k\Gamma(\mathfrak Y_g,R^nf_*(\mathcal F_k))
\]
est un isomorphisme~; mais on a (0\textsubscript{I}, 3.2.6)
\[
\varprojlim_k\Gamma(\mathfrak Y_g,R^nf_*(\mathcal F_k))
=\Gamma(\mathfrak Y_g,\varprojlim_kR^nf_*(\mathcal F_k))
\]

\oldpage[III]{122}
et comme le faisceau $R^nf_*(\mathcal F)$ est associé au préfaisceau
$\mathfrak Y_g\mapsto H^n(f^{-1}(\mathfrak Y_g),\mathcal F)$ sur les
$\mathfrak Y_g$ (0\textsubscript{I}, 3.2.1), on a bien montré que
l'homomorphisme (\hyperref[III.3.4.2.2-fr]{3.4.2.2}) est bijectif. Prouvons
ensuite que $R^nf_*(\mathcal F)$ est un $\mathcal O_{\mathfrak Y}$-Module
cohérent, et de façon plus précise que l'on a
\[
R^nf_*(\mathcal F)=(H^n(\mathfrak X,\mathcal F))^\Delta.
\tag{3.4.5.1}\label{III.3.4.5.1-fr}
\]
Avec les notations précédentes, on a, puisque $\mathcal F_k$ est un
$\mathcal O_{X_k}$-Module cohérent (\hyperref[III.1.4.13-fr]{1.4.13})
\[
\Gamma(\mathfrak Y_g,R^nf_*(\mathcal F_k))
=(\Gamma(\mathfrak Y,R^nf_*(\mathcal F_k)))_g
=(H^n(\mathfrak X,\mathcal F_k))_g.
\]
Or, les $H^n(\mathfrak X,\mathcal F_k)$ forment un système projectif
vérifiant (ML) et leur limite projective $H^n(\mathfrak X,\mathcal F)$ est
un $A$-module de type fini. On en conclut (0, 13.7.8) que l'on a
\[
\varprojlim_k(H^n(\mathfrak X,\mathcal F_k))_g
=H^n(\mathfrak X,\mathcal F)\otimes_AA_{\{g\}}
=\Gamma(\mathfrak Y_g,(H^n(\mathfrak X,\mathcal F))^\Delta)
\]
compte tenu de (I, 10.10.8) appliqué à $A$ et $A_{\{g\}}$~; ceci démontre
(\hyperref[III.3.4.5.1-fr]{3.4.5.1}) puisque
\[
\Gamma(\mathfrak Y_g,R^nf_*(\mathcal F))
=\varprojlim_k\Gamma(\mathfrak Y_g,R^nf_*(\mathcal F_k)).
\]
Comme (\hyperref[III.3.4.2.2-fr]{3.4.2.2}) est alors un isomorphisme de
$\mathcal O_{\mathfrak Y}$-Modules cohérents, c'est nécessairement un
isomorphisme \emph{topologique} (I, 10.11.6). Enfin, il résulte des relations
$R^nf_*(\mathcal F_k)=(H^n(\mathfrak X,\mathcal F_k))^\Delta$ que le système
projectif $(R^nf_*(\mathcal F_k))_{k\geq0}$ vérifie (ML) (I, 10.10.2).

Une fois (\hyperref[III.3.4.2-fr]{3.4.2}) et
(\hyperref[III.3.4.3-fr]{3.4.3}) démontrés dans le cas où le préschéma formel
$\mathfrak Y$ est affine noethérien, il est immédiat de passer de là au cas
général pour (\hyperref[III.3.4.2-fr]{3.4.2}) et la première assertion de
(\hyperref[III.3.4.3-fr]{3.4.3}), qui sont locales sur $\mathfrak Y$. Quant à
la seconde assertion de (\hyperref[III.3.4.3-fr]{3.4.3}), il suffit,
$\mathfrak Y$ étant noethérien, de le recouvrir par un nombre fini d'ouverts
affines noethériens $U_i$ et de remarquer que les restrictions du système
projectif $(R^qf_*(\mathcal F_k))$ à chacun des $U_i$ vérifient (ML).
\end{env}

Nous avons en outre démontré en cours de route~:

\begin{corollary}[3.4.6]
\label{III.3.4.6-fr}
Sous les hypothèses de (\hyperref[III.3.4.4-fr]{3.4.4}), l'homomorphisme
canonique
\[
H^q(\mathfrak X,\mathcal F)
\longrightarrow\Gamma(\mathfrak Y,R^qf_*(\mathcal F))
\tag{3.4.6.1}\label{III.3.4.6.1-fr}
\]
est bijectif.
\end{corollary}

\section{Le théorème fondamental des morphismes propres. Applications}
\label{section:III.4-fr}

\subsection{Le théorème fondamental}
\label{subsection:III.4.1-fr}

\begin{env}[4.1.1]
\label{III.4.1.1-fr}
Soient $X$, $Y$ deux préschémas usuels noethériens, $f:X\to Y$ un
morphisme propre, $Y'$ une partie fermée de $Y$, $X'$ son image réciproque
$f^{-1}(Y')$. Nous désignerons par $\widehat X$ et $\widehat Y$ les
préschémas formels $X_{/X'}$ et $Y_{/Y'}$, complétés de $X$ et $Y$ le long
de $X'$ et $Y'$ respectivement (I, 10.8.5), par $\widehat f$ le prolongement
de $f$ à ces complétés (I, 10.9.1) qui est un morphisme
$\widehat X\to\widehat Y$ de préschémas formels. Pour tout
$\mathcal O_X$-Module cohérent $\mathcal F$, nous

\oldpage[III]{123}

désignerons par $\widehat{\mathcal F}$ son complété $\mathcal F_{/X'}$ le
long de $X'$ (I, 10.8.4) qui est un $\mathcal O_{\widehat X}$-Module
cohérent (I, 10.8.8).
\end{env}

\begin{env}[4.1.2]
\label{III.4.1.2-fr}
Soit $\mathcal J$ un Idéal cohérent de $\mathcal O_Y$ tel que
$\operatorname{Supp}(\mathcal O_Y/\mathcal J)=Y'$ (I, 5.2.1)~; on sait
(I, 4.4.5) que
\[
\mathcal K=f^*(\mathcal J)\mathcal O_X
\]
est un Idéal cohérent de $\mathcal O_X$ tel que
\[
\operatorname{Supp}(\mathcal O_X/\mathcal K)=X'.
\]
Nous considérerons pour tout $k\geq0$ les $\mathcal O_X$-Modules cohérents
\[
\mathcal F_k
=\mathcal F\otimes_{\mathcal O_Y}(\mathcal O_Y/\mathcal J^{k+1})
=\mathcal F/\mathcal K^{k+1}\mathcal F.
\]
Les $\mathcal O_Y$-Modules $R^nf_*(\mathcal F)$ et
$R^nf_*(\mathcal F_k)$ sont cohérents pour tout $n$
(\hyperref[III.3.2.1-fr]{3.2.1}). Pour tout $k\geq0$ et tout $n$,
l'homomorphisme canonique $\mathcal F\to\mathcal F_k$ définit par
fonctorialité un homomorphisme
\[
R^nf_*(\mathcal F)\longrightarrow R^nf_*(\mathcal F_k).
\tag{4.1.2.1}\label{III.4.1.2.1-fr}
\]
En outre, comme $\mathcal F_k$ est un
$\mathcal O_X/\mathcal K^{k+1}$-Module, $R^nf_*(\mathcal F_k)$ est un
$\mathcal O_Y/\mathcal J^{k+1}$-Module (0, 12.2.1) et on déduit donc de
(\hyperref[III.4.1.2.1-fr]{4.1.2.1}) un homomorphisme
\[
R^nf_*(\mathcal F)\otimes_{\mathcal O_Y}
(\mathcal O_Y/\mathcal J^{k+1})
\longrightarrow R^nf_*(\mathcal F_k).
\tag{4.1.2.2}\label{III.4.1.2.2-fr}
\]
Les deux membres de (\hyperref[III.4.1.2.2-fr]{4.1.2.2}) forment deux
systèmes projectifs, et la limite projective du premier membre n'est autre
que le complété $(R^nf_*(\mathcal F))_{/Y'}$, que nous noterons
$\widehat{R^nf_*(\mathcal F)}$. En outre, il est immédiat que les
homomorphismes (\hyperref[III.4.1.2.2-fr]{4.1.2.2}) forment un système
projectif, d'où par passage à la limite un homomorphisme canonique
\[
\varphi_n:\widehat{R^nf_*(\mathcal F)}
\longrightarrow\varprojlim_kR^nf_*(\mathcal F_k).
\tag{4.1.2.3}\label{III.4.1.2.3-fr}
\]
D'ailleurs (\hyperref[III.4.1.2.2-fr]{4.1.2.2}) est un homomorphisme de
$(\mathcal O_Y/\mathcal J^{k+1})$-Modules, et par suite (I, 10.8.3) peut
être considéré comme un homomorphisme continu de
$\mathcal O_{\widehat Y}$-Modules topologiques pseudo-discrets
(0\textsubscript{I}, 3.8.1). L'homomorphisme $\varphi_n$ est par suite un
homomorphisme continu de $\mathcal O_{\widehat Y}$-Modules topologiques.
\end{env}

\begin{env}[4.1.3]
\label{III.4.1.3-fr}
Soit $i:\widehat X\to X$ le morphisme canonique d'espaces annelés défini
dans (I, 10.8.7), de sorte que l'on a le diagramme commutatif
\[
\xymatrix{
X_k\ar[r]^{h_k}\ar[dr]_{i_k}&\widehat X\ar[d]^i\\
&X
}
\tag{4.1.3.1}\label{III.4.1.3.1-fr}
\]
où $X_k$ est le sous-préschéma fermé de $X$ défini par l'Idéal
$\mathcal K^{k+1}$, $i_k$ l'injection canonique, $h_k$ le morphisme
d'espaces annelés correspondant à l'identité dans les espaces sous-jacents
et à l'homomorphisme canonique
$\mathcal O_{\widehat X}\to\mathcal O_X/\mathcal K^{k+1}$ (I, 10.5.2).
En outre, on a $\widehat{\mathcal F}=i^*(\mathcal F)$ (I, 10.8.8) à un
isomorphisme canonique près. On sait que l'on a
\[
H^n(X_k,i_k^*(\mathcal F_k))=H^n(X,\mathcal F_k)
\tag{4.1.3.2}\label{III.4.1.3.2-fr}
\]

\oldpage[III]{124}

à un isomorphisme canonique près, puisque
$\mathcal F_k=(i_k)_*(i_k^*(\mathcal F_k))$ (G, II, 4.9.1)~;
l'homomorphisme canonique
$H^n(\widehat X,\widehat{\mathcal F})\to
H^n(X_k,h_k^*(\widehat{\mathcal F}))$ (0, 12.1.3.5) s'écrit donc aussi
\[
H^n(\widehat X,\widehat{\mathcal F})\longrightarrow H^n(X,\mathcal F_k)
\tag{4.1.3.3}\label{III.4.1.3.3-fr}
\]
et ces homomorphismes forment évidemment un système projectif, d'où par
passage à la limite, un homomorphisme canonique
\[
\psi_{n,X}:H^n(\widehat X,\widehat{\mathcal F})
\longrightarrow\varprojlim_kH^n(X,\mathcal F_k).
\tag{4.1.3.4}\label{III.4.1.3.4-fr}
\]
Remplaçant $X$ par un ouvert de la forme $f^{-1}(V)$, où $V$ est un
ouvert affine de $Y$, on a, compte tenu de
(\hyperref[III.1.4.11-fr]{1.4.11}), des homomorphismes
\[
\psi_{n,V}:H^n(\widehat X\cap f^{-1}(V),\widehat{\mathcal F})
\longrightarrow\varprojlim_k\Gamma(V,R^nf_*(\mathcal F_k))~;
\tag{4.1.3.5}\label{III.4.1.3.5-fr}
\]
ces homomorphismes commutent de façon évidente à la restriction de $V$ à
un ouvert affine plus petit, donc définissent finalement un homomorphisme
canonique de faisceaux
\[
\psi_n:R^n\widehat f_*(\widehat{\mathcal F})
\longrightarrow\varprojlim_kR^nf_*(\mathcal F_k).
\tag{4.1.3.6}\label{III.4.1.3.6-fr}
\]
\end{env}

\begin{env}[4.1.4]
\label{III.4.1.4-fr}
Soit enfin $j:\widehat Y\to Y$ le morphisme canonique d'espaces annelés
(I, 10.8.7)~; comme $R^nf_*(\mathcal F)$ est un $\mathcal O_Y$-Module
cohérent (\hyperref[III.3.2.1-fr]{3.2.1}), on a
$j^*(R^nf_*(\mathcal F))=\widehat{R^nf_*(\mathcal F)}$ à un isomorphisme
canonique près (I, 10.8.8) et on a donc un homomorphisme canonique
\[
\rho_n:\widehat{R^nf_*(\mathcal F)}=j^*(R^nf_*(\mathcal F))
\longrightarrow R^n\widehat f_*(i^*(\mathcal F))
=R^n\widehat f_*(\widehat{\mathcal F})
\tag{4.1.4.1}\label{III.4.1.4.1-fr}
\]
défini de façon générale pour les espaces annelés (voir la démonstration
de (\hyperref[III.1.4.15-fr]{1.4.15})). Montrons que le diagramme
\[
\xymatrix{
\widehat{R^nf_*(\mathcal F)}\ar[rr]^{\rho_n}\ar[dr]_{\varphi_n}
&&R^n\widehat f_*(\widehat{\mathcal F})\ar[dl]^{\psi_n}\\
&\displaystyle\varprojlim_kR^nf_*(\mathcal F_k)&
}
\tag{4.1.4.2}\label{III.4.1.4.2-fr}
\]
est commutatif. Il suffit évidemment de prouver la commutativité du
diagramme correspondant d'homomorphismes de préfaisceaux, donc on peut se
borner au cas où $Y$ est affine, et tout revient à prouver que le diagramme
\[
\xymatrix{
\widehat{H^n(X,\mathcal F)}\ar[rr]^{\rho_n}\ar[dr]_{\varphi_n}
&&H^n(\widehat X,\widehat{\mathcal F})\ar[dl]^{\psi_{n,X}}\\
&\displaystyle\varprojlim_kH^n(X,\mathcal F_k)&
}
\tag{4.1.4.3}\label{III.4.1.4.3-fr}
\]
est commutatif. Mais la commutativité de
(\hyperref[III.4.1.3.1-fr]{4.1.3.1}) et les relations vues dans
(\hyperref[III.4.1.3-fr]{4.1.3}) entre les groupes de cohomologie donnent
aussitôt le diagramme commutatif

\oldpage[III]{125}

\[
\xymatrix{
H^n(X,\mathcal F)\ar[r]\ar[dr]&
H^n(\widehat X,\widehat{\mathcal F})
=H^n(\widehat X,i^*(\mathcal F))\ar[d]\\
&H^n(X,\mathcal F_k)=H^n(X_k,i_k^*(\mathcal F))
}
\]
d'où on déduit aussitôt la commutativité de
(\hyperref[III.4.1.4.3-fr]{4.1.4.3}).
\end{env}

\begin{theorem}[4.1.5]
\label{III.4.1.5-fr}
Soient $f:X\to Y$ un morphisme propre de préschémas noethériens, $Y'$ une
partie fermée de $Y$, $X'=f^{-1}(Y')$. Alors, pour tout
$\mathcal O_X$-Module cohérent $\mathcal F$,
$R^n\widehat f_*(\widehat{\mathcal F})$ est un
$\mathcal O_{\widehat X}$-Module cohérent et les homomorphismes
$\varphi_n$, $\psi_n$ et $\rho_n$ du diagramme
(\hyperref[III.4.1.4.2-fr]{4.1.4.2}) sont des isomorphismes topologiques.
\end{theorem}

\begin{proof}
Il suffira évidemment de prouver que $\varphi_n$ et $\psi_n$ sont des
isomorphismes~; comme $R^nf_*(\mathcal F)$ est cohérent
(\hyperref[III.3.2.1-fr]{3.2.1}), il en résultera que
$\widehat{R^nf_*(\mathcal F)}$ est cohérent (I, 10.8.8) et la bicontinuité
de $\varphi_n$, $\psi_n$ et $\rho_n$ est alors automatique (I, 10.11.6).
\end{proof}

\begin{env}[4.1.6]
\label{III.4.1.6-fr}
\emph{Remarques.}
\begin{enumerate}
\item[(i)] Si on pose
$\widehat{\mathcal F}_k=
\widehat{\mathcal F}/\widehat{\mathcal K}^{k+1}
\widehat{\mathcal F}$, il est immédiat que
$\mathcal F_k=i_*(\widehat{\mathcal F}_k)$, et l'homomorphisme canonique
(\hyperref[III.4.1.3.6-fr]{4.1.3.6}) n'est autre que l'homomorphisme déjà
défini en (\hyperref[III.3.4.2.2-fr]{3.4.2.2})
\[
R^n\widehat f_*(\widehat{\mathcal F})
\longrightarrow\varprojlim_kR^n\widehat f_*(\widehat{\mathcal F}_k)~;
\tag{4.1.6.1}\label{III.4.1.6.1-fr}
\]
par suite le fait que $\psi_n$ soit un isomorphisme est un cas particulier
de (\hyperref[III.3.4.3-fr]{3.4.3}). Mais nous en donnerons ci-dessous une
démonstration directe, évitant les considérations délicates sur les limites
projectives de suites spectrales (0, 13.7), sur lesquelles repose le
théorème général (\hyperref[III.3.4.3-fr]{3.4.3}).

\item[(ii)] Compte tenu du fait que les $\psi_n$ sont des isomorphismes, il
est équivalent de dire que les $\varphi_n$ ou les
$\rho_n=\psi_n^{-1}\circ\varphi_n$ sont des isomorphismes. Le th.
(\hyperref[III.4.1.5-fr]{4.1.5}) exprime entre autres que la formation des
$R^nf_*$ commute à la complétion et pourra s'appeler le premier théorème de
comparaison de la théorie «~algébrique~» à la théorie «~formelle~».
\end{enumerate}
Nous commencerons par établir la forme affine de
(\hyperref[III.4.1.5-fr]{4.1.5})~:
\end{env}

\begin{corollary}[4.1.7]
\label{III.4.1.7-fr}
Les hypothèses étant celles de (\hyperref[III.4.1.5-fr]{4.1.5}), supposons
en outre $Y=\operatorname{Spec}(A)$, où $A$ est noethérien, et
$\mathcal J=\widetilde{\mathfrak J}$, où $\mathfrak J$ est un idéal de
$A$, de sorte que
$\mathcal F_k=\mathcal F/\mathfrak J^{k+1}\mathcal F$.
L'homomorphisme canonique
\[
\varphi_n:\widehat{H^n(X,\mathcal F)}
\longrightarrow\varprojlim_kH^n(X,\mathcal F_k)
\tag{4.1.7.1}\label{III.4.1.7.1-fr}
\]
(où le premier membre est le séparé complété de $H^n(X,\mathcal F)$ pour
la topologie $\mathfrak J$-préadique) est un isomorphisme. Le système
projectif $(H^n(X,\mathcal F_k))_{k\geq0}$ vérifie pour tout $n$ la
condition (ML), et l'homomorphisme canonique
\[
\psi_n:H^n(\widehat X,\widehat{\mathcal F})
\longrightarrow\varprojlim_kH^n(X,\mathcal F_k)
\tag{4.1.7.2}\label{III.4.1.7.2-fr}
\]

\oldpage[III]{126}

est un isomorphisme. Enfin, la filtration sur $H^n(X,\mathcal F)$, définie
par les noyaux des homomorphismes canoniques
$H^n(X,\mathcal F)\to H^n(X,\mathcal F_k)$, est $\mathfrak J$-bonne
(0, 13.7.7) et $\varphi_n$ est un isomorphisme topologique.
\footnote{La démonstration qui suit, plus simple que la démonstration
initiale, et le complément relatif à la filtration de
$H^n(X,\mathcal F)$, nous ont été communiqués par J.-P. Serre.}
\end{corollary}

\begin{proof}
L'entier $n\geq0$ étant fixé dans cette démonstration, nous poserons pour
simplifier
\[
\mathrm H=H^n(X,\mathcal F),\qquad
\mathrm H_k=H^n(X,\mathcal F_k)
\tag{4.1.7.3}\label{III.4.1.7.3-fr}
\]
et
\[
\mathrm R_k=\operatorname{Ker}(\mathrm H\longrightarrow\mathrm H_k),
\qquad\text{sous-$A$-module de $\mathrm H$.}
\tag{4.1.7.4}\label{III.4.1.7.4-fr}
\]
La suite exacte de cohomologie
\[
H^n(X,\mathfrak J^{k+1}\mathcal F)\longrightarrow H^n(X,\mathcal F)
\longrightarrow H^n(X,\mathcal F_k)\xrightarrow{\partial}
H^{n+1}(X,\mathfrak J^{k+1}\mathcal F)
\longrightarrow H^{n+1}(X,\mathcal F)
\]
montre que l'on a aussi
\[
\mathrm R_k=\operatorname{Im}\bigl(
H^n(X,\mathfrak J^{k+1}\mathcal F)\longrightarrow H^n(X,\mathcal F)
\bigr)~;
\]
nous poserons
\[
\begin{aligned}
\mathrm Q_k
&=\operatorname{Ker}\bigl(H^{n+1}(X,\mathfrak J^{k+1}\mathcal F)
\longrightarrow H^{n+1}(X,\mathcal F)\bigr)\\
&=\operatorname{Im}\bigl(H^n(X,\mathcal F_k)
\longrightarrow H^{n+1}(X,\mathfrak J^{k+1}\mathcal F)\bigr).
\end{aligned}
\tag{4.1.7.5}\label{III.4.1.7.5-fr}
\]
On a donc la suite exacte
\[
0\longrightarrow\mathrm R_k\longrightarrow\mathrm H
\longrightarrow\mathrm H_k\longrightarrow\mathrm Q_k\longrightarrow0.
\tag{4.1.7.6}\label{III.4.1.7.6-fr}
\]

\begin{env}[4.1.7.7]
\label{III.4.1.7.7-fr}
Soit $x$ un élément de $\mathfrak J^m$ ($m\geq0$)~; la multiplication par
$x$ dans $\mathfrak J^k\mathcal F$ est un homomorphisme
$\mathfrak J^k\mathcal F\to\mathfrak J^{k+m}\mathcal F$, et donne lieu
par suite à un homomorphisme
\[
\mu_{x,m}:H^n(X,\mathfrak J^k\mathcal F)
\longrightarrow H^n(X,\mathfrak J^{k+m}\mathcal F).
\tag{4.1.7.8}\label{III.4.1.7.8-fr}
\]
Si on désigne par $S$ la $A$-algèbre graduée
$\bigoplus_{k\geq0}\mathfrak J^k$, on sait que les multiplications
$\mu_{x,m}$ définissent sur
$\mathrm E=\bigoplus_{k\geq0}H^n(X,\mathfrak J^k\mathcal F)$ une
structure de module gradué de type fini sur l'anneau gradué $S$
(\hyperref[III.3.3.2-fr]{3.3.2}), qui est noethérien (II, 2.1.5).
\end{env}

\begin{lemma}[4.1.7.9]
\label{III.4.1.7.9-fr}
Les sous-modules $\mathrm R_k$ de $\mathrm H$ définissent sur $\mathrm H$
une filtration $\mathfrak J$-bonne.
\end{lemma}

\begin{proof}
En premier lieu, montrons que l'on a
\[
\mathfrak J^m\mathrm R_k\subset\mathrm R_{k+m},
\tag{4.1.7.10}\label{III.4.1.7.10-fr}
\]
la multiplication dans $\mathrm H=H^n(X,\mathcal F)$ par un élément
$x\in\mathfrak J^m$ étant donc l'application $\mu_{x,0}$. Pour tout
$x\in\mathfrak J^m$, le diagramme
\[
\xymatrix{
\mathfrak J^{k+1}\mathcal F\ar[r]^x\ar[d]&
\mathfrak J^{k+m+1}\mathcal F\ar[d]\\
\mathcal F\ar[r]_x&\mathcal F
}
\tag{4.1.7.11}\label{III.4.1.7.11-fr}
\]

\oldpage[III]{127}

(où les flèches horizontales sont la multiplication par $x$, et les
flèches verticales les injections canoniques) est commutatif~; donc le
diagramme correspondant
\[
\xymatrix{
H^n(X,\mathfrak J^{k+1}\mathcal F)\ar[r]^{\mu_{x,m}}\ar[d]&
H^n(X,\mathfrak J^{k+m+1}\mathcal F)\ar[d]\\
H^n(X,\mathcal F)\ar[r]_{\mu_{x,0}}&H^n(X,\mathcal F)
}
\]
est commutatif, ce qui, compte tenu de l'interprétation de $\mathrm R_k$
comme image de
$H^n(X,\mathfrak J^{k+1}\mathcal F)\to H^n(X,\mathcal F)$, prouve
(\hyperref[III.4.1.7.10-fr]{4.1.7.10}) et montre en outre que le $S$-module
gradué $\mathrm R=\bigoplus_{k\geq0}\mathrm R_k$ est un quotient du
sous-$S$-module
$\mathrm M=\bigoplus_{k\geq0}H^n(X,\mathfrak J^{k+1}\mathcal F)$ de
$\mathrm E$~; la remarque faite plus haut montre donc que $\mathrm R$ est
un $S$-module de type fini, ce qui équivaut à l'assertion de
(\hyperref[III.4.1.7.9-fr]{4.1.7.9}) (Bourbaki, \emph{Alg. comm.},
chap. III, §~3, n\textsuperscript{o}~1, th.~1).
\end{proof}

\begin{env}[4.1.7.12]
\label{III.4.1.7.12-fr}
Considérons maintenant le $S$-module gradué
\[
\mathrm N=\bigoplus_{k\geq0}H^{n+1}(X,\mathfrak J^{k+1}\mathcal F)
\]
défini comme dans (4.7.1.8)~; c'est encore un $S$-module de type fini en
vertu de (\hyperref[III.3.3.2-fr]{3.3.2})~; on a
$\mathrm Q_k\subset\mathrm N_k$ pour tout $k$ par
(\hyperref[III.4.1.7.5-fr]{4.1.7.5}) et le diagramme
(\hyperref[III.4.1.7.11-fr]{4.1.7.11}) où on remplace $n$ par $n+1$,
montre que
$S_m\mathrm Q_k=\mathfrak J^m\mathrm Q_k\subset\mathrm Q_{k+m}$.
Autrement dit, $\mathrm Q$ est un sous-$S$-module gradué de $\mathrm N$,
et par suite est de type fini.
\end{env}

\begin{env}[4.1.7.13]
\label{III.4.1.7.13-fr}
Désignons par $\alpha_m$ l'injection canonique
$\mathfrak J^m\to A$, qu'on peut écrire $S_m\to S_0$. Comme
$\mathfrak J^{k+1}\mathcal F_k=0$, le $A$-module
$H^n(X,\mathcal F_k)$ est annulé par $\mathfrak J^{k+1}$~; comme
$\mathrm Q_k$ est l'image du $A$-homomorphisme
$H^n(X,\mathcal F_k)\to H^{n+1}(X,\mathfrak J^{k+1}\mathcal F)$,
$\mathrm Q_k$, en tant que $A$-module, est aussi annulé par
$\mathfrak J^{k+1}$~; cela signifie encore que, dans le $S$-module
$\mathrm Q$, on a
\[
\alpha_{k+1}(S_{k+1})\mathrm Q_k=0.
\tag{4.1.7.14}\label{III.4.1.7.14-fr}
\]
Comme $\mathrm Q$ est un $S$-module de type fini, il existe un entier
$k_0$ et un entier $h$ tels que
$\mathrm Q_{k+h}=S_h\mathrm Q_k$ pour $k\geq k_0$ (II, 2.1.6, (ii))~;
on déduit de cette relation et de
(\hyperref[III.4.1.7.14-fr]{4.1.7.14}) qu'il existe un entier $r>0$ tel
que
\[
\alpha_r(S_r)\mathrm Q=0.
\tag{4.1.7.15}\label{III.4.1.7.15-fr}
\]
\end{env}

\begin{env}[4.1.7.16]
\label{III.4.1.7.16-fr}
Notons maintenant que l'injection canonique
$\mathfrak J^{k+m}\mathcal F\to\mathfrak J^k\mathcal F$ donne en passant
à la cohomologie un $A$-homomorphisme
\[
\nu_m:H^{n+1}(X,\mathfrak J^{k+m}\mathcal F)
\longrightarrow H^{n+1}(X,\mathfrak J^k\mathcal F)
\tag{4.1.7.17}\label{III.4.1.7.17-fr}
\]
et, pour tout $x\in\mathfrak J^m$, on a évidemment la factorisation
\[
\xymatrix@C=4.5em{
H^{n+1}(X,\mathfrak J^k\mathcal F)\ar[r]^{\mu_{x,m}}&
H^{n+1}(X,\mathfrak J^{k+m}\mathcal F)\ar[r]^{\nu_m}&
H^{n+1}(X,\mathfrak J^k\mathcal F)
}
\tag{4.1.7.18}\label{III.4.1.7.18-fr}
\]
dont le composé est $\mu_{x,0}$.
\end{env}

\oldpage[III]{128}

d'où on conclut que, pour tout sous-$A$-module $P$ de
$H^{n+1}(X,\mathfrak J^k\mathcal F)$, on a, dans le $S$-module
$\mathrm N$,
\[
\nu_m(S_mP)=\alpha_m(S_m)P.
\tag{4.1.7.19}\label{III.4.1.7.19-fr}
\]

\begin{lemma}[4.1.7.20]
\label{III.4.1.7.20-fr}
Il existe un entier $m>0$ tel que $\nu_m(\mathrm Q_{k+m})=0$ pour tout
$k\geq k_0$.
\end{lemma}

\begin{proof}
Prenons en effet pour $m$ un multiple de $h$ qui soit $\geq r$~; comme
$\mathrm Q_{k+m}=S_m\mathrm Q_k$ pour $k\geq k_0$, on a en vertu de
(\hyperref[III.4.1.7.19-fr]{4.1.7.19}) et
(\hyperref[III.4.1.7.15-fr]{4.1.7.15})
\[
\nu_m(\mathrm Q_{k+m})=\alpha_m(S_m)\mathrm Q_k
\subset\alpha_r(S_r)\mathrm Q_k=0.
\]
\end{proof}

\begin{env}[4.1.7.21]
\label{III.4.1.7.21-fr}
Remarquons que du diagramme commutatif
\[
\xymatrix@C=2.5em{
H^n(X,\mathcal F)\ar[r]&H^n(X,\mathcal F_k)\ar[r]^-\partial&
H^{n+1}(X,\mathfrak J^{k+1}\mathcal F)\ar[r]&H^{n+1}(X,\mathcal F)\\
H^n(X,\mathcal F)\ar[r]\ar[u]&H^n(X,\mathcal F_{k+m})\ar[r]^-\partial\ar[u]&
H^{n+1}(X,\mathfrak J^{k+m+1}\mathcal F)\ar[r]\ar[u]&
H^{n+1}(X,\mathcal F)\ar[u]
}
\]
provenant lui-même du diagramme commutatif
\[
\xymatrix{
0\ar[r]&\mathfrak J^{k+1}\mathcal F\ar[r]&\mathcal F\ar[r]&
\mathcal F_k\ar[r]&0\\
0\ar[r]&\mathfrak J^{k+m+1}\mathcal F\ar[r]\ar[u]&
\mathcal F\ar[r]\ar[u]&\mathcal F_{k+m}\ar[r]\ar[u]&0
}
\]
où les flèches verticales sont les applications canoniques, on déduit un
diagramme commutatif
\[
\xymatrix@C=2.2em{
0\ar[r]&\mathrm R_k\ar[r]&\mathrm H\ar[r]&\mathrm H_k\ar[r]&
\mathrm Q_k\ar[r]&0\\
0\ar[r]&\mathrm R_{k+m}\ar[r]\ar[u]&\mathrm H\ar[r]\ar@{=}[u]&
\mathrm H_{k+m}\ar[r]\ar[u]&\mathrm Q_{k+m}\ar[r]\ar[u]^{\nu_m}&0
}
\]
où les lignes sont exactes. Comme la dernière flèche verticale est nulle
pour $k\geq k_0$ (\hyperref[III.4.1.7.20-fr]{4.1.7.20}), l'image de
$\mathrm H_{k+m}$ dans $\mathrm H_k$ est contenue dans
$\operatorname{Ker}(\mathrm H_k\to\mathrm Q_k)
=\operatorname{Im}(\mathrm H\to\mathrm H_k)$, mais par ailleurs elle
contient $\operatorname{Im}(\mathrm H\to\mathrm H_k)$ par la commutativité
du diagramme, donc elle lui est égale~; il en est donc de même des images
dans $\mathrm H_k$ des $\mathrm H_{k'}$ pour $k'\geq k+m$, ce qui démontre
la condition (ML) pour le système projectif $(\mathrm H_k)_{k\geq0}$. En
outre, pour tout ouvert affine $U$ de $X$, on a
$H^i(U,\mathcal F_k)=0$ pour $i>0$
(\hyperref[III.1.3.1-fr]{1.3.1}), et pour $m>0$, l'application
$H^0(U,\mathcal F_{k+m})\to H^0(U,\mathcal F_k)$ est surjective
(I, 1.3.9). On peut donc appli-

\oldpage[III]{129}

quer (0, 13.3.1), et l'homomorphisme canonique
$H^n(\widehat X,\widehat{\mathcal F})\to
\varprojlim_kH^n(X,\mathcal F_k)$ est bijectif pour tout $n\geq0$.

Comme le système projectif $(\mathrm H/\mathrm R_k)_{k\geq0}$ est strict,
on peut passer à la limite projective dans les suites exactes
\[
0\longrightarrow\mathrm H/\mathrm R_k\longrightarrow\mathrm H_k
\longrightarrow\mathrm Q_k\longrightarrow0
\tag{4.1.7.22}\label{III.4.1.7.22-fr}
\]
(0, 13.2.2)~; comme $\nu_m(\mathrm Q_{k+m})=0$, on a
$\varprojlim_k\mathrm Q_k=0$, d'où un isomorphisme topologique
\[
\varprojlim_k(\mathrm H/\mathrm R_k)\xrightarrow{\sim}
\varprojlim_k\mathrm H_k.
\]
Mais comme la filtration $(\mathrm R_k)$ de $\mathrm H$ est
$\mathfrak J$-bonne, elle définit sur $\mathrm H$ la topologie
$\mathfrak J$-préadique~; donc $\varprojlim_k(\mathrm H/\mathrm R_k)$ est
le séparé complété de $\mathrm H$ pour la topologie
$\mathfrak J$-préadique, ce qui achève de démontrer
(\hyperref[III.4.1.7-fr]{4.1.7}).
\end{env}
\end{proof}

\begin{env}[4.1.8]
\label{III.4.1.8-fr}
Passons enfin à la démonstration de (\hyperref[III.4.1.5-fr]{4.1.5})~:
pour tout ouvert affine $V$ de $Y$,
$\Gamma(V,\widehat{R^nf_*(\mathcal F)})$ est le séparé complété de
$\Gamma(V,R^nf_*(\mathcal F))$ pour la topologie $\mathfrak J$-préadique
(si $\mathcal J|V=\widetilde{\mathfrak J}$) puisque
$R^nf_*(\mathcal F)$ est un $\mathcal O_Y$-Module cohérent (I, 10.8.4),
et
\[
\Gamma\bigl(V,\varprojlim_kR^nf_*(\mathcal F_k)\bigr)
=\varprojlim_k\Gamma(V,R^nf_*(\mathcal F_k))
\]
(0\textsubscript{I}, 3.2.6)~; le fait que $\varphi_n$ soit un isomorphisme
topologique résulte alors de (\hyperref[III.4.1.7-fr]{4.1.7}) et de
(\hyperref[III.1.4.11-fr]{1.4.11}). D'autre part (toujours en vertu de
(\hyperref[III.1.4.11-fr]{1.4.11})), il résulte de
(\hyperref[III.4.1.7-fr]{4.1.7}) que l'homomorphisme $\psi_{n,V}$ de
(\hyperref[III.4.1.3.3-fr]{4.1.3.3}) est un isomorphisme, donc $\psi_n$
est un isomorphisme par définition de
$R^n\widehat f_*(\widehat{\mathcal F})$.
\end{env}

\begin{corollary}[4.1.9]
\label{III.4.1.9-fr}
Sous les hypothèses de (\hyperref[III.4.1.4-fr]{4.1.4}), pour tout ouvert
affine $V$ de $Y$, l'homomorphisme canonique
\[
H^n(\widehat X\cap f^{-1}(V),\widehat{\mathcal F})
\longrightarrow\Gamma(\widehat Y\cap V,
R^n\widehat f_*(\widehat{\mathcal F}))
\]
est bijectif.
\end{corollary}

\begin{env}[4.1.10]
\label{III.4.1.10-fr}
\emph{Remarque.}
Soit $f:X\to Y$ un morphisme de type fini de préschémas (usuels)
noethériens, et soit $\mathcal F$ un $\mathcal O_X$-Module cohérent dont le
support soit propre sur $Y$ (II, 5.4.10). On sait alors
(\hyperref[III.3.2.4-fr]{3.2.4}) que $R^nf_*(\mathcal F)$ est un
$\mathcal O_Y$-Module cohérent pour tout $n\geq0$. En outre, on peut
toujours supposer que $\mathcal F=u_*(\mathcal G)$, où
$\mathcal G=u^*(\mathcal F)$ est un $\mathcal O_Z$-Module cohérent, $Z$
désignant un sous-préschéma fermé convenable de $X$ dont l'espace
sous-jacent est $\operatorname{Supp}(\mathcal F)$, et $u:Z\to X$
l'injection canonique (I, 9.3.5). Si on pose
\[
\mathcal G_k=\mathcal G\otimes_{\mathcal O_Y}
(\mathcal O_Y/\mathcal J^{k+1}),
\]
on a $\mathcal G_k=u^*(\mathcal F_k)$,
$R^nf_*(\mathcal F_k)=R^n(f\circ u)_*(\mathcal G_k)$ et
$R^nf_*(\mathcal F)=R^n(f\circ u)_*(\mathcal G)$
(\hyperref[III.1.3.4-fr]{1.3.4}), et enfin, compte tenu de (I, 10.9.5),
\[
R^n\widehat f_*(\widehat{\mathcal F})
=R^n\widehat{(f\circ u)}_*(\widehat{\mathcal G}).
\]
On peut alors appliquer (\hyperref[III.4.1.5-fr]{4.1.5}) à $\mathcal G$ et
au morphisme propre $f\circ u$, et on en conclut que sous ces hypothèses,
les résultats de (\hyperref[III.4.1.5-fr]{4.1.5}) sont valables pour
$\mathcal F$ et $f$.
\end{env}

\subsection{Cas particuliers et variantes}
\label{subsection:III.4.2-fr}

La forme la plus utile du th. de comparaison
(\hyperref[III.4.1.5-fr]{4.1.5}) est la suivante~:

\begin{proposition}[4.2.1]
\label{III.4.2.1-fr}
Soient $Y$ un préschéma localement noethérien, $f:X\to Y$ un morphisme
propre, $\mathcal F$ un $\mathcal O_X$-Module cohérent. Alors, pour tout
$y\in Y$ et tout $p$, $(R^pf_*(\mathcal F))_y$ est

\oldpage[III]{130}

un $\mathcal O_y$-module de type fini, donc séparé pour la topologie
$\mathfrak m_y$-préadique, et on a un isomorphisme topologique canonique
\[
\widehat{(R^pf_*(\mathcal F))_y}
\xrightarrow{\sim}
\varprojlim_kH^p\bigl(f^{-1}(y),
\mathcal F\otimes_{\mathcal O_Y}(\mathcal O_y/\mathfrak m_y^k)\bigr)
\tag{4.2.1.1}\label{III.4.2.1.1-fr}
\]
où le premier membre est le complété de $(R^pf_*(\mathcal F))_y$ pour la
topologie $\mathfrak m_y$-préadique, et au second membre $f^{-1}(y)$ est
considéré, pour tout $k\geq0$, comme espace sous-jacent au préschéma
$X\times_Y\operatorname{Spec}(\mathcal O_y/\mathfrak m_y^k)$
(I, 3.6.1).
\end{proposition}

\begin{proof}
Comme $\mathcal O_y$ est un anneau local noethérien et
$(R^pf_*(\mathcal F))_y$ un $\mathcal O_y$-module de type fini
(\hyperref[III.3.2.1-fr]{3.2.1}), la topologie $\mathfrak m_y$-préadique
sur $(R^pf_*(\mathcal F))_y$ est séparée (0\textsubscript{I}, 7.3.5). Les
autres assertions sont conséquences de (\hyperref[III.4.1.7-fr]{4.1.7})
lorsque $Y$ est noethérien et le point $y$ fermé, en remplaçant $Y$ par un
voisinage affine de $y$ et prenant $Y'=\{y\}$, compte tenu de
(G, II, 4.9.1). Dans le cas général, posons
\[
Y_1=\operatorname{Spec}(\mathcal O_y),\qquad
X_1=X\times_YY_1,\qquad
\mathcal F_1=\mathcal F\otimes_{\mathcal O_Y}\mathcal O_{Y_1}
\]
et soit $f_1=f\times1_{Y_1}:X_1\to Y_1$~; $Y_1$ est noethérien et $f_1$
propre (II, 5.4.2, (iii)) et $\mathcal F_1$ est cohérent
(0\textsubscript{I}, 5.3.11). Soit $y_1$ l'unique point fermé de $Y_1$~;
la proposition est valable pour $f_1$, $\mathcal F_1$ et $y_1$~; on a
$\mathcal O_{y_1}=\mathcal O_y$, $f_1^{-1}(y_1)=f^{-1}(y)$ (I, 3.6.5),
les préschémas
$X\times_Y\operatorname{Spec}(\mathcal O_y/\mathfrak m_y^k)$ et
$X_1\times_{Y_1}\operatorname{Spec}(\mathcal O_y/\mathfrak m_y^k)$ étant
canoniquement identifiés (I, 3.3.9)~; en outre,
$\mathcal F_1\otimes_{\mathcal O_{Y_1}}
(\mathcal O_y/\mathfrak m_y^k)$ s'identifie à
$\mathcal F\otimes_{\mathcal O_Y}(\mathcal O_y/\mathfrak m_y^k)$
(I, 9.1.6). Il reste à voir que $R^pf_{1*}(\mathcal F_1)$ est
canoniquement isomorphe à
$R^pf_*(\mathcal F)\otimes_{\mathcal O_Y}\mathcal O_{Y_1}$, ce qui
résulte de (\hyperref[III.1.4.15-fr]{1.4.15}), le morphisme local
$\operatorname{Spec}(\mathcal O_y)\to Y$ étant plat
(0\textsubscript{I}, 6.7.1 et I, 2.4.2).
\end{proof}

Le corollaire suivant utilise la terminologie de la théorie de la dimension
(chap. IV) et ne sera pas appliqué avant le chap. IV.

\begin{corollary}[4.2.2]
\label{III.4.2.2-fr}
Soient $Y$ un préschéma localement noethérien, $f:X\to Y$ un morphisme
propre, $y$ un point de $Y$, $r$ la dimension de $f^{-1}(y)$. Alors, pour
tout $\mathcal O_X$-Module cohérent $\mathcal F$, les faisceaux
$R^pf_*(\mathcal F)$ sont nuls au voisinage de $y$ pour tout $p>r$.
\end{corollary}

\begin{proof}
En effet, on a alors
\[
H^p\bigl(f^{-1}(y),
\mathcal F\otimes_{\mathcal O_Y}(\mathcal O_y/\mathfrak m_y^k)\bigr)=0
\]
(G, II, 4.15.2) pour tout $k$, donc
(\hyperref[III.4.2.1-fr]{4.2.1}) le séparé complété de
$(R^pf_*(\mathcal F))_y$ pour la topologie $\mathfrak m_y$-préadique est
nul, et comme cette topologie est séparée, on a aussi
$(R^pf_*(\mathcal F))_y=0$~; d'où la conclusion, puisque
$R^pf_*(\mathcal F)$ est cohérent (0\textsubscript{I}, 5.2.2).
\end{proof}

\begin{env}[4.2.3]
\label{III.4.2.3-fr}
Le résultat (\hyperref[III.4.2.1-fr]{4.2.1}) est surtout employé pour
$p=0$~; on obtient donc le corollaire suivant~:
\end{env}

\begin{corollary}[4.2.4]
\label{III.4.2.4-fr}
Sous les hypothèses de (\hyperref[III.4.2.1-fr]{4.2.1}), on a un
isomorphisme canonique topologique
\[
\widehat{(f_*(\mathcal F))_y}
\xrightarrow{\sim}
\varprojlim_k\Gamma\bigl(f^{-1}(y),
\mathcal F_y/\mathfrak m_y^k\mathcal F_y\bigr).
\]
\end{corollary}

\subsection{Le théorème de connexion de Zariski}
\label{subsection:III.4.3-fr}
\label{III.4.3-fr}

Les résultats de ce numéro et du suivant généralisent des théorèmes bien
connus de Zariski, et peuvent tous se déduire de
(\hyperref[III.4.2.4-fr]{4.2.4}). Ils sont conséquences du th. suivant~:

\begin{theorem}[4.3.1]
\label{III.4.3.1-fr}
\textnormal{(théorème de connexion).}
Soient $Y$ un préschéma localement noethérien, $f:X\to Y$ un morphisme
propre. Alors $\mathcal A(X)=f_*(\mathcal O_X)$ est une
$\mathcal O_Y$-Algèbre cohérente.

\oldpage[III]{131}

Soit $Y'$ le $Y$-schéma fini au-dessus de $Y$ tel que
$\mathcal A(Y')=\mathcal A(X)$, qui est déterminé à un $Y$-isomorphisme
près (II, 1.3.1 et 6.1.3)~; si $f'=\mathcal A(e)$ est le $Y$-morphisme
$X\to Y'$ déduit de l'isomorphisme identique
$e:\mathcal A(Y')\to\mathcal A(X)$ (II, 1.2.7), $f'$ est propre,
$f'_*(\mathcal O_X)$ est isomorphe à $\mathcal O_{Y'}$, et les fibres
$f'^{-1}(y')$ du morphisme $f'$ sont connexes et non vides pour tout
$y'\in Y'$.
\end{theorem}

\begin{proof}
Soit $g:Y'\to Y$ le morphisme structural. Pour prouver que l'homomorphisme
$\theta:\mathcal O_{Y'}\to f'_*(\mathcal O_X)$ entrant dans la définition
du morphisme $f'$ est bijectif, il suffit, puisque $Y'$ est affine sur $Y$,
de prouver que
\[
g_*(\theta):g_*(\mathcal O_{Y'})\longrightarrow
g_*(f'_*(\mathcal O_X))=f_*(\mathcal O_X)
\]
est l'identité (II, 1.4.2)~; mais cela résulte des définitions puisque
$g_*(\mathcal O_{Y'})=\mathcal A(Y')$ et
$f_*(\mathcal O_X)=\mathcal A(X)$. Le fait que $\mathcal A(X)$ est
cohérent est un cas particulier du th. de finitude
(\hyperref[III.3.2.1-fr]{3.2.1}). Comme $f$ est propre et $g$ séparé,
$f'$ est propre (II, 5.4.3, (i))~; pour terminer la démonstration de
(\hyperref[III.4.3.1-fr]{4.3.1}), il suffit donc d'en prouver le
\end{proof}

\begin{corollary}[4.3.2]
\label{III.4.3.2-fr}
Sous les hypothèses de (\hyperref[III.4.3.1-fr]{4.3.1}), supposons de plus
que $f_*(\mathcal O_X)$ soit isomorphe à $\mathcal O_Y$. Alors les fibres
$f^{-1}(y)$ de $f$ sont connexes et non vides pour tout $y\in Y$.
\end{corollary}

\begin{proof}
L'hypothèse que $f_*(\mathcal O_X)$ est isomorphe à $\mathcal O_Y$
entraîne déjà que $f$ est dominant, donc surjectif puisque $f$ est une
application fermée. On peut se ramener, comme dans
(\hyperref[III.4.2.1-fr]{4.2.1}), au cas où $y$ est fermé dans $Y$~;
$f^{-1}(y)$ étant un espace noethérien, a un nombre fini de composantes
connexes, et c'est l'espace sous-jacent du complété $\widehat X$ le long de
$f^{-1}(y)$. Si $Z_i$ ($1\leq i\leq n$) sont ces composantes connexes, il
est clair que $\Gamma(\widehat X,\mathcal O_{\widehat X})$ est composé
direct des anneaux $\Gamma(Z_i,\mathcal O_{\widehat X})$, et chacun de ces
derniers n'est pas réduit à $0$, puisque la section unité est distincte de
$0$ en chaque point de $\widehat X$. Or, si on applique
(\hyperref[III.4.1.5-fr]{4.1.5}) à $\mathcal F=\mathcal O_X$, dont le
complété le long de $f^{-1}(y)$ est $\mathcal O_{\widehat X}$, on voit que
$\Gamma(\widehat X,\mathcal O_{\widehat X})$ est isomorphe au séparé
complété $\mathfrak m_y$-adique $\widehat{\mathcal O}_y$ de l'anneau local
$\mathcal O_y$~; c'est donc un anneau local qui ne peut être composé
direct de plusieurs anneaux non réduits à $0$ (sans quoi il aurait
plusieurs idéaux maximaux distincts). On a donc $n=1$, ce qui démontre le
corollaire.
\end{proof}

\begin{corollary}[4.3.3]
\label{III.4.3.3-fr}
Sous les hypothèses de (\hyperref[III.4.3.1-fr]{4.3.1}), pour tout $y\in Y$,
l'ensemble des composantes connexes de la fibre $f^{-1}(y)$ est en
correspondance biunivoque avec l'ensemble fini des points de la fibre
$g^{-1}(y)$, où $g:Y'\to Y$ est le morphisme structural (autrement dit,
l'ensemble des idéaux maximaux de $(f_*(\mathcal O_X))_y$).
\end{corollary}

\begin{proof}
Puisque $Y'$ est fini sur $Y$, on sait en effet que $g^{-1}(y)$ est un
espace fini discret (II, 6.1.7). Comme
$f^{-1}(y)=f'^{-1}(g^{-1}(y))$, le corollaire résulte de cette remarque et
de (\hyperref[III.4.3.1-fr]{4.3.1}).
\end{proof}

On a ainsi une interprétation remarquable du $Y$-préschéma $Y'$ défini
dans (\hyperref[III.4.3.1-fr]{4.3.1}). La factorisation $f=g\circ f'$ du
morphisme propre $f$ est analogue à la factorisation obtenue par K.~Stein
pour les applications holomorphes d'espaces analytiques, et nous
l'appellerons par la suite la \emph{factorisation de Stein} de $f$.

\begin{remark}[4.3.4]
\label{III.4.3.4-fr}
Soit $k$ une extension du corps $k(y)$~: si le préschéma
\[
f^{-1}(y)\otimes_{k(y)}k=X\times_Y\operatorname{Spec}(k)
\]
est connexe, il en est de même de $f^{-1}(y)$, qui en est l'image par un
morphisme de projection (I, 3.4.7). Nous dirons que, pour un morphisme
$f:X\to Y$ de préschémas et un point $y\in Y$, la fibre $f^{-1}(y)$ est
\emph{géométriquement connexe}, si pour toute extension $k$ de $k(y)$, le
préschéma
$f^{-1}(y)\otimes_{k(y)}k=X\times_Y\operatorname{Spec}(k)$ est connexe.

\oldpage[III]{132}

Sous les hypothèses de (\hyperref[III.4.3.2-fr]{4.3.2}), on peut alors en
renforcer la conclusion~: les fibres $f^{-1}(y)$ sont en effet
\emph{géométriquement connexes}. Pour le voir, observons que pour toute
extension $k$ de $k(y)$, il existe un anneau local noethérien $A$ et un
homomorphisme local $\varphi:\mathcal O_y\to A$ qui fait de $A$ un
$\mathcal O_y$-module \emph{plat} et tel que le corps résiduel de $A$ soit
$k(y)$-isomorphe à $k$ (0, 10.3.1). Soit alors
$Y_1=\operatorname{Spec}(A)$ et soit $h:Y_1\to Y$ le morphisme local
correspondant à $\varphi$, transformant l'unique point fermée $y_1$ de $Y_1$
en $y$ (I, 2.4.1)~; posons
$X_1=X\times_YY_1$, et $f_1=f\times1_{Y_1}$~; $f_1$ est propre
(II, 5.4.2, (iii)) et $f_1^{-1}(y_1)$ est un $k(y_1)$-préschéma isomorphe à
$X\times_Y\operatorname{Spec}(k)$. Tout revient donc à montrer que
$f_{1*}(\mathcal O_{X_1})=\mathcal O_{Y_1}$ pour pouvoir appliquer
(\hyperref[III.4.3.2-fr]{4.3.2}) à $f_1$. Or $g$ est un morphisme plat,
comme il résulte de (I, 2.4.2) et de
(\hyperref[III.1.4.15-fr]{1.4.15.5})~; on a donc
\[
f_{1*}(\mathcal O_{X_1})
=h^*(f_*(\mathcal O_X))
=h^*(\mathcal O_Y)
=\mathcal O_{Y_1}
\]
en vertu de (\hyperref[III.1.4.15-fr]{1.4.15}) appliqué pour $q=0$.

Dans le cas général (\hyperref[III.4.3.1-fr]{4.3.1}), le même raisonnement
montre que l'on a (avec les notations de
(\hyperref[III.4.3.1-fr]{4.3.1}))
\[
f_{1*}(\mathcal O_{X_1})=h^*(g_*(\mathcal O_{Y'})),
\]
et la factorisation de Stein $f_1=g_1\circ f'_1$ de $f_1$ est telle que
$g_1=g\times1_{Y_1}$ (II, 1.5.2), le $Y_1$-schéma fini correspondant étant
$Y'_1=Y'\times_YY_1$. Tenant compte de la transitivité des fibres
(I, 3.6.4), on voit donc que le nombre des composantes connexes de
$f_1^{-1}(y_1)$ est, en vertu de
(\hyperref[III.4.3.3-fr]{4.3.3}), égal au nombre d'éléments de
\[
g_1^{-1}(y_1)=g^{-1}(y)\otimes_{k(y)}k.
\]
Si on prend pour $k$ une extension algébriquement close de $k(y)$, ce nombre
est indépendant de l'extension algébriquement close considérée et égal au
\emph{nombre géométrique de points} de $g^{-1}(y)$ (I, 6.4.7), ou encore à
la \emph{somme des rangs séparables}
\[
\sum_i[k(y'_i):k(y)]_s
\]
où $y'_i$ parcourt l'ensemble fini $g^{-1}(y)$. On dit encore que ce nombre
est le \emph{nombre géométrique de composantes connexes} de $f^{-1}(y)$. On
notera que les $k(y'_i)$ ne sont autres que les corps résiduels de l'anneau
semi-local $(f_*(\mathcal O_X))_y$.
\end{remark}

\begin{proposition}[4.3.5]
\label{III.4.3.5-fr}
Soient $X$ et $Y$ deux préschémas localement noethériens intègres et
$f:X\to Y$ un morphisme propre dominant. Pour tout $y\in Y$, le nombre de
composantes connexes de $f^{-1}(y)$ est au plus égal au nombre des idéaux
maximaux de la fermeture intégrale $\mathcal O'_y$ de $\mathcal O_y$ dans le
corps des fonctions rationnelles $\mathrm R(X)$.
\end{proposition}

\begin{proof}
En effet, pour tout ouvert $U$ de $Y$,
\[
\Gamma(U,f_*(\mathcal O_X))=\Gamma(f^{-1}(U),\mathcal O_X)
\]
est l'intersection des anneaux locaux $\mathcal O_x$ tels que
$x\in f^{-1}(U)$ (I, 8.2.1.1). On en conclut aussitôt que la fibre
$(f_*(\mathcal O_X))_y$ est un sous-anneau de $\mathrm R(X)$ contenant
$\mathcal O_y$. En outre, comme $f_*(\mathcal O_X)$ est un
$\mathcal O_Y$-Module cohérent, $(f_*(\mathcal O_X))_y$ est un
$\mathcal O_y$-module de type fini, donc contenu dans $\mathcal O'_y$~; on
sait ([13], vol. I, p. 257 et 259) que tout idéal maximal d'un tel anneau
$A$ est l'intersection de $A$ et d'un idéal maximal de $\mathcal O'_y$,
d'où la proposition.
\end{proof}

\begin{definition}[4.3.6]
\label{III.4.3.6-fr}
On dit qu'un anneau local intègre est \emph{unibranche} si sa clôture
intégrale est un anneau local. On dit qu'un point $y$ d'un préschéma intègre
$Y$ est unibranche si l'anneau local $\mathcal O_y$ est unibranche (ce qui
est en particulier le cas lorsque $Y$ est normal au point $y$).
\end{definition}

Soit $A$ un anneau local intègre, et soit $K$ son corps des fractions~; pour
que $A$ soit unibranche, il faut et il suffit que tout sous-anneau $A_1$ de
$K$, contenant $A$ et qui est une $A$-algèbre finie, soit un anneau local.
En effet, soit $A'$ la clôture intégrale de $A$~; il résulte du premier
théorème de Cohen-Seidenberg (Bourbaki, \emph{Alg. comm.}, chap. V, \S~2,
n\textsuperscript{o}~1, th.~1) que tout idéal maximal de $A_1$ est la trace
d'un idéal maximal de $A'$, donc si $A'$ est local, il en est de même de
$A_1$. Réciproquement, $A'$ est limite inductive

\oldpage[III]{133}

de la famille filtrante croissante des sous-$A$-algèbres finies $A_\alpha$
de $A'$, et si chacune des $A_\alpha$ est un anneau local, l'idéal maximal
de $A_\alpha$ est la trace sur $A_\alpha$ de celui de $A_\beta$ pour
$A_\alpha\subset A_\beta$ par le même raisonnement que ci-dessus, donc $A'$
est un anneau local (0, 10.3.1.3).

On notera que si le complété d'un anneau local noethérien $A$ est intègre
(ce qu'on exprime en disant que $A$ est \emph{analytiquement intègre}), $A$
est unibranche. Soient en effet $\mathfrak m$ l'idéal maximal de $A$, $K$
son corps des fractions, $K'$ le corps des fractions de $\widehat A$~; on
a donc
\[
K'=K\otimes_A\widehat A.
\]
Soit $A'_F$ une sous-$A$-algèbre finie de $K$. Le sous-anneau $B_F$ de $K'$
engendré par $\widehat A$ et $A'_F$ est isomorphe à
$A'_F\otimes_A\widehat A$~; c'est un $\widehat A$-module de type fini,
complété de $A'_F$ pour la topologie $\mathfrak m$-adique
(0\textsubscript{I}, 7.3.3 et 7.3.6). Comme $A'_F$ est un anneau semi-local
(Bourbaki, \emph{Alg. comm.}, chap. IV, \S~2, n\textsuperscript{o}~5,
cor. 3 de la prop. 9) et que son complété est intègre, $A'_F$ ne peut avoir
qu'un seul idéal maximal $\mathfrak m'_F$ et on a
$\mathfrak m'_F\cap A=\mathfrak m$~; d'où notre assertion.

\begin{corollary}[4.3.7]
\label{III.4.3.7-fr}
Sous les hypothèses de (\hyperref[III.4.3.5-fr]{4.3.5}), supposons que la
fermeture algébrique de $\mathrm R(Y)$ dans $\mathrm R(X)$ soit de degré
séparable $n$ et que $y\in Y$ soit unibranche. Alors la fibre $f^{-1}(y)$ a
au plus $n$ composantes connexes. En particulier si la fermeture algébrique
de $\mathrm R(Y)$ dans $\mathrm R(X)$ est radicielle sur $\mathrm R(X)$,
$f^{-1}(y)$ est connexe.
\end{corollary}

\begin{proof}
En effet, soit $\mathcal O'_y$ la clôture intégrale de $\mathcal O_y$~; la
fermeture intégrale $\mathcal O''_y$ de $\mathcal O_y$ dans $\mathrm R(X)$
est aussi celle de $\mathcal O'_y$~; mais on sait que si $\mathcal O'_y$ est
un anneau local, $\mathcal O''_y$ est un anneau semi-local dont le nombre
d'idéaux maximaux est au plus égal à $n$ ([13], vol. I, p. 289, th. 22).

Ce corollaire est essentiellement la forme sous laquelle Zariski énonce son
«~théorème de connexion~» pour les schémas algébriques.
\end{proof}

\begin{remark}[4.3.8]
\label{III.4.3.8-fr}
Si on ajoute aux hypothèses de (\hyperref[III.4.3.7-fr]{4.3.7}) l'hypothèse
que $Y$ est normal au point $y$, la fibre $f^{-1}(y)$ est
\emph{géométriquement connexe}, puisque (avec les notations de
(\hyperref[III.4.3.4-fr]{4.3.4})) $g^{-1}(y)$ est réduit à un point $y'$ et
que $k(y')$ est radiciel sur $k(y)$.
\end{remark}

\begin{definition}[4.3.9]
\label{III.4.3.9-fr}
Étant donné un préschéma localement noethérien $Y$, on dit qu'un morphisme de
type fini $f:X\to Y$ est \emph{universellement ouvert} si, pour tout
préschéma irréductible localement noethérien $Y'$, et tout morphisme dominant
$g:Y'\to Y$, toute composante irréductible de $X'=X\times_YY'$ domine $Y'$.
\end{definition}

Si $Y$ est irréductible, cela revient à dire que si $\eta$, $\eta'$ sont les
points génériques de $Y$ et $Y'$ respectivement (de sorte que
$g(\eta')=\eta$), et si on pose $f'=f_{(Y')}$, toute composante irréductible
de $X'$ rencontre $f'^{-1}(\eta')$ (0\textsubscript{I}, 2.1.8)~; cela
implique donc que pour tout ouvert $U$ de $Y$, le morphisme
$f^{-1}(U)\to U$, restriction de $f$, est universellement ouvert.

\begin{corollary}[4.3.10]
\label{III.4.3.10-fr}
Soient $X$, $Y$ deux préschémas localement noethériens intègres.
$f:X\to Y$ un morphisme propre dominant et universellement ouvert. Si la
fermeture algébrique de $\mathrm R(Y)$ dans $\mathrm R(X)$ est radicielle
sur $\mathrm R(Y)$, toute fibre $f^{-1}(y)$ ($y\in Y$) est géométriquement
connexe.
\end{corollary}

\begin{proof}
On peut se borner au cas où $Y=\operatorname{Spec}(B)$, $B$ étant un anneau
intègre noethérien. Il résulte alors de (II, 7.1.7) qu'il existe un anneau
local noethérien intégralement clos $A$ qui domine $\mathcal O_y$ et a
$\mathrm R(Y)$ pour corps des fractions. Soit
$Y'=\operatorname{Spec}(A)$, et soit $h:Y'\to Y$ le morphisme correspondant
à l'injection canonique $B\to A$, qui est birationnel (donc dominant)~; en
outre, si $y'$ est l'unique point fermé de $Y'$, on a $h(y')=y$. Soient
$X'=X\times_YY'$, $f'=f\times1_{Y'}$~; désignons par $\eta$, $\eta'$,
$\xi$ les points génériques de $Y$, $Y'$ et $X$ respectivement, de sorte
que $f(\xi)=\eta$ et $h^{-1}(\eta)=\{\eta'\}$~; en outre,
$k(\eta)=k(\eta')=\mathrm R(Y)$, donc $f'^{-1}(\eta')$ est isomorphe à
$f^{-1}(\eta)$ (I, 3.6.4) et en particulier, puisque $\xi$ est le point
générique de $f^{-1}(\eta)$ (0\textsubscript{I}, 2.1.8),
$f'^{-1}(\eta')$ a un seul point générique. Mais par hypothèse, toute
composante irréductible de $X'$ a son point générique dans
$f'^{-1}(\eta')$, donc $X'$ est nécessairement irréductible, son point
générique $\xi'$ est le point générique de $f'^{-1}(\eta')$ et on a
$k(\xi')=k(\xi)$.

\oldpage[III]{134}

Posons $X''=X'_{\mathrm{red}}$, $f''=f'_{\mathrm{red}}$~; $X''$ est donc
intègre et noethérien, $f''$ est propre (II, 5.4.6) et les espaces
sous-jacents aux fibres $f'^{-1}(y')$ et $f''^{-1}(y')$ sont les mêmes~; en
outre, $\mathrm R(X'')=k(\xi')=\mathrm R(X)$, donc $f''$ vérifie les
hypothèses de (\hyperref[III.4.3.8-fr]{4.3.8}), et $f''^{-1}(y')$ est
géométriquement connexe. Soit maintenant $k$ une extension quelconque de
$k(y)$~; il existe une extension $k_1$ de $k(y)$ dont $k(y')$ et $k$ peuvent
être considérés comme des sous-extensions (Bourbaki, \emph{Alg.}, chap. V,
\S~4, prop. 2). Par hypothèse,
$f''^{-1}(y')\times_{Y'}\operatorname{Spec}(k_1)$ est connexe, et il a même
préschéma réduit que
$f'^{-1}(y')\times_{Y'}\operatorname{Spec}(k_1)$ (I, 5.1.8), donc ce dernier
est connexe, et comme il est isomorphe à
$f^{-1}(y)\times_Y\operatorname{Spec}(k_1)$ (I, 3.6.4), on en conclut que ce
dernier est connexe~; \emph{a fortiori}, il en est de même de
$f^{-1}(y)\times_Y\operatorname{Spec}(k)$ d'après la remarque du début de
(\hyperref[III.4.3.4-fr]{4.3.4}), ce qui achève la démonstration.
\end{proof}

\begin{env}[4.3.11]
\label{III.4.3.11-fr}
\emph{Remarques} (4.3.11). ---
\begin{enumerate}
\item[(i)] Le raisonnement précédent est dû en substance à Zariski [20], à
cela près qu'il peut prendre pour $A$ la clôture intégrale de
$\mathcal O_y$, celle-ci étant un anneau noethérien pour les anneaux locaux
de la géométrie algébrique classique. D'autre part, Zariski prouve que si
$Y$ est la variété de Chow d'un espace projectif $\mathbf P_k^r$ sur un
corps $k$, et si $X$ est la partie fermée de $\mathbf P_k^r\times_kY$ qui
définit la correspondance de Chow entre $\mathbf P_k^r$ et $Y$, alors la
projection $X\to Y$ est un morphisme universellement ouvert
(\emph{loc. cit.}, lemme de la p. 82). Il semble bien que ce soit la seule
propriété formelle des «~coordonnées de Chow~» dont on se soit servi dans
certaines applications~; il y a par suite intérêt dans une telle situation,
à substituer le langage~: fibres d'un morphisme propre (éventuellement
supposé universellement ouvert ou soumis à d'autres restrictions analogues
de régularité locale) au langage~: spécialisation de cycles dans l'espace
projectif.

\item[(ii)] Au chap. IV, nous verrons qu'un morphisme universellement ouvert
$f:X\to Y$ peut encore être défini de la façon suivante (qui justifie la
terminologie)~: pour tout morphisme $Z\to Y$, le morphisme
$f_{(Z)}:X_{(Z)}\to Z$ est ouvert. On peut montrer en outre que si $f$
vérifie les hypothèses de (\hyperref[III.4.3.10-fr]{4.3.10}), alors, si
$y,y'$ sont deux points de $Y$ tels que $y$ soit une spécialisation de $y'$,
le nombre géométrique de composantes connexes de $f^{-1}(y)$ est au plus
égal à celui des composantes connexes de $f^{-1}(y')$.
\end{enumerate}
\end{env}

\begin{corollary}[4.3.12]
\label{III.4.3.12-fr}
Sous les hypothèses de (\hyperref[III.4.3.5-fr]{4.3.5}), supposons en outre
$\mathrm R(Y)$ algébriquement fermé dans $\mathrm R(X)$, et soit $y$ un
point normal de $Y$. Alors $f^{-1}(y)$ est géométriquement connexe, et il
existe un voisinage ouvert $U$ de $y$ dans $Y$ tel que
$f_*(\mathcal O_X|f^{-1}(U))$ soit isomorphe à $\mathcal O_Y|U$. Si plus
particulièrement, on suppose $Y$ normal (et $\mathrm R(Y)$ algébriquement
fermé dans $\mathrm R(X)$), $f_*(\mathcal O_X)$ est isomorphe à
$\mathcal O_Y$.
\end{corollary}

\begin{proof}
La première assertion relative à $f^{-1}(y)$ est un cas particulier de
(\hyperref[III.4.3.8-fr]{4.3.8}). On en

\oldpage[III]{135}

déduit que si $X\xrightarrow{f'}Y'\xrightarrow{g}Y$ est la factorisation
de Stein de $f$ (\hyperref[III.4.3.3-fr]{4.3.3}), $g^{-1}(y)$ se réduit à
un seul point $y'$~; en outre, on a
\[
\mathcal O_y\subset\mathcal O_{y'}=(f_*(\mathcal O_X))_y
\subset\mathrm R(X),
\]
et comme $\mathcal O_{y'}$ est fini sur $\mathcal O_y$ (et \emph{a
fortiori} sur $\mathrm R(Y)$), il est contenu dans $\mathrm R(Y)$ en vertu
de l'hypothèse~; comme $y$ est normal, on a nécessairement
$\mathcal O_{y'}=\mathcal O_y$~; on en conclut que $g$ est un isomorphisme
local au point $y'$ (I, 6.5.4), ce qui achève de prouver la première partie
du corollaire. La seconde résulte de la première, car l'hypothèse
supplémentaire entraîne que $g$ est bijective et un isomorphisme local au
voisinage de tout point de $Y'$, donc un isomorphisme.
\end{proof}

Le fait que (\hyperref[III.4.3.7-fr]{4.3.7}) soit établi dans le cadre des
schémas permet des applications telles que la suivante~:

\begin{proposition}[4.3.13]
\label{III.4.3.13-fr}
Soient $A$ un anneau local noethérien unibranche, $\mathfrak a$ un idéal de
définition de $A$, $A_0=A/\mathfrak a$,
$S=\operatorname{gr}_{\mathfrak a}(A)$ l'anneau gradué associé à $A$ pour
la filtration $\mathfrak a$-préadique~; $S$ est une $A_0$-algèbre graduée
engendrée par $S_1$, $S_1$ étant un $A_0$-module de type fini. Alors
$\operatorname{Proj}(S)$ est un $A_0$-schéma connexe.
\end{proposition}

\begin{proof}
Soit $\mathfrak m$ l'idéal maximal de $A$~; $Y=\operatorname{Spec}(A)$ est
un schéma intègre dont le point $y$ correspondant à $\mathfrak m$ est
l'unique point fermé. Par hypothèse, on a
$\mathfrak m^p\subset\mathfrak a\subset\mathfrak m$ pour un entier $p$,
donc $V(\mathfrak a)=\{\mathfrak m\}$. Soit
$S'=\bigoplus_{n\geq0}\mathfrak a^n$, et soit
$X=\operatorname{Proj}(S')$, qui est le $Y$-schéma obtenu en faisant éclater
l'idéal $\mathfrak a$~; $X$ est intègre et le morphisme structural
$f:X\to Y$ est birationnel (II, 8.1.4) et évidemment projectif. Par suite,
(\hyperref[III.4.3.7-fr]{4.3.7}) est applicable et montre que $f^{-1}(y)$
est connexe~; mais l'espace $f^{-1}(y)$ est sous-jacent à
$\operatorname{Proj}(S'\otimes_AA_0)$ (I, 3.6.1 et II, 2.8.10)~; comme
$S'\otimes_AA_0=S$ par définition, la proposition est démontrée.
\end{proof}


\subsection{Le «~main theorem~» de Zariski}
\label{subsection:III.4.4-fr}

\begin{proposition}[4.4.1]
\label{III.4.4.1-fr}
Soient $Y$ un préschéma localement noethérien, $f:X\to Y$ un morphisme
propre. Soit $X'$ l'ensemble des points $x\in X$ qui sont isolés dans leur
fibre $f^{-1}(f(x))$. Alors l'ensemble $X'$ est ouvert dans $X$, et si
$f=gf'$ est la factorisation de Stein de $f$
(\hyperref[III.4.3.3-fr]{4.3.3}), la restriction de $f'$ à $X'$ est un
isomorphisme de $X'$ sur un sous-préschéma induit sur un ouvert $U$ de $Y'$,
et on a $X'=f'^{-1}(U)$.
\end{proposition}

\begin{proof}
Comme $g^{-1}(f(x))$ est fini et discret
(\hyperref[III.4.3.3-fr]{4.3.3} et II, 6.1.7), pour que $x$ soit isolé dans
$f^{-1}(f(x))$, il faut et il suffit qu'il le soit dans
$f'^{-1}(f'(x))$~; on peut donc se borner au cas où $f'=f$, donc
$f_*(\mathcal O_X)=\mathcal O_Y$. Alors, si $x\in X'$,
$f^{-1}(f(x))$, qui est connexe (\hyperref[III.4.3.2-fr]{4.3.2}), est
nécessairement réduit au point $x$. Comme $f$ est fermé, pour tout voisinage
ouvert $V$ de $x$ dans $X$, $f(X-V)$ est fermé dans $Y$ et ne contient pas
$y=f(x)$, puisque $f^{-1}(y)=\{x\}$~; si $U$ est le complémentaire de
$f(X-V)$ dans $Y$, on a $f^{-1}(U)\subset V$, et on en conclut que les
images réciproques par $f$ d'un système fondamental de voisinages ouverts de
$y$ forment un système fondamental de voisinages ouverts de $x$. L'hypothèse
$f_*(\mathcal O_X)=\mathcal O_Y$ et la définition de l'image directe d'un
faisceau (0\textsubscript{I}, 3.4.1 et 4.2.1) entraînent alors que, si
$f=(\psi,\theta)$, l'homomorphisme
$\theta_y^\#:\mathcal O_y\to\mathcal O_x$ est un isomorphisme. On en conclut
qu'il existe un voisinage ouvert $V$ de $x$ et un voisinage ouvert $U$ de
$y$ tels que la restriction de $f$ à $V$ soit un isomorphisme de $V$ sur
$U$ (I, 6.5.4)~; en outre, d'après ce qu'on vient de voir,

\oldpage[III]{136}

on peut supposer que $f^{-1}(U)=V$, d'où on conclut aussitôt, par définition,
que $V\subset X'$, ce qui achève la démonstration.
\end{proof}

La proposition suivante a été démontrée par Chevalley dans le cas des schémas
algébriques~:

\begin{proposition}[4.4.2]
\label{III.4.4.2-fr}
Soient $Y$ un préschéma localement noethérien, $f:X\to Y$ un morphisme. Les
conditions suivantes sont équivalentes~:
\begin{enumerate}
\item[a)] $f$ est fini.
\item[b)] $f$ est affine et propre.
\item[c)] $f$ est propre et, pour tout $y\in Y$, $f^{-1}(y)$ est un ensemble
fini.
\end{enumerate}
\end{proposition}

\begin{proof}
On sait que a) entraîne b) (II, 6.1.2 et 6.1.11). Si $f$ est propre et affine,
il en est de même du morphisme
$f^{-1}(y)\to\operatorname{Spec}(k(y))$ (II, 1.6.2, (iii) et 5.4.2, (iii)),
et le th. de finitude (\hyperref[III.3.2.1-fr]{3.2.1}) appliqué au faisceau
structural de $f^{-1}(y)$, montre que $f^{-1}(y)=\operatorname{Spec}(A)$, où
$A$ est une $k(y)$-algèbre finie~; donc $f^{-1}(y)$ est un ensemble fini
(II, 6.1.7), et on voit que b) entraîne c). Enfin, comme $f^{-1}(y)$ est un
préschéma algébrique sur $k(y)$, l'hypothèse que l'ensemble $f^{-1}(y)$ est
fini entraîne que l'espace $f^{-1}(y)$ est discret (I, 6.4.4). Avec les
notations de (\hyperref[III.4.4.1-fr]{4.4.1}), on a donc $X'=X$, et
$f':X\to Y'$ est un isomorphisme~; comme $g$ est un morphisme fini, on voit
que c) entraîne a).
\end{proof}

\begin{theorem}[4.4.3]
\label{III.4.4.3-fr}
\textnormal{(«~Main theorem~» de Zariski).}
Soient $Y$ un préschéma noethérien, $f:X\to Y$ un morphisme quasi-projectif,
$X'$ l'ensemble des points $x\in X$ qui sont isolés dans leur fibre
$f^{-1}(f(x))$. Alors $X'$ est une partie ouverte de $X$, et le
sous-préschéma induit $X'$ est isomorphe à un préschéma induit sur une partie
ouverte d'un $Y$-préschéma $Y'$ fini sur $Y$.
\end{theorem}

\begin{proof}
L'hypothèse entraîne qu'il existe un $Y$-préschéma projectif $Z$ tel que $X$
soit $Y$-isomorphe à un sous-préschéma induit sur un ouvert de $Z$
(II, 5.3.2 et 5.5.1). On est donc ramené à démontrer le théorème lorsque $f$
est un morphisme projectif, donc propre (II, 5.5.3), et il résulte alors
aussitôt de (\hyperref[III.4.4.1-fr]{4.4.1}).
\end{proof}

\begin{remark}[4.4.4]
\label{III.4.4.4-fr}
Si $X$ est réduit (resp. irréductible et $X'$ non vide), on peut supposer,
dans l'énoncé de (\hyperref[III.4.4.3-fr]{4.4.3}), que $Y'$ est réduit (resp.
irréductible). En effet, on peut toujours remplacer $Y'$ par le
sous-préschéma adhérence $\overline{X'}$ de $X'$ dans $Y'$ (I, 9.5.11 et
II, 6.1.5, (i) et (ii)), et on sait que si $X'$ est réduit, il en est de même
de $\overline{X'}$ (I, 9.5.9, (i))~; par ailleurs, si $X'$ n'est pas vide, il
est irréductible si $X$ l'est, et $\overline{X'}$ est alors aussi
irréductible.
\end{remark}

\begin{corollary}[4.4.5]
\label{III.4.4.5-fr}
Soient $Y$ un schéma localement noethérien, $f:X\to Y$ un morphisme de type
fini, $x$ un point de $X$ isolé dans sa fibre $f^{-1}(f(x))$. Alors il existe
un voisinage ouvert de $x$ dans $X$ qui est isomorphe à une partie ouverte
d'un $Y$-préschéma fini sur $Y$.
\end{corollary}

\begin{proof}
Soient en effet $y=f(x)$, $U$ un voisinage ouvert affine de $y$ dans $Y$,
$V$ un voisinage ouvert affine de $x$ dans $X$, contenu dans $f^{-1}(U)$.
Comme $Y$ est séparé, l'injection $U\to Y$ est affine (II, 1.6.3), et comme
$V$ est affine sur $U$ (\emph{ibid.}), la restriction de $f$ à $V$ est un
morphisme affine $V\to Y$ (II, 1.6.2, (ii))~; \emph{a fortiori}, cette
restriction est un morphisme quasi-projectif puisqu'il est de type fini
(I, 6.3.5 et II, 5.3.4, (i)). Il suffit alors d'appliquer à cette restriction
le th. (\hyperref[III.4.4.3-fr]{4.4.3}).
\end{proof}

Le cor. (\hyperref[III.4.4.5-fr]{4.4.5}) s'énonce dans le langage de
l'algèbre commutative~:

\oldpage[III]{137}

\begin{corollary}[4.4.6]
\label{III.4.4.6-fr}
Soient $A$ un anneau noethérien, $B$ une $A$-algèbre de type fini,
$\mathfrak q$ un idéal premier de $B$, $\mathfrak p$ son image réciproque
dans $A$. On suppose que $\mathfrak q$ soit à la fois maximal et minimal dans
l'ensemble des idéaux premiers de $B$ dont l'image réciproque est
$\mathfrak p$. Alors il existe $g\in B-\mathfrak q$, une $A$-algèbre finie
$A'$ et un élément $f'\in A'$ tels que les $A$-algèbres $B_g$ et $A'_{f'}$
soient isomorphes.
\end{corollary}

\begin{proof}
Il suffit en effet d'appliquer (\hyperref[III.4.4.5-fr]{4.4.5}) à
$Y=\operatorname{Spec}(A)$ et $X=\operatorname{Spec}(B)$, l'hypothèse sur
$\mathfrak q$ signifiant exactement qu'il est isolé dans sa fibre
(I, 1.1.7).
\end{proof}

On en déduit le résultat suivant, moins général en apparence~:

\begin{corollary}[4.4.7]
\label{III.4.4.7-fr}
Soient $A$ un anneau local noethérien, $B$ une $A$-algèbre de type fini,
$\mathfrak n$ un idéal premier de $B$ dont l'image réciproque dans $A$ est
l'idéal maximal $\mathfrak m$. On suppose que $\mathfrak n$ est maximal dans
$B$ et est minimal dans l'ensemble des idéaux premiers de $B$ dont l'image
réciproque est $\mathfrak m$ (ce qui signifie aussi que
$B\mathfrak m$ est primaire pour $\mathfrak n$). Alors il existe une
$A$-algèbre finie $A'$ et un idéal maximal $\mathfrak m'$ de $A'$ (dont
$\mathfrak m$ est l'image réciproque dans $A$) tels que $B_{\mathfrak n}$
soit isomorphe à la $A$-algèbre $A'_{\mathfrak m'}$.
\end{corollary}

Le cas particulier suivant de (\hyperref[III.4.4.7-fr]{4.4.7}) est aussi
parfois appelé «~Main Theorem~»~:

\begin{corollary}[4.4.8]
\label{III.4.4.8-fr}
Sous les conditions de (\hyperref[III.4.4.7-fr]{4.4.7}), supposons en outre
$A$ et $B$ intègres et ayant même corps des fractions $K$. Alors, si $A$ est
intégralement clos, on a $B=A$.
\end{corollary}

\begin{proof}
En effet, la Remarque (\hyperref[III.4.4.4-fr]{4.4.4}) montre que l'on peut
supposer, dans l'application de (\hyperref[III.4.4.7-fr]{4.4.7}) que $A'$ est
intègre et a $K$ pour corps des fractions~; l'hypothèse sur $A$ entraîne alors
$A'=A$, donc $B_{\mathfrak n}=A$~; comme on a
$A\subset B\subset B_{\mathfrak n}$, on en conclut bien $A=B$.
\end{proof}

L'énoncé (\hyperref[III.4.4.8-fr]{4.4.8}) est la forme donnée par Zariski à
son «~Main theorem~» (étendu aux anneaux locaux noethériens intègres
quelconques).

Les corollaires précédents étaient des variantes de nature locale de
(\hyperref[III.4.4.3-fr]{4.4.3}), qui est un résultat global. Voici une autre
conséquence de nature globale~:

\begin{corollary}[4.4.9]
\label{III.4.4.9-fr}
Soient $Y$ un préschéma intègre localement noethérien, $f:X\to Y$ un
morphisme séparé, de type fini et birationnel. Supposons en outre $Y$ normal
et toutes les fibres $f^{-1}(y)$ finies pour $y\in Y$. Alors $f$ est une
immersion ouverte~; si en outre $f$ est fermé (et en particulier si $f$ est
propre), $f$ est un isomorphisme.
\end{corollary}

\begin{proof}
Soit en effet $x\in X$, et posons $y=f(x)$. Comme $f^{-1}(y)$ est un schéma
algébrique sur $k(y)$, l'hypothèse qu'il est fini entraîne qu'il est discret
(I, 6.4.4)~; en outre $\mathcal O_y$ est intégralement clos et
$\mathcal O_x$ et $\mathcal O_y$ ont même corps des fractions (I, 7.1.5).
On peut donc appliquer (\hyperref[III.4.4.8-fr]{4.4.8}), et si
$f=(\psi,\theta)$, l'homomorphisme
$\theta_y^\#:\mathcal O_y\to\mathcal O_x$ est bijectif~; on en conclut
(I, 6.5.4) que $f$ est un isomorphisme local. Mais comme $f$ est séparé et
$X$ intègre, $f$ est une immersion ouverte (I, 8.2.8). La dernière assertion
résulte de ce que $f$ est dominant.
\end{proof}

\begin{proposition}[4.4.10]
\label{III.4.4.10-fr}
Soient $Y$ un préschéma localement noethérien, $f:X\to Y$ un morphisme
localement de type fini. L'ensemble $X'$ des $x\in X$ isolés dans leur fibre
$f^{-1}(f(x))$ est ouvert dans $X$.
\end{proposition}

\begin{proof}
La question étant locale sur $X$ et $Y$, on peut supposer $X$ et $Y$ affines
noethériens et $f$ de type fini~; $f$ est alors un morphisme affine de type
fini, donc quasi-projectif (II, 5.3.4, (i)), et il suffit d'appliquer
(\hyperref[III.4.4.3-fr]{4.4.3}).
\end{proof}

\oldpage[III]{138}

\begin{corollary}[4.4.11]
\label{III.4.4.11-fr}
Soient $Y$ un préschéma localement noethérien, $f:X\to Y$ un morphisme
propre. L'ensemble $U$ des points $y\in Y$ tels que $f^{-1}(y)$ soit discret
est ouvert dans $Y$, et le morphisme $f^{-1}(U)\to U$ restriction de $f$ est
fini. En particulier, un morphisme propre et quasi-fini $X\to Y$ est fini.
\end{corollary}

\begin{proof}
En effet, le complémentaire de $U$ dans $Y$ est l'image par $f$ de $X-X'$ qui
est fermé dans $X$ en vertu de (\hyperref[III.4.4.10-fr]{4.4.10})~; comme $f$
est une application fermée, $U$ est ouvert. En outre, il résulte de
(II, 6.2.2) que $f^{-1}(y)$ est fini pour tout $y\in U$~; comme le morphisme
$f^{-1}(U)\to U$, restriction de $f$ est propre (II, 5.4.1), il est fini en
vertu de (\hyperref[III.4.4.2-fr]{4.4.2}).
\end{proof}

\begin{env}[4.4.12]
\label{III.4.4.12-fr}
\emph{Remarques} (4.4.12). ---
\begin{enumerate}
\item[(i)] Comme on l'a annoncé dans (II, 6.2.7), nous montrerons au chap. V
que si $Y$ est localement noethérien, tout morphisme quasi-fini et séparé
$f:X\to Y$ est quasi-affine, donc quasi-projectif. Il s'ensuivra que, dans le
Main Theorem (\hyperref[III.4.4.3-fr]{4.4.3}) la conclusion reste valable
lorsqu'on suppose seulement $f$ séparé et de type fini. En effet, il résulte
de (\hyperref[III.4.4.10-fr]{4.4.10}) que $X'$ est ouvert dans $X$, et comme
$X$ est localement noethérien, la restriction de $f$ à $X'$ est encore de
type fini (I, 6.3.5), donc quasi-fini par définition de $X'$, et évidemment
séparé~; on peut donc appliquer (\hyperref[III.4.4.3-fr]{4.4.3}) à cette
restriction, d'où la conclusion.

\item[(ii)] Nous donnerons au chap. IV une démonstration plus élémentaire de
(\hyperref[III.4.4.10-fr]{4.4.10}), utilisant la théorie de la dimension.
\end{enumerate}
\end{env}

\subsection{Complétés de modules d'homomorphismes}
\label{subsection:III.4.5-fr}

\begin{proposition}[4.5.1]
\label{III.4.5.1-fr}
Soient $A$ un anneau noethérien, $\mathfrak J$ un idéal de $A$, $X$ un
$A$-préschéma de type fini, $\mathcal F$, $\mathcal G$ deux
$\mathcal O_X$-Modules cohérents dont les supports ont une intersection
propre sur $Y=\operatorname{Spec}(A)$ (II, 5.4.10). Alors, pour tout entier
$n\geq0$, $\operatorname{Ext}^n_{\mathcal O_X}(X;\mathcal F,\mathcal G)$ est un $A$-module
de type fini, et son séparé complété pour la topologie $\mathfrak J$-préadique
s'identifie canoniquement (avec les notations de
(\hyperref[III.4.1.7-fr]{4.1.7})) à
$\operatorname{Ext}^n_{\mathcal O_{\widehat X}}(\widehat X;\widehat{\mathcal F},
\widehat{\mathcal G})$.
\end{proposition}

\begin{proof}
On sait (T, 4.2) qu'il existe une suite spectrale birégulière
$\mathrm E(\mathcal F,\mathcal G)$ dont l'aboutissement est
$\operatorname{Ext}_{\mathcal O_X}(X;\mathcal F,\mathcal G)$ et dont les termes
$\mathrm E_2$ sont donnés par
\[
\mathrm E_2^{pq}=\mathrm H^p\bigl(X,
\mathcal{E}\!xt^q_{\mathcal O_X}(\mathcal F,\mathcal G)\bigr).
\]
On sait que $\mathcal{E}\!xt^q_{\mathcal O_X}(\mathcal F,\mathcal G)$ est un
$\mathcal O_X$-Module cohérent (0, 12.3.3) dont le support est contenu dans
l'intersection de ceux de $\mathcal F$ et $\mathcal G$ (T, 4.2.2), et est
par suite propre sur $Y$ (II, 5.4.10). On conclut de
(\hyperref[III.3.2.4-fr]{3.2.4}) que les $\mathrm E_2^{pq}$ sont des
$A$-modules de type fini, et par suite (0, 11.1.8) il en est de même de tous
les termes $\mathrm E_r^{pq}$ de la suite spectrale et de son aboutissement.
D'autre part, si $i:\widehat X\to X$ est le morphisme canonique,
$\widehat{\mathcal F}$ et $\widehat{\mathcal G}$ s'identifient canoniquement
à $i^*(\mathcal F)$ et $i^*(\mathcal G)$, et $i$ est plat (I, 10.8.8 et
10.8.9). On sait alors (0, 12.3.4) qu'il existe, pour tout $q\geq0$, un
$i$-morphisme canonique
\[
u_q:\mathcal{E}\!xt^q_{\mathcal O_X}(\mathcal F,\mathcal G)
\longrightarrow
\mathcal{E}\!xt^q_{\mathcal O_{\widehat X}}
(\widehat{\mathcal F},\widehat{\mathcal G}),
\]
et que le $\mathcal O_{\widehat X}$-homomorphisme correspondant
\[
u_q^\#:i^*\bigl(\mathcal{E}\!xt^q_{\mathcal O_X}(\mathcal F,\mathcal G)\bigr)
\longrightarrow
\mathcal{E}\!xt^q_{\mathcal O_{\widehat X}}
(\widehat{\mathcal F},\widehat{\mathcal G})
\]
est un isomorphisme (0, 12.3.5)~; autrement dit (I, 10.8.8),
$\mathcal{E}\!xt^q_{\mathcal O_{\widehat X}}
(\widehat{\mathcal F},\widehat{\mathcal G})$ s'identifie canoniquement au
complété $(\mathcal{E}\!xt^q_{\mathcal O_X}(\mathcal F,\mathcal G))^\wedge$ (avec les
notations de (\hyperref[III.4.1.7-fr]{4.1.7})). On conclut alors du th. de
comparaison (\hyperref[III.4.1.10-fr]{4.1.10}) que pour tout $p\geq0$,
\[
\mathrm H^p\bigl(\widehat X,
\mathcal{E}\!xt^q_{\mathcal O_{\widehat X}}
(\widehat{\mathcal F},\widehat{\mathcal G})\bigr)
\]
s'identifie canoniquement au séparé complété de
$\mathrm H^p\bigl(X,\mathcal{E}\!xt^q_{\mathcal O_X}
(\mathcal F,\mathcal G)\bigr)$ pour la topologie $\mathfrak J$-préadique.

\oldpage[III]{139}

Si on désigne par $\mathrm E(\widehat{\mathcal F},
\widehat{\mathcal G})$ la suite spectrale birégulière définie dans (T, 4.2)
relative à $\widehat{\mathcal F}$ et $\widehat{\mathcal G}$, on voit donc que
si $\widehat A$ désigne le séparé complété de $A$ pour la topologie
$\mathfrak J$-préadique, on a, à un isomorphisme canonique près,
\[
\mathrm E_2^{pq}(\widehat{\mathcal F},\widehat{\mathcal G})
=\mathrm E_2^{pq}(\mathcal F,\mathcal G)\otimes_A\widehat A
\qquad (0\textsubscript{I}, 7.3.3).
\]

Cela étant, on sait que la donnée du morphisme plat $i$ définit un
homomorphisme canonique de suites spectrales
\[
\varphi:\mathrm E(\mathcal F,\mathcal G)\longrightarrow
\mathrm E(\widehat{\mathcal F},\widehat{\mathcal G})
=\mathrm E(i^*(\mathcal F),i^*(\mathcal G))
\]
qui, pour les termes $\mathrm E_2$ (resp. l'aboutissement) se réduit à
l'homomorphisme
\[
\varphi_2^{pq}:\mathrm H^p\bigl(X,
\mathcal{E}\!xt^q_{\mathcal O_X}(\mathcal F,\mathcal G)\bigr)
\longrightarrow
\mathrm H^p\bigl(\widehat X,
\mathcal{E}\!xt^q_{\mathcal O_{\widehat X}}
(\widehat{\mathcal F},\widehat{\mathcal G})\bigr)
\]
\[
\text{(resp. }\varphi^n:\operatorname{Ext}^n_{\mathcal O_X}
(X;\mathcal F,\mathcal G)\longrightarrow
\operatorname{Ext}^n_{\mathcal O_{\widehat X}}
(\widehat X;\widehat{\mathcal F},\widehat{\mathcal G})\text{)}
\]
déduit de $u_q$ (resp. $u_0$) par fonctorialité (0, 12.3.4). Par
tensorisation avec $\widehat A$, les $\varphi_r^{pq}$ et $\varphi^n$ donnent
des homomorphismes de $\widehat A$-modules
\[
\psi_r^{pq}:\mathrm E_r^{pq}(\mathcal F,\mathcal G)
\otimes_A\widehat A\longrightarrow
\mathrm E_r^{pq}(\widehat{\mathcal F},\widehat{\mathcal G})
\]
\[
\psi^n:\operatorname{Ext}^n_{\mathcal O_X}(X;\mathcal F,\mathcal G)
\otimes_A\widehat A\longrightarrow
\operatorname{Ext}^n_{\mathcal O_{\widehat X}}
(\widehat X;\widehat{\mathcal F},\widehat{\mathcal G}).
\]
Comme $\widehat A$ est un $A$-module plat (0\textsubscript{I}, 7.3.3), les
$\widehat A$-modules
$\mathrm E_r^{pq}(\mathcal F,\mathcal G)\otimes_A\widehat A$ forment une
suite spectrale birégulière d'aboutissement les
$\operatorname{Ext}^n_{\mathcal O_X}(X;\mathcal F,\mathcal G)\otimes_A\widehat A$, et les
$\psi_r^{pq}$ et $\psi^n$ un morphisme de suites spectrales. Comme les
$\psi_2^{pq}$ sont des isomorphismes, il en est de même des $\psi^n$
(0, 11.1.5).
\end{proof}

\begin{corollary}[4.5.2]
\label{III.4.5.2-fr}
Sous les hypothèses de (\hyperref[III.4.5.1-fr]{4.5.1}), supposons en outre
que $A$ soit un anneau $\mathfrak J$-adique noethérien. Alors, pour tout
entier $n\geq0$, $\operatorname{Ext}^n_{\mathcal O_X}(X;\mathcal F,\mathcal G)$
s'identifie canoniquement à
$\operatorname{Ext}^n_{\mathcal O_{\widehat X}}
(\widehat X;\widehat{\mathcal F},\widehat{\mathcal G})$.
\end{corollary}

\begin{proof}
Il suffit de remarquer que $\operatorname{Ext}^n_{\mathcal O_X}
(X;\mathcal F,\mathcal G)$, étant un $A$-module de type fini, est séparé et
complet pour la topologie $\mathfrak J$-préadique (0\textsubscript{I}, 7.3.6).
\end{proof}

Le cas particulier $n=0$ de (\hyperref[III.4.5.1-fr]{4.5.1}) s'énonce de la
façon suivante~:

\begin{corollary}[4.5.3]
\label{III.4.5.3-fr}
Sous les hypothèses de (\hyperref[III.4.5.1-fr]{4.5.1}), pour tout
homomorphisme $u:\mathcal F\to\mathcal G$, désignons par
$\widehat u:\widehat{\mathcal F}\to\widehat{\mathcal G}$ l'homomorphisme
complété (I, 10.8.4). Alors on a un isomorphisme canonique
\[
\label{III.4.5.3.1-fr}
\left(\operatorname{Hom}_{\mathcal O_X}(\mathcal F,\mathcal G)\right)^\wedge
\xrightarrow{\sim}
\operatorname{Hom}_{\mathcal O_{\widehat X}}
(\widehat{\mathcal F},\widehat{\mathcal G}),
\tag{4.5.3.1}
\]
où le premier membre est le séparé complété pour la topologie
$\mathfrak J$-préadique du $A$-module
$\operatorname{Hom}_{\mathcal O_X}(\mathcal F,\mathcal G)$, cet isomorphisme étant obtenu
par passage aux séparés complétés à partir de l'homomorphisme
$u\mapsto\widehat u$.
\end{corollary}

\subsection{Relations entre morphismes formels et morphismes usuels}

\begin{proposition}[4.6.1]
\label{III.4.6.1-fr}
Soient $Y$ un préschéma localement noethérien, $f:X\to Y$ un morphisme
propre, $\mathcal F$ un $\mathcal O_X$-Module cohérent et $f$-plat, $y$ un
point de $Y$. Supposons que pour un entier $n$, on ait
\[
H^n\bigl(f^{-1}(y),\mathcal F\otimes_{\mathcal O_Y}k(y)\bigr)=0.
\]
Alors il existe un voisinage $U$ de $y$ dans $Y$ tel que
$R^nf_*(\mathcal F)|U=0$, et pour tout entier $p\geq0$, l'homomorphisme
canonique
\[
\bigl(R^{n-1}f_*(\mathcal F)\bigr)_y
\longrightarrow
H^{n-1}\bigl(f^{-1}(y),
\mathcal F\otimes_{\mathcal O_Y}
(\mathcal O_y/\mathfrak m_y^{p+1})\bigr)
\]
de (\hyperref[III.4.2.1.1-fr]{4.2.1.1}) est surjectif.
\end{proposition}

\begin{proof}
Comme $R^nf_*(\mathcal F)$ est un $\mathcal O_Y$-Module cohérent
(\hyperref[III.3.2.1-fr]{3.2.1}), la première assertion de la proposition
sera établie si l'on prouve que l'on a
$(R^nf_*(\mathcal F))_y=0$ (0\textsubscript{I}, 5.2.2)~; en vertu de
(\hyperref[III.4.2.1-fr]{4.2.1}), il suffira de prouver que
\[
H^n\bigl(f^{-1}(y),\mathcal F\otimes_{\mathcal O_Y}
(\mathcal O_y/\mathfrak m_y^{p+1})\bigr)=0
\]
pour tout $p$. C'est vrai par hypothèse pour $p=0$~; nous allons le démontrer
par récurrence sur $p$. Posons
\[
X_p=X\times_Y\operatorname{Spec}
(\mathcal O_y/\mathfrak m_y^{p+1}),
\]
de sorte que $X_{p-1}$ est un sous-préschéma fermé de $X_p$, ayant même
espace sous-jacent (I, 3.6.1)~; l'hypothèse de récurrence
\[
H^n\bigl(X_{p-1},\mathcal F\otimes_{\mathcal O_Y}
(\mathcal O_y/\mathfrak m_y^p)\bigr)=0
\]
entraîne donc
\[
H^n\bigl(X_p,\mathcal F\otimes_{\mathcal O_Y}
(\mathcal O_y/\mathfrak m_y^p)\bigr)=0~;
\]
d'autre part, la suite exacte de cohomologie donne, à partir de la suite
exacte
\[
0\longrightarrow
\mathfrak m_y^p\mathcal F/\mathfrak m_y^{p+1}\mathcal F
\longrightarrow
\mathcal F/\mathfrak m_y^{p+1}\mathcal F
\longrightarrow
\mathcal F/\mathfrak m_y^p\mathcal F
\longrightarrow0
\]
de $\mathcal O_{X_p}$-Modules, la suite exacte
\[
H^n\bigl(X_p,
\mathfrak m_y^p\mathcal F/\mathfrak m_y^{p+1}\mathcal F\bigr)
\longrightarrow
H^n\bigl(X_p,\mathcal F/\mathfrak m_y^{p+1}\mathcal F\bigr)
\longrightarrow
H^n\bigl(X_p,\mathcal F/\mathfrak m_y^p\mathcal F\bigr)
\]
et il suffira de montrer que l'on a
\[
\label{III.4.6.1.1-fr}
H^n\bigl(X_p,
\mathfrak m_y^p\mathcal F/\mathfrak m_y^{p+1}\mathcal F\bigr)=0.
\tag{4.6.1.1}
\]
car alors $H^n(X_p,\mathcal F/\mathfrak m_y^{p+1}\mathcal F)$ sera un
sous-module de $H^n(X_p,\mathcal F/\mathfrak m_y^p\mathcal F)$, donc $0$ en
vertu de l'hypothèse de récurrence.

\oldpage[III]{140}

Notons maintenant que la fibre
$Z=f^{-1}(y)=X\times_Y\operatorname{Spec}(k(y))$ est un sous-préschéma fermé
de $X_p$, et que
$\mathfrak m_y^p\mathcal F/\mathfrak m_y^{p+1}\mathcal F$ est annulé par
$\mathfrak m_y$, donc peut être considéré comme un $\mathcal O_Z$-Module, de
sorte que
\[
H^n\bigl(Z,\mathfrak m_y^p\mathcal F/
\mathfrak m_y^{p+1}\mathcal F\bigr)
=H^n\bigl(X_p,\mathfrak m_y^p\mathcal F/
\mathfrak m_y^{p+1}\mathcal F\bigr).
\]
Cela étant, nous allons montrer que le $\mathcal O_Z$-homomorphisme canonique
\[
\label{III.4.6.1.2-fr}
(\mathcal F/\mathfrak m_y\mathcal F)
\otimes_{k(y)}(\mathfrak m_y^p/\mathfrak m_y^{p+1})
\longrightarrow
\mathfrak m_y^p\mathcal F/\mathfrak m_y^{p+1}\mathcal F
\tag{4.6.1.2}
\]
est bijectif~; cela établi, il en résultera, puisque
$\mathfrak m_y^p/\mathfrak m_y^{p+1}$ est un $k(y)$-module libre, que l'on a
\[
\begin{split}
H^n\bigl(Z,\mathfrak m_y^p\mathcal F/
\mathfrak m_y^{p+1}\mathcal F\bigr)
&=H^n\bigl(Z,\mathcal F/\mathfrak m_y\mathcal F\bigr)
\otimes_{k(y)}(\mathfrak m_y^p/\mathfrak m_y^{p+1})\\
&=0
\end{split}
\]
(0, 12.2.3), puisque
$H^n(Z,\mathcal F/\mathfrak m_y\mathcal F)=0$ par hypothèse, d'où
(\hyperref[III.4.6.1.1-fr]{4.6.1.1}). Pour établir la première assertion, il
reste donc à prouver que (\hyperref[III.4.6.1.2-fr]{4.6.1.2}) est bijectif~;
comme la question est ponctuelle sur $X$ et que $\mathcal F_x$ est un
$\mathcal O_y$-module plat par hypothèse pour tout $x\in f^{-1}(y)$, il
suffit d'appliquer (0, 10.2.1, c)), car
$\mathcal F_x/\mathfrak m_y\mathcal F_x$ est un module plat sur le corps
$k(y)=\mathcal O_y/\mathfrak m_y$.

Pour démontrer la seconde assertion de (\hyperref[III.4.6.1-fr]{4.6.1}), on
se ramène aussitôt, comme dans (\hyperref[III.4.2.1-fr]{4.2.1}), au cas où
$Y$ est affine et $y$ fermé. Notons que
(\hyperref[III.4.6.1.1-fr]{4.6.1.1}) donne, par un raisonnement analogue,
pour tout $k>0$, la relation
\[
\label{III.4.6.1.3-fr}
H^n\bigl(X_{k+1},
\mathfrak m_y^k\mathcal F/
\mathfrak m_y^{k+p+1}\mathcal F\bigr)=0.
\tag{4.6.1.3}
\]

\oldpage[III]{141}

d'où on déduit, par (\hyperref[III.4.2.1-fr]{4.2.1}), que l'on a aussi
\[
\label{III.4.6.1.4-fr}
\bigl(R^nf_*(\mathfrak m_y^k\mathcal F)\bigr)_y=0.
\tag{4.6.1.4}
\]
Cela étant, on tire de la suite exacte de cohomologie l'exactitude de la
suite
\[
\bigl(R^{n-1}f_*(\mathcal F)\bigr)_y
\longrightarrow
\bigl(R^{n-1}f_*(\mathcal F/\mathfrak m_y^p\mathcal F)\bigr)_y
\longrightarrow
\bigl(R^nf_*(\mathfrak m_y^p\mathcal F)\bigr)_y=0
\]
et comme $y$ est fermé et $Y$ affine, on a (\hyperref[III.1.4.11-fr]{1.4.11})
\[
R^{n-1}f_*(\mathcal F/\mathfrak m_y^p\mathcal F)
=\bigl(H^{n-1}(X,\mathcal F/\mathfrak m_y^p\mathcal F)\bigr)^\sim
=\bigl(H^{n-1}(f^{-1}(y),
\mathcal F/\mathfrak m_y^p\mathcal F)\bigr)^\sim
\]
(G, II, 4.9.1)~; or
$H^{n-1}(f^{-1}(y),\mathcal F/\mathfrak m_y^p\mathcal F)$ est un
$(\mathcal O_y/\mathfrak m_y^p)$-module, d'où
\[
\bigl(R^{n-1}f_*(\mathcal F/\mathfrak m_y^p\mathcal F)\bigr)_y
=H^{n-1}(f^{-1}(y),\mathcal F/\mathfrak m_y^p\mathcal F)
\]
et cela achève de prouver (\hyperref[III.4.6.1-fr]{4.6.1}).
\end{proof}

\begin{corollary}[4.6.2]
\label{III.4.6.2-fr}
Soient $Y$ un préschéma localement noethérien, $f:X\to Y$ un morphisme propre
et plat, $\mathcal F$, $\mathcal G$ deux $\mathcal O_X$-Modules localement
libres, $y$ un point de $Y$. Posons
\[
X_y=f^{-1}(y)=X\otimes_Yk(y),\qquad
\mathcal F_y=\mathcal F\otimes_{\mathcal O_Y}k(y),\qquad
\mathcal G_y=\mathcal G\otimes_{\mathcal O_Y}k(y),
\]
et supposons que l'on ait
\[
\label{III.4.6.2.1-fr}
H^1\bigl(X_y,
\mathcal{H}\!om_{\mathcal O_{X_y}}(\mathcal F_y,\mathcal G_y)\bigr)=0.
\tag{4.6.2.1}
\]
Alors, pour tout homomorphisme $u_0:\mathcal F_y\to\mathcal G_y$, il existe
un voisinage ouvert $U$ de $y$ et un homomorphisme
\[
u:\mathcal F|f^{-1}(U)\longrightarrow\mathcal G|f^{-1}(U)
\]
tels que $u_0$ soit égal à l'homomorphisme $u\otimes1$.
\end{corollary}

\begin{proof}
En effet, l'hypothèse permet d'appliquer (\hyperref[III.4.6.1-fr]{4.6.1}) au
$\mathcal O_X$-Module cohérent
$\mathcal H=\mathcal{H}\!om_{\mathcal O_X}(\mathcal F,\mathcal G)$ pour
$n=1$ et $p=0$, car $\mathcal H$ est localement libre et \emph{a fortiori}
$f$-plat, et le $\mathcal O_{X_y}$-Module
$\mathcal H\otimes_{\mathcal O_Y}k(y)$ s'identifie alors à
$\mathcal{H}\!om_{\mathcal O_{X_y}}(\mathcal F_y,\mathcal G_y)$
(0\textsubscript{I}, 6.2.2). On peut supposer $Y=\operatorname{Spec}(A)$
affine, et alors (\hyperref[III.1.4.11-fr]{1.4.11})
\[
R^0f_*(\mathcal H)=
\bigl(\operatorname{Hom}_{\mathcal O_X}(\mathcal F,\mathcal G)\bigr)^\sim,
\]
donc
\[
\bigl(R^0f_*(\mathcal H)\bigr)_y
=\operatorname{Hom}_{\mathcal O_X}(\mathcal F,\mathcal G)
\otimes_A\mathcal O_y~;
\]
l'homomorphisme canonique
\[
\operatorname{Hom}_{\mathcal O_X}(\mathcal F,\mathcal G)
\otimes_A\mathcal O_y
\longrightarrow
\operatorname{Hom}_{\mathcal O_{X_y}}(\mathcal F_y,\mathcal G_y)
\]
étant surjectif par (\hyperref[III.4.6.1-fr]{4.6.1}), cela établit le
corollaire, tout élément de
$\operatorname{Hom}_{\mathcal O_X}(\mathcal F,\mathcal G)
\otimes_A\mathcal O_y$ pouvant toujours se mettre sous la forme
$u\otimes(1/s)$, où $s\notin\mathfrak m_y$ est un élément de $A$.
\end{proof}

Ce corollaire peut se compléter par le suivant~:

\begin{corollary}[4.6.3]
\label{III.4.6.3-fr}
Sous les hypothèses de (\hyperref[III.4.6.2-fr]{4.6.2}), si $u_0$ est
injectif (resp. surjectif, bijectif), on peut supposer qu'il en est de même de
$u$.
\end{corollary}

\begin{proof}
On peut se borner au cas où $U=Y$. Il suffit de prouver que si $u_0$ est
injectif (resp. surjectif), $\operatorname{Ker}u_x=0$ (resp.
$\operatorname{Coker}u_x=0$) pour tout $x\in f^{-1}(y)$~: en effet,
$\operatorname{Ker}u$ et $\operatorname{Coker}u$ sont des
$\mathcal O_X$-Modules cohérents (0\textsubscript{I}, 5.3.4), donc il
existera un voisinage $V$ de $f^{-1}(y)$ dans $X$ tel que la restriction de
$\operatorname{Ker}u$ (resp. $\operatorname{Coker}u$) à $V$ soit $0$
(0\textsubscript{I}, 5.2.2)~; comme $f$ est fermé, il existera un voisinage
$U'\subset U$ de $y$ tel que $f^{-1}(U')\subset V$, et
(\hyperref[III.4.6.3-fr]{4.6.3}) sera démontré. Par hypothèse,
\[
u_x\otimes1:
\mathcal F_x\otimes_{\mathcal O_y}k(y)
\longrightarrow
\mathcal G_x\otimes_{\mathcal O_y}k(y)
\]
est injectif (resp. surjectif), $\mathcal F_x$ et $\mathcal G_x$ sont des
$\mathcal O_x$-modules libres de type fini et $\mathcal O_x$ est un
$\mathcal O_y$-module plat. Lorsque l'on suppose $u_x\otimes1$ injectif, le
fait que $u_x$ soit injectif résulte de (0, 10.2.4). Lorsque l'on suppose
$u_x\otimes1$ surjectif, \emph{a fortiori} l'homomorphisme
\[
\mathcal F_x/\mathfrak m_x\mathcal F_x
\longrightarrow
\mathcal G_x/\mathfrak m_x\mathcal G_x,
\]
qui s'en déduit par passage aux quotients, est surjectif~; comme
$\mathcal G_x$ est un $\mathcal O_x$-module de type fini et que
$\mathcal O_x$ est un anneau local d'idéal maximal $\mathfrak m_x$, la
conclusion résulte du lemme de Nakayama (Bourbaki, \emph{Alg.}, chap. VIII,
\S\,6, no~3, cor. 4 de la prop. 6).
\end{proof}

On déduit en particulier de (\hyperref[III.4.6.3-fr]{4.6.3})~:

\begin{corollary}[4.6.4]
\label{III.4.6.4-fr}
Soient $Y$ un préschéma localement noethérien, $f:X\to Y$ un morphisme propre
et plat, $y$ un point de $Y$, $X_y=X\otimes_Yk(y)$. Soit $\mathcal E_0$ un
$\mathcal O_{X_y}$-Module localement libre tel que
\[
\label{III.4.6.4.1-fr}
H^1\bigl(X_y,
\mathcal{H}\!om_{\mathcal O_{X_y}}(\mathcal E_0,\mathcal E_0)\bigr)=0.
\tag{4.6.4.1}
\]
Soient $\mathcal F$, $\mathcal G$ deux $\mathcal O_X$-Modules localement
libres tels que $\mathcal F_y$ et $\mathcal G_y$ (avec les notations de
(\hyperref[III.4.6.2-fr]{4.6.2})) soient isomorphes à $\mathcal E_0$. Alors
il existe un voisinage ouvert $U$ de $y$ tel que $\mathcal F|f^{-1}(U)$ et
$\mathcal G|f^{-1}(U)$ soient isomorphes.
\end{corollary}

Plus particulièrement~:

\begin{corollary}[4.6.5]
\label{III.4.6.5-fr}
Sous les hypothèses de (\hyperref[III.4.6.4-fr]{4.6.4}) sur $f$, $X$, $Y$,
supposons que $H^1(X_y,\mathcal O_{X_y})=0$. Si $\mathcal F$ et $\mathcal G$
sont deux $\mathcal O_X$-Modules inversibles tels que $\mathcal F_y$ et
$\mathcal G_y$ soient isomorphes, il existe un voisinage ouvert $U$ de $y$
tel que $\mathcal F|f^{-1}(U)$ et $\mathcal G|f^{-1}(U)$ soient isomorphes.
\end{corollary}

\begin{proof}
Il suffit d'appliquer (\hyperref[III.4.6.4-fr]{4.6.4}) aux modules
$\mathcal F^{-1}\otimes\mathcal G$ et $\mathcal O_X$.
\end{proof}

\begin{remarks}[4.6.6]
\label{III.4.6.6-fr}
\begin{enumerate}
\item[(i)] En utilisant (\hyperref[III.4.6.5-fr]{4.6.5}), nous établirons au
chap. V la classification des faisceaux inversibles sur un fibré projectif,
annoncée dans (II, 4.2.7).
\item[(ii)] Le résultat de (\hyperref[III.4.6.1-fr]{4.6.1}) apparaîtra au
\S\,7 comme conséquence de propositions plus générales.
\end{enumerate}
\end{remarks}

\begin{proposition}[4.6.7]
\label{III.4.6.7-fr}
Soient $Z$ un préschéma localement noethérien, $X$, $Y$ deux $Z$-préschémas
tels que les morphismes structuraux $g:X\to Z$, $h:Y\to Z$ soient propres.
Soient $f:X\to Y$ un $Z$-morphisme, $z$ un point de $Z$, et soit
\[
f_z=f\times_Z1:
X\otimes_Zk(z)\longrightarrow Y\otimes_Zk(z).
\]
\begin{enumerate}
\item[(i)] Si $f_z$ est un morphisme fini (resp. une immersion fermée), il
existe un voisinage ouvert $U$ de $z$ tel que le morphisme
$g^{-1}(U)\to h^{-1}(U)$, restriction de $f$, soit un morphisme fini (resp.
une immersion fermée).
\item[(ii)] On suppose de plus que $g$ soit un morphisme plat. Alors, si
$f_z$ est un isomorphisme, il existe un voisinage ouvert $U$ de $z$ tel que
le morphisme $g^{-1}(U)\to h^{-1}(U)$, restriction de $f$, soit un
isomorphisme.
\end{enumerate}
\end{proposition}

\begin{proof}
Dans les deux cas, il suffira de prouver que pour tout $y\in h^{-1}(z)$ il
existe un voisinage $V_y$ de $y$ tel que la restriction
$f^{-1}(V_y)\to V_y$ de $f$ soit un morphisme fini (resp. une immersion
fermée, un isomorphisme)~; il en résultera alors en effet que si $V$ est la
réunion des $V_y$, la restriction $f^{-1}(V)\to V$ de $f$ est un morphisme
fini (resp. une immersion fermée, un isomorphisme) (II, 6.1.1 et I, 4.2.4).
Comme $h$ est un morphisme fermé, il existera un voisinage ouvert $U$ de $z$
tel que $h^{-1}(U)\subset V$, et la proposition sera démontrée.

\noindent\textit{(i)} Notons tout d'abord que $f$ est un morphisme propre
(II, 5.4.3)~; si on suppose $f_z$ fini, l'existence pour tout
$y\in h^{-1}(z)$ d'un voisinage $V_y$ tel que $f^{-1}(V_y)\to V_y$ soit fini
résulte de (\hyperref[III.4.4.11-fr]{4.4.11}).

\oldpage[III]{143}

Pour traiter le cas où $f_z$ est une immersion fermée, on peut donc déjà
supposer que le morphisme $f$ est fini, donc que
$X=\operatorname{Spec}(\mathcal B)$, où $\mathcal B$ est une
$\mathcal O_Y$-Algèbre cohérente, le morphisme $f$ correspondant (II, 1.2.7)
à l'homomorphisme canonique $u:\mathcal O_Y\to\mathcal B$. Si l'on prouve
que pour tout $y\in h^{-1}(z)$, l'homomorphisme
$u_y:\mathcal O_y\to\mathcal B_y$ est surjectif, il en résultera que pour un
voisinage $V_y$ de $y$, $u|V_y$ sera surjectif, le faisceau
$\operatorname{Coker}u$ étant cohérent (0\textsubscript{I}, 5.3.4 et 5.2.2).
Cela étant, le morphisme fini $f_z$ correspond à l'homomorphisme
\[
v=u\otimes1:
\mathcal O_Y\otimes_{\mathcal O_Z}k(z)
\longrightarrow
\mathcal B\otimes_{\mathcal O_Z}k(z),
\]
et l'hypothèse que $f_z$ est une immersion fermée entraîne que
l'homomorphisme
\[
u_y\otimes1:
\mathcal O_y\otimes_{\mathcal O_z}k(z)
\longrightarrow
\mathcal B_y\otimes_{\mathcal O_z}k(z)
\]
est surjectif. Comme $\mathcal B_y$ est un $\mathcal O_y$-module de type
fini et $\mathcal O_y$ un anneau local noethérien, la conclusion résulte
comme dans (\hyperref[III.4.6.3-fr]{4.6.3}) du lemme de Nakayama.

\noindent\textit{(ii)} Le même raisonnement que ci-dessus montre qu'il
suffit cette fois de prouver que $u_y$ est bijectif, sachant que
$u_y\otimes1$ est bijectif.

Cela résultera du lemme suivant~:
\end{proof}

\begin{lemma}[4.6.7.1]
\label{III.4.6.7.1-fr}
Soient $A$, $B$ deux anneaux locaux noethériens,
$\rho:A\to B$ un homomorphisme local, $u:N\to M$ un homomorphisme de
$B$-modules. On suppose que $M$ est un $A$-module plat, $N$ un $B$-module
de type fini et que
\[
u\otimes1:N\otimes_Ak\longrightarrow M\otimes_Ak
\]
(où $k$ est le corps résiduel de $A$) est injectif. Alors $N$ est un
$A$-module plat et $u$ est injectif.
\end{lemma}

\begin{proof}
Pour établir la première assertion, il faut montrer que pour tout couple de
$A$-modules de type fini $P$, $Q$ et tout $A$-homomorphisme injectif
$v:P\to Q$,
$1_N\otimes v:N\otimes_AP\to N\otimes_AQ$ est injectif. Or, on a le
diagramme commutatif
\[
\xymatrix{
N\otimes_AP\ar[r]^{1_N\otimes v}\ar[d]_{u\otimes1_P}
  &N\otimes_AQ\ar[d]^{u\otimes1_Q}\\
M\otimes_AP\ar[r]_{1_M\otimes v}
  &M\otimes_AQ
}
\]
et comme $1_M\otimes v$ est injectif par hypothèse, il suffira de prouver
qu'il en est de même de $u\otimes1_P$. Soit $\mathfrak m$ l'idéal maximal
de $A$~; la filtration $\mathfrak m$-adique sur le $A$-module
$N\otimes_AP$ est aussi sa filtration $\mathfrak mB$-adique en tant que
$B$-module~; la topologie définie par cette filtration est donc séparée,
puisque $B$ est noethérien, que $\mathfrak mB$ est contenu dans le radical de
$B$, et que $N\otimes_AP$ est un $B$-module de type fini, $N$ étant un
$B$-module de type fini et $P$ un $A$-module de type fini
(0\textsubscript{I}, 7.3.5). Il suffit donc de prouver que l'homomorphisme
\[
\operatorname{gr}(u\otimes1_P):
\operatorname{gr}_\bullet(N\otimes_AP)
\longrightarrow
\operatorname{gr}_\bullet(M\otimes_AP)
\]
(où les modules gradués sont relatifs aux filtrations $\mathfrak m$-adiques)
est injectif (Bourbaki, \emph{Alg. comm.}, chap. III, \S\,2, no~8, cor. 1 du
th. 1). Notons maintenant que puisque $M$ est un $A$-module plat, les
homomorphismes
\[
M\otimes_A(\mathfrak m^nP)
\longrightarrow
\mathfrak m^n(M\otimes_AP)
\]
sont bijectifs~; il en est donc de même de l'homomorphisme canonique
\[
\varphi_M:
\operatorname{gr}_0(M)\otimes_A\operatorname{gr}_\bullet(P)
\longrightarrow
\operatorname{gr}_\bullet(M\otimes_AP).
\]

\oldpage[III]{144}

Or, on a un diagramme commutatif
\[
\xymatrix{
\operatorname{gr}_0(N)\otimes_A\operatorname{gr}_\bullet(P)
  \ar[r]^{\operatorname{gr}_0(u)\otimes1}\ar[d]_{\varphi_N}
  &\operatorname{gr}_0(M)\otimes_A\operatorname{gr}_\bullet(P)
  \ar[d]^{\varphi_M}\\
\operatorname{gr}_\bullet(N\otimes_AP)
  \ar[r]_{\operatorname{gr}(u\otimes1)}
  &\operatorname{gr}_\bullet(M\otimes_AP)
}
\]
dans lequel $\varphi_M$ est bijectif, $\varphi_N$ surjectif~; en outre,
$\operatorname{gr}_0(u)$ est injectif par hypothèse, et comme
\[
\begin{split}
\operatorname{gr}_0(N)\otimes_A\operatorname{gr}_\bullet(P)
&=\operatorname{gr}_0(N)\otimes_k\operatorname{gr}_\bullet(P),\\
\operatorname{gr}_0(M)\otimes_A\operatorname{gr}_\bullet(P)
&=\operatorname{gr}_0(M)\otimes_k\operatorname{gr}_\bullet(P),
\end{split}
\]
$\operatorname{gr}_0(u)\otimes1$ est aussi injectif. On en conclut que
$\operatorname{gr}(u\otimes1)$ est injectif, ce qui achève de démontrer la
première assertion. La seconde se déduit du raisonnement précédent en faisant
$P=A$.
\end{proof}

\begin{proposition}[4.6.8]
\label{III.4.6.8-fr}
Soient $Z$ un préschéma localement noethérien, $X$, $Y$, deux $Z$-préschémas
tels que les morphismes structuraux $g:X\to Z$, $h:Y\to Z$ soient propres,
$Z'$ une partie fermée de $Z$, $X'=g^{-1}(Z')$, $Y'=h^{-1}(Z')$ ses images
réciproques,
\[
\widehat X=X_{/X'},\qquad
\widehat Y=Y_{/Y'},\qquad
\widehat Z=Z_{/Z'}
\]
les complétés formels de $X$, $Y$, $Z$ le long de ces parties fermées,
$f:X\to Y$ un $Z$-morphisme, $\widehat f:\widehat X\to\widehat Y$ son
prolongement aux complétés. Pour que $\widehat f$ soit un isomorphisme (resp.
une immersion fermée), il faut et il suffit qu'il existe un voisinage ouvert
$U$ de $Z'$ tel que le morphisme $g^{-1}(U)\to h^{-1}(U)$, restriction de
$f$, soit un isomorphisme (resp. une immersion fermée).
\end{proposition}

\begin{proof}
La suffisance de la condition est immédiate (I, 10.14.7). Pour en montrer la
nécessité, il suffit encore de prouver que pour tout $y\in Y'$, il existe un
voisinage ouvert $V_y$ de $y$ tel que la restriction
$f^{-1}(V_y)\to V_y$ de $f$ soit un isomorphisme (resp. une immersion
fermée), par le même raisonnement que dans
(\hyperref[III.4.6.7-fr]{4.6.7}). On est ainsi ramené au cas où $Y=Z$,
$Y=\operatorname{Spec}(A)$ étant affine noethérien. Par hypothèse
(I, 10.9.1 et 10.14.2) la fibre $f^{-1}(y)$ est réduite à un point pour
$y\in Y'$, donc comme $f$ est propre (II, 5.4.3), il existe un voisinage
ouvert $U$ de $y$ tel que la restriction $f^{-1}(U)\to U$ de $f$ soit un
morphisme fini (\hyperref[III.4.4.11-fr]{4.4.11}). On peut donc déjà supposer
que $f$ soit un morphisme fini, donc $X=\operatorname{Spec}(B)$, où $B$ est
une $A$-algèbre finie sur $A$. Si $Y'=V(\mathfrak J)$, on a alors
\[
\widehat Y=\operatorname{Spf}(\widehat A),\qquad
\widehat X=\operatorname{Spf}(\widehat B),
\]
$\widehat A$ étant le séparé complété de $A$ pour la topologie
$\mathfrak J$-préadique, $\widehat B$ le séparé complété de $B$ pour la
topologie $\mathfrak JB$-préadique, ou (ce qui revient au même), le séparé
complété du $A$-module $B$ pour la topologie $\mathfrak J$-préadique~; en
outre, $\widehat f$ est le morphisme de schémas formels affines correspondant
au prolongement continu $\widehat\varphi:\widehat A\to\widehat B$ de
l'homomorphisme canonique d'anneaux $\varphi:A\to B$, et l'hypothèse est que
$\widehat\varphi$ est surjectif (resp. bijectif) (I, 10.14.2). Or,
$\widehat\varphi$ est aussi le prolongement continu de $\varphi$ considéré
comme homomorphisme de $A$-modules~; on sait alors (I, 10.8.14) qu'il existe
un voisinage ouvert $U$ de $Y'$ tel que la restriction à $U$ de
l'homomorphisme $\widetilde\varphi:\widetilde A\to\widetilde B$ de
$\mathcal O_Y$-Modules soit surjectif (resp. bijectif), ce qui achève la
démonstration.
\end{proof}

\oldpage[III]{145}

\subsection{Un critère d'amplitude}

\begin{theorem}[4.7.1]
\label{III.4.7.1-fr}
Soient $Y$ un préschéma localement noethérien, $f:X\to Y$ un morphisme
propre, $\mathcal L$ un $\mathcal O_X$-Module inversible, $y$ un point de
$Y$, $X_y=X\otimes_Yk(y)=f^{-1}(y)$, $g$ la projection de $X_y$ dans $X$.
Si $\mathcal L_y=g^*(\mathcal L)=\mathcal L\otimes_{\mathcal O_Y}k(y)$ est
ample sur $X_y$, il existe un voisinage ouvert $U$ de $y$ dans $Y$ tel que
$\mathcal L|f^{-1}(U)$ soit ample pour la restriction de $f$ à $f^{-1}(U)$.
\end{theorem}

\begin{proof}
\noindent\textit{I)} Posons
\[
Y'=\operatorname{Spec}(\mathcal O_y),\qquad
X'=X\times_YY',\qquad
\mathcal L'=\mathcal L\otimes_{\mathcal O_Y}\mathcal O_y~;
\]
nous allons d'abord prouver que $\mathcal L'$ est ample pour
$f'=f_{(Y')}$. On a le diagramme commutatif
\[
\xymatrix{
X\ar[d]_f & X'\ar[l]\ar[d]_{f'} & X_y\ar[l]\ar[d]_{f_y}\\
Y & Y'\ar[l] & \operatorname{Spec}(k(y))\ar[l]
}
\]
Comme $f'$ est propre (II, 5.4.2, (iii)) et $\mathcal O_y$ noethérien, on
voit qu'on peut se borner au cas où
$Y=Y'=\operatorname{Spec}(\mathcal O_y)$, donc $X=X'$, supposer que
$\mathcal L_y$ est ample pour $f_y$ et prouver que $\mathcal L$ est ample
pour $f$ (II, 4.6.6). Nous allons appliquer le critère
(\hyperref[III.2.6.1-fr]{2.6.1, c)}) et montrer en fait que pour tout
$\mathcal O_X$-Module cohérent $\mathcal F$, il existe un entier $N$ tel que
\[
H^1(X,\mathcal F(n))=0\quad\text{pour tout }n\geq N,
\qquad
\mathcal F(n)=\mathcal F\otimes\mathcal L^{\otimes n}.
\]
Notons que $y$ est un point fermé de $Y$ correspondant à l'idéal maximal
$\mathfrak m$ de $\mathcal O_y$~; $X_y$ est donc un sous-préschéma fermé de
$X$ défini par l'Idéal cohérent
$\mathcal J=f^*(\widetilde{\mathfrak m})\mathcal O_X=
\mathfrak m\mathcal O_X$ de $\mathcal O_X$ (I, 4.4.5), et $g$ l'injection
canonique. Considérons alors la $k(y)$-algèbre graduée
\[
S=\operatorname{gr}(\mathcal O_y)
=\bigoplus_{j\geq0}\mathfrak m^j/\mathfrak m^{j+1},
\]
qui est de type fini puisque $\mathcal O_y$ est noethérien~; la
$\mathcal O_X$-Algèbre $\mathcal S=f^*(\widetilde S)$ est donc
quasi-cohérente et de type fini, et elle est évidemment annulée par
$\mathcal J$, donc si on pose $\mathcal S_y=g^*(\mathcal S)$,
$\mathcal S_y$ est une $\mathcal O_{X_y}$-Algèbre quasi-cohérente de type
fini, et $\mathcal S=g_*(\mathcal S_y)$. Posons d'autre part
\[
\mathcal M_j=\mathfrak m^j\mathcal F/\mathfrak m^{j+1}\mathcal F,
\qquad
\mathcal M=\bigoplus_{j\geq0}\mathcal M_j
=\operatorname{gr}(\mathcal F)~;
\]
comme $\mathcal F$ est cohérent, $\mathcal M$ est un $\mathcal S$-Module
quasi-cohérent de type fini (0, 10.1.1) qui est aussi annulé par
$\mathcal J$, de sorte que si l'on pose
\[
\mathcal M'_j=g^*(\mathcal M_j),\qquad
\mathcal M'=g^*(\mathcal M)=\bigoplus_{j\geq0}\mathcal M'_j,
\]
$\mathcal M'$ est un $\mathcal S_y$-Module gradué quasi-cohérent de type
fini tel que $\mathcal M=g_*(\mathcal M')$. En outre, si on pose
$\mathcal M'_j(n)=\mathcal M'_j\otimes\mathcal L_y^{\otimes n}$, on a
$\mathcal M'_j(n)=g^*(\mathcal M_j(n))$. Cela étant, $f_y$ est propre
(II, 5.4.2, (iii)) et $\mathcal L_y$ est ample, donc $f_y$ est projectif
(II, 5.5.4 et 4.6.11), et on peut appliquer à
$\operatorname{Spec}(k(y))$, $f_y$, $\mathcal S_y$, $\mathcal L_y$ et
$\mathcal M'$ le théorème (\hyperref[III.2.4.1-fr]{2.4.1, (ii)})~: il
existe un entier $N$ tel que pour $n\geq N$, on ait
$H^q(X_y,\mathcal M'_j(n))=0$ pour tout $q>0$ et tout $j$~; par suite, on a
aussi $H^q(X,\mathcal M_j(n))=0$ pour tout $q>0$ et tout $j$
(G, II, 4.9.1). Posons alors
\[
\mathcal F(n)_j=\mathcal F(n)/\mathfrak m^{j+1}\mathcal F(n),
\]
de sorte que $\mathcal F(n)_{j-1}=\mathcal F(n)_j/\mathcal M_j(n)$ pour
$j\geq1$ et $\mathcal F(n)_0=\mathcal M_0(n)$. On a
$H^1(X,\mathcal F(n)_0)=0$, et, par la suite exacte de cohomologie,
$H^1(X,\mathcal F(n)_j)=H^1(X,\mathcal F(n)_{j-1})$ pour tout $j\geq1$,
donc $H^1(X,\mathcal F(n)_j)=0$ pour tout $j\geq0$. On conclut donc de
(\hyperref[III.4.2.1-fr]{4.2.1}) que $H^1(X,\mathcal F(n))=0$, ce qui
achève de prouver notre assertion.

\oldpage[III]{146}

\noindent\textit{II)} Revenons aux notations du début de la démonstration,
et remarquons qu'on peut toujours supposer $Y=\operatorname{Spec}(A)$
affine~; comme $f'$ est de type fini et $\mathcal L'$ ample pour $f'$, il
existe un entier $m>0$ tel que $\mathcal L'^{\otimes m}$ soit très ample pour
$f'$ (II, 4.6.11)~; remplaçant au besoin $\mathcal L$ par
$\mathcal L^{\otimes m}$, on peut se borner à considérer le cas où
$\mathcal L'$ est très ample pour $f'$, et à prouver que
$\mathcal L|f^{-1}(U)$ est alors très ample pour $f$. Comme $f'$ est propre,
il existe alors une $Y'$-immersion fermée
$j:X'\to P=\mathbf P_{Y'}^r$ pour un entier $r>0$ convenable, telle que
$\mathcal L'$ soit isomorphe à $j^*(\mathcal O_P(1))$ (II, 5.5.4, (ii))~;
cette immersion correspond canoniquement à un $\mathcal O_{X'}$-homomorphisme
surjectif
\[
u:\mathcal O_{X'}^{r+1}\longrightarrow\mathcal L'
\]
(II, 4.2.3). Ce dernier correspond (0\textsubscript{I}, 5.1.1) à la donnée
de $r+1$ sections $s'_i$ ($0\leq i\leq r$) de $\mathcal L'$ au-dessus de
$X'$ qui engendrent ce $\mathcal O_{X'}$-Module. Ces sections sont aussi par
définition des sections de $f'_*(\mathcal L')$ au-dessus de $Y'$~; on a
\[
f'_*(\mathcal L')=f_*(\mathcal L)\otimes_{\mathcal O_Y}\mathcal O_{Y'}
\]
(0\textsubscript{I}, 5.4.10), $Y$ est affine et $\mathcal O_y$ est l'anneau
local en l'idéal premier $\mathfrak j_y$ de $A$, donc on a
$s'_i=s''_i/t_i$, où les $s''_i$ sont des sections de $f_*(\mathcal L)$
au-dessus de $Y$ et les $t_i$ des éléments de $A$ n'appartenant pas à
$\mathfrak j_y$~; on en conclut qu'il existe un voisinage ouvert affine $V$
de $y$ dans $Y$ et des sections $s_i$ de $f_*(\mathcal L)|V$ telles que
$s'_i=s_i/1$ (on rappelle que l'espace $Y'$ est contenu dans $V$, cf. I,
2.4.2). Les $s_i$ sont alors des sections de $\mathcal L$ au-dessus de
$f^{-1}(V)$, définissant donc un homomorphisme
\[
v:\bigl(\mathcal O_X|f^{-1}(V)\bigr)^{r+1}
\longrightarrow\mathcal L|f^{-1}(V)
\]
qui, par hypothèse, est surjectif en tous les points de $f^{-1}(y)$~; comme
$\operatorname{Coker}(v)$ est cohérent (0\textsubscript{I}, 5.3.4), son
support est fermé (0\textsubscript{I}, 5.2.2) et par suite il existe un
voisinage ouvert $W\subset f^{-1}(V)$ de $f^{-1}(y)$ tel que la restriction
de $v$ à $W$ soit un homomorphisme surjectif. Puisque le morphisme $f$ est
fermé, on peut supposer que $W$ est de la forme $f^{-1}(U)$, où $U$ est un
voisinage ouvert de $y$, et la conclusion résulte alors de (II, 4.2.3).
\end{proof}

\subsection{Morphismes finis de préschémas formels}

\begin{proposition}[4.8.1]
\label{III.4.8.1-fr}
Soient $\mathfrak Y$ un préschéma formel localement noethérien, $\mathcal K$ un
Idéal de définition de $\mathfrak Y$, $f:\mathfrak X\to\mathfrak Y$ un
morphisme de préschémas formels. Les conditions suivantes sont équivalentes~:
\begin{enumerate}
\item[a)] $\mathfrak X$ est localement noethérien, $f$ est un morphisme adique
  (I, 10.12.1) et si l'on pose
  $\mathcal J=f^*(\mathcal K)\mathcal O_{\mathfrak X}$, le morphisme
  \[
  f_0:(\mathfrak X,\mathcal O_{\mathfrak X}/\mathcal J)
  \longrightarrow
  (\mathfrak Y,\mathcal O_{\mathfrak Y}/\mathcal K)
  \]
  déduit de $f$ est fini.
\item[b)] $\mathfrak X$ est localement noethérien et est limite inductive d'un
  $(Y_n)$-système inductif adique $(X_n)$ tel que le morphisme
  $X_0\to Y_0$ soit fini.
\item[c)] Tout point de $\mathfrak Y$ possède un voisinage ouvert formel affine
  noethérien $V$ tel que $f^{-1}(V)$ soit un ouvert formel affine et que
  $\Gamma(f^{-1}(V),\mathcal O_{\mathfrak X})$ soit un
  $\Gamma(V,\mathcal O_{\mathfrak Y})$-module de type fini.
\end{enumerate}
\end{proposition}

\begin{proof}
Il est immédiat que a) entraîne b) en vertu de (I, 10.12.3). Pour voir que b)
entraîne c), on peut supposer que
$\mathfrak Y=\operatorname{Spf}(B)$, où $B$ est adique noethérien et
$\mathcal K=\mathfrak R^\Delta$, où $\mathfrak R$ est un idéal de définition
de $B$. Par hypothèse, $X_0$ est un schéma affine dont l'anneau $A_0$ est un
$B/\mathfrak R$-module de type fini (II, 6.1.3). En vertu de (I, 5.1.9), chacun
des $X_n$ est un schéma affine, et si $A_n$ est son anneau, l'hypothèse b)
entraîne que pour $m\leq n$, $A_m$ est isomorphe à
$A_n/\mathfrak R^{m+1}A_n$. On en déduit que $\mathfrak X$ est isomorphe à
$\operatorname{Spf}(A)$, où $A=\varprojlim A_n$~; on conclut en vertu de
(0\textsubscript{I}, 7.2.9). Enfin, pour prouver que c)

\oldpage[III]{147}

entraîne a), on peut se borner encore au cas où
$\mathfrak Y=\operatorname{Spf}(B)$, $\mathfrak X=\operatorname{Spf}(A)$,
$A$ étant une $B$-algèbre finie~; comme $A/\mathfrak RA$ est alors une
$B/\mathfrak R$-algèbre finie, il résulte de (I, 10.10.9) que les conditions
de a) sont satisfaites.
\end{proof}

\begin{definition}[4.8.2]
\label{III.4.8.2-fr}
Lorsque les propriétés équivalentes a), b), c) de
(\hyperref[III.4.8.1-fr]{4.8.1}) sont vérifiées, on dit que le morphisme $f$
est \emph{fini}, ou que $\mathfrak X$ est un \emph{$\mathfrak Y$-préschéma
formel fini}, ou un préschéma formel \emph{fini au-dessus de $\mathfrak Y$}.
\end{definition}

\begin{proposition}[4.8.3]
\label{III.4.8.3-fr}
\begin{enumerate}
\item[(i)] Une immersion fermée de préschémas formels localement noethériens
  est un morphisme fini.
\item[(ii)] Le composé de deux morphismes finis de préschémas formels
  localement noethériens est un morphisme fini.
\item[(iii)] Soient $\mathfrak X$, $\mathfrak Y$, $\mathfrak S$ trois
  préschémas formels localement noethériens, $f:\mathfrak X\to\mathfrak S$ un
  morphisme fini, $g:\mathfrak Y\to\mathfrak S$ un morphisme~; alors le
  morphisme
  $\mathfrak X\times_{\mathfrak S}\mathfrak Y\to\mathfrak Y$ est fini.
\item[(iv)] Soient $\mathfrak S$ un préschéma formel localement noethérien,
  $\mathfrak X'$, $\mathfrak Y'$ deux préschémas formels localement
  noethériens tels que $\mathfrak X'\times_{\mathfrak S}\mathfrak Y'$ soit
  localement noethérien. Si $\mathfrak X$, $\mathfrak Y$ sont des
  $\mathfrak S$-préschémas formels localement noethériens,
  $f:\mathfrak X\to\mathfrak X'$, $g:\mathfrak Y\to\mathfrak Y'$ deux
  $\mathfrak S$-morphismes finis, alors $f\times_{\mathfrak S}g$ est un
  morphisme fini.
\item[(v)] Soient $f:\mathfrak X\to\mathfrak Y$,
  $g:\mathfrak Y\to\mathfrak Z$ deux morphismes de préschémas formels
  localement noethériens tels que $g$ soit de type fini et séparé~; alors, si
  $g\circ f$ est un morphisme fini, $f$ est un morphisme fini.
\end{enumerate}
\end{proposition}

\begin{proof}
(i) est trivial, et les autres assertions se ramènent aussitôt aux
propositions correspondantes pour les morphismes de préschémas usuels
(II, 6.1.5) à l'aide du critère a) de
(\hyperref[III.4.8.1-fr]{4.8.1})~; nous laissons les détails au lecteur, sur le
modèle de (I, 10.13.5).
\end{proof}

\begin{corollary}[4.8.4]
\label{III.4.8.4-fr}
Sous les hypothèses de (I, 10.9.9), si $f$ est un morphisme fini, il en est de
même de son prolongement $\widehat f$ aux complétés.
\end{corollary}

\begin{corollary}[4.8.5]
\label{III.4.8.5-fr}
Si $\mathfrak X$ est un préschéma formel fini au-dessus de $\mathfrak Y$,
$f:\mathfrak X\to\mathfrak Y$ le morphisme structural, alors, pour tout ouvert
$U\subset\mathfrak Y$, $f^{-1}(U)$ est fini au-dessus de $U$.
\end{corollary}

\begin{proposition}[4.8.6]
\label{III.4.8.6-fr}
Si $f:\mathfrak X\to\mathfrak Y$ est un morphisme fini de préschémas formels
localement noethériens, $f_*(\mathcal O_{\mathfrak X})$ est une
$\mathcal O_{\mathfrak Y}$-Algèbre cohérente.
\end{proposition}

\begin{proof}
On peut considérer $f$ comme limite inductive d'un système inductif $(f_n)$ de
morphismes $f_n:X_n\to Y_n$~; montrons que les $f_n$ sont des morphismes finis
et que $f_*(\mathcal O_{\mathfrak X})$ est isomorphe à la limite projective des
$(f_n)_*(\mathcal O_{X_n})$, ce qui établira notre assertion (I, 10.10.5). Il
suffit de se borner au cas où $\mathfrak Y=\operatorname{Spf}(B)$,
$\mathfrak X=\operatorname{Spf}(A)$, et de remarquer que si $\mathfrak R$ est
un idéal de définition de $B$ et $A$ un $B$-module de type fini,
$A/\mathfrak R^{n+1}A$ est un module de type fini sur
$B/\mathfrak R^{n+1}B$, et que $A$ est limite projective des
$A/\mathfrak R^{n+1}A$.
\end{proof}

Réciproquement~:

\begin{proposition}[4.8.7]
\label{III.4.8.7-fr}
Soient $\mathfrak Y$ un préschéma formel localement noethérien, $\mathcal A$
une $\mathcal O_{\mathfrak Y}$-Algèbre cohérente. Il existe un préschéma
formel $\mathfrak X$ fini au-dessus de $\mathfrak Y$, défini à un
$\mathfrak Y$-isomorphisme unique près, et tel que
$f_*(\mathcal O_{\mathfrak X})=\mathcal A$, $f:\mathfrak X\to\mathfrak Y$
étant le morphisme structural.
\end{proposition}

\begin{proof}
Soit $\mathcal K$ un Idéal de définition de $\mathfrak Y$, et posons
\[
Y_n=(\mathfrak Y,\mathcal O_{\mathfrak Y}/\mathcal K^{n+1}),\qquad
\mathcal A_n=\mathcal A/\mathcal K^{n+1}\mathcal A~;
\]
il est clair que $\mathcal A_n$ est une $\mathcal O_{Y_n}$-Algèbre finie et
définit donc un

\oldpage[III]{148}

$Y_n$-préschéma fini $X_n=\operatorname{Spec}(\mathcal A_n)$ (II, 6.1.3)~;
pour $m\leq n$, l'homomorphisme canonique surjectif
$h_{mn}:\mathcal A_n\to\mathcal A_m$ définit un morphisme
$u_{mn}:X_m\to X_n$ tel que le diagramme
\[
\xymatrix{
X_n\ar[d]_{f_n} & X_m\ar[l]_{u_{mn}}\ar[d]^{f_m}\\
Y_n & Y_m\ar[l]
}
\]
($f_n$ étant le morphisme structural) soit commutatif et identifie $X_m$ au
produit $X_n\times_{Y_n}Y_m$, comme on le voit aussitôt (II, 1.4.6). Le
préschéma formel $\mathfrak X$, limite inductive du système inductif $(X_n)$
est alors localement noethérien et tel que le morphisme structural
$f:\mathfrak X\to\mathfrak Y$, limite inductive du système $(f_n)$, soit fini
(\hyperref[III.4.8.1-fr]{4.8.1} et II, 10.12.3.1)~; on a vu en outre dans la
démonstration de (\hyperref[III.4.8.6-fr]{4.8.6}) que
$f_*(\mathcal O_{\mathfrak X})$ est limite projective des $\mathcal A_n$, donc
égale à $\mathcal A$ (I, 10.10.6). Quant à l'assertion d'unicité, elle est
conséquence du résultat plus général suivant~:
\end{proof}

\begin{proposition}[4.8.8]
\label{III.4.8.8-fr}
Soient $\mathfrak Y$ un préschéma formel localement noethérien,
$\mathfrak X$, $\mathfrak X'$ deux $\mathfrak Y$-préschémas formels finis
au-dessus de $\mathfrak Y$, $f:\mathfrak X\to\mathfrak Y$,
$f':\mathfrak X'\to\mathfrak Y$ les morphismes structuraux. Il existe une
bijection canonique de
\[
\operatorname{Hom}_{\mathfrak Y}(\mathfrak X,\mathfrak X')
\quad\hbox{sur}\quad
\operatorname{Hom}_{\mathcal O_{\mathfrak Y}}
\bigl(f'_*(\mathcal O_{\mathfrak X'}),f_*(\mathcal O_{\mathfrak X})\bigr)
\footnotemark.
\]
\footnotetext{La dernière expression désigne l'ensemble des homomorphismes de
$\mathcal O_{\mathfrak Y}$-Algèbres
$f'_*(\mathcal O_{\mathfrak X'})\to f_*(\mathcal O_{\mathfrak X})$.}
\end{proposition}

\begin{proof}
La définition de cette application $h\to\mathcal A(h)$ est la même que dans
(II, 1.1.2), et pour voir qu'elle est bijective, on est aussitôt ramené au cas
où $\mathfrak Y=\operatorname{Spf}(B)$ est un schéma formel affine noethérien.
Mais alors $\mathfrak X=\operatorname{Spf}(A)$,
$\mathfrak X'=\operatorname{Spf}(A')$, où $A$ et $A'$ sont deux $B$-algèbres
finies et $f_*(\mathcal O_{\mathfrak X})=A^\Delta$,
$f'_*(\mathcal O_{\mathfrak X'})=A'^\Delta$. La conclusion résulte alors de
la correspondance biunivoque, d'une part entre les $\mathfrak Y$-morphismes
$\mathfrak X\to\mathfrak X'$ et les $B$-homomorphismes (nécessairement
continus) $A'\to A$ qui sont des homomorphismes d'algèbres (I, 10.2.2), et
d'autre part entre les homomorphismes de $B$-modules $A'\to A$ et les
homomorphismes de $\mathcal O_{\mathfrak Y}$-Modules
$A'^\Delta\to A^\Delta$ (I, 10.10.2.3).
\end{proof}

\begin{corollary}[4.8.9]
\label{III.4.8.9-fr}
Dans la correspondance biunivoque canonique définie dans
(\hyperref[III.4.8.8-fr]{4.8.8}), les immersions fermées
$\mathfrak X\to\mathfrak X'$ correspondent aux homomorphismes surjectifs de
$\mathcal O_{\mathfrak Y}$-Algèbres
$f'_*(\mathcal O_{\mathfrak X'})\to f_*(\mathcal O_{\mathfrak X})$.
\end{corollary}

\begin{proof}
La question étant encore locale sur $\mathfrak Y$, on est ramené à la
définition des immersions fermées de préschémas formels localement noethériens
(I, 10.14.2).
\end{proof}

\begin{corollary}[4.8.10]
\label{III.4.8.10-fr}
Les notations et hypothèses étant celles de
(\hyperref[III.4.8.1-fr]{4.8.1}), pour qu'un morphisme adique $f$ soit une
immersion fermée, il faut et il suffit que $f_0$ soit une immersion fermée (de
préschémas usuels).
\end{corollary}

\begin{proof}
Cela résulte aussitôt de (\hyperref[III.4.8.9-fr]{4.8.9}) et de la condition de
surjectivité pour un homomorphisme de $\mathcal O_{\mathfrak Y}$-Modules
cohérents (I, 10.11.5).
\end{proof}

\begin{proposition}[4.8.11]
\label{III.4.8.11-fr}
Pour qu'un morphisme $f:\mathfrak X\to\mathfrak Y$ de préschémas formels localement

\oldpage[III]{149}

noethériens soit fini, il faut et il suffit qu'il soit propre et ait ses fibres
$f^{-1}(y)$ finies (pour tout $y\in\mathfrak Y$).
\end{proposition}

\begin{proof}
Grâce aux définitions (3.4.1 et 4.8.2), on est aussitôt ramené à la même proposition
pour $f_0$ (notations de (\hyperref[III.4.8.1-fr]{4.8.1})), ce qui n'est autre que
(4.4.2).
\end{proof}

\section{Un théorème d'existence de faisceaux algébriques cohérents}
\label{section:III.5-fr}

\subsection{Énoncé du théorème}
\label{subsection:III.5.1-fr}

\begin{env}[5.1.1]
\label{III.5.1.1-fr}
Soient $A$ un anneau adique noethérien, $\mathfrak J$ un idéal de définition de
$A$, de sorte que $A$ est séparé et complet pour la topologie $\mathfrak J$-adique.
Si $Y=\operatorname{Spec}(A)$, le schéma formel affine $\operatorname{Spf}(A)$
s'identifie au complété $\widehat Y$ de $Y$ le long de la partie fermée
$Y'=\mathrm V(\mathfrak J)$ (I, 10.10.1). Soient $X$ un $Y$-préschéma (usuel) de
type fini, $f:X\to Y$ le morphisme structural~; nous désignerons par $\widehat X$
le complété de $X$ le long de la partie fermée $X'=f^{-1}(Y')$, ou encore le
$\widehat Y$-préschéma formel $X\times_Y\widehat Y$~; par
$\widehat f:\widehat X\to\widehat Y$ le prolongement de $f$ aux complétés~; enfin,
pour tout $\mathcal O_X$-Module cohérent $\mathcal F$, nous noterons
$\widehat{\mathcal F}$ son complété, $\mathcal F_{/X'}$, qui est un
$\mathcal O_{\widehat X}$-Module cohérent.
\end{env}

\begin{proposition}[5.1.2]
\label{III.5.1.2-fr}
Les hypothèses et notations étant celles de
(\hyperref[III.5.1.1-fr]{5.1.1}), soit $\mathcal F$ un $\mathcal O_X$-Module
cohérent dont le support est propre sur $Y$ (II, 5.4.10). Les homomorphismes
canoniques (4.1.4)
\[
\rho_i:\mathrm H^i(X,\mathcal F)\longrightarrow
\mathrm H^i(\widehat X,\widehat{\mathcal F})
\]
sont alors des isomorphismes.
\end{proposition}

\begin{proof}
Comme $\mathrm H^i(X,\mathcal F)$ est un $A$-module de type fini (3.2.4), donc
identique à son séparé complété pour la topologie $\mathfrak J$-préadique
($0_{\mathrm I}$, 7.3.6), la proposition n'est qu'un cas particulier de (4.1.10).
\end{proof}

Rappelons que les isomorphismes canoniques $\rho_i$ commutent aux cobords pour
toute suite exacte
$0\to\mathcal F'\to\mathcal F\to\mathcal F''\to0$ de $\mathcal O_X$-Modules
cohérents (($0$, 12.1.6) et (I, 10.8.9)).

\begin{corollary}[5.1.3]
\label{III.5.1.3-fr}
Soient $\mathcal F$, $\mathcal G$ deux $\mathcal O_X$-Modules cohérents tels que
l'intersection de leurs supports soit propre sur $Y$. Alors l'homomorphisme canonique
\begin{equation}
\operatorname{Hom}_{\mathcal O_X}(\mathcal F,\mathcal G)
\longrightarrow
\operatorname{Hom}_{\mathcal O_{\widehat X}}
(\widehat{\mathcal F},\widehat{\mathcal G})
\tag{5.1.3.1}
\label{III.5.1.3.1.eq-fr}
\end{equation}
qui, à tout homomorphisme $u:\mathcal F\to\mathcal G$, fait correspondre son
complété $\widehat u:\widehat{\mathcal F}\to\widehat{\mathcal G}$, est un
isomorphisme. De plus, lorsque le morphisme $f$ est fermé, pour que $\widehat u$ soit
injectif (resp. surjectif), il faut et il suffit que $u$ le soit.
\end{corollary}

\begin{proof}
La première assertion est un cas particulier de (4.5.3), dû encore au fait que le
premier membre de (\ref{III.5.1.3.1.eq-fr}) est un $A$-module de type fini, donc
identique à son séparé complété. Pour démontrer la seconde, notons en vertu de
(I, 10.8.14) que $\widehat u$ est injectif (resp. surjectif) si et seulement s'il existe
un voisinage de $X'$ dans lequel $u$ soit injectif (resp. surjectif). La conclusion
résulte donc du lemme suivant~:
\end{proof}

\oldpage[III]{150}

\begin{lemma}[5.1.3.1]
\label{III.5.1.3.1-fr}
Sous les hypothèses de (\hyperref[III.5.1.1-fr]{5.1.1}), si on suppose en outre le
morphisme $f$ fermé, tout voisinage de $X'$ dans $X$ est identique à $X$.
\end{lemma}

\begin{proof}
Tout d'abord, on peut se ramener au cas où $f(X)=Y$. En effet, par hypothèse,
$f(X)$ est une partie fermée $Z$ de $Y$~; on peut en outre remplacer $f$ par
$f_{\mathrm{red}}$ (I, 6.3.4), et supposer par suite $X$ et $Y$ réduits~; on peut alors
remplacer $Y$ par le sous-préschéma fermé réduit de $Y$ ayant $Z$ pour espace
sous-jacent (I, 5.2.2), car tout idéal de $A$ est fermé, et tout anneau quotient de
$A$ est donc adique et noethérien. On a alors $f(X')=Y'$~; si $V$ est un voisinage
ouvert de $X'$ dans $X$, $f(X-V)$ est fermé dans $Y$ par hypothèse, et ne rencontre
pas $Y'$~; mais cela est impossible à moins que $X-V$ ne soit vide, puisque
$\mathfrak J$ est contenu dans le radical de $A$ (I, 1.1.15 et $0_{\mathrm I}$,
7.1.10), d'où la conclusion.
\end{proof}

Lorsqu'on se borne aux $\mathcal O_X$-Modules cohérents dont le support est propre
sur $Y$, (\hyperref[III.5.1.3-fr]{5.1.3}) peut s'énoncer, dans le langage des
catégories, en disant que le foncteur
$\mathcal F\mapsto\widehat{\mathcal F}$ est pleinement fidèle de la catégorie des
$\mathcal O_X$-Modules du type précédent, dans la catégorie des
$\mathcal O_{\widehat X}$-Modules cohérents, et établit par suite une équivalence de
la première de ces catégories avec une sous-catégorie pleine de la seconde
($0$, 8.1.6). Le théorème d'existence va prouver que lorsque $X$ est propre sur $Y$,
cette sous-catégorie est en fait la catégorie de tous les
$\mathcal O_{\widehat X}$-Modules cohérents. De façon précise~:

\begin{theorem}[5.1.4]
\label{III.5.1.4-fr}
Soient $A$ un anneau adique noethérien, $Y=\operatorname{Spec}(A)$,
$\mathfrak J$ un idéal de définition de $A$, $Y'=\mathrm V(\mathfrak J)$,
$f:X\to Y$ un morphisme séparé de type fini, $X'=f^{-1}(Y')$. Soient
$\widehat Y=Y_{/Y'}=\operatorname{Spf}(A)$, $\widehat X=X_{/X'}$ les complétés de
$Y$ et $X$ le long de $Y'$ et $X'$, $\widehat f:\widehat X\to\widehat Y$ le
prolongement de $f$ aux complétés~; alors, le foncteur
$\mathcal F\mapsto\mathcal F_{/X'}=\widehat{\mathcal F}$ est une équivalence de la
catégorie des $\mathcal O_X$-Modules cohérents de support propre sur
$\operatorname{Spec}(A)$, avec la catégorie des $\mathcal O_{\widehat X}$-Modules
cohérents de support propre sur $\operatorname{Spf}(A)$.
\end{theorem}

En d'autres termes, compte tenu de (\hyperref[III.5.1.3-fr]{5.1.3})~:

\begin{corollary}[5.1.5]
\label{III.5.1.5-fr}
Pour qu'un $\mathcal O_{\widehat X}$-Module soit isomorphe au complété d'un
$\mathcal O_X$-Module cohérent et de support propre sur $\operatorname{Spec}(A)$,
il faut et il suffit qu'il soit cohérent et de support propre sur
$\operatorname{Spf}(A)$.
\end{corollary}

Le cas le plus important est le

\begin{corollary}[5.1.6]
\label{III.5.1.6-fr}
Supposons $X$ propre sur $Y=\operatorname{Spec}(A)$. Alors le foncteur
$\mathcal F\mapsto\widehat{\mathcal F}$ est une équivalence de la catégorie des
$\mathcal O_X$-Modules cohérents et de la catégorie des
$\mathcal O_{\widehat X}$-Modules cohérents.
\end{corollary}

\begin{env}[5.1.7]
\label{III.5.1.7-fr}
\emph{Scholie.}
Si on tient compte de la caractérisation des faisceaux cohérents sur les préschémas
formels (I, 10.11.3), on voit que sous les conditions de
(\hyperref[III.5.1.1-fr]{5.1.1}), la donnée d'un $\mathcal O_X$-Module cohérent de
support propre sur $\operatorname{Spec}(A)$ équivaut (en posant
$Y_n=\operatorname{Spec}(A/\mathfrak J^{n+1})$ et $X_n=X\times_Y Y_n$) à la donnée
d'un système projectif de $\mathcal O_{X_n}$-Modules cohérents $(\mathcal F_n)$ tel
que pour $m\leq n$ on ait
\[
\mathcal F_m=\mathcal F_n\otimes_{\mathcal O_{Y_n}}\mathcal O_{Y_m}
\quad\text{(ou encore }\mathcal F_m=\mathcal F_n/\mathfrak J^{m+1}\mathcal F_n\text{)}
\]
et que le support de $\mathcal F_0$ soit une partie de $X_0$ propre sur $Y_0$. Au
moyen de (I, 10.11.4), on interprète de même les homomorphismes de
$\mathcal O_X$-Modules cohérents comme des homomorphismes de systèmes projectifs de
$\mathcal O_{X_n}$-Modules cohérents.

\oldpage[III]{151}

Dans tous les cas d'application connus, $A$ est en fait un anneau adique local
noethérien, donc les $Y_n$ sont des spectres d'anneaux artiniens locaux, et les
résultats de ce paragraphe et des précédents réduisent dans une large mesure la
géométrie algébrique sur un anneau adique local noethérien, à la géométrie algébrique
sur des anneaux locaux artiniens.
\end{env}

\begin{corollary}[5.1.8]
\label{III.5.1.8-fr}
Sous les conditions de (\hyperref[III.5.1.4-fr]{5.1.4}), l'application
$Z\mapsto\widehat Z=Z_{/(Z\cap X')}$ est une bijection de l'ensemble des
sous-préschémas fermés $Z$ de $X$, propres sur $Y$, sur l'ensemble des
sous-préschémas formels fermés de $\widehat X$, propres sur $\widehat Y$.
\end{corollary}

\begin{proof}
En effet, un sous-préschéma formel fermé de $\widehat X$ est de la forme
$(T,(\mathcal O_{\widehat X}/\mathcal A)|T)$, où $\mathcal A$ est un Idéal
cohérent de $\mathcal O_{\widehat X}$ (I, 10.14.2)~; si $T$ est propre sur
$\widehat Y$, il résulte de (\hyperref[III.5.1.4-fr]{5.1.4}) que
$\mathcal O_{\widehat X}/\mathcal A$ est isomorphe à un
$\mathcal O_{\widehat X}$-Module de la forme $\widehat{\mathcal F}$, où
$\mathcal F$ est un $\mathcal O_X$-Module cohérent de support propre sur $Y$~; en
outre, il résulte de (\hyperref[III.5.1.3-fr]{5.1.3}) que l'homomorphisme canonique
$\mathcal O_{\widehat X}\to\mathcal O_{\widehat X}/\mathcal A$ est de la forme
$\widehat u$, où $u:\mathcal O_X\to\mathcal F$ est un homomorphisme surjectif de
$\mathcal O_X$-Modules. Donc $\mathcal F$ est de la forme
$\mathcal O_X/\mathcal N$, où $\mathcal N$ est un Idéal cohérent de
$\mathcal O_X$, et $\mathcal A=\widehat{\mathcal N}$ (I, 10.8.8), d'où la
conclusion (I, 10.14.7).
\end{proof}

\subsection{Démonstration du théorème d'existence : cas projectif et quasi-projectif}
\label{subsection:III.5.2-fr}

\begin{env}[5.2.1]
\label{III.5.2.1-fr}
Sous les conditions de (\hyperref[III.5.1.4-fr]{5.1.4}), nous dirons provisoirement
qu'un $\mathcal O_{\widehat X}$-Module cohérent est algébrisable s'il est isomorphe
à un complété $\widehat{\mathcal F}$ d'un $\mathcal O_X$-Module cohérent
$\mathcal F$ de support propre sur $Y$.
\end{env}

\begin{lemma}[5.2.2]
\label{III.5.2.2-fr}
Soient $\mathcal F'$, $\mathcal G'$ deux $\mathcal O_{\widehat X}$-Modules
algébrisables. Pour tout homomorphisme $u:\mathcal F'\to\mathcal G'$,
$\operatorname{Ker}(u)$, $\operatorname{Im}(u)$ et $\operatorname{Coker}(u)$ sont
algébrisables.
\end{lemma}

\begin{proof}
En effet, si $\mathcal F'=\widehat{\mathcal F}$,
$\mathcal G'=\widehat{\mathcal G}$, où $\mathcal F$ et $\mathcal G$ sont des
$\mathcal O_X$-Modules cohérents de supports propres sur $Y$, on a
$u=\widehat v$, où $v:\mathcal F\to\mathcal G$ est un homomorphisme
(\hyperref[III.5.1.3-fr]{5.1.3}). En vertu de l'exactitude du foncteur
$\mathcal F\mapsto\widehat{\mathcal F}$,
$\operatorname{Ker}(\widehat v)$ est isomorphe à
$\widehat{\operatorname{Ker}(v)}$ et comme le support de
$\operatorname{Ker}(v)$ est contenu dans celui de $\mathcal F$, on voit que
$\operatorname{Ker}(u)$ est algébrisable~; démonstration analogue pour
$\operatorname{Im}(u)$ et $\operatorname{Coker}(u)$.
\end{proof}

\begin{proposition}[5.2.3]
\label{III.5.2.3-fr}
Soient $A$ un anneau adique noethérien, $\mathfrak J$ un idéal de définition de
$A$, $\mathfrak Y=\operatorname{Spf}(A)$,
$f:\mathfrak X\to\mathfrak Y$ un morphisme propre de préschémas formels. On pose
$Y_k=\operatorname{Spec}(A/\mathfrak J^{k+1})$,
$X_k=\mathfrak X\times_{\mathfrak Y}Y_k$, et pour tout
$\mathcal O_{\mathfrak X}$-Module $\mathcal F$,
\[
\mathcal F_k=\mathcal F\otimes_{\mathcal O_{\mathfrak X}}\mathcal O_{X_k}
=\mathcal F/\mathfrak J^{k+1}\mathcal F.
\]
Soit $\mathcal L$ un $\mathcal O_{\mathfrak X}$-Module inversible, et supposons que
$\mathcal L_0=\mathcal L/\mathfrak J\mathcal L$ soit un
$\mathcal O_{X_0}$-Module ample~; pour tout $\mathcal O_{\mathfrak X}$-Module
$\mathcal F$ et tout entier $n$, posons
$\mathcal F(n)=\mathcal F\otimes\mathcal L^{\otimes n}$. Alors, pour tout
$\mathcal O_{\mathfrak X}$-Module cohérent $\mathcal F$, il existe un entier $n_0$
tel que, pour tout $n\geq n_0$, les propriétés suivantes aient lieu~:
\begin{enumerate}
\item[(i)] L'homomorphisme canonique
$\mathrm H^0(\mathfrak X,\mathcal F(n))\to
\mathrm H^0(\mathfrak X,\mathcal F_k(n))$ est surjectif pour tout $k\geq0$.
\item[(ii)] On a $\mathrm H^q(\mathfrak X,\mathcal F(n))=0$ pour tout $q>0$.
\end{enumerate}
\end{proposition}

\begin{proof}
On sait que les espaces sous-jacents à $\mathfrak X$ et à $X_0$ sont les mêmes~;
les faisceaux
\[
\mathcal M_k=\mathfrak J^k\mathcal F/\mathfrak J^{k+1}\mathcal F
\]
étant annulés par $\mathfrak J$, peuvent être considérés comme des
$\mathcal O_{X_0}$-Modules cohérents ($0_{\mathrm I}$, 5.3.10)~; en outre, si on
pose $\mathcal M_k(n)=\mathcal M_k\otimes_{\mathcal O_{X_0}}
\mathcal L_0^{\otimes n}$, on voit aussitôt que
\[
\mathcal M_k(n)=\mathfrak J^k\mathcal F(n)/
\mathfrak J^{k+1}\mathcal F(n).
\]
Notons que, puisque $\mathcal L_0$ est ample pour $f_0:X_0\to Y_0$, et

\oldpage[III]{152}

que $f_0$ est propre, on en conclut que $f_0$ est projectif (II, 5.5.4). Soit
\[
S=\bigoplus_{k\geq0}\mathfrak J^k/\mathfrak J^{k+1}
\]
la $A$-algèbre graduée associée à la filtration $\mathfrak J$-adique de $A$, qui est
de type fini puisque $A$ est noethérien~; si on pose
$\mathcal S'=f_0^*(\widetilde S)$,
$\mathcal M=\bigoplus_{k\geq0}\mathcal M_k$, $\mathcal S'$ est une
$\mathcal O_{X_0}$-Algèbre quasi-cohérente de type fini, et $\mathcal M$ un
$\mathcal S'$-Module gradué quasi-cohérent et de type fini (puisque $\mathcal F_0$
est cohérent et engendre le $\mathcal S'$-Module $\mathcal M$). On est donc dans les
conditions d'application du théorème (2.4.1, (ii)), et on en conclut qu'il existe
$n_0$ tel que, pour $n\geq n_0$ et pour tout $k$, on ait
\begin{equation}
\mathrm H^q(X_0,\mathcal M_k(n))=0\qquad\text{pour tout }q>0.
\tag{5.2.3.1}
\label{III.5.2.3.1-fr}
\end{equation}
On a donc aussi $\mathrm H^q(\mathfrak X,\mathcal M_k(n))=0$ pour $q>0$ et
$n\geq n_0$, $\mathcal M_k(n)$ étant cette fois considéré comme
$\mathcal O_{\mathfrak X}$-Module. Appliquant la suite exacte de cohomologie à
\[
0\to
\frac{\mathfrak J^k\mathcal F(n)}{\mathfrak J^{k+1}\mathcal F(n)}
\to
\frac{\mathfrak J^h\mathcal F(n)}{\mathfrak J^{k+1}\mathcal F(n)}
\to
\frac{\mathfrak J^h\mathcal F(n)}{\mathfrak J^k\mathcal F(n)}
\to0,
\]
on en déduit tout d'abord que pour $0\leq h<k$, $n\geq n_0$ et $q>0$, on a par
récurrence sur $k-h$
\begin{equation}
\mathrm H^q\left(\mathfrak X,
\frac{\mathfrak J^h\mathcal F(n)}{\mathfrak J^k\mathcal F(n)}\right)=0
\tag{5.2.3.2}
\label{III.5.2.3.2-fr}
\end{equation}
et en particulier pour $h=0$
\begin{equation}
\mathrm H^q(\mathfrak X,\mathcal F_k(n))=0
\qquad\text{pour }n\geq n_0,\ k\geq0\text{ et }q>0.
\tag{5.2.3.3}
\label{III.5.2.3.3-fr}
\end{equation}
Une autre portion de la suite exacte de cohomologie, pour $h=0$, donne la suite
exacte
\begin{equation}
\mathrm H^0(\mathfrak X,\mathcal F_{k+1}(n))\to
\mathrm H^0(\mathfrak X,\mathcal F_k(n))\to
\mathrm H^1\left(\mathfrak X,
\frac{\mathfrak J^k\mathcal F(n)}{\mathfrak J^{k+1}\mathcal F(n)}\right)=0
\tag{5.2.3.4}
\label{III.5.2.3.4-fr}
\end{equation}
d'où on déduit que pour $h\leq k$, l'application canonique
\begin{equation}
\mathrm H^0(\mathfrak X,\mathcal F_k(n))\longrightarrow
\mathrm H^0(\mathfrak X,\mathcal F_h(n))
\tag{5.2.3.5}
\label{III.5.2.3.5-fr}
\end{equation}
est surjective. Pour tout $q$, le système projectif
$(\mathrm H^q(\mathfrak X,\mathcal F_k(n)))_{k\geq0}$ vérifie donc la condition
(ML) pour $n\geq n_0$. Par ailleurs, tout ouvert formel affine $U$ de $\mathfrak X$
est aussi un ouvert affine dans chacun des $X_k$ (I, 10.5.2), donc on a
$\mathrm H^q(U,\mathcal F_k(n))=0$ pour tout $q>0$ (I, 1.3.1), et
$\mathrm H^0(U,\mathcal F_k(n))\to\mathrm H^0(U,\mathcal F_h(n))$ est surjective
pour $h\leq k$ (I, 1.3.9). Les conditions d'application de ($0$, 13.3.1) sont par
suite remplies, et on en conclut que, pour $n\geq n_0$~:
\begin{enumerate}
\item[$1^\circ$] Pour tout $q>0$,
\[
\mathrm H^q(\mathfrak X,\mathcal F(n))\longrightarrow
\varprojlim_k\mathrm H^q(\mathfrak X,\mathcal F_k(n))
\]
est bijectif, donc, en vertu de (\ref{III.5.2.3.3-fr}),
$\mathrm H^q(\mathfrak X,\mathcal F(n))=0$.
\item[$2^\circ$] L'homomorphisme
\[
\mathrm H^0(\mathfrak X,\mathcal F(n))\longrightarrow
\varprojlim_k\mathrm H^0(\mathfrak X,\mathcal F_k(n))
\]
est bijectif~; par ailleurs, comme les homomorphismes
(\ref{III.5.2.3.5-fr}) sont surjectifs, il en est de même de chacun des
homomorphismes
\[
\varprojlim_k\mathrm H^0(\mathfrak X,\mathcal F_k(n))
\longrightarrow\mathrm H^0(\mathfrak X,\mathcal F_h(n)),
\]
\end{enumerate}
ce qui achève la démonstration.
\end{proof}

\begin{corollary}[5.2.4]
\label{III.5.2.4-fr}
Les hypothèses étant celles de (\hyperref[III.5.2.3-fr]{5.2.3}), pour tout
$\mathcal O_{\mathfrak X}$-Module cohérent $\mathcal F$, il existe un entier $N$
tel que pour $n\geq N$, $\mathcal F(n)$ soit engendré par ses sections au-dessus

\oldpage[III]{153}

de $\mathfrak X$~; en d'autres termes, $\mathcal F$ est isomorphe au quotient d'un
$\mathcal O_{\mathfrak X}$-Module de la forme
$(\mathcal O_{\mathfrak X}(-n))^k$.
\end{corollary}

\begin{proof}
Comme $X_0$ est noethérien, il résulte de l'hypothèse sur $\mathcal L_0$ et de
(II, 4.5.5) qu'il existe $n_0$ tel que, pour $n\geq n_0$, $\mathcal F_0(n)$ soit
engendré par ses sections au-dessus de $\mathfrak X$~; par ailleurs, on peut supposer
$n_0$ pris assez grand pour que l'homomorphisme
\[
\Gamma(\mathfrak X,\mathcal F(n))\longrightarrow
\Gamma(\mathfrak X,\mathcal F_0(n))
\]
soit surjectif pour $n\geq n_0$ (\hyperref[III.5.2.3-fr]{5.2.3}). Il existe donc
un nombre fini de sections $s_i\in\Gamma(\mathfrak X,\mathcal F(n))$ dont les
images dans $\Gamma(\mathfrak X,\mathcal F_0(n))$ engendrent $\mathcal F_0(n)$
($0_{\mathrm I}$, 5.2.3). Comme $\mathfrak J$ est contenu dans l'idéal maximal de
l'anneau local en tout point de $\mathfrak X$, il résulte du lemme de Nakayama,
appliqué à ces anneaux locaux, que les $s_i$ engendrent $\mathcal F(n)$
($0_{\mathrm I}$, 5.1.1).
\end{proof}

\begin{env}[5.2.5]
\label{III.5.2.5-fr}
\emph{Démonstration du théorème d'existence : cas projectif.}
Les notations étant celles de (\hyperref[III.5.1.4-fr]{5.1.4}), supposons $f$
projectif, de sorte qu'il existe un $\mathcal O_X$-Module ample $\mathcal L$
(I, 5.5.4). Par définition, $X_n=\widehat X\times_{\widehat Y}Y_n$ est égal au
sous-préschéma fermé $X\times_Y Y_n=f^{-1}(Y_n)$ de $X$~; si
$\mathcal L'=\mathcal L\otimes_{\mathcal O_X}\mathcal O_{\widehat X}$ est le
complété de $\mathcal L$, on a donc
$\mathcal L'_0=\mathcal L/\mathfrak J\mathcal L$, considéré comme
$\mathcal O_{X_0}$-Module~; on sait que $\mathcal L'_0$ est ample
(II, 4.6.13, (i bis)). On peut donc appliquer à $\mathcal L'$ et à tout
$\mathcal O_{\widehat X}$-Module cohérent $\mathcal F$ le cor.
(\hyperref[III.5.2.4-fr]{5.2.4})~; on voit donc que $\mathcal F$ est isomorphe à
un quotient de
\[
\mathcal G=(\mathcal L'^{\otimes(-n)})^k
\]
pour des entiers $n>0$ et $k>0$ convenables. Or, il est clair que $\mathcal G$ est
le complété de $(\mathcal L^{\otimes(-n)})^k$ (I, 10.8.10), donc est
algébrisable. Considérons ensuite l'homomorphisme canonique surjectif
$u:\mathcal G\to\mathcal F$, et soit $\mathcal H=\operatorname{Ker}(u)$, qui est
un $\mathcal O_{\widehat X}$-Module cohérent ($0_{\mathrm I}$, 5.3.4). On voit de
même qu'il existe un $\mathcal O_{\widehat X}$-Module algébrisable $\mathcal K$ et
un homomorphisme $v:\mathcal K\to\mathcal G$ tel que
$\mathcal H=\operatorname{Im}(v)$. On a alors
$\mathcal F=\operatorname{Coker}(v)$, et $\mathcal F$ est algébrisable en vertu de
(\hyperref[III.5.2.2-fr]{5.2.2}).
\end{env}

\begin{env}[5.2.6]
\label{III.5.2.6-fr}
\emph{Démonstration du théorème d'existence; cas quasi-projectif.}
Les notations étant toujours celles de (\hyperref[III.5.1.4-fr]{5.1.4}), supposons
maintenant que $f$ soit quasi-projectif. Il existe alors un morphisme projectif
$g:Z\to Y$ tel que $X$ s'identifie au $Y$-préschéma induit sur un ouvert de $Z$
(II, 5.3.2)~; si on pose $Z'=g^{-1}(Y')$, on a $X'=X\cap Z'$. Par suite, le
complété $\widehat X=X_{/X'}$ s'identifie au préschéma formel induit par le complété
$\widehat Z=Z_{/Z'}$ sur l'ouvert $X\cap Z'$ de $\widehat Z$ (I, 10.8.5). Soit
$\mathcal F'$ un $\mathcal O_{\widehat X}$-Module cohérent, dont le support $T'$ est
propre sur $\widehat Y$~; cela signifie par définition qu'il existe un sous-préschéma
fermé de $X'$, ayant $T'\subset X'$ comme espace sous-jacent, tel que la restriction
$T'\to Y'$ de $f$ soit propre~; on en conclut que $T'$ est propre sur $Y$, donc
fermé dans $Z'$ (II, 5.4.10). Il en résulte que $\mathcal F'$ est le faisceau induit
sur $\widehat X$ par le $\mathcal O_{\widehat Z}$-Module $\mathcal G'$ obtenu par
recollement de $\mathcal F'$ (défini sur l'ouvert $\widehat X$ de $\widehat Z$) et
du faisceau $0$ sur l'ouvert $\widehat Z-T'$ de $\widehat Z$, ces deux faisceaux
coïncidant dans l'ouvert intersection $\widehat X-T'$. Il est clair que
$\mathcal G'$ est cohérent~; en vertu de
(\hyperref[III.5.2.5-fr]{5.2.5}), il existe un $\mathcal O_Z$-Module cohérent
$\mathcal G$ tel que $\mathcal G'=\widehat{\mathcal G}$~; soit $T$ le support de
$\mathcal G$, de sorte que $T'=T\cap Z'$ (I, 10.8.12). Si $h$ est la restriction
de $g$ au sous-préschéma fermé réduit de $Z$ ayant $T$ comme espace sous-jacent, on
a donc $T'=h^{-1}(Y')=T\cap g^{-1}(Y')$, et par suite $X\cap T$ est un ouvert de
$T$ contenant $T'$.

\oldpage[III]{154}

Comme $h$ est propre (II, 5.4.2), donc fermé, il résulte de
(\hyperref[III.5.1.3.1-fr]{5.1.3.1}) que $X\cap T=T$, autrement dit $T\subset X$,
et comme $T$ est fermé dans $Z$, $T$ est propre sur $Y$. Si $\mathcal F$ est le
faisceau induit sur $X$ par $\mathcal G$, son complété $\widehat{\mathcal F}$ est
induit sur $\widehat X$ par $\widehat{\mathcal G}$ (I, 10.8.4), donc est égal à
$\mathcal F'$, ce qui achève la démonstration.
\end{env}

\subsection{Démonstration du théorème d'existence : cas général}
\label{subsection:III.5.3-fr}

\begin{lemma}[5.3.1]
\label{III.5.3.1-fr}
Sous les conditions de (\hyperref[III.5.1.4-fr]{5.1.4}), si
\[
0\to\mathcal H\to\mathcal F\to\mathcal G\to0
\]
est une suite exacte de $\mathcal O_{\widehat X}$-Modules cohérents telle que
$\mathcal G$ et $\mathcal H$ soient algébrisables, alors $\mathcal F$ est
algébrisable.
\end{lemma}

\begin{proof}
En effet, supposons que $\mathcal G=\widehat{\mathcal B}$,
$\mathcal H=\widehat{\mathcal C}$, $\mathcal B$ et $\mathcal C$ étant des
$\mathcal O_X$-Modules cohérents à supports propres sur $Y$~; $\mathcal F$ définit
canoniquement un élément du $A$-module
\[
\operatorname{Ext}^1_{\mathcal O_{\widehat X}}
(\widehat X;\widehat{\mathcal B},\widehat{\mathcal C})
\]
($0$, 12.3.2), et les hypothèses entraînent que ce $A$-module est canoniquement
isomorphe à
\[
\operatorname{Ext}^1_{\mathcal O_X}(X;\mathcal B,\mathcal C)
\]
(4.5.2)~; il existe donc une suite exacte
$0\to\mathcal C\to\mathcal A\to\mathcal B\to0$ de $\mathcal O_X$-Modules
cohérents telle que l'image canonique de l'élément de
$\operatorname{Ext}^1_{\mathcal O_X}(X;\mathcal B,\mathcal C)$ correspondant à
$\mathcal A$ soit l'élément de
$\operatorname{Ext}^1_{\mathcal O_{\widehat X}}
(\widehat X;\widehat{\mathcal B},\widehat{\mathcal C})$ correspondant à
$\mathcal F$. Mais par définition (compte tenu de (I, 10.8.8, (iii))), cela signifie
que $\mathcal F$ est isomorphe à $\widehat{\mathcal A}$, d'où le lemme, car
$\operatorname{Supp}(\mathcal A)$ est contenu dans la réunion de
$\operatorname{Supp}(\mathcal B)$ et $\operatorname{Supp}(\mathcal C)$, donc est
propre sur $Y$.
\end{proof}

\begin{corollary}[5.3.2]
\label{III.5.3.2-fr}
Sous les conditions de (\hyperref[III.5.1.1-fr]{5.1.1}), soit
$u:\mathcal F\to\mathcal G$ un homomorphisme de
$\mathcal O_{\widehat X}$-Modules cohérents~; si $\mathcal G$,
$\operatorname{Ker}(u)$ et $\operatorname{Coker}(u)$ sont algébrisables, il en est
de même de $\mathcal F$.
\end{corollary}

\begin{proof}
Le lemme (\hyperref[III.5.2.2-fr]{5.2.2}) appliqué à l'homomorphisme
$\mathcal G\to\operatorname{Coker}(u)$ montre en effet que
$\operatorname{Im}(u)$ est algébrisable, et il suffit ensuite d'appliquer le lemme
(\hyperref[III.5.3.1-fr]{5.3.1}) à la suite exacte
\[
0\to\operatorname{Ker}(u)\to\mathcal F\to\operatorname{Im}(u)\to0.
\]
\end{proof}

\begin{lemma}[5.3.3]
\label{III.5.3.3-fr}
Sous les conditions de (\hyperref[III.5.1.1-fr]{5.1.1}), soient $h:Z\to Y$ un
morphisme de type fini, $\widehat Z$ le complété de $Z$ le long de
$Z'=h^{-1}(Y')$, $g:Z\to X$ un $Y$-morphisme propre,
$\widehat g:\widehat Z\to\widehat X$ son prolongement aux complétés. Pour tout
$\mathcal O_{\widehat Z}$-Module algébrisable $\mathcal F'$,
$\widehat g_*(\mathcal F')$ est un $\mathcal O_{\widehat X}$-Module algébrisable.
\end{lemma}

\begin{proof}
En effet, si $\mathcal F$ est un $\mathcal O_Z$-Module cohérent tel que
$\mathcal F'=\widehat{\mathcal F}$, il résulte du premier th. de comparaison
(4.1.5) que $\widehat g_*(\widehat{\mathcal F})$ est isomorphe au complété de
$g_*(\mathcal F)$.
\end{proof}

\begin{lemma}[5.3.4]
\label{III.5.3.4-fr}
Soient $X$ un schéma (usuel) noethérien, $X'$ une partie fermée de $X$,
$f:Z\to X$ un morphisme propre, $Z'=f^{-1}(X')$,
$\widehat X=X_{/X'}$, $\widehat Z=Z_{/Z'}$,
$\widehat f:\widehat Z\to\widehat X$ le prolongement de $f$ aux complétés. Soit
$\mathcal M$ un Idéal cohérent de $\mathcal O_X$ tel que, si
$U=X-\operatorname{Supp}(\mathcal O_X/\mathcal M)$, la restriction
$f^{-1}(U)\to U$ de $f$ soit un isomorphisme. Alors, pour tout
$\mathcal O_{\widehat X}$-Module cohérent $\mathcal F$, il existe un entier $n>0$
tel que le noyau et le conoyau de l'homomorphisme canonique
\[
\rho_{\mathcal F}:\mathcal F\longrightarrow
\widehat f_*\bigl(\widehat f^{,*}(\mathcal F)\bigr)
\]
($0_{\mathrm I}$, 4.4.3) soient annulés par $\widehat{\mathcal M}^{,n}$.
\end{lemma}

\begin{proof}
On peut se borner au cas où $X=\operatorname{Spec}(B)$ où $B$ est un anneau
noethérien, donc $X'=\mathrm V(\mathfrak K)$, où $\mathfrak K$ est un idéal de
$B$. Nous allons voir qu'on peut se ramener au cas où $B$ est un anneau adique
noethérien et $\mathfrak K$ un idéal de définition de $B$. Soit en effet $B_1$ le
séparé complété de $B$ pour la topologie $\mathfrak K$-préadique~; si
$\mathfrak K_1=\mathfrak K B_1$, $B_1$ est donc un

\oldpage[III]{155}

anneau adique noethérien dont $\mathfrak K_1$ est un idéal de définition. Posons
$X_1=\operatorname{Spec}(B_1)$ et soit $h:X_1\to X$ le morphisme correspondant à
l'homomorphisme canonique $B\to B_1$~; si $X'_1=h^{-1}(X')$, on a donc
$X'_1=\mathrm V(\mathfrak K_1)$. Posons enfin
$Z_1=Z\times_X X_1=Z_{(X_1)}$, $f_1=f_{(X_1)}:Z_1\to X_1$, qui est un morphisme
propre (II, 5.4.2), et désignons par $\widehat X_1$ le complété de $X_1$ le long
de $X'_1$, par
$\widehat Z_1=Z_1\times_{X_1}\widehat X_1$ le complété de $Z_1$ le long de
$Z'_1=f_1^{-1}(X'_1)$, par $\widehat f_1$ le prolongement de $f_1$ aux complétés.
Il est immédiat que le prolongement $\widehat h:\widehat X_1\to\widehat X$ de $h$
aux complétés est un isomorphisme, correspondant à l'application identique de $B_1$
(I, 10.9.1)~; on en conclut que l'homomorphisme correspondant
$\widehat Z_1\to\widehat Z$ est aussi un isomorphisme, ces isomorphismes identifiant
$\widehat f_1$ et $\widehat f$. Enfin,
$\mathcal M_1=h^*(\mathcal M)$ est un Idéal cohérent de $\mathcal O_{X_1}$ et
\[
\operatorname{Supp}(\mathcal O_{X_1}/\mathcal M_1)
=h^{-1}(\operatorname{Supp}(\mathcal O_X/\mathcal M))
\]
(I, 9.1.13), donc, si
$U_1=X_1-\operatorname{Supp}(\mathcal O_{X_1}/\mathcal M_1)$, on a
$U_1=h^{-1}(U)$, d'où résulte aussitôt que la restriction
$f_1^{-1}(U_1)\to U_1$ de $f_1$ est un isomorphisme (I, 3.2.7)~; en outre, les
complétés $\widehat{\mathcal M}$ et $\widehat{\mathcal M}_1$ s'identifient par
$\widehat h$ (I, 10.9.5). Toutes les hypothèses de
(\hyperref[III.5.3.4-fr]{5.3.4}) sont donc remplies par $X_1$, $X'_1$, $f_1$ et
$\mathcal M_1$, et on peut donc désormais supposer $B$ adique noethérien et
$\mathfrak K$ un idéal de définition de $B$. On a alors
$\widehat X=\operatorname{Spf}(B)$, et $\mathcal F=N^\Delta$, où $N$ est un
$B$-module de type fini, d'où $\mathcal F=\widehat{\mathcal G}$, où $\mathcal G$
est le $\mathcal O_X$-Module cohérent $\widetilde N$ (I, 10.10.5), et par suite
\[
\widehat f^{,*}(\mathcal F)=\widehat{f^*(\mathcal G)}
\]
(I, 10.9.5). En outre, en vertu du premier théorème de comparaison (4.1.5),
\[
\widehat f_*\bigl(\widehat{f^*(\mathcal G)}\bigr)
\simeq\widehat{f_*\bigl(f^*(\mathcal G)\bigr)},
\]
et l'homomorphisme canonique $\rho_{\mathcal F}$ n'est autre que
$\widehat{\rho}_{\mathcal G}$ en vertu de
(\hyperref[III.5.1.3-fr]{5.1.3}). Or, le noyau $\mathcal P$ et le conoyau
$\mathcal R$ de
\[
\rho_{\mathcal G}:\mathcal G\longrightarrow f_*\bigl(f^*(\mathcal G)\bigr)
\]
sont des $\mathcal O_X$-Modules cohérents, et par hypothèse leurs restrictions à
$U$ sont évidemment nulles. Il existe par suite un entier $n>0$ tel que
$\mathcal M^n\mathcal P=\mathcal M^n\mathcal R=0$ (I, 9.3.4)~; on en conclut que
$\widehat{\mathcal M}^{,n}\widehat{\mathcal P}=
\widehat{\mathcal M}^{,n}\widehat{\mathcal R}=0$
(I, 10.8.8 et 10.8.10).
\end{proof}

\begin{env}[5.3.5]
\label{III.5.3.5-fr}
\emph{Fin de la démonstration du théorème d'existence.}
Les hypothèses étant celles de (\hyperref[III.5.1.4-fr]{5.1.4}), nous allons
utiliser le principe de récurrence noethérienne ($0_{\mathrm I}$, 2.2.2) en supposant
donc le théorème vrai pour tout sous-préschéma fermé $T$ de $X$ dont l'espace
sous-jacent est distinct de $X$ (le complété $\widehat T$ étant bien entendu le
complété de $T$ le long de $T\cap X'$). On peut supposer $X$ non vide. Comme $f$
est séparé et de type fini, on peut appliquer le lemme de Chow (II, 5.6.1)~: il
existe donc un $Y$-schéma $Z$ et un $Y$-morphisme $g:Z\to X$ tels que le morphisme
structural $h:Z\to Y$ soit quasi-projectif, le morphisme $g$ projectif et surjectif,
et en outre un ouvert non vide $U$ de $X$ tel que la restriction
$g^{-1}(U)\to U$ soit un isomorphisme. Soit $\mathcal M$ un Idéal cohérent de
$\mathcal O_X$ définissant un sous-préschéma fermé d'espace sous-jacent $X-U$
(I, 5.2.2), et soit $\mathcal F$ un $\mathcal O_{\widehat X}$-Module cohérent dont
le support $E$ est propre sur $Y$~; désignons par $\widehat Z$ le complété de $Z$
le long de $h^{-1}(Y')$, par $\widehat g:\widehat Z\to\widehat X$ le prolongement
de $g$ aux complétés. Alors $\widehat g^{,*}(\mathcal F)$ est un
$\mathcal O_{\widehat Z}$-Module cohérent dont le support est contenu dans
$g^{-1}(E)$ et est par suite propre sur $Y$, puisque $g$ est projectif, donc propre
(II, 5.4.6). Comme $h$ est quasi-projectif, $\widehat g^{,*}(\mathcal F)$ est
algébrisable en vertu de (\hyperref[III.5.2.6-fr]{5.2.6}). On en conclut

\oldpage[III]{156}

que $\widehat g_*\bigl(\widehat g^{,*}(\mathcal F)\bigr)$ est un
$\mathcal O_{\widehat X}$-Module algébrisable
(\hyperref[III.5.3.3-fr]{5.3.3}) puisque $g$ est propre. On peut maintenant
appliquer à $\mathcal F$ et à $g$ le résultat de
(\hyperref[III.5.3.4-fr]{5.3.4})~: le noyau $\mathcal P$ et le conoyau
$\mathcal R$ de l'homomorphisme
$\rho_{\mathcal F}:\mathcal F\to
\widehat g_*\bigl(\widehat g^{,*}(\mathcal F)\bigr)$ sont annulés par une
puissance $\widehat{\mathcal M}^{,n}$~; soient $T$ le sous-préschéma fermé de $X$
défini par $\mathcal M^n$, ayant $X-U$ pour espace sous-jacent,
$j:T\to X$ l'injection canonique, de sorte que le prolongement aux complétés
$\widehat j:\widehat T\to\widehat X$ est l'injection canonique (I, 10.14.7). On
peut donc écrire
$\mathcal P=\widehat j_*\bigl(\widehat j^{,*}(\mathcal P)\bigr)$ et
$\mathcal R=\widehat j_*\bigl(\widehat j^{,*}(\mathcal R)\bigr)$ et comme $U$
n'est pas vide, il résulte de l'hypothèse de récurrence que
$\widehat j^{,*}(\mathcal P)$ et $\widehat j^{,*}(\mathcal R)$ sont des
$\mathcal O_{\widehat T}$-Modules algébrisables~; en vertu de
(\hyperref[III.5.3.3-fr]{5.3.3}), $\mathcal P$ et $\mathcal R$ sont
algébrisables, et on peut alors appliquer (\hyperref[III.5.3.2-fr]{5.3.2}), qui
prouve finalement que $\mathcal F$ est algébrisable.

\hfill C.Q.F.D.
\end{env}

\subsection{Application : comparaison de morphismes de schémas usuels et de morphismes de schémas formels. Schémas formels algébrisables}
\label{subsection:III.5.4-fr}

\begin{theorem}[5.4.1]
\label{III.5.4.1-fr}
Soient $A$ un anneau adique noethérien, $\mathfrak J$ un idéal de définition de
$A$, $S=\operatorname{Spec}(A)$, $S'=\mathrm V(\mathfrak J)$. Soient
$u:X\to S$ un morphisme propre, $v:Y\to S$ un morphisme séparé de type fini, et
soient $\widehat S$, $\widehat X$, $\widehat Y$ les complétés de $S$, $X$, $Y$ le
long de $S'$, $u^{-1}(S')$, $v^{-1}(S')$ respectivement. Si, pour tout
$S$-morphisme $f:X\to Y$, $\widehat f:\widehat X\to\widehat Y$ est le
prolongement de $f$ aux complétés, l'application $f\mapsto\widehat f$ est une
bijection
\[
\operatorname{Hom}_S(X,Y)\xrightarrow{\ \sim\ }
\operatorname{Hom}_{\widehat S}(\widehat X,\widehat Y).
\]
\end{theorem}

\begin{proof}
Montrons d'abord que $f\mapsto\widehat f$ est injective. Supposons en effet que deux
$S$-morphismes $f$, $g$ de $X$ dans $Y$ soient tels que $\widehat f=\widehat g$.
On sait alors (I, 10.9.4) qu'il existe un voisinage ouvert $V$ de
$X'=u^{-1}(S')$ dans lequel $f$ et $g$ coïncident. Or, comme $u$ est une application
fermée, on a $V=X$ (\hyperref[III.5.1.3.1-fr]{5.1.3.1}), d'où $f=g$.

Prouvons maintenant que $f\mapsto\widehat f$ est surjective, et soit donc $h$ un
$\widehat S$-morphisme $\widehat X\to\widehat Y$. Soit $Z=X\times_S Y$, et
désignons par $p:Z\to X$ et $q:Z\to Y$ les projections canoniques~; $Z$ est de
type fini sur $S$ (I, 6.3.4), donc noethérien~; désignons par $\widehat Z$ son
complété le long de $Z'=p^{-1}(u^{-1}(S'))$~; on sait que $\widehat Z$ s'identifie
canoniquement à $\widehat X\times_{\widehat S}\widehat Y$, les projections
$\widehat Z\to\widehat X$ et $\widehat Z\to\widehat Y$ s'identifiant aux
prolongements $\widehat p$ et $\widehat q$ (I, 10.9.7). Comme $Y$ est séparé sur
$S$, $\widehat Y$ est séparé sur $\widehat S$ (I, 10.15.7), donc le morphisme
graphe
\[
\Gamma_h=(1_{\widehat X},h):\widehat X\to\widehat Z
\]
est une immersion fermée (I, 10.15.4). Soient $\mathfrak T$ le sous-préschéma
formel fermé de $\widehat Z$ associé à cette immersion, et
$j:\mathfrak T\to\widehat Z$ l'injection canonique, de sorte que
$\Gamma_h=j\circ w$, où $w:\widehat X\to\mathfrak T$ est un isomorphisme
(I, 10.14.3) dont l'isomorphisme réciproque est $\widehat p\circ j$~; en outre,
$\mathfrak T$ est évidemment propre sur $\widehat S$, puisque $\widehat X$ l'est~;
on en conclut (\hyperref[III.5.1.8-fr]{5.1.8}) qu'il existe un sous-préschéma
fermé $T$ de $Z$ tel que
$\mathfrak T=\widehat T=T_{/(T\cap Z')}$, et que $j=\widehat i$, où $i$ est
l'injection canonique $T\to Z$ (I, 10.14.7). Alors $p\circ i:T\to X$ est un
isomorphisme, car il en est ainsi de
$\widehat{p\circ i}=\widehat p\circ\widehat i$ par hypothèse, et il suffit
d'appliquer

\oldpage[III]{157}

(4.6.8) en notant comme ci-dessus que $S$ est le seul voisinage de $S'$ dans $S$.
Soit $g:X\to T$ l'isomorphisme réciproque de $p\circ i$, et posons
$f=q\circ i\circ g$, qui est un morphisme $X\to Y$ dont par définition le graphe
$\Gamma_f=i\circ g$. Comme $\widehat g$ est l'isomorphisme réciproque de
$\widehat{p\circ i}=w$, on a
\[
\widehat{\Gamma_f}=\widehat i\circ\widehat g=j\circ w=\Gamma_h.
\]
Mais on sait que $\widehat{\Gamma_f}=\Gamma_{\widehat f}$ (I, 10.9.8), d'où
finalement $h=\widehat f$, ce qui achève la démonstration.
\end{proof}

On peut donc dire, dans le langage des catégories, que le foncteur
$X\mapsto\widehat X$ est pleinement fidèle ($0$, 8.1.6) de la catégorie des
schémas propres sur $\operatorname{Spec}(A)$ dans la catégorie des schémas formels
propres sur $\operatorname{Spf}(A)$, pour tout anneau adique noethérien $A$~; il
établit par suite une équivalence entre la première de ces catégories et une
sous-catégorie de la seconde~; les objets de cette dernière seront appelés
\emph{schémas formels algébrisables}. Pour un tel schéma $\mathfrak X$, il existe
un schéma usuel $X$, propre sur $\operatorname{Spec}(A)$, déterminé à isomorphisme
unique près, tel que $\mathfrak X$ soit isomorphe à $\widehat X$.

\begin{env}[5.4.2]
\label{III.5.4.2-fr}
\emph{Scholie.}
Avec les notations de (\hyperref[III.5.4.1-fr]{5.4.1}), posons
$S_n=\operatorname{Spec}(A/\mathfrak J^{n+1})$,
$X_n=X\times_S S_n$, $Y_n=Y\times_S S_n$. Il résulte de
(\hyperref[III.5.4.1-fr]{5.4.1}) et de (I, 10.12.3) que se donner un
$S$-morphisme $f:X\to Y$ équivaut à se donner un $(S_n)$-système inductif adique
(I, 10.12.2) de $S_n$-morphismes $f_n:X_n\to Y_n$.
\end{env}

\begin{remark}[5.4.3]
\label{III.5.4.3-fr}
Contrairement à ce que pourrait suggérer le th. d'existence
(\hyperref[III.5.1.6-fr]{5.1.6}), il y a des schémas formels propres sur
$\operatorname{Spf}(A)$ et qui ne sont pas algébrisables (tout comme il y a des
espaces analytiques compacts qui ne proviennent pas de variétés algébriques
complexes). Nous rencontrerons plus tard de tels schémas dans la «~théorie des
modules~», qui traite précisément (lorsque le corps de base est $\mathbb C$) des
variations infinitésimales de la structure complexe d'une variété algébrique
complète, et on sait que de telles variations peuvent donner naissance à des
variétés analytiques qui ne sont pas algébriques.
\end{remark}

\begin{proposition}[5.4.4]
\label{III.5.4.4-fr}
Soient $A$ un anneau adique noethérien, $\mathfrak S=\operatorname{Spf}(A)$,
$g:\mathfrak X\to\mathfrak S$, $h:\mathfrak Y\to\mathfrak S$ deux morphismes
propres de schémas formels, $f:\mathfrak X\to\mathfrak Y$ un
$\mathfrak S$-morphisme. Si $f$ est fini et si $\mathfrak Y$ est algébrisable,
alors $\mathfrak X$ est algébrisable.
\end{proposition}

\begin{proof}
On notera que les hypothèses sur $g$ et $h$ entraînent déjà que $f$ est propre
(3.4.1), et pour que $f$ soit fini, il suffit que pour tout $y\in\mathfrak Y$, la
fibre $f^{-1}(y)$ soit finie (\hyperref[III.4.8.11-fr]{4.8.11}). L'hypothèse
entraîne que $\mathcal B=f_*(\mathcal O_{\mathfrak X})$ est une
$\mathcal O_{\mathfrak Y}$-Algèbre cohérente
(\hyperref[III.4.8.6-fr]{4.8.6}), donc il résulte du th. d'existence que, si
$\mathfrak Y=\widehat Y$ et $h=\widehat w$, où
$w:Y\to\operatorname{Spec}(A)$ est un morphisme propre de schémas usuels, il
existe une $\mathcal O_Y$-Algèbre cohérente $\mathcal C$ telle que
$\mathcal B=\widehat{\mathcal C}$. Soient $X=\operatorname{Spec}(\mathcal C)$, et
$u:X\to Y$ le morphisme structural~; alors, il résulte aussitôt de la définition de
$\mathfrak X$ à partir de $\mathcal B$
(\hyperref[III.4.8.7-fr]{4.8.7}) que $\mathfrak X$ est canoniquement isomorphe à
$\widehat X$ et que $f$ s'identifie à $\widehat u$ (il suffit pour le voir de
considérer le cas où $Y$ est affine).

On notera que (\hyperref[III.5.1.8-fr]{5.1.8}) est un cas particulier de
(\hyperref[III.5.4.4-fr]{5.4.4}).
\end{proof}

\begin{theorem}[5.4.5]
\label{III.5.4.5-fr}
Soient $A$ un anneau adique noethérien, $\mathfrak J$ un idéal de définition de
$A$, $S=\operatorname{Spec}(A)$,
$\mathfrak S=\widehat S=\operatorname{Spf}(A)$,
$f:\mathfrak X\to\mathfrak S$ un morphisme propre de schémas formels. On pose
$S_k=\operatorname{Spec}(A/\mathfrak J^{k+1})$,
$X_k=\mathfrak X\times_{\mathfrak S}S_k$, et pour tout
$\mathcal O_{\mathfrak X}$-Module $\mathcal F$,
$\mathcal F_k=\mathcal F\otimes_{\mathcal O_{\mathfrak X}}\mathcal O_{X_k}
=\mathcal F/\mathfrak J^{k+1}\mathcal F$.

\oldpage[III]{158}

Soit $\mathcal L$ un $\mathcal O_{\mathfrak X}$-Module inversible, et supposons
que $\mathcal L_0=\mathcal L/\mathfrak J\mathcal L$ soit un
$\mathcal O_{X_0}$-Module ample. Alors $\mathfrak X$ est algébrisable, et si
$X$ est un $S$-schéma propre tel que $\mathfrak X$ soit isomorphe à
$\widehat X$, il existe un $\mathcal O_X$-Module ample $\mathcal M$ tel que
$\mathcal L$ soit isomorphe à $\widehat{\mathcal M}$ (ce qui entraîne que
$X$ est projectif sur $S$).
\end{theorem}

\begin{proof}
Appliquons (\hyperref[III.5.2.3-fr]{5.2.3}) à
$\mathcal F=\mathcal O_{\mathfrak X}$~: il existe donc un entier $n_0$ tel que
pour $n\geq n_0$, l'homomorphisme canonique
$\Gamma(\mathfrak X,\mathcal L^{\otimes n})\to
\Gamma(X_0,\mathcal L_0^{\otimes n})$ soit surjectif. On peut supposer
$n\geq n_0$ choisi assez grand pour que $\mathcal L_0^{\otimes n}$ soit très
ample pour $S_0$ (II, 4.5.10). Comme le morphisme $f_0:X_0\to S_0$ est propre,
$\Gamma(X_0,\mathcal L_0^{\otimes n})$ est un $A$-module de type fini (3.2.1),
donc il existe un sous-$A$-module de type fini $E$ de
$\Gamma(\mathfrak X,\mathcal L^{\otimes n})$ dont l'image dans
$\Gamma(X_0,\mathcal L_0^{\otimes n})$ soit ce dernier module tout entier.
Cela étant, pour tout $k\geq0$, considérons l'homomorphisme
\[
u_k:E/\mathfrak J^{k+1}E\longrightarrow
\Gamma(X_k,\mathcal L_k^{\otimes n})
\]
déduit de l'injection canonique
$E\to\Gamma(\mathfrak X,\mathcal L^{\otimes n})$. Notons que
$(f_k)_*(\mathcal L_k^{\otimes n})$ est quasi-cohérent, et comme
$\Gamma(S_k,(f_k)_*(\mathcal L_k^{\otimes n}))
=\Gamma(X_k,\mathcal L_k^{\otimes n})$, $u_k$ définit un homomorphisme
\[
\widetilde u_k:(E/\mathfrak J^{k+1}E)^\sim
\longrightarrow(f_k)_*(\mathcal L_k^{\otimes n}),
\]
et par suite aussi un homomorphisme
\[
\widetilde u_k^\sharp:
f_k^*((E/\mathfrak J^{k+1}E)^\sim)\longrightarrow
\mathcal L_k^{\otimes n}.
\]
D'ailleurs, si on pose
$\mathcal G_k=f_k^*((E/\mathfrak J^{k+1}E)^\sim)$, on a
$\mathcal G_0=\mathcal G_k/\mathfrak J\mathcal G_k$ (I, 9.1.5), donc
$\widetilde u_0^\sharp:\mathcal G_k/\mathfrak J\mathcal G_k
\to\mathcal L_k^{\otimes n}/\mathfrak J\mathcal L_k^{\otimes n}$ se déduit
de $\widetilde u_k^\sharp$ par passage aux quotients. Or, par définition de
$E$, $\widetilde u_0^\sharp$ n'est autre que l'homomorphisme canonique
\[
\sigma:f_0^*((f_0)_*(\mathcal L_0^{\otimes n}))\longrightarrow
\mathcal L_0^{\otimes n},
\]
et l'hypothèse que $\mathcal L_0^{\otimes n}$ est très ample entraîne que
$\widetilde u_0^\sharp$ est surjectif (II, 4.4.3)~; on déduit alors du lemme
de Nakayama que chacun des $\widetilde u_k^\sharp$ est aussi surjectif.
Chacun des $\widetilde u_k^\sharp$ définit donc (II, 4.2.2) un
$S_k$-morphisme
\[
g_k:X_k\longrightarrow P_k=\mathbf P(E/\mathfrak J^{k+1}E),
\]
et comme $P_h=P_k\times_{S_k}S_h$ pour $h\leq k$ en vertu de (II, 4.1.3),
$(g_k)$ est un $(S_k)$-système inductif adique (I, 10.12.2) en vertu des
relations entre les $\widetilde u_k^\sharp$ et de (II, 4.2.10). Les $g_k$
définissent donc un $\mathfrak S$-morphisme de schémas formels
$g:\mathfrak X\to\mathfrak P$, où $\mathfrak P$ est la limite inductive du
système $(P_k)$, ou encore le complété $\widehat P$, où $P=\mathbf P(E)$. De
plus, l'hypothèse que $\mathcal L_0^{\otimes n}$ est très ample entraîne que
$g_0$ est une immersion fermée (II, 4.4.3)~; on en conclut que $g$ est une
immersion fermée de schémas formels (4.8.10), donc $\mathfrak X$ est
algébrisable (5.1.8). Le fait que $\mathcal L$ soit isomorphe au complété
$\widehat{\mathcal M}$ d'un $\mathcal O_X$-Module inversible résulte alors du
th. d'existence (5.1.6). En outre, $\mathcal L^{\otimes n}$ est alors le
complété de $\mathcal M^{\otimes n}$ (I, 10.8.10), et les homomorphismes
$\widetilde u_k^\sharp$ définissent un homomorphisme bien déterminé
\[
v:f^*(\widetilde E)\longrightarrow\mathcal M^{\otimes n}
\]
(5.1.7)~; d'ailleurs, comme $\widetilde u_k^\sharp$ est surjectif, il en est de
même de $\widehat v$ (I, 10.11.5), donc de $v$ (5.1.3)~; en outre, le morphisme
$r:X\to P$ défini par $v$ (II, 4.2.2) a pour prolongement aux complétés $g$,
et comme $g$ est une immersion fermée, il en est de même de $r$, par (5.1.8)
et (5.4.1)~; on en conclut que $\mathcal M^{\otimes n}$ est très ample
(II, 4.4.6) et $\mathcal M$ est ample (II, 4.5.10).
\end{proof}

\begin{remark}[5.4.6]
\label{III.5.4.6-fr}
Soient $A$ un anneau adique noethérien,
$\mathfrak S=\operatorname{Spf}(A)$,
$f:\mathfrak X\to\mathfrak S$ un morphisme propre de schémas formels. Soit
$\mathcal N$ l'Idéal cohérent de $\mathcal O_{\mathfrak X}$ tel que pour tout
ouvert formel affine $U$ de $\mathfrak X$, $\Gamma(U,\mathcal N)$ soit le
nilradical de $\Gamma(U,\mathcal O_{\mathfrak X})$~; l'existence de cet Idéal
résulte facilement de (I, 10.10.2) et du fait que tout homomorphisme d'anneaux
$B\to C$ applique le nilradical de $B$ dans celui de $C$. Soit
$\mathfrak X'$ le sous-schéma formel fermé de $\mathfrak X$ défini par
$\mathcal N$ (I, 10.14.2)~; il serait intéressant de savoir si, lorsque
$\mathfrak X'$ est algébrisable, $\mathfrak X$ lui-même est algébrisable. On
parviendrait sans doute à une solution de ce problème si on savait classifier
(par exemple au moyen d'invariants de nature

\oldpage[III]{159}

cohomologique), les \emph{extensions} d'un faisceau structural
$\mathcal O_{\mathfrak X}$ (pour un préschéma usuel ou un préschéma formel) par
un Idéal de carré nul, autrement dit les $\mathcal O_{\mathfrak X}$-Algèbres
$\mathcal A$ telles que $\mathcal O_{\mathfrak X}$ soit isomorphe à
$\mathcal A/\mathcal J$, où $\mathcal J$ est un Idéal de carré nul de
$\mathcal A$.
\end{remark}

\subsection{Une décomposition de certains schémas}
\label{subsection:III.5.5-fr}

\begin{proposition}[5.5.1]
\label{III.5.5.1-fr}
Soient $A$ un anneau adique noethérien, $\mathfrak J$ un idéal de définition
de $A$, $Y=\operatorname{Spec}(A)$. Soit $f:X\to Y$ un morphisme séparé de
type fini~; on pose $Y_0=\operatorname{Spec}(A/\mathfrak J)$,
$X_0=X\times_Y Y_0=f^{-1}(Y_0)$. Soit $Z_0$ une partie ouverte de $X_0$,
propre sur $Y_0$~; il existe alors dans $X$ une partie ouverte et fermée $Z$,
propre sur $Y$ et telle que $Z\cap X_0=Z_0$.
\end{proposition}

\begin{proof}
Par hypothèse, il y a un ouvert $T$ de $X$ tel que $T\cap X_0=Z_0$~; soit
$\widehat T$ le complété le long de $Z_0$ du schéma induit par $X$ sur
l'ouvert $T$~; le support de $\mathcal O_{\widehat T}$ étant $Z_0$, qui est
propre sur $Y_0$, $\widehat T$ est propre sur
$\widehat Y=\operatorname{Spf}(A)$ (3.4.1). Il résulte de (5.1.8) qu'il existe
un sous-schéma fermé $Z$ de $T$ propre sur $Y$, tel que, si $i:Z\to T$ est
l'injection canonique, $\widehat i:\widehat Z\to\widehat T$ soit un
isomorphisme ($\widehat Z$ étant le complété de $Z$ le long de $Z_0$). On en
conclut (4.6.8) qu'il existe dans $T$ un voisinage ouvert $V$ de $Z_0$ tel que
la restriction $i^{-1}(V)\to V$ de $i$ soit un isomorphisme. Mais
$i^{-1}(V)$ est un voisinage de $Z_0$ dans $Z$, donc est nécessairement
identique à $Z$ (5.1.3.1). On en conclut que $Z$ est ouvert dans $T$, donc
dans $X$, ce qui achève la démonstration.
\end{proof}

\begin{corollary}[5.5.2]
\label{III.5.5.2-fr}
Si $X_0$ est propre sur $Y_0$, $X$ est réunion de deux parties ouvertes
disjointes $Z$ et $Z'$ telles que $Z$ soit propre sur $Y$ et contienne $X_0$~;
en outre, toute partie fermée $P$ de $X$, propre sur $Y$, est contenue dans
$Z$.
\end{corollary}

\begin{proof}
La dernière assertion résulte de ce que $P\cap Z'$ étant fermé dans $P$ est
propre sur $Y$~; si $P\cap Z'$ n'était pas vide, $f(P\cap Z')$ serait fermé
non vide dans $Y$, donc rencontrerait $Y_0$ (5.1.3.1), ce qui contredit la
définition de $Z$.
\end{proof}

\begin{flushright}
\emph{(A suivre.)}
\end{flushright}

% La page imprimée 504 est blanche.
\clearpage
\thispagestyle{empty}
\null
\clearpage
\thispagestyle{plain}

\oldpage[III]{161}
\section*{BIBLIOGRAPHIE (suite)}

\begin{flushleft}
[27] P. Gabriel, \emph{Des catégories abéliennes}, thèse, Paris (1961).

\medskip
[28] J. E. Roos, \emph{Sur les foncteurs dérivés de $\varprojlim$.
Applications}, C. R. Acad. Sci. Paris, t. CCLII, 1961, p. 3702--3704.

\medskip
[29] H. Cartan, \emph{Séminaire de l'École Normale Supérieure},
13\textsuperscript{e} année (1960--1961), exposé n\textsuperscript{o} 11.
\end{flushleft}

\clearpage
\oldpage[III]{162}
\section*{INDEX DES NOTATIONS}

\begin{flushleft}
$\mathbf{Ens}$ : 0, 8.1.1.

$h_X(Y)$, $h_X(u)$ ($X$, $Y$ objets d'une catégorie $C$, $u$ morphisme de
$C$) : 0, 8.1.1.

$h_w(u)$ ($u$, $w$ morphismes de $C$) : 0, 8.1.2.

$Z_k(E_r^{pq})$, $B_k(E_r^{pq})$ ($k\geq r$) : 0, 11.1.1.

$Z_\infty(E_2^{pq})$, $B_\infty(E_2^{pq})$ : 0, 11.1.1.

$K^\bullet$, $K_\bullet$ (complexes) : 0, 11.2.1.

$K^{\bullet\bullet}$, $K_{\bullet\bullet}$ (bicomplexes) : 0, 11.2.1.

$B^i(K^\bullet)$, $Z^i(K^\bullet)$, $H^i(K^\bullet)$,
$H^\bullet(K^\bullet)$ : 0, 11.2.1.

$B_i(K_\bullet)$, $Z_i(K_\bullet)$, $H_i(K_\bullet)$,
$H_\bullet(K_\bullet)$ : 0, 11.2.1.

$T(K^\bullet)$, $T(K_\bullet)$ ($T$ foncteur) : 0, 11.2.1.

$E(K^\bullet)$ ($K^\bullet$ complexe filtré) : 0, 11.2.2.

$K^{i,\bullet}$, $K^{\bullet,j}$, $Z_{\mathrm{II}}^p(K^{i,\bullet})$,
$B_{\mathrm{II}}^p(K^{i,\bullet})$, $H_{\mathrm{II}}^p(K^{i,\bullet})$,
$Z_{\mathrm I}^p(K^{\bullet,j})$, $B_{\mathrm I}^p(K^{\bullet,j})$,
$H_{\mathrm I}^p(K^{\bullet,j})$ : 0, 11.3.1.

$H_{\mathrm{II}}^p(K^{\bullet\bullet})$,
$Z_{\mathrm I}^q(H_{\mathrm{II}}^p(K^{\bullet\bullet}))$,
$B_{\mathrm I}^q(H_{\mathrm{II}}^p(K^{\bullet\bullet}))$,
$H_{\mathrm I}^q(H_{\mathrm{II}}^p(K^{\bullet\bullet}))$,
$H^n(K^{\bullet\bullet})$ : 0, 11.3.1.

$F_{\mathrm I}^p(K^{\bullet\bullet})$,
$F_{\mathrm{II}}^p(K^{\bullet\bullet})$, ${}'E(K^{\bullet\bullet})$,
${}''E(K^{\bullet\bullet})$ : 0, 11.3.2.

$K_{i,\bullet}$, $K_{\bullet,j}$, $H_p^{\mathrm{II}}(K_{i,\bullet})$,
$H_p^{\mathrm I}(K_{\bullet,j})$, $H_p^{\mathrm{II}}(K_{\bullet\bullet})$,
$H_p^{\mathrm I}(K_{\bullet\bullet})$,
$H_q^{\mathrm I}(H_p^{\mathrm{II}}(K_{\bullet\bullet}))$,
$H_q^{\mathrm{II}}(H_p^{\mathrm I}(K_{\bullet\bullet}))$,
$H_n(K_{\bullet\bullet})$ : 0, 11.3.5.

$\mathrm R^\bullet T(K^\bullet)$ ($T$ foncteur covariant additif) : 0, 11.4.3.

$\mathrm R^\bullet T(K^\bullet,K'^{\bullet})$ ($T$ bifoncteur covariant
additif) : 0, 11.4.6.

$\mathrm L_\bullet T(K_\bullet)$ ($T$ foncteur covariant additif) : 0, 11.6.2.

$\mathrm L_\bullet T(K_\bullet,K'_\bullet)$ ($T$ bifoncteur covariant
additif) : 0, 11.6.6.

$\mathrm L_\bullet T(K_{\bullet\bullet})$ ($T$ foncteur covariant additif) :
0, 11.7.2.

$|\sigma|$, $\Sigma(A)$, $\mathrm C_\bullet(A)$, $\mathrm L_\bullet(A)$ ($A$ ensemble fini,
$\sigma$ simplexe de $A$) : 0, 11.8.1.

$\mathrm C^\bullet(A,B;\mathscr S)$, $\mathrm L^\bullet(A,B;\mathscr S)$ ($A$, $B$ ensembles
finis, $\mathscr S$ système de coefficients) : 0, 11.8.4.

$\mathrm P^\bullet(A,B,\mathscr S)$ ($A$, $B$ ensembles finis, $\mathscr S$ système
de coefficients) : 0, 11.8.6.

$\mathrm C^\bullet(A,B;\mathscr S^\bullet)$,
$\mathrm L^\bullet(A,B;\mathscr S^\bullet)$ ($A$, $B$ ensembles finis,
$\mathscr S^\bullet$ complexe de systèmes de coefficients) : 0, 11.8.8.

$\mathrm P^\bullet(A,B;\mathscr S^\bullet)$ ($A$, $B$ ensembles finis,
$\mathscr S^\bullet$ complexe de systèmes de coefficients) : 0, 11.8.10.

$\chi(M^\bullet)$ ($M^\bullet$ complexe fini de modules de longueur finie) :
0, 11.10.1.

$\check C^\bullet(\mathfrak U,\mathcal F)$ ($\mathcal F$ faisceau de groupes abéliens
sur $X$, $\mathfrak U$ recouvrement ouvert de $X$) : 0, 12.1.2.

$\mathrm R^pf_*(\mathcal F)$ ($f:X\to Y$ morphisme d'espaces annelés,
$\mathcal F$ faisceau de $\mathcal O_X$-Modules) : 0, 12.2.1.

$\mathcal H^p(f,\mathcal K^\bullet)$,
$\mathcal H_f^p(\mathcal K^\bullet)$, ${}'E(f,\mathcal K^\bullet)$,
${}''E(f,\mathcal K^\bullet)$ ($f:X\to Y$ morphisme d'espaces
annelés, $\mathcal K^\bullet$ complexe de $\mathcal O_X$-Modules) : 0, 12.4.1.

$\mathbf H^\bullet(X,\mathcal K^\bullet)$,
$\mathbf H^\bullet(V,\mathcal K^\bullet)$ ($X$ espace annelé,
$\mathcal K^\bullet$ complexe de $\mathcal O_X$-Modules, $V$ ouvert de $X$) :
0, 12.4.2.

$\operatorname{gr}^\bullet(A)$ ($A$ système projectif vérifiant (ML)) :
0, 13.4.2.

$E(A)$ ($A$ objet filtré) : 0, 13.6.4.

$K_\bullet(\mathbf f)$, $i_{\mathbf f}$ ($\mathbf f$ famille finie d'éléments
d'un anneau) : III, 1.1.1.

$K^\bullet(\mathbf f,M)$, $K_\bullet(\mathbf f,M)$ ($\mathbf f$ famille
finie d'éléments d'un anneau $A$, $M$ $A$-module) : III, 1.1.2.

$H^\bullet(\mathbf f,M)$, $H_\bullet(\mathbf f,M)$ ($\mathbf f$ famille
finie d'éléments d'un anneau $A$, $M$ $A$-module) : III, 1.1.3.

$C^\bullet((\mathbf f),M)$, $H^\bullet((\mathbf f),M)$ ($\mathbf f$ famille
finie d'éléments d'un anneau $A$, $M$ $A$-module) : III, 1.1.6.

$H^\bullet(U,\mathcal F^{(*)})$,
$H^\bullet(\mathfrak U,\mathcal F^{(*)})$ ($\mathcal F$
$\mathcal O_X$-Module quasi-cohérent, $U$ ouvert de $X$, $\mathfrak U$
recouvrement ouvert de $X$) : III, 2.1.1.
\end{flushleft}

\clearpage
\oldpage[III]{163}
\section*{INDEX TERMINOLOGIQUE}

\begin{flushleft}
Aboutissement d'une suite spectrale : 0, 11.1.1.

Algébrisable ($\mathcal O_X$-Module) : III, 5.2.1.

Algébrisable (schéma formel) : III, 5.4.2.

Analytiquement intègre (anneau) : III, 4.3.6.

Application quasi-compacte : 0, 9.1.1.

Augmentation d'une résolution : 0, 11.4.1.

Caractéristique d'Euler-Poincaré d'un complexe : 0, 11.10.1.

Caractéristique d'Euler-Poincaré d'un $\mathcal O_X$-Module cohérent ($X$
projectif sur $\operatorname{Spec}(A)$, $A$ anneau artinien) : III, 2.5.1.

Cochaîne bi-alternée : 0, 11.8.4.

Complexe défini par un bicomplexe : 0, 11.3.1.

Complexe de l'algèbre extérieure : III, 1.1.1.

Condition (ML) : 0, 13.1.2.

Constructible (partie, ensemble) : 0, 9.1.2.

Constructible (fonction) : 0, 9.3.1.

Cup-produit : 0, 12.1.2.

Dihomomorphisme pour les structures algébriques sur les catégories : 0, 8.2.1.

Exact (sous-ensemble) dans une catégorie abélienne : III, 3.1.1.

Essentiellement constant (système projectif) : 0, 13.4.2.

Factorisation de Stein : III, 4.3.3.

Filtration co-discrète, co-séparée, discrète, exhaustive, finie, séparée :
0, 11.1.3.

Final (objet) d'une catégorie : 0, 8.1.10.

Fini (morphisme) de préschémas formels : III, 4.8.2.

Fini (préschéma formel) au-dessus de $\mathfrak Y$,
$\mathfrak Y$-fini (préschéma formel) : III, 4.8.2.

Foncteur covariant canonique $C\to\operatorname{Hom}(C^o,\mathbf{Ens})$ :
0, 8.1.2.

Genre arithmétique : III, 2.5.1.

Géométriquement connexe (fibre) : III, 4.3.4.

Homomorphisme de structures algébriques sur les catégories : 0, 8.2.1.

Hypercohomologie d'un foncteur par rapport à un complexe $K^\bullet$ :
0, 11.4.3.

Hypercohomologie d'un bifoncteur par rapport à deux complexes
$K^\bullet$, $K'^{\bullet}$ : 0, 11.4.6.

Hyperhomologie d'un foncteur par rapport à un complexe $K_\bullet$ :
0, 11.6.2.

Hyperhomologie d'un bifoncteur par rapport à deux complexes $K_\bullet$,
$K'_\bullet$ : 0, 11.6.6.

Hyperhomologie d'un foncteur par rapport à un bicomplexe : 0, 11.7.2.

Localement constructible (ensemble) : 0, 9.1.11.

Loi de composition externe (interne) sur les objets d'une catégorie :
0, 8.2.1.

\oldpage[III]{164}

$S$-$C$-module filtré, $\operatorname{gr}^\bullet(S)$-$C$-module gradué,
$\operatorname{gr}^{\bullet\bullet}(S)$-$C$-module bigradué : 0, 13.6.6.

Morphisme de suites spectrales : 0, 11.1.2.

Nombre géométrique de composantes connexes d'une fibre : III, 4.3.4.

$C$-objet en groupes, $C$-groupe, $C$-anneau, $C$-module : 0, 8.2.3.

Objet des bords, objet des cobords, objet des cocycles, objet des cycles :
0, 11.2.1.

Objet des images universelles : 0, 13.1.1.

Objet gradué associé à un système projectif satisfaisant à (ML) : 0, 13.4.2.

Objet gradué limité inférieurement (supérieurement) : 0, 11.2.1.

Partie propre (sur $\mathfrak Y$) d'un préschéma formel : III, 3.4.1.

Pleine (sous-catégorie) : 0, 8.1.5.

Pleinement fidèle (foncteur) : 0, 8.1.5.

Polynôme de Hilbert : III, 2.5.3.

Propre (morphisme) de préschémas formels : III, 3.4.1.

Représentable (foncteur) : 0, 8.1.8.

Résolution cohomologique, droite, gauche, homologique, injective, libre,
plate, projective : 0, 11.4.1.

Résolution de Cartan-Eilenberg droite, injective : 0, 11.4.2.

Résolution de Cartan-Eilenberg gauche, projective : 0, 11.6.1.

Rétrocompact (ensemble) : 0, 9.1.1.

Strict (système projectif) : 0, 13.4.2.

Suite spectrale dans une catégorie abélienne : 0, 11.1.1.

Suite spectrale faiblement convergente, régulière, corégulière, birégulière :
0, 11.1.3.

Suite spectrale dégénérée : 0, 11.1.6.

Suite spectrale d'un foncteur relativement à un objet filtré : 0, 13.6.4.

Suites spectrales d'un bicomplexe : 0, 11.3.2 et 11.3.5.

Suites spectrales d'hypercohomologie d'un foncteur par rapport à un complexe :
0, 11.4.3.

Suites spectrales d'hyperhomologie d'un foncteur par rapport à un complexe :
0, 11.6.2.

Suites spectrales d'hyperhomologie d'un foncteur par rapport à un bicomplexe :
0, 11.7.2.

Système de coefficients : 0, 11.8.4.

Unibranche (anneau, point) : III, 4.3.6.

Universellement ouvert (morphisme) : III, 4.3.9.
\end{flushleft}

\clearpage
\oldpage[III]{165}
\section*{TABLE DES MATIÈRES}

\begin{flushleft}
\textsc{Chapitre 0.} --- \textbf{Préliminaires (suite)} \dotfill 5

\medskip
\S~8. Foncteurs représentables \dotfill 5

8.1. Foncteurs représentables \dotfill 5

8.2. Structures algébriques dans les catégories \dotfill 9

\medskip
\S~9. Ensembles constructibles \dotfill 12

9.1. Ensembles constructibles \dotfill 12

9.2. Ensembles constructibles dans les espaces noethériens \dotfill 14

9.3. Fonctions constructibles \dotfill 16

\medskip
\S~10. Compléments sur les modules plats \dotfill 17

10.1. Relations entre modules plats et modules libres \dotfill 17

10.2. Critères locaux de platitude \dotfill 18

10.3. Existence d'extensions plates d'anneaux locaux \dotfill 20

\medskip
\S~11. Compléments d'algèbre homologique \dotfill 23

11.1. Rappels sur les suites spectrales \dotfill 23

11.2. La suite spectrale d'un complexe filtré \dotfill 27

11.3. Les suites spectrales d'un bicomplexe \dotfill 29

11.4. Hypercohomologie d'un foncteur par rapport à un complexe $K^\bullet$
\dotfill 32

11.5. Passage à la limite inductive dans l'hypercohomologie \dotfill 35

11.6. Hyperhomologie d'un foncteur par rapport à un complexe $K_\bullet$
\dotfill 39

11.7. Hyperhomologie d'un foncteur par rapport à un bicomplexe
$K_{\bullet\bullet}$ \dotfill 41

11.8. Compléments sur la cohomologie des complexes simpliciaux \dotfill 43

11.9. Un lemme sur les complexes de type fini \dotfill 46

11.10. Caractéristique d'Euler-Poincaré d'un complexe de modules de longueur
finie \dotfill 48

\medskip
\S~12. Compléments sur la cohomologie des faisceaux \dotfill 49

12.1. Cohomologie des faisceaux de modules sur les espaces annelés \dotfill 49

12.2. Images directes supérieures \dotfill 57

12.3. Compléments sur les foncteurs Ext de faisceaux \dotfill 60

12.4. Hypercohomologie du foncteur image directe \dotfill 62

\oldpage[III]{166}

\S~13. Limites projectives en algèbre homologique \dotfill 64

13.1. La condition de Mittag-Leffler \dotfill 64

13.2. La condition de Mittag-Leffler pour les groupes abéliens \dotfill 65

13.3. Application~: cohomologie d'une limite projective de faisceaux \dotfill 68

13.4. Condition de Mittag-Leffler et objets gradués associés aux systèmes
projectifs \dotfill 69

13.5. Limites projectives de suites spectrales de complexes filtrés \dotfill 71

13.6. Suite spectrale d'un foncteur relative à un objet muni d'une filtration
finie \dotfill 73

13.7. Foncteurs dérivés d'une limite projective d'arguments \dotfill 75

\medskip
\textsc{Chapitre III.} --- \textbf{Étude cohomologique des faisceaux
cohérents} \dotfill 81

\medskip
\S~1. Cohomologie des schémas affines \dotfill 82

1.1. Rappels sur le complexe de l'algèbre extérieure \dotfill 82

1.2. Cohomologie de Čech d'un recouvrement ouvert \dotfill 85

1.3. Cohomologie d'un schéma affine \dotfill 88

1.4. Application à la cohomologie des préschémas quelconques \dotfill 89

\medskip
\S~2. Étude cohomologique des morphismes projectifs \dotfill 95

2.1. Calculs explicites de certains groupes de cohomologie \dotfill 95

2.2. Le théorème fondamental des morphismes projectifs \dotfill 100

2.3. Application aux faisceaux gradués d'algèbres et de modules \dotfill 102

2.4. Une généralisation du théorème fondamental \dotfill 107

2.5. Caractéristique d'Euler-Poincaré et polynôme de Hilbert \dotfill 109

2.6. Application~: critères d'amplitude \dotfill 111

\medskip
\S~3. Le théorème de finitude pour les morphismes propres \dotfill 115

3.1. Le lemme de dévissage \dotfill 115

3.2. Le théorème de finitude~: cas des schémas usuels \dotfill 116

3.3. Généralisation du théorème de finitude (schémas usuels) \dotfill 118

3.4. Le théorème de finitude~: cas des schémas formels \dotfill 119

\medskip
\S~4. Le théorème fondamental des morphismes propres. Applications \dotfill 122

4.1. Le théorème fondamental \dotfill 122

4.2. Cas particuliers et variantes \dotfill 129

4.3. Le théorème de connexion de Zariski \dotfill 130

4.4. Le «~main theorem~» de Zariski \dotfill 135

4.5. Complétés de modules d'homomorphismes \dotfill 138

4.6. Relations entre morphismes formels et morphismes usuels \dotfill 139

4.7. Un critère d'amplitude \dotfill 145

4.8. Morphismes finis de préschémas formels \dotfill 146

\oldpage[III]{167}

\S~5. Un théorème d'existence de faisceaux algébriques cohérents \dotfill 149

5.1. Énoncé du théorème \dotfill 149

5.2. Démonstration du théorème d'existence~: cas projectif et quasi-projectif
\dotfill 151

5.3. Démonstration du théorème d'existence~: cas général \dotfill 154

5.4. Application~: comparaison de morphismes de schémas usuels et de
morphismes de schémas formels. Schémas formels algébrisables \dotfill 156

5.5. Une décomposition de certains schémas \dotfill 159

\medskip
\textsc{Bibliographie (suite)} \dotfill 161

\textsc{Index des notations} \dotfill 162

\textsc{Index terminologique} \dotfill 163

\begin{flushright}
\emph{Reçu le 28 juin 1961.}
\end{flushright}
\end{flushleft}
